\documentclass[11pt,oneside]{book}

\usepackage[a4paper,margin=1in]{geometry}

\usepackage[T1]{fontenc}

\usepackage[utf8]{inputenc}

\usepackage{amsmath,amssymb}

\usepackage{booktabs,longtable,array}

\usepackage{xcolor}

\usepackage{titlesec}

\usepackage[most]{tcolorbox}

\usepackage{hyperref}

\hypersetup{colorlinks=true,linkcolor=blue!55!black,urlcolor=blue!55!black,citecolor=blue!55!black}

\pagestyle{plain}

\emergencystretch=3em

\definecolor{quantumgreen}{HTML}{3F6F5A}

\definecolor{softgray}{HTML}{F5F6F2}

\definecolor{ink}{HTML}{1C1E1A}

\titleformat{\chapter}[display]{\normalfont\huge\bfseries\color{ink}\raggedright}{\chaptertitlename\ \thechapter}{14pt}{\Huge}

\newtcolorbox{claimbox}{colback=softgray,colframe=quantumgreen!70!black,arc=2pt,boxrule=0.7pt,left=8pt,right=8pt,top=6pt,bottom=6pt}

\newenvironment{centeredalign}{\[\begin{gathered}[c]}{\end{gathered}\]}

\newcommand{\ket}[1]{\lvert #1\rangle}

\newcommand{\bra}[1]{\langle #1\rvert}

\begin{document}

\frontmatter

\hypersetup{pageanchor=false}

\begin{titlepage}

\centering

\vspace*{1.7cm}

{\Huge\bfseries Quantum Theory\par}

\vspace{0.12cm}

{\Huge\bfseries As A Mechanism Tree\par}

\vspace{0.3cm}

\vspace{1.2cm}

{\large Synthetix Institute\par}

{\large Edition dated 2026-10-15 07:24 UTC\par}

\vfill

\begin{claimbox}

\noindent Interactions generate phases, correlations and transitions. Their observable consequences connect the preparation of a quantum system to its later behaviour. This book develops those mechanisms through equations and follows their realization in different physical systems.

\end{claimbox}

\vspace{0.5cm}

\begin{minipage}{0.94\linewidth}\small\raggedright

\textbf{Scope and authorship.} This book develops established quantum theory through its physical mechanisms. The mechanism tree and its cross-topic connections are the authors' organization of those mechanisms. References identify the particular relations used in each treatment.

\end{minipage}

\end{titlepage}

\hypersetup{pageanchor=true}

\tableofcontents

\chapter*{About This Edition}

This book develops quantum theory through worked calculations. Each topic begins with a physical question, derives the relation that answers it, and follows that relation to a measurable consequence. The opening part introduces the ingredients that every prediction needs, shows how a fixed background becomes part of the system, and follows what is left behind when a state is carried around a loop or a law is carried into another system. The grouping of topics into chapters is the authors' choice; the physics in each treatment is established theory, and the calculations are written for this volume.

The book contains 125 worked treatments together with 17 historical, bibliographic or interpretive entries and 4 alternative names. 

References follow the calculations they support. Links headed Relations In The Original Papers identify particular displays in those articles; Equation Sources identifies relations aligned to corpus records.

\mainmatter

\part{Mechanisms And Predictions}

\chapter{Quantum Mechanisms And Their Predictions}\label{chap:mechanisms-predictions}

A spin prepared perpendicular to a magnetic field precesses, and its
transverse polarization oscillates. Coupling it to a second spin changes
the problem: the field acting on the first now depends on the second.
The measured polarization can fall and return even though the pair is
isolated and no energy is dissipated. A controlled experiment uses this
interaction to generate correlations; an experiment that leaves the
neighbour uncontrolled sees a loss of coherence. Both observations follow
from the same joint evolution.

Consider two spin-1/2 systems with Hamiltonian \(H=\hbar gZ_1Z_2\), where
\(g\) is an angular frequency. The dimensionless Pauli operators
\(X_j,Y_j,Z_j\) represent the components of spin \(j\), in units of
\(\hbar/2\). Each \(Z_j\) has eigenvalues \(+1\) and \(-1\), and
\([Z_j,X_j]=2iY_j\). When spin 2 is in a \(Z_2\) eigenstate, spin 1
experiences the Hamiltonian \(+\hbar gZ_1\) or \(-\hbar gZ_1\).
It therefore precesses in opposite directions in the two cases. A
superposition of those alternatives can entangle the spins.

Let the apparatus measure the transverse polarization of spin 1,
\(x(t)=\operatorname{Tr}[\rho(t)X_1]\), where \(\rho\) is the joint
density operator. Its time derivative follows from the Heisenberg
equation \(d\langle O\rangle/dt=(i/\hbar)\langle[H,O]\rangle\).
The commutator contains the joint operator \(Y_1Z_2\), so a correlation
\(c(t)=\operatorname{Tr}[\rho(t)Y_1Z_2]\) enters the measured motion.
Taking its derivative returns the original polarization:
\begin{centeredalign}
\frac{i}{\hbar}[H,X_1]=-2gY_1Z_2,\qquad
\frac{i}{\hbar}[H,Y_1Z_2]=2gX_1,\\
\dot x=-2gc,\qquad \dot c=2gx,\\
x(t)=x(0)\cos(2gt)-c(0)\sin(2gt).
\end{centeredalign}
The two expectation values form a closed pair of equations: their
derivatives require no further operator. This is a smaller description
than the full two-spin density matrix, yet it determines this
polarization exactly. For an initial state with \(x(0)=1\) and
\(c(0)=0\), the apparatus observes \(x(t)=\cos(2gt)\). At
\(t=\pi/(4g)\), for \(g>0\), the transverse polarization is zero
and the correlation has reached one. The vanishing signal represents
a transfer into a joint correlation, rather than disappearance of the
joint state. The polarization later returns.

Measuring \(Z_1\) would give a constant expectation, because
\([H,Z_1]=0\). Thus the same Hamiltonian generates visible motion
for one observable and no motion for another. The interaction and
the measured quantity together determine which part of the state
must be followed. In a larger spin network, repeated commutators
can introduce further correlations; the number required by a
particular measurement is a physical property of the interactions,
not the number of topics used to describe them.

The initial correlation matters independently of the initial
polarization. Let \(I_4\) denote the identity on the two-spin space
and choose real \(x_0,c_0\) with \(x_0^2+c_0^2\leq1\). The two
preparations
\begin{centeredalign}
\rho_\pm=\tfrac14(I_4+x_0X_1\pm c_0Y_1Z_2),\qquad
\operatorname{Tr}_2\rho_\pm=\tfrac12(I_2+x_0X_1)
\end{centeredalign}
give the same reduced density operator for spin 1. Every immediate
measurement on that spin consequently has identical statistics.
The states are physical: \(X_1\) and \(Y_1Z_2\) anticommute and
square to the identity, giving eigenvalues
\((1\pm\sqrt{x_0^2+c_0^2})/4\). After the interaction begins,
their polarizations differ by \(-2c_0\sin(2gt)\). In particular,
\begin{centeredalign}
x(0)=x_0\ \text{for both},\qquad \dot x(0)=\mp2gc_0.
\end{centeredalign}
No evolution law depending only on \(x_0\) can reproduce both initial slopes.
The two preparations are distinguished by a joint quantity that the
interaction converts into a local signal. Predicting that signal therefore
requires \(c_0\), even though the apparatus measures only \(x\).

Integrating \(\dot c=2gx\) gives
\(c(t)=c_0+2g\int_0^t x(s)\,\mathrm ds\). Substituting this into
the polarization equation eliminates the correlation without
losing its effect:
\begin{centeredalign}
\dot x(t)=-2gc_0-4g^2\int_0^t x(s)\,\mathrm ds.
\end{centeredalign}
The first term carries the correlation present at preparation.
The integral carries the correlation produced during the motion.
Setting \(c_0=0\) removes the former, but leaves the history
dependence. The two first-order equations and this memory equation
predict the same signal because both retain the physical influence
of the second expectation value.

These calculations explain what is meant here by a mechanism: the
interaction through which a specified preparation produces a
measurable change. Its description has dependent parts. The joint
spin space makes \(Z_1Z_2\) and \(Y_1Z_2\) meaningful operators;
the Hamiltonian generates the motion; positivity and normalization
make the preparation physical; the measured polarization selects
the required expectations. Preparing the state, applying the
interaction for a given duration and choosing the measurement
complete the experiment. A particular pair of spins supplies the
coupling \(g\), whose value fixes the time scale.

The dependence of these parts is the reason for organizing quantum
theory as a mechanism tree. An operator acts on a chosen space of
states. Its spectrum and evolution depend on its domain, including
boundary conditions. A prediction uses that evolution together
with a preparation and a measured operator. Material parameters
then specify an experimental realization of this problem. The
tree records these mathematical and physical dependencies; it
does not prescribe a chronological order in which a researcher
must choose them. A change in an outer part can require a change
within: replacing a measuring apparatus that observes \(X_1\)
by one that observes \(Z_1\) changes which equations are needed,
while adding a new interaction can enlarge their closed set.

The following chapters use this dependence to connect familiar
subjects. Confinement fixes an operator domain and changes energy
levels. Quantizing a field introduces states whose fluctuations
act on an atom. Eliminating an environment retains its influence
as decoherence or memory. Changing representation moves states
and observables together while preserving predictions. These
relations explain how different physical systems can realize the
same equations and where a changed interaction, boundary or
preparation requires another term. The individual topic
treatments develop the spectra, interference patterns and
correlations through which those changes are measured.



\chapter{What A Quantum Prediction Requires}\label{chap:physical-identity}
\begin{claimbox}
A quantum prediction is fixed only when the possible states, the law of motion, the conditions that make that law
well defined, the measured quantity and the order of operations are all specified. Changing any one of them can change
the predicted signal while the others stay the same.
\end{claimbox}

\section{One Equation, Two Spectra}
A particle of mass \(m\) confined to an interval of length \(L\) obeys the same differential equation,
\(-(\hbar^2/2m)\,\psi''=E\psi\), whether its ends are hard walls or ends at which the wave function has zero slope.
The two situations nevertheless have different energies. Hard walls require \(\psi(0)=\psi(L)=0\) and admit only sine
waves with a whole number of half-wavelengths, whereas zero-slope ends require \(\psi'(0)=\psi'(L)=0\) and admit cosine
waves, including the constant function:
\begin{centeredalign}
E_n^{\mathrm{wall}}=\frac{n^2\pi^2\hbar^2}{2mL^2},\quad n=1,2,\ldots\qquad
E_n^{\mathrm{slope}}=\frac{n^2\pi^2\hbar^2}{2mL^2},\quad n=0,1,2,\ldots
\end{centeredalign}
A weak probe coupled to position can excite the lowest state to the next
one in either case. The energy absorbed is
\(3\pi^2\hbar^2/(2mL^2)\) for hard walls and
\(\pi^2\hbar^2/(2mL^2)\) for zero-slope ends. The transition energies
differ by a factor of three because the second problem admits the
constant, zero-energy state. The differential expression is the same,
but the set of wave functions on which it acts, its domain, is different.
A boundary condition therefore enters the law of motion itself, and
spectroscopy distinguishes the two choices.

\section{What Each Part Of A Prediction Contributes}
The two coupled spins of Chapter~\ref{chap:mechanisms-predictions} show the same dependence on every ingredient of a
calculation. The states are density operators of two spins, and the law of motion is generated by \(H=\hbar gZ_1Z_2\).
Positivity and unit trace restrict which operators describe physical preparations, just as the boundary condition
restricted the wave functions of the confined particle. The measured quantity matters independently: the transverse
polarization \(\langle X_1\rangle\) oscillates at the angular frequency \(2g\), while \(\langle Z_1\rangle\) does not
change at all, because \(Z_1\) commutes with \(H\). The interaction is the same in both experiments; only the question
put to the system differs.

Preparing spin 1 along \(+z\) and applying two quarter-turn pulses
gives another experimentally distinguishable change. Define
\(U_x=e^{-i\pi X_1/4}\) and \(U_z=e^{-i\pi Z_1/4}\). The first
sequence \(U_zU_x\) turns the spin toward \(+x\), whereas
\(U_xU_z\) turns it toward \(-y\). A final measurement along \(x\)
gives
\begin{centeredalign}
P(+x\mid U_zU_x)=1,\qquad P(+x\mid U_xU_z)=\tfrac12.
\end{centeredalign}
The initial state, pulse angles and measured component are identical;
the different probabilities come from the order of the two rotations.
Their specification, including the duration of any intervening
evolution, is the experimental protocol. Numerical values such as
\(g\), \(m\) and \(L\) then fix frequencies or energies within the
chosen physical model. Opening a box onto the free line changes more
than a value: it replaces the domain and its discrete spectrum by a
different physical problem.

\section{When Two Descriptions Describe The Same Physics}
Two descriptions of a system are equivalent when translating a state from one to the other and then evolving it gives
the same result as evolving first and translating afterwards. A mechanical oscillator and an electrical LC circuit
provide the standard example. Mapping the position \(x\) to the charge \(Q\), the momentum \(p\) to the magnetic flux
\(\Phi\), the mass \(m\) to the inductance \(L\) and the spring constant to \(1/C\) carries Hamilton's equations of one
system into those of the other. The quantum spectra then agree term by term, and the zero-point fluctuations of the
charge follow from those of the position:
\begin{centeredalign}
E_n=\hbar\omega\left(n+\tfrac12\right),\qquad \omega=\frac{1}{\sqrt{LC}},\qquad
\langle Q^2\rangle_0=\frac{\hbar}{2L\omega}=\frac{\hbar\omega C}{2}.
\end{centeredalign}
The translation fails when the second description cannot follow the first. Chapter~\ref{chap:mechanisms-predictions}
gave an example: describing the two spins by the polarization \(x\) of spin 1 alone, two preparations with the same
\(x_0\) but opposite correlations \(c_0\) start with opposite slopes \(\mp2gc_0\). Evolving and then translating
differs from translating and then evolving, and the difference is exactly the correlation that the reduced
description discarded. Restoring it as a second variable makes the translation exact again. The failure of a
translation, when it has a definite form, points to the physical quantity that a smaller description is missing; the
chapters that follow use this observation repeatedly.



\chapter{When A Fixed Background Becomes Part Of The System}\label{chap:background}
\begin{claimbox}
An atom can exchange energy with a cavity, and a spin can exchange
correlations with its surroundings. Treating the field or environment
as fixed discards those changes. Their dynamics becomes part of the
model when it affects the measurement being predicted.
\end{claimbox}

\section{A Test For Missing Variables}
A description of a quantum system is complete for a given set of experiments when equal present states always lead to
equal future measurement statistics. If two histories \(h_1\) and \(h_2\) prepare the same state \(q\) but a later
measurement of an outcome \(y\) distinguishes them,
\begin{centeredalign}
q(h_1)=q(h_2)\qquad\text{but}\qquad p(y,t\mid h_1)\ne p(y,t\mid h_2),
\end{centeredalign}
then the state has discarded something that still acts. The two-spin preparations \(\rho_\pm\) of
Chapter~\ref{chap:mechanisms-predictions} are the simplest case: both give spin 1 the same reduced state, yet their
later polarizations differ by \(-2c_0\sin(2gt)\). The missing variable there is a correlation with a second spin. The
sections below show the same test at work for a field, an environment, a boundary and a sequence of pulses.

\section{A Classical Field Becomes A Quantized Mode}
An atom with transition frequency \(\omega_0\) driven on resonance by a classical field of Rabi frequency \(\Omega_R\)
oscillates between its ground and excited states as \(P_e(t)=\sin^2(\Omega_Rt/2)\). A field of zero amplitude does
nothing to it. When the field is a single mode of a cavity, coupled to the atom with strength \(g\), the Jaynes--Cummings
Hamiltonian
\begin{centeredalign}
H=\hbar\omega a^\dagger a+\tfrac12\hbar\omega_0\sigma_z+\hbar g\,(a\sigma_++a^\dagger\sigma_-)
\end{centeredalign}
couples the excited atom with \(n\) photons only to the ground atom with \(n+1\) photons. On resonance the excited
population then oscillates at a frequency that depends on the photon number,
\begin{centeredalign}
P_e(t)=\cos^2\!\left(g\sqrt{n+1}\,t\right),
\end{centeredalign}
so that an excited atom in an empty cavity exchanges its energy with the vacuum at the angular frequency \(2g\). A
classical field of zero amplitude predicts no such exchange. The vacuum oscillation, and the collapse and revival that a
coherent field produces because its photon numbers oscillate at different frequencies, were observed with Rydberg atoms
crossing a microwave cavity (Brune et al., Phys. Rev. Lett. 76, 1800, 1996); the treatment of quantum optics works them
out. The field had to become part of the system because atoms prepared identically but sent through cavities holding
different photon-number distributions behaved differently.

\section{An Environment Becomes A Source Of Decay And Of Memory}
Let spin 1 of Chapter~\ref{chap:mechanisms-predictions} interact with several unpolarized neighbours through
\(H=\hbar Z_1\sum_kg_kZ_k\). Each neighbour, found up or down with equal probability, shifts the precession of spin 1 in
one sense or the other, and averaging over their orientations gives the transverse polarization
\begin{centeredalign}
x(t)=x_0\prod_k\cos(2g_kt)\;\approx\;x_0\exp\!\Big(-2t^2\sum_kg_k^2\Big).
\end{centeredalign}
With one neighbour, \(x\) changes sign at \(t=\pi/(2|g_1|)\)
and returns to its initial signed value at \(t=\pi/|g_1|\).
The magnitude of the transverse polarization has already recovered
at the first time. With several unequal couplings these revivals can
occur at different times. For \(|g_kt|\ll1\), expanding each cosine
gives the Gaussian in the second expression and the decay time
\(T_2^*=(2\sum_kg_k^2)^{-1/2}\). This short-time decay resembles
free-induction decay in nuclear magnetic resonance. The exact finite
product still permits recurrences; it does not establish irreversible
loss. The neighbours retain the conditional phases even when their
average contribution to the measured polarization is small.

\section{A Pulse Sequence Becomes Part Of The Dynamics}
The decay above is not irreversible. A pulse that flips spin 1 halfway through the evolution reverses the sense of
every conditional precession, so that the phases accumulated before the pulse are undone after it, and the full
polarization returns at twice the pulse time. This is the spin echo. The final signal depends on when the pulse acts:
the static couplings have not changed, but the sequence of operations has, and with it the effective evolution. In
quantum computing, error correction and periodically driven matter, the order and timing of operations are in the
same way part of the physics implemented in an experiment.

\section{A Boundary Condition Carries A Field}
An electron confined to a ring of radius \(R\) that encloses a magnetic flux \(\Phi\) experiences no magnetic field
along its path. Its energies nevertheless depend on the flux,
\begin{centeredalign}
E_n=\frac{\hbar^2}{2mR^2}\left(n+\frac{\Phi}{\Phi_0}\right)^2,\qquad \Phi_0=\frac{h}{e},\qquad n\in\mathbb Z,
\end{centeredalign}
because a gauge transformation that removes the vector potential from the ring moves it into the condition
\(\psi(\varphi+2\pi)=e^{2\pi i\Phi/\Phi_0}\psi(\varphi)\) imposed on the wave function. The spectrum repeats with
period \(\Phi_0\), and the conductance of small metal rings oscillates with the same period (Webb et al., Phys. Rev.
Lett. 54, 2696, 1985). The enclosed flux, a background quantity for a classical particle, enters the quantum
description through a boundary condition. The chapters on boundaries develop this dependence of spectra on domains for
tunnelling, scattering and confined particles.



\chapter{Loops: Geometric Phase, Frustration And Memory}\label{chap:loops}
\begin{claimbox}
Going around a closed loop can leave a trace: a phase acquired by a transported state, a constraint that cannot be
satisfied by every bond of a magnet, or a dependence of a reduced description on the path by which it was reached.
Each trace is a definite physical quantity that a purely local description misses.
\end{claimbox}

\section{A Spin Carried Around A Cone Acquires A Phase}
A spin-1/2 aligned with a magnetic field follows the field adiabatically if the direction changes slowly. When the
field direction is carried once around a cone of half-angle \(\theta\) and returns to its start, the spin returns to
its initial state up to a phase. Besides the dynamical phase set by the energy, the state acquires a geometric phase
equal to minus half the solid angle enclosed by the path:
\begin{centeredalign}
\gamma=-\tfrac12\,\Omega_{\mathrm{solid}}=-\pi(1-\cos\theta).
\end{centeredalign}
The phase depends only on the geometry of the loop, not on how fast it was traversed, and it is measured by letting the
transported state interfere with a reference (Berry, Proc. R. Soc. A 392, 45, 1984). In gauge theories the same
structure appears in space-time: the phase acquired around a small loop is set by the field strength enclosed by it,
\begin{centeredalign}
W(\gamma)=\operatorname{Tr}\,\mathcal P\exp\!\left(-ig\oint_\gamma A_\mu\,dx^\mu\right),\qquad
\mathcal P\exp\!\left(-ig\oint_{\square}A\right)=I-igF_{\mu\nu}\ell^2+O(\ell^3),
\end{centeredalign}
for a square of side \(\ell\) in the \(\mu\nu\) plane, with the covariant derivative \(D_\mu=\partial_\mu+igA_\mu\). The
Aharonov--Bohm ring of the previous chapter is the case of an abelian field in which the enclosed flux, not the local
field, fixes the phase.

\section{Three Spins That Cannot All Be Satisfied}
Three Ising spins \(s_i=\pm1\) on the corners of a triangle, coupled antiferromagnetically by
\(H=J(s_1s_2+s_2s_3+s_3s_1)\) with \(J>0\), would each like to point opposite to both neighbours. No assignment achieves
this for all three bonds at once: going around the triangle, three sign reversals cannot return to the starting sign.
The best configurations satisfy two bonds and leave one unsatisfied, giving the ground-state energy \(-J\) instead of
\(-3J\), and six of the eight configurations share it. In general, for bonds of either sign in
\(H=-\sum_{\langle ij\rangle}J_{ij}s_is_j\), a loop is frustrated when
\begin{centeredalign}
\prod_{\langle ij\rangle\in\gamma}\operatorname{sign}J_{ij}=-1,
\end{centeredalign}
and no relabelling of individual spins can remove this sign. In the
triangle, choosing which bond remains unsatisfied gives three possibilities,
each with its globally reversed partner. Extended frustrated lattices can
similarly have many minimizing configurations. Whether their degeneracy
grows extensively, and hence leaves a nonzero residual entropy per spin,
depends on the lattice and its interactions as well as the frustrated loops.

\section{A Reduced Description Remembers Its Path}
The polarization of spin 1 in Chapter~\ref{chap:mechanisms-predictions},
for a preparation without correlation, evolves as
\(x(t)=x_0\cos(2gt)\). At \(t=\pi/(4g)\), for \(g>0\), all
these preparations give \(x=0\), although their correlations are
\(c=x_0\). The subsequent slope is \(-2gx_0\). Knowing the present
polarization alone therefore does not determine its subsequent motion.
The reduced map for this preparation family is multiplication by
\(\cos(2gt)\), and its composition gives
\begin{centeredalign}
V_{t+s}\ne V_tV_s,\qquad
\cos[2g(t+s)]\ne\cos(2gt)\cos(2gs)
\end{centeredalign}
at general times. The larger rotation of \((x,c)\) does compose:
its final state contains the correlation required to start the next
interval. Eliminating \(c\) instead leaves the memory integral and
initial-correlation term derived in the opening chapter.

Order can also be recorded by applying two different operations to a
state. If \(W_BW_A\) and \(W_AW_B\) produce different states, a
later measurement can distinguish their order only to the extent that
the subsequent relaxation retains that difference and the measured
operator is sensitive to it. Noncommuting operations can create a
record; a relaxation that erases the difference, or a measurement
orthogonal to it, removes the record from the observed signal.

\section{What The Three Loops Share}
In each case local information fails to determine a global outcome. A connection fixes how a state is carried from point
to point, and its curvature makes the result depend on the path. Signed bonds fix local preferences whose product around
a loop can forbid satisfying all of them. A reduced state fixes present expectation values but not the correlations that
will feed back on them. The three objects are mathematically distinct, and the comparison is a guide for where to look
for missing quantities rather than an identity between them.


\chapter{Equivalence Across Quantum Descriptions}


An energy spectrum predicts both the results of an energy measurement and
the phases acquired during evolution. Changing from a moving state to a
moving observable leaves those predictions unchanged. By contrast,
discarding a correlated subsystem or changing a boundary can change the
measurement itself. Quantum theory connects its descriptions through these
particular mathematical relations. The following calculations explain
which relations are identities, which retain only specified observations,
and which define a changed physical problem.


\section{One operator, two roles: generator and observable}


Preparing a system in an energy eigenstate gives a definite energy, but its
state vector still evolves. For a time-independent, self-adjoint Hamiltonian
\(H\), the eigenvalue equation \(H|n\rangle=E_n|n\rangle\) gives
\(U(t)|n\rangle=e^{-iE_nt/\hbar}|n\rangle\). The phase of an individual
eigenstate leaves every measurement probability unchanged. A superposition of
two energies accumulates a relative phase, which an observable connecting the
two states can reveal.

For a two-level system with \(H=\hbar\omega Z/2\), prepare
\(|+x\rangle=(|+z\rangle+|-z\rangle)/\sqrt2\). An energy measurement gives
\(\pm\hbar\omega/2\) with equal probabilities. A later measurement of \(X\)
instead detects the relative phase:
\begin{centeredalign}
U(t)=e^{-iHt/\hbar},\qquad
\langle X\rangle_t=\cos(\omega t),\\
H=\int_{\sigma(H)}E\,dP_H(E),\qquad
U(t)=\int_{\sigma(H)}e^{-iEt/\hbar}\,dP_H(E).
\end{centeredalign}
Here \(P_H\) is the projection-valued spectral measure of \(H\); it assigns
a projector to each energy interval. The probability of an energy in an
interval \(B\) is \(\operatorname{Tr}[\rho P_H(B)]\). The same projectors
weight the phases in the evolution operator. Energy as a measured quantity
and energy as the generator of time evolution are therefore two uses of the
same operator, rather than two mechanisms inferred from similar notation.

For a particle in a box, the boundary conditions enter the domain of \(H\)
and hence both integrals. Changing those conditions changes the allowed
energies and the evolution together. The treatments of
\hyperref[topic:hamiltonian-quantum-mechanics]{the Hamiltonian} and
\hyperref[topic:spectral-theory]{spectral theory} develop the two uses in
greater detail.


\section{Schrödinger, Heisenberg, and path-integral representations}


The oscillation in the preceding section can be calculated by evolving the
prepared spin state or by evolving its measured operator. These are the
Schr\"odinger and Heisenberg pictures. Starting from the same density operator
\(\rho_0\) and unitary evolution \(U(t)\), they give
\begin{centeredalign}
\rho_S(t)=U(t)\rho_0U(t)^\dagger,\qquad
A_H(t)=U(t)^\dagger A_SU(t),\\
\operatorname{Tr}[\rho_S(t)A_S]
=\operatorname{Tr}[\rho_0A_H(t)].
\end{centeredalign}
The equality follows by moving the final \(U^\dagger\) around the trace.
For \(H=\hbar\omega Z/2\), the Heisenberg operator is
\(X_H(t)=X\cos(\omega t)-Y\sin(\omega t)\). Its expectation in
\(|+x\rangle\) gives the same cosine as the moving state. The measured
polarization has not changed; only the location of time dependence has.

For a particle, a third representation follows by splitting the evolution
time into short intervals and inserting complete position states between
them. Each intermediate position is summed over. In the limit, the
propagator is expressed as an integral over paths joining the prescribed
endpoints:
\begin{centeredalign}
\langle q_f|U(t_f-t_i)|q_i\rangle
=\int_{q(t_i)=q_i}^{q(t_f)=q_f}\mathcal Dq\,e^{iS[q]/\hbar}.
\end{centeredalign}
The measure includes the normalization inherited from the short-time
propagators. For a free particle of mass \(m\), evaluating the Gaussian
integrals gives, for \(T>0\),
\begin{centeredalign}
\langle q_f|U(T)|q_i\rangle
=\sqrt{\frac{m}{2\pi i\hbar T}}
\exp\!\left[\frac{im(q_f-q_i)^2}{2\hbar T}\right].
\end{centeredalign}
This kernel also solves the free Schr\"odinger equation and tends to a delta
function as \(T\) approaches zero. Boundary conditions or constraints must
enter the short-time kernels before the limit is taken. Thus the path
integral represents a specified Hamiltonian evolution; an action written
without its measure and domain does not yet specify the same problem.
\hyperref[topic:path-integral-formulation]{The path-integral treatment}
explains the construction for a potential and for imaginary time.


\section{The choice of subsystems defines locality}


When a detector acts on one spin of a pair, its operator has the form
\(O_A\otimes I_B\). The tensor-product decomposition identifies the two
subsystems and therefore what the detector can measure locally. For a joint
state \(\rho_{AB}\), all such measurements are reproduced by its reduced
state:
\begin{centeredalign}
\mathcal H=\mathcal H_A\otimes\mathcal H_B,\qquad
\rho_A=\operatorname{Tr}_B\rho_{AB},\\
\operatorname{Tr}[\rho_{AB}(O_A\otimes I_B)]
=\operatorname{Tr}(\rho_AO_A).
\end{centeredalign}
The singlet \((|01\rangle-|10\rangle)/\sqrt2\) gives \(\rho_A=I/2\).
Every direction of spin measurement is equally likely on either spin,
although measurements on the pair are perfectly anticorrelated along a
common axis. The reduced state has retained local probabilities and
discarded joint information. During an interaction, that joint information
can return to a local measurement, as the polarization--correlation pair of
Chapter~\ref{chap:mechanisms-predictions} shows.

A change of basis preserves the full experiment when it transforms both
states and operators. For a unitary \(V\), the transformed objects are
\(V\rho V^\dagger\) and \(VOV^\dagger\), and their trace is unchanged.
An entangling \(V\) generally turns a local operator into a joint operator.
If only the state is transformed while the same physical detector is kept,
the experiment has changed. This is why an arbitrary change of tensor
factors can change which states are called entangled without changing the
abstract dimension of the total Hilbert space.

In a spatial experiment the interactions and accessible detectors determine
the relevant subsystem algebras. In a quantum field theory, locality is
expressed by observable algebras associated with regions; a simple tensor
factorization for every region need not exist. The common physical question
is which measurements belong to each subsystem and which correlations
their combination can expose. The chapters on
\hyperref[topic:quantum-entanglement]{entanglement} and
\hyperref[topic:quantum-field-theory]{fields} develop these distinct settings.


\section{Constraints define physical observables}


A redundant coordinate description assigns several mathematical states to
one physical state. Gauge constraints remove this redundancy. In a Dirac
description, constraint operators \(\widehat{\mathcal C}_a\) annihilate
physical states. An operator representing a physical measurement must act
within that state space:
\begin{centeredalign}
\widehat{\mathcal C}_a|\psi_{\rm phys}\rangle=0,\qquad
[\widehat O_{\rm phys},\widehat{\mathcal C}_a]
|\psi_{\rm phys}\rangle=0.
\end{centeredalign}
The second relation follows by applying the first to
\(\widehat O_{\rm phys}|\psi_{\rm phys}\rangle\). It is a condition on
physical states; it need not be an operator identity throughout a larger
auxiliary space. In continuum gauge theories, defining the physical inner
product and the constraint domains is also part of the construction.

A finite example makes the restriction explicit. In a two-spin space impose
\(\mathcal C=Z_1Z_2-I\). The allowed subspace is spanned by
\(|00\rangle\) and \(|11\rangle\). The single-spin flip \(X_1\) sends
either state outside this subspace. The paired flip \(X_1X_2\) interchanges
the two allowed states and commutes with the constraint. Consequently a
Hamiltonian \(H=\hbar gX_1X_2\) gives physical motion inside the subspace:
\begin{centeredalign}
e^{-iHt/\hbar}|00\rangle
=\cos(gt)|00\rangle-i\sin(gt)|11\rangle,\qquad
P(11,t)=\sin^2(gt).
\end{centeredalign}
The example is an algebraic constrained system, not a model of an
electromagnetic gauge field. In an actual gauge theory, Gauss's law similarly
restricts the allowed matter and field states, and observables such as
electric fields or closed-loop quantities respect the restriction.

The measured probability is defined only after this choice of physical
states. For a positive measurement effect \(E_y\), its restriction to the
physical space must remain positive and the complete set of effects must
sum to the identity on that space. Constraints, evolution and measurement
therefore enter one calculation. Their roles in electromagnetic fields and
in gravitational theories are developed under
\hyperref[topic:gauge-theory]{gauge theory} and
\hyperref[topic:quantum-gravity]{quantum gravity}.


\section{Boundary conditions control spectra}


The energies of a confined particle and the transmission of a wave through
a barrier both follow from the Schr\"odinger equation. What distinguishes
them is how solutions are joined to the boundary or to incoming and
outgoing waves. The differential expression alone describes neither
experiment. Writing
\begin{centeredalign}
H_D=-\frac{\hbar^2}{2m}\Delta_D+V
\end{centeredalign}
uses the subscript \(D\) to retain the operator domain, including its
boundary conditions. The wall and zero-slope spectra of
Chapter~\ref{chap:physical-identity} are different operators in this sense.
Their common bulk expression does not make their level spacings equal.

For scattering on the line, a stationary wave at energy \(E\) is specified
by an incident amplitude and reflected and transmitted amplitudes. In
regions of constant potential its probability current is
\begin{centeredalign}
j=\frac{\hbar}{m}\operatorname{Im}(\psi^*\partial_x\psi),\qquad
\mathcal R=|r|^2,\qquad
\mathcal T=\frac{k_{\rm out}}{k_{\rm in}}|t|^2,\qquad
\mathcal R+\mathcal T=1.
\end{centeredalign}
Here \(r\) and \(t\) are amplitudes relative to the incident wave, and
\(k_{\rm in}\) and \(k_{\rm out}\) are positive wavenumbers in propagating
channels of equal mass. The velocity factor in \(\mathcal T\) is needed
because equal amplitudes need not carry equal current. For a real,
time-independent potential with no absorption, matching the wave function
and its appropriate derivative conserves this current.

Collecting all propagating channels gives a scattering map
\begin{centeredalign}
S:\mathcal H_{\rm in}\to\mathcal H_{\rm out},\qquad S^\dagger S=I,
\end{centeredalign}
when amplitudes are normalized to unit flux and all open channels are
included. Transmission probabilities then follow from its matrix elements.
A confining wall, a finite barrier and an absorbing boundary impose
different physical conditions. Changing between them constructs a different
spectral or scattering problem; it is not a change of notation. The
\hyperref[topic:particle-in-a-box]{box},
\hyperref[topic:quantum-tunnelling]{tunnelling} and
\hyperref[topic:scattering]{scattering} treatments calculate these outcomes.


\section{Measurement, channels and error correction share one description}


A measurement does two things: it assigns a probability to each recorded
outcome and it changes the state conditioned on that outcome. An instrument
\(\mathcal I_i\) keeps both pieces. If its Kraus operators are
\(K_{i\alpha}\), the outcome probability and the conditional state are
\begin{centeredalign}
\mathcal I_i(\rho)=\sum_\alpha K_{i\alpha}\rho K_{i\alpha}^\dagger,
\qquad p_i=\operatorname{Tr}\mathcal I_i(\rho),\\
\rho_i=\frac{\mathcal I_i(\rho)}{p_i}\quad(p_i>0),\qquad
\sum_{i,\alpha}K_{i\alpha}^\dagger K_{i\alpha}=I.
\end{centeredalign}
If the outcome is discarded, the state is instead
\(\mathcal E(\rho)=\sum_i\mathcal I_i(\rho)\). For a projective
measurement of \(Z\), the two operators are \((I\pm Z)/2\). Keeping
the outcome prepares the corresponding spin direction. Discarding it
removes coherence between those directions: an input \(|+x\rangle\)
becomes \(I/2\). The same coupling to a measuring apparatus has produced
different available information depending on whether its record is kept.

Error correction uses a record without measuring the unknown logical
amplitudes. In the three-qubit repetition code, the logical state is
\(a|000\rangle+b|111\rangle\). Measuring \(Z_1Z_2\) and \(Z_2Z_3\)
gives a syndrome that distinguishes no error from a flip of any one spin,
while leaving \(a\) and \(b\) undetermined. The syndrome selects a
recovery operation \(\mathcal R_i\). The complete corrected channel is
\begin{centeredalign}
\mathcal E=\sum_i\mathcal I_i,\qquad
\mathcal R_i\circ\mathcal I_i,\qquad
\rho_{\rm corrected}=\sum_i(\mathcal R_i\circ\mathcal I_i)(\rho_{\rm noisy}).
\end{centeredalign}
For a mixture of no error and single-spin flips, the recovered logical
state equals the original one. The code does not correct a phase flip:
that error acts inside the code space and gives no changed syndrome.

Complete positivity ensures that these maps also give physical states
when the encoded system is entangled with a reference. For successful
quantum correction the recovered correlations with that reference must
be preserved, not merely the populations of two chosen basis states.
\hyperref[topic:quantum-error-correction]{Quantum error correction}
derives the general condition for this preservation.


\section{Duality and simulation require a map that follows the dynamics}


A quantum simulator represents the states and interactions of one model
in another physical device. An encoding \(V\) sends a target state into
a chosen subspace of the device. When \(V^\dagger V=I\), the encoding
preserves norms and inner products. Its Hamiltonian must also preserve
the encoded evolution:
\begin{centeredalign}
VH_{\rm target}=H_{\rm carrier}V,\qquad
VO_{\rm target}=O_{\rm carrier}V.
\end{centeredalign}
The first relation implies
\(Ve^{-iH_{\rm target}t/\hbar}=e^{-iH_{\rm carrier}t/\hbar}V\).
The second identifies which device measurement represents a target
observable. Together they preserve the predicted expectation for every
encoded preparation.

For example, encode one spin by
\(V|0\rangle=|00\rangle\) and \(V|1\rangle=|11\rangle\).
The target Hamiltonian \(H_{\rm target}=\hbar\omega Z/2\) is
represented by \(H_{\rm carrier}=\hbar\omega Z_1/2\) within the
code space, and the target \(X\) is represented by \(X_1X_2\).
Starting from the encoded \(|+x\rangle\), the pair therefore gives
\begin{centeredalign}
\langle X_1X_2\rangle_t=\cos(\omega t)=\langle X\rangle_t.
\end{centeredalign}
Measuring \(X_1\) alone would give a different answer: its action
leaves the encoded subspace. The observable correspondence is as
necessary as the Hamiltonian correspondence.

For an approximate simulation write
\(D=H_{\rm carrier}V-VH_{\rm target}\), so that
\(VH_{\rm target}\simeq H_{\rm carrier}V\). Integrating the derivative
of the two unitary evolutions gives, for bounded operators and \(t\geq0\),
\begin{centeredalign}
\big\|e^{-iH_{\rm carrier}t/\hbar}V
-Ve^{-iH_{\rm target}t/\hbar}\big\|
\leq\frac{t}{\hbar}\|D\|.
\end{centeredalign}
Similarly, \(VO_{\rm target}\simeq O_{\rm carrier}V\) introduces
a separate measurement error. A small discrepancy at preparation can
thus grow during evolution, and its scale determines the duration over
which the simulation predicts accurately. Dualities and field-theory
correspondences require their own domains and observable maps; the
finite-spin encoding is an explicit example of the calculation, not
a proof of those broader correspondences.



\chapter{Carrying A Law Into Another System}\label{chap:transfer}
\begin{claimbox}
A coordinate transformation, an encoding or a reduction relates two
physical descriptions. Substituting that relation into their equations
shows which evolution is preserved and which additional contribution
is needed to reproduce the original motion.
\end{claimbox}

\section{Exact Transfer And Its Failure}
A mechanical oscillator and an LC circuit have different constituents,
but charge and magnetic flux can reproduce the equations for position
and momentum. The comparison becomes a calculation when the
correspondence is specified. Let \(\alpha\) be a linear map from source
states to target states, and let \(L_s,L_t\) generate their linear
evolutions on compatible domains. Evolving and then mapping agrees
with mapping and then evolving when
\begin{centeredalign}
e^{tL_t}\alpha=\alpha e^{tL_s},\qquad
L_t\alpha=\alpha L_s,\qquad \Delta=L_t\alpha-\alpha L_s.
\end{centeredalign}
The generator relation gives the infinitesimal version of the same
correspondence. The measured quantities must also agree after the map.
The oscillator--circuit correspondence of
Chapter~\ref{chap:physical-identity} meets both requirements. A nonlinear
change of a stochastic coordinate instead acts by composition on test
functions; applying its generator includes the second derivative of
that composition. This is the origin of the It\^o correction developed
in the calculation companion.

For a proposed target description, \(\Delta\) gives the missing rate
of change on mapped source states. A correction \(\mathcal K\) satisfies
\((L_t+\mathcal K)\alpha=\alpha L_s\) when \(\mathcal K\alpha=-\Delta\).
In the two-spin example, mapping \(\rho\) to the single number
\(x\) leaves a missing rate \(-2gc\). Different correlations
can have the same \(x\), so no function of \(x\) alone supplies
that rate for every preparation. Retaining \(c\), or its initial
value and the memory integral, restores the prediction. The
calculation determines what has been lost; the physical source
equations determine whether it is a changed local term or a
degree of freedom that must be retained.

A particle confined to a thin curved surface gives a classic example. Carrying the free Schr\"odinger equation from
flat space onto the surface by keeping only the motion along it leaves a difference that does not vanish as the
confinement tightens. It takes the form of a potential set by the curvature of the surface (da Costa, Phys. Rev. A 23,
1982, 1981),
\begin{centeredalign}
V_{\mathrm{geom}}=-\frac{\hbar^2}{2m}\left(M^2-K\right),
\end{centeredalign}
with \(M\) the mean and \(K\) the Gaussian curvature. If
\(k_1,k_2\) are the principal curvatures, then
\(M^2-K=(k_1-k_2)^2/4\). The potential consequently vanishes
on a sphere, although the sphere is curved, and is
\(-\hbar^2/(8mR^2)\) on a cylinder of radius \(R\).
Spatial variations in this difference of curvatures can trap
states. The term follows from taking the thin-confinement limit
of the three-dimensional dynamics; replacing derivatives by
surface derivatives alone would omit it. Its physical content
is the influence of transverse confinement on the retained
surface motion.



\section{How Fluctuations Select The Phase Of An Oscillator}
A parametric drive can make an oscillator settle into either of two
states whose phases differ by \(\pi\). A weak field at the oscillator
frequency favours one state, but fluctuations can select the other.
The measurable quantity is the fraction of repeated preparations that
end on the favoured side. In a classical noisy oscillator swept
through a symmetric instability, this fraction depends on the bias,
the sweep rate and the fluctuations present during growth. A quantum
oscillator raises a precise additional question: do vacuum and thermal
fluctuations enter the same calculation, or does the quantum dynamics
change how the phase is selected?

Take a single bosonic mode with annihilation operator \(a\),
\([a,a^\dagger]=1\), in a frame rotating at half the pump frequency.
In units with \(\hbar=1\), let \(K\) be the Kerr coefficient,
\(\varepsilon_2(t)\) the two-photon drive and \(h\) a real
one-photon bias. The mode loses photons at rate \(\kappa\) to a
bath with occupation \(\bar n\). Its Hamiltonian and density-operator
evolution are
\begin{centeredalign}
H=-Ka^{\dagger2}a^2
+\tfrac{i}{2}\varepsilon_2(t)(a^{\dagger2}-a^2)
+ih(a^\dagger-a),\\
\dot\rho=-i[H,\rho]
+\kappa(\bar n+1)\mathcal D[a]\rho
+\kappa\bar n\mathcal D[a^\dagger]\rho,\\
\mathcal D[b]\rho=b\rho b^\dagger
-\tfrac12(b^\dagger b\rho+\rho b^\dagger b).
\end{centeredalign}
The two-photon term amplifies the quadrature
\(\widehat x=(a+a^\dagger)/2\), while the bias pushes it toward
positive values. The linear gain changes sign at
\(\varepsilon_2=\kappa/2\). Above this threshold, the Kerr term
limits the growth and produces opposite-phase states. The measured
selection probability is the probability that \(\widehat x>0\)
at the end of the specified ramp and hold.

Before the nonlinear term becomes important, the Hamiltonian is
quadratic and the dissipation is linear. The Wigner function then
obeys a Fokker--Planck equation with a positive Gaussian marginal
for this quadrature. That marginal is reproduced by a classical
random variable \(x\) satisfying
\begin{centeredalign}
dx=[(\varepsilon_2(t)-\kappa/2)x+h]\,dt+\sqrt{2D}\,dW_t,
\qquad D=\frac{\kappa(2\bar n+1)}{8}.
\end{centeredalign}
Here \(W_t\) is a Wiener process. The random variable represents
quadrature statistics, whereas \(\langle\widehat x\rangle\)
is its mean and contains no random increment. The diffusion
coefficient is fixed by loss and bath temperature. At zero bath
temperature it remains nonzero because the vacuum quadrature
variance is \(1/4\).

For a linear ramp \(\varepsilon_2=rt\), define
\(\varphi(t)=rt^2/2-\kappa t/2\). Multiplication by
\(e^{-\varphi(t)}\) removes the gain from the stochastic
equation. Starting from a centred thermal preparation with
variance \(\sigma_0^2=(2\bar n_0+1)/4\), its mean and variance
at time \(t\) are
\begin{centeredalign}
\mu(t)=h e^{\varphi(t)}I_1(t),\qquad
\sigma^2(t)=e^{2\varphi(t)}[\sigma_0^2+2DI_2(t)],\\
I_k(t)=\int_0^t e^{-k\varphi(s)}\,ds,\qquad
P(t)=\Phi\!\left(\frac{hI_1(t)}{\sqrt{\sigma_0^2+2DI_2(t)}}\right).
\end{centeredalign}
The function \(\Phi\) is the cumulative distribution function
of a standard normal random variable. The mean and standard
deviation grow by the same gain factor, so that factor cancels
from their ratio. The probability depends instead on the bias
accumulated during the ramp and on the initial and continuously
supplied fluctuations.

For an indefinitely continued ramp, the convergent integrals are
\(I_k=\int_0^\infty e^{-k\varphi(s)}\,ds\). When
\(\kappa\gg\sqrt r\), bath fluctuations overwhelm the initial
variance before the threshold is crossed. Evaluating the
Gaussian integrals yields
\begin{centeredalign}
P\simeq\Phi\!\left(\frac{\pi^{1/4}h}{D^{1/2}r^{1/4}}\right).
\end{centeredalign}
This is the classical selection law of Kondepudi and Nelson
(Nature 314, 438, 1985), now with \(D\) determined by the
quantum bath. In the opposite limit \(\kappa\to0\), an
initial vacuum supplies the variance and
\(I_1=\sqrt{\pi/(2r)}\). The result becomes
\(P=\Phi(h\sqrt{2\pi/r})\): the rate has exponent
\(-1/2\) rather than \(-1/4\). A different initial thermal
variance changes this limit. The distinction is whether the
choice is seeded by the prepared state or by fluctuations
supplied during the ramp. Related bias-selection calculations
for a step in gain are given by Gu et al. (Phys. Rev. Research
7, L022056, 2025), alongside the oscillator experiment of
Roques-Carmes et al. (Science 381, 205, 2023).

\section{Following The Selection Into The Nonlinear Regime}
The Gaussian probability is exact for the quadratic model.
To compare it with a Kerr oscillator, the finite ramp, hold
and nonlinear evolution must be followed. In the FieldBridge
example, \(\kappa=1\), \(K=0.02\), \(\bar n=0\), and
the drive rises from zero to \(1.2\) at rate \(r=0.1\),
then remains fixed for ten loss-time units. A bias
\(h=0.10691\) favours positive \(x\). The finite-time
Gaussian law gives \(P=0.75047\), and density-operator
evolution gives \(0.75045\), a difference of approximately
\(1.4\times10^{-5}\). The final states contain about
27 photons. In this calculation nonlinear growth largely
preserves the sign selected during the early stage.

The reported rate scan at the same final photon number gives
differences below \(6\times10^{-5}\) at rates \(0.1,1,10\),
but about \(6\times10^{-4}\) at \(r=0.01\). The latter
is the beginning of a nonlinear correction, rather than
agreement to \(6\times10^{-5}\) at every rate. The
\href{https://github.com/synthetix-institute/fieldbridge/blob/main/docs/tutorial/30_quantum_write.md}{quantum-oscillator calculation}
specifies the Hamiltonians, ramps and numerical comparisons.

Increasing \(K\) to one, with the other settings unchanged,
reduces the final mean-field amplitude to about half a photon.
The computed probability is then approximately \(0.684\),
against the Gaussian value \(0.750\); the biased steady state
gives about \(0.686\). The phase distribution can change
after the initial growth because transitions between the
two weakly populated alternatives remain appreciable. In
the reported few-photon scans, deficits reach \(0.22\),
and sufficiently slow sweeps approach the steady-state
probability within \(0.006\). These are numerical results
for the specified oscillators, not exact consequences of
the Gaussian law.

If a separated, slowly decaying mode of the Lindblad
generator governs the exchange of population, a two-state
reduction writes
\begin{centeredalign}
\frac{dp}{dt}=\Gamma(\varepsilon_2(t))
[P_{\mathrm{eq}}(\varepsilon_2(t))-p].
\end{centeredalign}
Here \(p\) is the positive-side population,
\(P_{\mathrm{eq}}\) its instantaneous steady-state value,
and \(\Gamma\) the decay rate of the relevant switching
mode. The gap determines that rate only when this mode
is the slowest relaxing one and other modes have decayed.
For the reported 1.5-photon example the reduction reproduces
the density-operator result within \(0.01\) over two decades
of sweep rate. The sweep time near threshold is of order
\(1/\sqrt r\), so \(\Gamma\gg\sqrt r\) permits
relaxation during that interval, whereas a smaller
\(\Gamma\) can retain the earlier selection. The
coefficient \(\pi^{1/4}\) comes from the preceding
Gaussian integrals, not from this time-scale comparison.

\section{What The Comparison Establishes}
The change from a classical noise parameter to vacuum and
thermal quadrature fluctuations is an explicit transfer:
the quantum generator supplies the coefficients of the
same Gaussian selection equation. Following the full
oscillator then shows when that equation predicts the
final measurement and when later dynamics alters the
selection. The proposed switching equation represents
that later change through quantities calculated from
the Lindblad generator. It is an approximation whose
accuracy depends on a separated switching mode.

The curved-surface potential and the two-spin memory
equation illustrate the same method with different
mathematics. In each case the full source equations
determine a contribution omitted by a smaller model.
Their consequences are different: a shifted surface
energy, a history-dependent polarization, or a changed
oscillator-selection probability. Preserving a common
operation means retaining those contributions when the
target experiment makes them relevant.

\section{Open Calculations}
For few-photon oscillators, the biased steady state and
switching mode are currently obtained numerically.
An analytic reduction would have to derive their
dependence on drive and Kerr strength and determine
where other decay modes can be eliminated. A second
question concerns a state lost at a turning point,
where escape may begin before the deterministic state
disappears; that change requires a different growth
and fluctuation calculation. The information carried
by a selected quantum phase and the work expended in
preparing it also remain to be related. The classical
Gaussian-channel bound \(\tfrac12\ln(1+W)\) is relevant
only when \(W\) is its specified signal-to-noise ratio
and the same input and noise assumptions are retained.
It supplies no universal quantum work--information law.


\part{State Spaces And Dynamical Laws}

\chapter{State space, domain, and representation}

The space of states and the domain of each operator decide which preparations and questions exist before any Hamiltonian is chosen.

The topics of this chapter concern that space itself: Hilbert space, the mathematical formulation of the theory, the change between representations, and the old quantum theory that preceded them.

\section*{Topics}

\begin{longtable}{p{0.92\linewidth}}
\toprule
Page \\
\midrule
\endhead
\phantomsection\label{page:transformation-theory-quantum-mechanics}Transformation theory (quantum mechanics) \\
\phantomsection\label{page:mathematical-formulation-of-quantum-mechanics}Mathematical formulation of quantum mechanics \\
\phantomsection\label{page:hilbert-space}Hilbert space \\
\phantomsection\label{page:quantum-differential-calculus}Quantum differential calculus \\
\phantomsection\label{page:quantum-mechanics}Quantum mechanics \\
\phantomsection\label{page:old-quantum-theory}Old quantum theory \\
\phantomsection\label{page:fourier-transform}Fourier transform \\
\bottomrule
\end{longtable}

\clearpage

\section{Transformation theory (quantum mechanics)}

\label{topic:transformation-theory-quantum-mechanics}

Cold atoms released into a harmonic trap form, after a quarter of an oscillation period, an image of their initial momentum distribution. The explanation is a calculation of the kind introduced by Dirac and Jordan in 1927. Their transformation theory treats every prediction as a product of amplitudes between eigenstates of different observables.

The amplitude \(\langle a|b\rangle\) is the transformation function between eigenstates of two observables. Its squared modulus is the probability of finding \(a\) after preparing \(b\). Amplitudes compose by summing over a complete set of intermediate eigenstates. Wave mechanics and matrix mechanics then differ only in which basis the same amplitudes are written:

\begin{centeredalign}
\langle a|c\rangle=\int db\,\langle a|b\rangle\langle b|c\rangle,\qquad
\psi(q)=\langle q|\psi\rangle,\qquad \widetilde\psi(p)=\langle p|\psi\rangle.
\end{centeredalign}

The central transformation function connects position and momentum. Because \([Q,P]=i\hbar\), momentum generates translations, and in the position basis it acts as \(-i\hbar\,\partial_q\). Applied to a momentum eigenstate this gives a first-order differential equation. Normalization to \(\langle p|p'\rangle=\delta(p-p')\) fixes the constant:

\begin{centeredalign}
-i\hbar\,\partial_q\langle q|p\rangle=p\,\langle q|p\rangle
\quad\Longrightarrow\quad
\langle q|p\rangle=\frac{e^{ipq/\hbar}}{\sqrt{2\pi\hbar}}.
\end{centeredalign}

The modulus is the same for every \(q\) and \(p\). A state of sharp momentum therefore has a uniform position distribution, and the converse also holds. The spin version is \(|\langle{\pm}x|{\pm}z\rangle|^2=\tfrac12\): complete knowledge of \(S_z\) leaves the outcomes of \(S_x\) equally likely.

Time evolution is also a transformation, from positions at one time to positions at another. For a particle of mass \(m\) in a harmonic trap of angular frequency \(\omega\), the energy eigenfunctions are Hermite functions \(h_n\), each acquiring the phase \(e^{-i(n+1/2)\omega t}\). At the quarter period \(t=\pi/(2\omega)\) this phase equals \(e^{-i\pi/4}(-i)^n\). The Hermite functions are also eigenfunctions of the Fourier transform with eigenvalues \((-i)^n\). The quarter-period evolution is therefore the position-to-momentum transformation itself:

\begin{centeredalign}
\psi\!\left(x,\frac{\pi}{2\omega}\right)=e^{-i\pi/4}\sqrt{m\omega}\;\widetilde\psi(m\omega x),
\qquad |\psi|^2dx=|\widetilde\psi(p)|^2dp\ \ \text{at}\ p=m\omega x.
\end{centeredalign}

For rubidium-87 in a trap with a frequency of 10 hertz, the quarter period is 25 milliseconds. A momentum of \(\hbar\) times one inverse micrometre maps to a displacement of 11.6 micrometres. Experiments with cold gases use this mapping to measure momentum distributions. The image is faithful when interactions are negligible during the quarter period and the trap is harmonic over the size of the cloud; anharmonicity blurs it.

The same composition rule underlies the Fourier transform between representations and the far-field imaging of a freely expanding cloud. Dirac also noted that the transformation function for a short time interval behaves like \(\exp(iS/\hbar)\), with \(S\) the classical action. Feynman developed that observation into the path integral. Von Neumann later recast the theory in Hilbert space without delta-normalized states; rigged Hilbert spaces restore them as generalized eigenvectors.

\subsection*{References For The Physical Derivation}
\begin{itemize}
\item \href{https://arxiv.org/abs/1408.4680}{P. A. Murthy and colleagues, Matter wave Fourier optics with a strongly interacting two-dimensional Fermi gas; imaging momentum distributions by evolution in a harmonic potential.}
\item \href{https://arxiv.org/abs/quant-ph/9907069}{F. Gieres, Mathematical surprises and Dirac's formalism in quantum mechanics; equivalence of wave mechanics, matrix mechanics and Dirac's formalism, and the status of delta-normalized states.}
\end{itemize}





\clearpage

\section{Mathematical formulation of quantum mechanics}

\label{topic:mathematical-formulation-of-quantum-mechanics}

A neutron spin rotated through a full turn points in its original direction again, but its state vector changes sign. In 1975 the groups of Rauch and of Werner observed this sign with neutron interferometers, in which one of two beams crossed a magnetic field. Predicting the measured intensity uses each rule of the standard formulation once.

The outcomes of a measurement correspond to mutually orthogonal projectors \(P_a\) on a complex Hilbert space, with \(\sum_aP_a=I\). Probabilities must add over orthogonal projectors, and the probability of \(P_a\) must not depend on which other projectors complete the measurement. In dimension three or more, Gleason's theorem shows that every such assignment has the trace form, so the state is a density operator \(\rho\):

\begin{centeredalign}
p(a)=\operatorname{Tr}(\rho P_a),\qquad \rho\geq0,\qquad \operatorname{Tr}\rho=1.
\end{centeredalign}

In dimension two, additional assumptions are needed for the same conclusion. Reversible evolution must preserve probabilities and compose in time. If it is continuous, Stone's theorem gives a self-adjoint generator \(H\), and the Schrödinger equation follows:

\begin{centeredalign}
U(t+s)=U(t)U(s),\qquad U(t)=e^{-iHt/\hbar},\qquad
\rho(t)=U(t)\rho\,U(t)^\dagger.
\end{centeredalign}

A measurement that would repeat its result immediately changes the state to \(P_a\rho P_a\) divided by its trace, the Lüders rule. If the outcome is not recorded, the state becomes the sum \(\sum_aP_a\rho P_a\). Composite systems are described on tensor products of the component spaces.

In the interferometer a silicon crystal splits each neutron into two paths, \(|1\rangle\) and \(|2\rangle\), and later recombines them. The state space is the tensor product of this path space with the spin space \(\mathbb C^2\). In path 2 a field \(B\) along \(z\) acts over a length \(L\). For a neutron of speed \(v\) and gyromagnetic ratio \(\gamma_n\), the field Hamiltonian generates a spin rotation by the angle \(\alpha=\gamma_nBL/v\):

\begin{centeredalign}
U(\alpha)=e^{-i\alpha\sigma_z/2},\qquad U(2\pi)=-I,\qquad U(4\pi)=I.
\end{centeredalign}

Let \(\rho_s\) be the incoming spin state and \(\chi\) the phase difference between the paths without field. The trace rule, with the projector onto one output beam, gives the fraction of neutrons leaving that beam. For an unpolarized beam, \(\rho_s=I/2\):

\begin{centeredalign}
I(\alpha)=\frac12\Bigl[1+\operatorname{Re}\bigl(e^{i\chi}\operatorname{Tr}\rho_sU(\alpha)\bigr)\Bigr]
=\frac12\Bigl[1+\cos\chi\,\cos\frac{\alpha}{2}\Bigr].
\end{centeredalign}

The interference term reverses its sign at \(\alpha=2\pi\) and returns only at \(\alpha=4\pi\). A single beam cannot show the sign, because \(-I\) multiplies every spin state by the same factor. The unrotated path supplies the reference. Thermal neutrons of wavelength 0.18 nanometres move at about 2200 metres per second. A full turn then needs \(BL=2\pi v/\gamma_n\approx7.5\times10^{-5}\) tesla metre, for example 7.5 millitesla over one centimetre. Both 1975 experiments found the intensity periodic in the field with this 4\(\pi\) period. A spread of neutron speeds spreads \(\alpha\) in proportion to its size, which lowers the contrast after many turns.

A measurement of the spin in path 2 changes the prediction. Its outcomes correspond to the projectors \(|2\rangle\langle2|\otimes P_\pm\), completed by \(|1\rangle\langle1|\otimes I\) for neutrons in path 1. The Lüders rule removes every term of the density operator that connects the two paths, so the interference term vanishes. The intensity becomes one half for every field and every \(\chi\). The record shows which path was taken, whether or not anyone reads it.

The formulation does not state which interactions count as measurements; decoherence and the measurement problem address that physical question. Within its rules, it gives the probability of any sequence of outcomes. The half-angle in \(U(\alpha)\) comes from the spin-1/2 representation of rotations described under the Pauli matrices.

\subsection*{References For The Physical Derivation}
\begin{itemize}
\item \href{https://arxiv.org/abs/1508.06951}{V. Moretti, Mathematical Foundations of Quantum Mechanics: An Advanced Short Course; Gleason's theorem, Stone's theorem and the Lüders rule.}
\item \href{https://arxiv.org/abs/1407.2526}{J. Klepp, S. Sponar and Y. Hasegawa, Fundamental phenomena of quantum mechanics explored with neutron interferometers; spinor rotation, spin-path states and measurements in neutron interferometry.}
\end{itemize}





\clearpage

\section{Hilbert space}

\label{topic:hilbert-space}

A particle between two hard walls a distance \(L\) apart is prepared with the same amplitude at every point of the box. A measurement of its energy then returns the lowest level in a definite fraction of runs. Computing that fraction uses the two properties that define a Hilbert space: an inner product and completeness.

The states of the particle are square-integrable functions on the interval. The inner product of two of them gives an amplitude, and its squared modulus gives the probability of finding one prepared state in another. The energy eigenfunctions \(\phi_n(x)=\sqrt{2/L}\,\sin(n\pi x/L)\) are orthonormal, and every state is a convergent sum of them:

\begin{centeredalign}
\langle\phi|\psi\rangle=\int_0^L\phi^*(x)\psi(x)\,dx,\qquad
\psi=\sum_{n\geq1}c_n\phi_n,\\
c_n=\langle\phi_n|\psi\rangle,\qquad
\sum_{n\geq1}|c_n|^2=\langle\psi|\psi\rangle=1.
\end{centeredalign}

For the flat state \(\psi=1/\sqrt L\) the coefficients are \(c_n=\sqrt2\,[1-(-1)^n]/(n\pi)\). The even terms vanish because the flat state is symmetric about the centre of the box, while the even eigenfunctions are antisymmetric. The odd terms give

\begin{centeredalign}
P(E_1)=|c_1|^2=\frac{8}{\pi^2}\approx0.811,\qquad
P(E_n)=\frac{8}{n^2\pi^2}\quad(n\ \text{odd}).
\end{centeredalign}

The probabilities add up to one because the sum of \(1/n^2\) over odd \(n\) equals \(\pi^2/8\). This is Parseval's identity. It holds because the space is complete: a sequence of finite superpositions whose members approach each other in mean square converges to a state inside the space. The probability missing after the terms up to \(n=N\) falls roughly as \(4/(\pi^2N)\) for large \(N\). The partial sums vanish at the walls themselves and overshoot the flat value by 18 percent next to them, however many terms are kept. Convergence in the norm is nevertheless the physically relevant sense, because the norm controls every probability.

The same state has infinite mean energy. With \(E_n=n^2\pi^2\hbar^2/(2mL^2)\), every odd term of the energy sum contributes the same amount:

\begin{centeredalign}
\langle H\rangle=\sum_{n\ \text{odd}}\frac{8}{n^2\pi^2}\,
\frac{n^2\pi^2\hbar^2}{2mL^2}
=\sum_{n\ \text{odd}}\frac{4\hbar^2}{mL^2}=\infty.
\end{centeredalign}

The flat wave function belongs to the Hilbert space, but not to the domain of the Hamiltonian. It jumps at the walls, where every eigenfunction vanishes, and an abrupt step needs arbitrarily large momentum components. The divergence comes from the step itself, so walls of finite height do not remove it. A preparation with finite energy must round the edges over some length, and that length cuts off the sum.

The Hilbert space is not physical space. For one particle in three dimensions it is a space of wave functions of three coordinates; for \(N\) particles the functions depend on \(3N\) coordinates. A single spin needs only the two-dimensional space \(\mathbb C^2\), and \(N\) spins need dimension \(2^N\). A finite set of box modes suffices for questions about low energies, but not for the mean energy of a state with sharp edges. The distinction between the space and an operator's domain recurs in spectral theory and in the treatment of self-adjoint operators.

\subsection*{References For The Physical Derivation}
\begin{itemize}
\item \href{https://arxiv.org/abs/1508.06951}{V. Moretti, Mathematical Foundations of Quantum Mechanics: An Advanced Short Course; Hilbert spaces, operator domains and spectral decompositions.}
\end{itemize}





\clearpage

\section{Quantum differential calculus}

\label{topic:quantum-differential-calculus}

An electron hopping between neighbouring atoms of a crystal changes its position in steps of the lattice constant \(a\). Its kinetic energy then involves differences of the wave function between sites, not derivatives. On such a lattice the product rule of ordinary calculus fails. Repairing it gives the simplest example of a quantum differential calculus.

The name comes from quantum groups and noncommutative geometry. It denotes a consistent algebra of functions and differentials, not the quantization of a mechanical system. Take the forward difference \(\partial f(x)=[f(x+a)-f(x)]/a\). Expanding the difference of a product gives an extra term:

\begin{centeredalign}
\partial(fg)=(\partial f)\,g+f\,(\partial g)+a\,(\partial f)(\partial g).
\end{centeredalign}

The extra term vanishes as \(a\to0\). It can be removed exactly at finite \(a\) if the differential \(dx\) no longer commutes with functions. Define \(df=(\partial f)\,dx\) and require the ordinary Leibniz rule \(d(fg)=(df)\,g+f\,dg\). Both are consistent only if moving a function across \(dx\) shifts its argument:

\begin{centeredalign}
dx\,f(x)=f(x+a)\,dx,\qquad dx\,x-x\,dx=a\,dx.
\end{centeredalign}

Functions on the lattice commute with each other, but differentials do not commute with them. This is a noncommutative differential calculus on a commutative algebra. Replacing the translation by a dilation gives Jackson's q-derivative, \(D_qf(x)=[f(qx)-f(x)]/[(q-1)x]\), with \(dx\,f(x)=f(qx)\,dx\). On powers it produces the q-numbers:

\begin{centeredalign}
D_qx^n=[n]_qx^{n-1},\qquad [n]_q=\frac{q^n-1}{q-1}=1+q+\cdots+q^{n-1}.
\end{centeredalign}

The lattice calculus has a measurable consequence. The kinetic energy built from a forward and a backward difference has plane waves \(e^{ikx}\) as eigenfunctions, with

\begin{centeredalign}
E(k)=\frac{2\hbar^2}{ma^2}\sin^2\frac{ka}{2}
=\frac{\hbar^2k^2}{2m}-\frac{\hbar^2a^2k^4}{24m}+\cdots
\end{centeredalign}

The group velocity \((\hbar/ma)\sin ka\) vanishes at the zone edge \(k=\pi/a\). A constant force \(F\) advances \(k\) uniformly in time, so the velocity oscillates instead of growing. The period of this Bloch oscillation is \(2\pi\hbar/(Fa)\). The same structure describes the lowest band of a deep optical lattice, with a hopping energy in place of \(\hbar^2/(2ma^2)\). The period does not depend on the hopping. In a vertical lattice of spacing 266 nanometres, gravity gives strontium-88 atoms a Bloch frequency \(mga/h\) of about 574 hertz. Bloch oscillations of strontium-88 in a vertical optical lattice have been followed for several seconds and used to measure gravity at micrometre scales. The single-band description holds while the force is too weak to drive transitions to the next band.

The q-calculus has a less direct physical status. For a q-deformed oscillator with \(aa^\dagger-qa^\dagger a=1\), the operator \(a^\dagger a\) has eigenvalues \([n]_q\), so its levels are unequally spaced. To first order in \(q-1\) the spacing changes linearly with \(n\), like the anharmonic term of a molecular vibrational ladder. Such fits are phenomenological. It is established that these calculi are consistent algebraic structures. Quantum-group symmetry also appears exactly in integrable models, such as the anisotropic XXZ spin chain with suitable boundary terms. Whether spacetime itself has a deformed differential structure at short distances is open, and no experiment has shown it.

\subsection*{References For The Physical Derivation}
\begin{itemize}
\item \href{https://arxiv.org/abs/hep-th/9404112}{A. Dimakis and F. Müller-Hoissen, Discrete differential calculus, graphs, topologies and gauge theory; differential calculus on discrete sets with noncommuting differentials.}
\item \href{https://arxiv.org/abs/physics/0605018}{G. Ferrari, N. Poli, F. Sorrentino and G. M. Tino, Long-lived Bloch oscillations with bosonic Sr atoms and application to gravity measurement at micrometer scale; lattice dispersion and the Bloch period in gravity.}
\end{itemize}





\clearpage

\section{Quantum mechanics}

\label{topic:quantum-mechanics}

A beam of silver atoms passing an inhomogeneous magnet splits into two spots instead of a continuous smear. Stern and Gerlach observed this in 1922. Sequences of such magnets display most of the structure of quantum mechanics in one experiment.

The magnetic moment of a silver atom comes from one unpaired electron spin. Its spin state is a unit vector \(|\psi\rangle=\alpha|\uparrow\rangle+\beta|\downarrow\rangle\) in a two-dimensional complex space. A magnet with its field gradient along \(z\) sends the atom up or down with probabilities \(|\alpha|^2\) and \(|\beta|^2\). Blocking the lower beam prepares \(|\uparrow\rangle\): every atom of the upper beam goes up again in a second \(z\) magnet.

A magnet tilted by an angle \(\theta\) from \(z\) in the \(xz\) plane has its own pair of output states. The probability of each output is the squared modulus of an overlap, the Born rule:

\begin{centeredalign}
|{+}\mathbf n\rangle=\cos\tfrac\theta2|\uparrow\rangle+\sin\tfrac\theta2|\downarrow\rangle,\qquad
p({+}\mathbf n\,|\uparrow)=|\langle{+}\mathbf n|\uparrow\rangle|^2=\cos^2\tfrac\theta2.
\end{centeredalign}

At \(\theta=60^\circ\) three quarters of the atoms take the upper path; at \(90^\circ\) half do. Single atoms arrive at random, and only the relative counts approach these values.

Now place a third magnet along \(z\) after the tilted one. Suppose the path through the tilted magnet is recorded. The probabilities of the two paths then add, and atoms that entered as \(|\uparrow\rangle\) leave as \(|\downarrow\rangle\) with probability \(\tfrac12\sin^2\theta\). If the two paths are recombined without any record, the amplitudes add instead. Completeness of the two output states then gives zero:

\begin{centeredalign}
p_{\rm recorded}(\downarrow)=\sum_\pm|\langle\downarrow|{\pm}\mathbf n\rangle|^2|\langle{\pm}\mathbf n|\uparrow\rangle|^2
=\tfrac12\sin^2\theta,\\
p_{\rm recombined}(\downarrow)=\Bigl|\sum_\pm\langle\downarrow|{\pm}\mathbf n\rangle\langle{\pm}\mathbf n|\uparrow\rangle\Bigr|^2=0.
\end{centeredalign}

The difference between the two lines is set by whether any physical system keeps information about the path. Coherent recombination is demanding, because the two paths must overlap again in position and momentum. A full-loop Stern--Gerlach interferometer driven by the precise field of an atom chip has achieved it.

The spin components \(S_z\) and \(S_x\) are represented by noncommuting operators, \([S_z,S_x]=i\hbar S_y\). No preparation gives definite values to both. After the \(x\) magnet selects its upper beam, a \(z\) magnet again splits the atoms evenly. The commutator and the uncertainty principle treat this incompatibility quantitatively.

Between magnets the state evolves by the Schrödinger equation. In a uniform field \(B\) along \(z\), take \(H=\omega S_z\) with \(\omega\) proportional to \(B\). A spin prepared along \(x\) precesses about \(z\), and the probability of finding it along \(x\) oscillates:

\begin{centeredalign}
i\hbar\frac{d}{dt}|\psi\rangle=H|\psi\rangle,\qquad
p_{+x}(t)=\cos^2\frac{\omega t}{2}.
\end{centeredalign}

For an electron spin, \(\omega/2\pi\) equals 28.0 gigahertz for each tesla of field. The same rules extend to any system. States are unit vectors or density operators, and outcomes are eigenvalues of self-adjoint operators. Probabilities are squared overlaps, evolution between measurements is unitary, and composite systems use tensor products. The two spots appear only when the gradient separates the paths by more than the beam width. With weaker gradients the spots overlap, and the magnet performs an unsharp measurement of the kind discussed under observables.

\subsection*{References For The Physical Derivation}
\begin{itemize}
\item \href{https://arxiv.org/abs/1801.02708}{Y. Margalit and colleagues, Realization of a complete Stern-Gerlach interferometer; coherent recombination of the two spin paths with an atom-chip field.}
\end{itemize}





\clearpage

\section{Old quantum theory}

\label{topic:old-quantum-theory}

Hydrogen emits light at sharp wavelengths, among them the red Balmer line at 656 nanometres. Between 1913 and 1925 these lines were explained by selecting particular classical orbits, before wave mechanics existed. The selection rule of that period, Bohr--Sommerfeld quantization, worked for hydrogen and for the oscillator, and its failures show what the later theory had to change.

The rule applies to motion that separates into periodic coordinates \(q_k\) with conjugate momenta \(p_k\). The action integral over one period of each coordinate must be an integer multiple of Planck's constant \(h\):

\begin{centeredalign}
J_k=\oint p_k\,dq_k=n_k h,\qquad n_k=0,1,2,\ldots
\end{centeredalign}

For an oscillator of mass \(m\) and angular frequency \(\omega\), the orbit at energy \(E\) is an ellipse in phase space with area \(2\pi E/\omega\). The rule gives \(E_n=n\hbar\omega\), with the spacing \(\hbar\omega\) seen in molecular vibrational ladders. The ground energy is wrong, however. Wave mechanics gives \((n+\tfrac12)\hbar\omega\), and isotope shifts in molecular band spectra, analysed by Mulliken in 1924 and 1925, already required the half quantum. The modern semiclassical rule replaces \(n\) by \(n+\tfrac12\) for motion between two smooth turning points. That correction comes from the phase of the wave at the turning points, which an orbit cannot supply.

For hydrogen, write \(k=e^2/(4\pi\varepsilon_0)\) for the Coulomb strength and use circular orbits. Quantized angular momentum \(mvr=n\hbar\) together with the force balance \(mv^2/r=k/r^2\) gives

\begin{centeredalign}
r_n=\frac{n^2\hbar^2}{mk}=n^2a_0,\qquad
E_n=-\frac{mk^2}{2\hbar^2n^2}=-\frac{13.6\ \text{eV}}{n^2}.
\end{centeredalign}

Here \(a_0=52.9\) picometres is the Bohr radius. The transition from \(n=3\) to \(n=2\), with the reduced-mass correction, gives a vacuum wavelength of 656.5 nanometres. Sommerfeld's elliptic orbits add a radial action. The sum of the radial and angular actions depends only on the energy, so the levels depend only on \(n=n_r+n_\varphi\). This reproduces the degeneracy of the Kepler problem:

\begin{centeredalign}
J_r+J_\varphi=2\pi k\sqrt{\frac{m}{-2E}}=(n_r+n_\varphi)h.
\end{centeredalign}

For large \(n\) the spacing of neighbouring levels, divided by \(h\), approaches the classical orbital frequency \(mk^2/(2\pi\hbar^3n^3)\). Bohr used this correspondence in 1913 to fix the constant in his rule, requiring large orbits to radiate at their classical frequency. Sommerfeld's relativistic treatment of the elliptic orbits gave a fine-structure formula that agrees with the later Dirac result. Semiclassical analysis of the Dirac equation, including the precession of the spin, helps explain that agreement.

The failures are specific. The rule requires \(n_\varphi\geq1\), because an orbit with zero angular momentum would pass through the nucleus. The predicted ground state therefore carries angular momentum \(\hbar\), whereas the hydrogen ground state has none; wave mechanics gives \(L^2=\ell(\ell+1)\hbar^2\) with \(\ell=0,\ldots,n-1\). Einstein pointed out in 1917 that the rule can be stated without special coordinates, as conditions on loops around invariant tori in phase space. Motion without such tori cannot be quantized this way. The three-body motion of helium has no such structure, and orbit models failed to reproduce its ionization energy of 24.6 eV. The anomalous Zeeman effect demanded half-integer quantum numbers, later traced to electron spin. Line intensities lay outside the rule, and their calculation led Heisenberg to matrix mechanics in 1925.

The corrected rule survives as the leading term of semiclassical WKB theory. It remains accurate for highly excited Rydberg states and for the vibrational levels of molecules. The quantum harmonic oscillator and the angular momentum operator give the exact results that replaced the orbit picture.

\subsection*{References For The Physical Derivation}
\begin{itemize}
\item \href{https://arxiv.org/abs/quant-ph/0212082}{S. Keppeler, Semiclassical quantisation rules for the Dirac and Pauli equations; Maslov corrections and the relativistic Kepler problem behind Sommerfeld's fine-structure formula.}
\item \href{https://arxiv.org/abs/2112.14194}{A. Afriat, Einstein and the old quantum theory; Einstein's coordinate-independent quantum conditions on invariant tori.}
\end{itemize}





\clearpage

\section{Fourier transform}

\label{topic:fourier-transform}

A free particle propagates most simply in momentum space, where each momentum component accumulates its own phase. A localized detector, however, asks for a position probability. The Fourier transform connects these descriptions while preserving the physical state. On the real line, for a square-integrable wave function \(\psi\) and momentum \(p\), take

\begin{centeredalign}
\widetilde\psi(p)=\frac{1}{\sqrt{2\pi\hbar}}\int_{\mathbb R}
e^{-ipx/\hbar}\psi(x)\,dx,\qquad
\int|\widetilde\psi(p)|^2dp=\int|\psi(x)|^2dx.
\end{centeredalign}

For wave functions on a common suitable domain, integration by parts carries the position-space momentum operator \(-i\hbar\partial_x\) into multiplication by \(p\). Multiplication by position becomes \(i\hbar\partial_p\). The kinetic equation and the detector must both be transformed:

\begin{centeredalign}
\mathcal F(-i\hbar\partial_x)\mathcal F^{-1}=p,\qquad
\mathcal F x\mathcal F^{-1}=i\hbar\partial_p,\qquad
\widetilde H=\frac{p^2}{2m}+V(i\hbar\partial_p).
\end{centeredalign}

The last expression is immediately useful for polynomial potentials; for a general potential, multiplication in position space is represented by an integral kernel in momentum space. A local interaction can therefore become nonlocal in the new coordinates without changing the theory. Similarity of displayed formulas is neither necessary nor sufficient for physical equivalence.

For a normalized Gaussian packet with initial position variance \(\sigma_x^2\), zero position-momentum covariance and minimum uncertainty, the free evolution multiplies each momentum amplitude by \(\exp[-ip^2t/(2m\hbar)]\). Transforming back gives

\begin{centeredalign}
\operatorname{Var}x(t)=\sigma_x^2+
\frac{\hbar^2t^2}{4m^2\sigma_x^2}.
\end{centeredalign}

The packet broadens because its different momentum components separate during propagation. The momentum distribution itself remains unchanged. This provides a direct comparison between descriptions: one representation makes the conservation of momentum probabilities explicit, while the other displays the spatial spreading produced by those same components.

Boundaries change this correspondence. Periodic functions on a finite interval have a discrete Fourier series; Dirichlet eigenfunctions lead instead to a sine expansion. Surface terms must vanish under the actual boundary conditions before the integration-by-parts identity defines the required operator map. The operation of Fourier transformation is transferable, but the domain and measure determine which transform and which spectrum describe the physical system.

\subsection*{References For The Physical Derivation}
\begin{itemize}
\item \href{https://arxiv.org/abs/1508.06951}{V. Moretti, Mathematical Foundations of Quantum Mechanics: An Advanced Short Course; states, operators and symmetry.}
\end{itemize}

\subsection*{Relations In The Original Papers}

\href{https://arxiv.org/html/1508.06951#S2.EGx7}{arXiv:1508.06951, S2.EGx7}. The Fourier transformation expresses the position wavefunction in wavevector coordinates. This source uses wavevector k; the chapter uses momentum p = hbar k with the corresponding normalization.



\chapter{Quantum states and subsystem structure}

A quantum state assigns probabilities to every possible measurement, and for a composite system it also records correlations between the parts.

Wave functions, density matrices, superpositions and entangled states are different ways of writing this information; decoherence describes how it spreads into an environment.

\section*{Topics}

\begin{longtable}{p{0.92\linewidth}}
\toprule
Page \\
\midrule
\endhead
\phantomsection\label{page:density-matrix}Density matrix \\
\phantomsection\label{page:quantum-superposition}Quantum superposition \\
\phantomsection\label{page:quantum-decoherence}Quantum decoherence \\
\phantomsection\label{page:coherence-physics}Coherence (physics) \\
\phantomsection\label{page:superposition-principle}Superposition principle \emph{(see Quantum Superposition)} \\
\phantomsection\label{page:wave-function}Wave function \\
\phantomsection\label{page:quantum-state}Quantum state \\
\phantomsection\label{page:two-state-quantum-system}Two-state quantum system \\
\phantomsection\label{page:qubit}Qubit \\
\phantomsection\label{page:quantum-number}Quantum number \\
\phantomsection\label{page:spin-physics}Spin (physics) \\
\phantomsection\label{page:schr-dinger-s-cat}Schrödinger's cat \\
\phantomsection\label{page:quantum-entanglement}Quantum entanglement \\
\phantomsection\label{page:quantum-biology}Quantum biology \\
\phantomsection\label{page:macroscopic-quantum-phenomena}Macroscopic quantum phenomena \\
\phantomsection\label{page:quantum-fluctuation}Quantum fluctuation \\
\bottomrule
\end{longtable}

\clearpage

\section{Density matrix}

\label{topic:density-matrix}

Nuclear spins in a slightly inhomogeneous magnet are tipped by a radio-frequency pulse into the plane perpendicular to the field. The induced signal decays quickly, yet a second pulse can bring it back. Over this time each spin evolves unitarily, precessing at its own local frequency. The decay and the return belong to the ensemble, and the density matrix is the state that describes it.

If a source prepares the pure state \(\ket{\psi_a}\) with probability \(p_a\), every measurement probability is linear in these weights. Collecting the weights in one operator gives

\begin{centeredalign}
\rho=\sum_a p_a\ket{\psi_a}\bra{\psi_a},\qquad \rho\geq0,\qquad \operatorname{Tr}\rho=1,
\qquad p(i)=\operatorname{Tr}(\rho P_i),
\end{centeredalign}

where \(P_i\) projects onto outcome \(i\). Positivity and unit trace make every such probability nonnegative and their sum one. In a chosen basis the diagonal elements are populations and the off-diagonal elements are coherences; only the coherences produce interference between the basis states. Different ensembles with the same \(\rho\) give identical statistics, as described for the quantum state. The purity separates a single pure preparation from a mixture, and for an isolated system \(\rho\) obeys the von Neumann equation:

\begin{centeredalign}
\operatorname{Tr}\rho^2=\sum_{a,b}p_ap_b|\langle\psi_a|\psi_b\rangle|^2\leq1,
\qquad i\hbar\frac{d\rho}{dt}=[H,\rho].
\end{centeredalign}

The purity equals one only when all prepared states coincide up to a phase. In the energy basis the von Neumann equation keeps the populations fixed and rotates each coherence at its Bohr frequency, \(\rho_{mn}(t)=\rho_{mn}(0)e^{-i(E_m-E_n)t/\hbar}\). Unitary evolution of the ensemble therefore preserves its purity.

Now take spin-one-half nuclei that the pulse has left polarized along \(x\) with polarization \(P\), so that each has \(\rho_{01}=P/2\) in the \(\sigma_z\) basis. In thermal nuclear magnetic resonance \(P\) is small, but it only scales the signal. A nucleus whose local precession frequency differs from the reference by \(\delta\) has Hamiltonian \(\hbar\delta\sigma_z/2\) in the rotating frame, and its coherence becomes \(\tfrac P2e^{-i\delta t}\). If \(\delta\) is Gaussian distributed across the sample with standard deviation \(\sigma\), the ensemble average is

\begin{centeredalign}
\rho_{01}(t)=\tfrac P2\langle e^{-i\delta t}\rangle=\tfrac P2e^{-\sigma^2t^2/2},\qquad
\operatorname{Tr}\rho^2=\tfrac12\left(1+P^2e^{-\sigma^2t^2}\right).
\end{centeredalign}

The transverse magnetization is proportional to \(|\rho_{01}|\), so the signal decays on the time \(\sqrt2/\sigma\). A frequency spread \(\sigma/2\pi=1\) kHz gives 0.23 ms. The ensemble loses purity although every member evolves unitarily. A \(\pi\) pulse about \(x\) at time \(\tau\) exchanges \(\rho_{01}\) and \(\rho_{10}\), so the phase acquired afterwards cancels the phase acquired before. At \(2\tau\) every member returns to \(\rho_{01}=P/2\) whatever its \(\delta\). This is the Hahn spin echo.

The echo is complete only if each detuning stays constant over the interval \(2\tau\). Diffusion of the spins through the field gradient, or coupling to other degrees of freedom, leaves an irreversible decay of the echo. That remaining loss is quantum decoherence: the phase has moved into correlations with an environment. A mixed state of the same kind arises whenever a subsystem is entangled with an unmeasured partner, and then \(\rho_A=\operatorname{Tr}_B\rho_{AB}\). At a single time the two origins can give the same density matrix. Only a further operation, such as the echo pulse, separates them.

\subsection*{References For The Physical Derivation}
\begin{itemize}
\item \href{https://arxiv.org/abs/1106.1445}{M. M. Wilde, From Classical to Quantum Shannon Theory; density operators, ensembles, purity and partial trace.}
\item \href{https://arxiv.org/abs/1508.06951}{V. Moretti, Mathematical Foundations of Quantum Mechanics: An Advanced Short Course; states as positive trace-class operators and the trace rule.}
\end{itemize}





\clearpage

\section{Quantum superposition}

\label{topic:quantum-superposition}

The Super-Kamiokande detector in Japan counts muon neutrinos produced by cosmic rays in the atmosphere. In 1998 the collaboration reported a deficit of muon neutrinos of a few GeV arriving from below, after crossing up to 13000 km of the Earth. They were about half as numerous as those arriving from above, which travel some 15 km. Electron neutrinos showed no such deficit. The missing muon neutrinos had changed flavour in flight. That change requires each neutrino to travel as a superposition of states with different masses.

The Schrödinger equation is linear. If \(\ket a\) and \(\ket b\) are possible states, so is every normalized combination of them, and each component evolves with its own phase. For an equal superposition with relative phase \(\varphi\), a measurement that projects onto \(\ket\pm=(\ket a\pm\ket b)/\sqrt2\) gives

\begin{centeredalign}
\ket\psi=\frac{1}{\sqrt2}\left(\ket a+e^{i\varphi}\ket b\right),\qquad
P_\pm=|\langle\pm|\psi\rangle|^2=\frac12\left(1\pm\cos\varphi\right).
\end{centeredalign}

In a two-path interferometer \(\ket a\) and \(\ket b\) are the paths and the two outputs realize \(\ket\pm\). A mixture with equal probability of either state gives a different prediction:

\begin{centeredalign}
\rho_{\rm mix}=\frac12\left(\ket a\bra a+\ket b\bra b\right),\qquad
P_\pm=\operatorname{Tr}\left(\rho_{\rm mix}\ket\pm\bra\pm\right)=\frac12.
\end{centeredalign}

Both descriptions assign probability one half to each of \(\ket a\) and \(\ket b\), so a measurement in that basis cannot tell them apart. Only the superposition responds to \(\varphi\). The term \(\cos\varphi\) comes from the off-diagonal elements of the density matrix and is the measurable signature of the superposition. Whether a state counts as a superposition depends on the basis: \(\ket+\) is a superposition of \(\ket a\) and \(\ket b\) but an eigenstate of the output measurement.

A neutrino is created by the weak interaction in a flavour state, here \(\nu_\mu\) together with a muon. In a two-flavour model the flavour states are rotations of the mass eigenstates \(\nu_2\) and \(\nu_3\) by a mixing angle \(\theta\):

\begin{centeredalign}
\ket{\nu_\mu}=\cos\theta\ket{\nu_2}+\sin\theta\ket{\nu_3},\qquad
\ket{\nu_\tau}=-\sin\theta\ket{\nu_2}+\cos\theta\ket{\nu_3}.
\end{centeredalign}

In natural units a mass state of momentum \(p\gg m_i\) has energy \(E_i\approx p+m_i^2/2E\). After a distance \(L\approx t\) the two components have therefore acquired the relative phase \(\Delta m^2L/2E\), with \(\Delta m^2=m_3^2-m_2^2\). Projecting the evolved state onto \(\ket{\nu_\tau}\) gives

\begin{centeredalign}
P(\nu_\mu\to\nu_\tau)=\sin^22\theta\,\sin^2\frac{\Delta m^2L}{4E}
=\sin^22\theta\,\sin^2\!\left(1.27\,\frac{\Delta m^2[\mathrm{eV}^2]\,L[\mathrm{km}]}{E[\mathrm{GeV}]}\right).
\end{centeredalign}

The conversion probability returns to zero after the oscillation length \(4\pi E/\Delta m^2\). For \(\Delta m^2=2.5\times10^{-3}\ \mathrm{eV}^2\) and \(E=1\) GeV this is about 990 km, with the first maximum near 500 km. Downward neutrinos, after 15 km, have converted with probability below one percent. Upward neutrinos have passed through several oscillation lengths. Averaged over their spread of energies and path lengths, the second factor gives one half. With nearly maximal mixing, \(\sin^22\theta\approx1\), about half of them arrive as \(\nu_\tau\), which the detector did not record as muon neutrinos.

Suppose instead that the source emitted an incoherent mixture, \(\nu_2\) with probability \(\cos^2\theta\) and \(\nu_3\) with probability \(\sin^2\theta\). Each mass state is stationary and keeps its flavour content, so the conversion probability is \(\cos^2\theta\sin^2\theta+\sin^2\theta\cos^2\theta=\tfrac12\sin^22\theta\) at every distance. That value equals the average of the oscillating formula but does not depend on \(L\). It would remove half of the downward muon neutrinos as well, which the data exclude. Only the coherent superposition starts as a pure \(\nu_\mu\) and develops the conversion with distance.

The coherent description requires the wave packets of the two mass components to overlap at the detector. Over very long distances their different group velocities separate them, and the oscillating term averages away, as described under quantum decoherence. The two-flavour formula neglects the third neutrino; matter in the Earth changes electron-neutrino oscillations but hardly affects the conversion of \(\nu_\mu\) into \(\nu_\tau\). Superpositions of states with different electric charge are never observed, an example of a superselection rule. Superpositions of macroscopically distinct states are the subject of Schrödinger's cat.

\subsection*{References For The Physical Derivation}
\begin{itemize}
\item \href{https://arxiv.org/abs/hep-ex/9807003}{Super-Kamiokande Collaboration (Y. Fukuda and colleagues), Evidence for Oscillation of Atmospheric Neutrinos; up-down asymmetry of muon neutrinos.}
\item \href{https://arxiv.org/abs/hep-ph/0202058}{M. C. Gonzalez-Garcia and Y. Nir, Neutrino Masses and Mixing: Evidence and Implications; two-flavour oscillation formula and atmospheric parameters.}
\item \href{https://arxiv.org/abs/hep-ph/0001264}{E. Kh. Akhmedov, Neutrino Physics; coherence of mass eigenstates, oscillation length and wave-packet separation.}
\end{itemize}





\clearpage

\section{Quantum decoherence}

\label{topic:quantum-decoherence}

Fullerene molecules sent through a grating interferometer produce clear fringes. When gas is admitted to the vacuum chamber, or when the molecules are heated before they enter, the contrast falls, although every molecule still passes the gratings. Collisions with gas molecules and thermally emitted photons leave a record of the path in the surroundings. Decoherence is the loss of interference that follows from such records.

Consider a particle in an equal superposition of two paths, \(\ket L\) and \(\ket R\), and an environment initially in \(\ket E\). A scattering event that depends on the path leaves the environment in \(\ket{E_L}\) or \(\ket{E_R}\). The joint evolution is unitary. Tracing out the environment gives the reduced density matrix of the particle:

\begin{centeredalign}
\tfrac{1}{\sqrt2}(\ket L+\ket R)\ket E\;\longrightarrow\;
\tfrac{1}{\sqrt2}\left(\ket L\ket{E_L}+\ket R\ket{E_R}\right),\\
\rho_S=\tfrac12\left(\ket L\bra L+\ket R\bra R+\langle E_R|E_L\rangle\ket L\bra R
+\langle E_L|E_R\rangle\ket R\bra L\right).
\end{centeredalign}

At the screen the path amplitudes \(\psi_L(x)\) and \(\psi_R(x)\) overlap, and the interference term is multiplied by \(\langle E_R|E_L\rangle\). For equal intensities the fringe visibility equals \(|\langle E_R|E_L\rangle|\). Visibility is lost to the extent that the environment could in principle identify the path; nobody needs to examine the scattered particle.

Independent collisions multiply their overlaps. Suppose collisions occur at random with mean rate \(\Gamma\), and each leaves a real overlap \(\chi\) between zero and one. Averaging \(\chi^N\) over the Poisson distribution of the number \(N\) of collisions in time \(t\) gives

\begin{centeredalign}
V(t)=\sum_{N=0}^\infty e^{-\Gamma t}\frac{(\Gamma t)^N}{N!}\,\chi^N=e^{-\Gamma(1-\chi)t}.
\end{centeredalign}

The decoherence rate \(\Gamma(1-\chi)\) depends on how well one collision resolves the path. A scattered wave whose wavelength \(\lambda\) greatly exceeds the path separation \(d\) gives \(\chi\) close to one, with \(1-\chi\) growing as \((d/\lambda)^2\). When \(d\gg\lambda\), a single event suffices and the rate saturates at \(\Gamma\). Since \(\Gamma\) is proportional to the gas pressure, the visibility falls exponentially with pressure. This dependence was measured for \(\mathrm{C}_{70}\) molecules in a Talbot--Lau interferometer. Heated molecules emit more thermal photons and lose visibility in the same way.

Over times long compared with one collision, the same decay follows from a Lindblad master equation with a dephasing operator. With \(\sigma_z=\ket L\bra L-\ket R\bra R\) and rate \(\gamma=\Gamma(1-\chi)\),

\begin{centeredalign}
\frac{d\rho}{dt}=-\frac{i}{\hbar}[H,\rho]+\frac{\gamma}{2}\left(\sigma_z\rho\sigma_z-\rho\right),
\qquad \frac{d\rho_{LR}}{dt}\Big|_{\rm env}=-\gamma\rho_{LR}.
\end{centeredalign}

The populations of the two paths are unchanged; only their coherence decays. Coherence is lost in the basis in which the interaction is diagonal, here position. Decoherence therefore explains why records form in that basis. It does not determine on which path a given molecule is found; that question belongs to the measurement problem.

The master equation assumes that each scattered particle leaves and never returns, so that records made at different times are independent. A small environment violates this assumption. For one environmental spin, initially up along its \(z\) axis and coupled through \(\hbar g\sigma_z\otimes\sigma_x^E\), the overlap is \(\cos2gt\), and the visibility \(|\cos2gt|\) returns to one at \(t=\pi/2g\). Such revivals, and the restored fringes of the quantum eraser experiment, show that the phase information has moved into system-environment correlations rather than being destroyed. A Markovian master equation for the particle alone cannot describe them.

\subsection*{References For The Physical Derivation}
\begin{itemize}
\item \href{https://arxiv.org/abs/quant-ph/0612118}{K. Hornberger, Introduction to Decoherence Theory; which-path records, reduced dynamics and master equations.}
\item \href{https://arxiv.org/abs/quant-ph/0303093}{K. Hornberger and colleagues, Collisional Decoherence Observed in Matter Wave Interferometry; exponential loss of fullerene fringe visibility with gas pressure.}
\item \href{https://arxiv.org/abs/quant-ph/0105127}{W. H. Zurek, Decoherence, Einselection, and the Quantum Origins of the Classical; pointer basis and the role of the environment.}
\end{itemize}



\subsection*{Equation Sources}
\begin{itemize}
\item \href{https://arxiv.org/abs/quant-ph/0306087}{arXiv:quant-ph/0306087}
\item \href{https://arxiv.org/abs/quant-ph/0008131}{arXiv:quant-ph/0008131}
\item \href{https://arxiv.org/abs/quant-ph/0211096}{arXiv:quant-ph/0211096}
\end{itemize}

\clearpage

\section{Coherence (physics)}

\label{topic:coherence-physics}

A Michelson interferometer illuminated by a filtered lamp shows fringes only while its two arms have nearly equal lengths. With a single-frequency laser, fringes persist over path differences of many metres. The fringe contrast as a function of the delay between the arms is a measurable property of the light, its first-order coherence.

Let \(E(t)\) be the complex field amplitude at the detector, let one arm delay it by \(\tau\), and let a balanced beam splitter recombine the two parts. Averaged over the fluctuations of a stationary source of intensity \(I_0\), the detected intensity is

\begin{centeredalign}
I(\tau)=\frac{I_0}{2}\left[1+\operatorname{Re}g^{(1)}(\tau)\right],\qquad
g^{(1)}(\tau)=\frac{\langle E^*(t)E(t+\tau)\rangle}{\langle|E(t)|^2\rangle}.
\end{centeredalign}

As \(\tau\) changes over one optical period, \(I\) oscillates between maxima and minima with visibility \(V=|g^{(1)}(\tau)|\). By the Wiener--Khinchin theorem, \(g^{(1)}\) is the Fourier transform of the normalized power spectrum \(S(\omega)\). A Lorentzian line centred at \(\omega_0\) with full width \(\Gamma\) at half maximum gives

\begin{centeredalign}
g^{(1)}(\tau)=\int S(\omega)\,e^{-i\omega\tau}d\omega=e^{-i\omega_0\tau-\Gamma|\tau|/2}.
\end{centeredalign}

Fringes therefore fade once the delay exceeds a time of order \(1/\Gamma\). In terms of path difference, light of wavelength \(\lambda\) and spectral width \(\Delta\lambda\) has a coherence length of order \(\lambda^2/\Delta\lambda\). A filter 1 nm wide at 600 nm gives about 0.36 mm, and a laser line 1 MHz wide gives about 300 m. The same quantity describes one photon in a two-path interferometer. There the visibility equals \(2|\rho_{LR}|\), twice the coherence between the paths in the density matrix, and quantum decoherence reduces it.

Intensity correlations distinguish sources with identical spectra. Hanbury Brown and Twiss correlated the photocurrents of two detectors and found correlated intensity fluctuations in thermal light. Chaotic light is the sum of fields from very many independent emitters, so the field has Gaussian statistics. Its fourth moment then factorizes into pairs of second moments, which gives the Siegert relation:

\begin{centeredalign}
g^{(2)}(\tau)=\frac{\langle I(t)I(t+\tau)\rangle}{\langle I\rangle^2}=1+|g^{(1)}(\tau)|^2,
\qquad g^{(2)}(0)=\frac{\langle n(n-1)\rangle}{\langle n\rangle^2}=2.
\end{centeredalign}

At zero delay the intensity correlation is twice the uncorrelated value; in photon terms, thermal light arrives in bunches. The same value 2 follows from the geometric distribution of photon numbers \(n\) in one thermal mode. The excess lasts about one coherence time, and a detector much slower than that averages most of it away. Hanbury Brown and Twiss nevertheless measured how the excess declined with the separation of two telescopes, and obtained the angular diameter of Sirius in 1956.

A classical intensity satisfies \(\langle I^2\rangle\geq\langle I\rangle^2\), which forces \(g^{(2)}(0)\geq1\), and an ideal laser beam gives exactly one. Values below one, which no fluctuating classical wave can produce, come from single-photon sources treated under photon. First-order coherence cannot separate these cases, because sources with the same spectrum give the same fringes.

\subsection*{References For The Physical Derivation}
\begin{itemize}
\item \href{https://arxiv.org/abs/nucl-th/9804026}{G. Baym, The Physics of Hanbury Brown--Twiss Intensity Interferometry: From Stars to Nuclear Collisions; first- and second-order coherence and stellar interferometry.}
\end{itemize}





\clearpage
\section{Superposition principle}
\label{topic:superposition-principle}
Superposition principle is an alternative name for Quantum Superposition.
The physical derivation is given in \hyperref[topic:quantum-superposition]{Quantum Superposition}.

\clearpage

\section{Wave function}

\label{topic:wave-function}

Electrons accelerated through 54 volts and directed onto a nickel crystal leave its surface preferentially at about \(50^\circ\) from the incident beam. Davisson and Germer observed this peak in 1927. Point particles scattered by the surface atoms give no reason for such a sharp preferred angle, whereas a wave of definite wavelength does. The wave function supplies that wave and also the probability of detecting each electron at a given place.

For one spinless particle, the state \(\ket\psi\) is represented in the position basis by the complex function \(\psi(\mathbf x)=\langle\mathbf x|\psi\rangle\). The Born rule gives the probability that a detector covering the region \(\Delta\) responds:

\begin{centeredalign}
\Pr(\mathbf x\in\Delta)=\int_\Delta|\psi(\mathbf x)|^2d^3x,\qquad
\int|\psi|^2d^3x=1,\qquad \psi\sim e^{i\alpha}\psi.
\end{centeredalign}

A constant overall phase \(e^{i\alpha}\) changes no probability, so \(\psi\) and \(e^{i\alpha}\psi\) describe the same state. Phase differences between regions are physical. The Schrödinger equation \(i\hbar\partial_t\psi=-(\hbar^2/2m)\nabla^2\psi+V\psi\) conserves probability locally, and the current through a surface depends on the phase gradient. Writing \(\psi=R\,e^{iS/\hbar}\) with real \(R\) and \(S\) gives

\begin{centeredalign}
\partial_t|\psi|^2+\nabla\cdot\mathbf j=0,\qquad
\mathbf j=\frac{\hbar}{m}\operatorname{Im}(\psi^*\nabla\psi)=\frac{R^2\,\nabla S}{m}.
\end{centeredalign}

The phase gradient divided by the mass acts as a local velocity. A plane wave \(e^{i\mathbf k\cdot\mathbf x}\) carries a uniform current, is an eigenfunction of the momentum operator \(-i\hbar\nabla\) with eigenvalue \(\hbar\mathbf k\), and has wavelength \(\lambda=2\pi/k=h/p\). For an electron of kinetic energy \(E\) this gives

\begin{centeredalign}
\lambda=\frac{h}{\sqrt{2m_eE}}=0.167\ \mathrm{nm}\ \ (E=54\ \mathrm{eV}),\qquad
D\sin\theta=\lambda.
\end{centeredalign}

Waves scattered from parallel rows of surface atoms with spacing \(D\) add in phase when their path difference \(D\sin\theta\) equals one wavelength. With \(D=0.215\) nm for nickel, the first-order angle is \(50.9^\circ\), close to the observed peak. The electron speed is about 1.5 percent of the speed of light, so the nonrelativistic relation suffices. The angle arises from the relative phases of the waves from different atoms; adding the probabilities \(|\psi|^2\) contributed by each atom would give no peak.

For several particles the wave function lives on configuration space. Two particles require \(\psi(\mathbf x_1,\mathbf x_2)\), a function of six coordinates, not two waves in ordinary space. Because \(\psi\) need not factor into a function of \(\mathbf x_1\) times one of \(\mathbf x_2\), quantum entanglement is possible. Spin adds discrete components, so a particle of spin \(s\) has a wave function with \(2s+1\) components. For identical particles \(\psi\) must be symmetric or antisymmetric under exchange of their coordinates. The wave function is not recorded by a single detector. Its modulus follows from position statistics of identically prepared particles, and its phase from interference or from momentum statistics. The Fourier transform gives the equivalent momentum-space wave function, and a statistical preparation requires a density matrix.

\subsection*{References For The Physical Derivation}
\begin{itemize}
\item \href{https://arxiv.org/abs/1508.06951}{V. Moretti, Mathematical Foundations of Quantum Mechanics: An Advanced Short Course; wave functions as elements of L2, rays and the position measure.}
\end{itemize}





\clearpage

\section{Quantum state}

\label{topic:quantum-state}

Preparing a spin along one axis fixes probabilities for measurements along every other axis. The information needed for those predictions can be collected in a density operator \(\rho\). For a two-level system let \(\boldsymbol\sigma=(\sigma_x,\sigma_y,\sigma_z)\) be the Pauli matrices and \(\mathbf r\) the vector of their expectation values. The three components of \(\mathbf r\) determine the state:

\begin{centeredalign}
\rho=\frac12(I+\mathbf r\cdot\boldsymbol\sigma),\qquad |\mathbf r|\leq1.
\end{centeredalign}

The bound follows from positivity: the two eigenvalues are \((1\pm|\mathbf r|)/2\). A unit vector describes a pure preparation. An interior point describes a mixed state, which may arise from uncontrolled preparation or from ignoring a correlated partner. Different ensembles can produce the same density operator. A later experiment acting only on this spin cannot distinguish those ensembles unless additional information about the preparation is accessible.

A Stern--Gerlach measurement along a unit vector \(\mathbf n\) has projectors \(E_\pm=(I\pm\mathbf n\cdot\boldsymbol\sigma)/2\). The probability of each outcome is consequently

\begin{centeredalign}
p_\pm=\operatorname{Tr}(\rho E_\pm)=\frac12(1\pm\mathbf r\cdot\mathbf n).
\end{centeredalign}

The same state gives certainty along its preparation axis and an equal distribution along a perpendicular axis. The uncertainty therefore belongs to the relation between preparation and measurement. It is not evidence that a pure state is an unspecified classical direction. Measurements along three independent axes reconstruct the Bloch vector, while measurements along only one axis leave a disk of compatible states.

For a composite system the density operator also contains correlations. The singlet and an equal mixture of oppositely aligned spins both give \(\rho_A=I/2\) for either individual spin, but they predict different joint measurements. A partial trace retains every expectation of an operator acting on subsystem \(A\) alone. It need not retain the information required for subsequent interacting evolution, because the interaction can convert a joint correlation into a local expectation value.

\begin{centeredalign}
\langle O_A\rangle=\operatorname{Tr}_{AB}[\rho_{AB}(O_A\otimes I_B)]
=\operatorname{Tr}_A(\rho_A O_A),\qquad \rho_A=\operatorname{Tr}_B\rho_{AB}.
\end{centeredalign}

Thus the state required by a mechanism is fixed by the available measurements and subsequent interactions together. Local spin probabilities require \(\rho_A\) at the present time; a prediction after coupling to \(B\) generally requires the joint state or an independently specified preparation of \(B\). This distinction is the physical reason to keep the state space, interaction and preparation connected in the description.

\subsection*{References For The Physical Derivation}
\begin{itemize}
\item \href{https://arxiv.org/abs/1508.06951}{V. Moretti, Mathematical Foundations of Quantum Mechanics: An Advanced Short Course; states, operators and symmetry.}
\end{itemize}

\subsection*{Relations In The Original Papers}

\href{https://arxiv.org/html/1508.06951#S3.Ex131}{arXiv:1508.06951, S3.Ex131}. The spectral probability of an observable is the trace of the density operator with its spectral projector. Density operators and projection-valued measurements; the Bloch-state example is derived in the chapter.



\clearpage

\section{Two-state quantum system}

\label{topic:two-state-quantum-system}

An avoided crossing couples two states that would otherwise evolve independently. In their basis, a constant energy offset contributes only an overall phase, whereas the remaining Hamiltonian determines both transition probabilities and phase evolution. Let \(\Delta\) be the angular-frequency detuning between the uncoupled levels and \(\Omega\) a real coupling frequency. With \(|0\rangle\) and \(|1\rangle\) the eigenstates of \(\sigma_z\), choose

\begin{centeredalign}
H=\frac{\hbar}{2}(\Delta\sigma_z+\Omega\sigma_x),\qquad \omega_R=\sqrt{\Delta^2+\Omega^2}.
\end{centeredalign}

The Pauli algebra gives \((\Delta\sigma_z+\Omega\sigma_x)^2=\omega_R^2I\). Every higher power in the exponential therefore reduces to either the identity or the Hamiltonian. The evolution follows without fitting a transition curve:

\begin{centeredalign}
U(t)=\cos\frac{\omega_Rt}{2}\,I-i\sin\frac{\omega_Rt}{2}\,
\frac{\Delta\sigma_z+\Omega\sigma_x}{\omega_R}.
\end{centeredalign}

For the preparation \(|0\rangle\), the off-diagonal matrix element of \(U\) gives the transition probability

\begin{centeredalign}
P_{0\to1}(t)=\frac{\Omega^2}{\Delta^2+\Omega^2}
\sin^2\!\left(\frac{t}{2}\sqrt{\Delta^2+\Omega^2}\right).
\end{centeredalign}

On resonance the coupling can transfer the entire population; a pulse of duration \(\pi/|\Omega|\) interchanges the levels. Detuning increases the oscillation frequency but reduces the greatest achievable population transfer. These two consequences distinguish a changed level splitting from a changed coupling even when a short segment of an oscillation looks similar.

The Bloch vector rotates about the effective axis \((\Omega,0,\Delta)\). A second pulse about another axis can convert a relative phase into a population difference. Reversing two such pulses generally changes the result because their Hamiltonians do not commute. The state space is the same in both experiments; the difference lies in the sequence of interactions, and the measured population exposes that difference.

This calculation applies to any isolated pair of levels governed by the stated Hamiltonian, including two coupled modes or selected atomic states. Its use for a larger physical system requires that other levels remain unpopulated and environmental relaxation be negligible over the pulse duration. If the Hamiltonian was obtained in a rotating frame, the rotating-wave approximation and the transformation of the measured operator are additional parts of the correspondence. Specifying a two-dimensional matrix does not by itself establish those conditions in a device.

\subsection*{References For The Physical Derivation}
\begin{itemize}
\item \href{https://arxiv.org/abs/1508.06951}{V. Moretti, Mathematical Foundations of Quantum Mechanics: An Advanced Short Course; states, operators and symmetry.}
\end{itemize}





\clearpage

\section{Qubit}

\label{topic:qubit}

A trapped ion, the polarization of a photon and the lowest two levels of a superconducting circuit each provide two states that can be prepared, rotated and distinguished. Whether such a pair works as a qubit depends on how strictly the dynamics stays within the two states. It also depends on how long their relative phase survives. Within the chosen pair a pure state is

\begin{centeredalign}
\ket\psi=\alpha\ket0+\beta\ket1,\qquad |\alpha|^2+|\beta|^2=1,\qquad
p(0)=|\langle0|\psi\rangle|^2=|\alpha|^2,\qquad p(1)=|\beta|^2.
\end{centeredalign}

A global phase has no observable effect, so a pure state is fixed by two real numbers: the population ratio and the relative phase. These coordinates form the Bloch sphere described for the quantum state. A measurement in the basis \(\{\ket0,\ket1\}\) returns one of two outcomes and removes the relative phase. The continuous amplitudes therefore cannot be read off in one shot. If a bit is encoded in one of two pure states with equal prior probability, any measurement errs with probability at least the Helstrom bound:

\begin{centeredalign}
P_{\rm err}=\frac12\left(1-\sqrt{1-|\langle\psi_0|\psi_1\rangle|^2}\right).
\end{centeredalign}

Orthogonal states give zero error, while an overlap \(|\langle\psi_0|\psi_1\rangle|=\cos(\pi/8)\) gives 0.31. More generally, Holevo's theorem limits the classical information retrievable from one qubit to one bit, despite the continuum of states. A previously shared entangled pair raises this to two bits in superdense coding.

A superconducting transmon shows what restricting a system to two levels involves. It is a weakly anharmonic oscillator whose lowest transition frequency \(\omega_{01}/2\pi\) typically lies between 4 and 6 GHz. The next transition is lower by an anharmonicity \(|\eta|/2\pi\) of about 200 to 300 MHz. A control pulse at \(\omega_{01}\) also drives the transition from \(\ket1\) to \(\ket2\), detuned by \(\eta\). By the off-resonant Rabi formula of the two-state quantum system, this leakage stays small only when the drive strength and the pulse bandwidth are well below \(|\eta|\). Single-qubit gates therefore last tens of nanoseconds.

Relaxation sets the second limit. Energy decay from \(\ket1\) to \(\ket0\) at rate \(1/T_1\), together with additional pure dephasing at rate \(\Gamma_\phi\), gives in the Markovian regime at zero temperature

\begin{centeredalign}
\rho_{11}(t)=\rho_{11}(0)e^{-t/T_1},\qquad |\rho_{01}(t)|=|\rho_{01}(0)|e^{-t/T_2},
\qquad \frac{1}{T_2}=\frac{1}{2T_1}+\Gamma_\phi.
\end{centeredalign}

The factor one half follows from the decay itself. The amplitude of \(\ket1\) falls as \(e^{-t/2T_1}\), and the coherence \(\rho_{01}\) is proportional to that amplitude. Hence \(T_2\leq2T_1\), with equality when energy decay is the only process. An echo experiment measures \(T_2\) from the decay of the recovered phase coherence. Comparing it with \(T_1\) from a population decay separates loss of energy from loss of phase. Quantum error correction encodes one logical qubit in several physical ones so that errors of both kinds can be detected and reversed.

\subsection*{References For The Physical Derivation}
\begin{itemize}
\item \href{https://arxiv.org/abs/1904.06560}{P. Krantz and colleagues, A Quantum Engineer's Guide to Superconducting Qubits; transmon anharmonicity, leakage, gate times, T1 and T2.}
\item \href{https://arxiv.org/abs/cond-mat/0703002}{J. Koch and colleagues, Charge-Insensitive Qubit Design Derived from the Cooper Pair Box; the transmon and its level structure.}
\item \href{https://arxiv.org/abs/quant-ph/0002077}{D. P. DiVincenzo, The Physical Implementation of Quantum Computation; requirements on a physical two-level system used as a qubit.}
\end{itemize}





\clearpage

\section{Quantum number}

\label{topic:quantum-number}

The visible lines of hydrogen, such as the red Balmer line near 656 nm, are labelled by integers. With higher resolution each line splits into components, and a magnetic field splits them further. Quantum numbers name the stationary states between which these transitions occur. Which integers are good labels depends on which interactions the Hamiltonian contains.

A quantum number labels an eigenvalue of an observable that commutes with the Hamiltonian, so that its value is conserved by the evolution. Several such observables that also commute with one another can be diagonalized together. For the nonrelativistic hydrogen atom with a Coulomb potential, \(H\), \(L^2\), \(L_z\) and \(S_z\) form such a set:

\begin{centeredalign}
[H,L^2]=[H,L_z]=[L^2,L_z]=0,\qquad H\ket{n\,l\,m_l\,m_s}=E_n\ket{n\,l\,m_l\,m_s},\\
E_n=-\frac{R_y}{n^2},\qquad R_y=\frac{m_ee^4}{2(4\pi\varepsilon_0)^2\hbar^2}=13.6\ \mathrm{eV}.
\end{centeredalign}

Here \(n=1,2,\ldots\); for each \(n\) the orbital number \(l\) runs from \(0\) to \(n-1\), \(m_l\) from \(-l\) to \(l\), and \(m_s=\pm1/2\). The ranges of \(l\) and \(m_l\) follow from the algebra of the angular momentum operator, and the bound \(l\leq n-1\) from the radial equation. The energy depends on \(n\) alone, so the number of states in a level is

\begin{centeredalign}
g_n=\sum_{l=0}^{n-1}2(2l+1)=2n^2.
\end{centeredalign}

The counts 2, 8, 18 and 32 match the shell capacities of the periodic table. In atoms with many electrons, screening by inner electrons removes the degeneracy in \(l\), and the filling order departs from the shell order. The Balmer line connects \(n=3\) to \(n=2\), with \(1/\lambda=(R_y/hc)(1/4-1/9)\) up to a small reduced-mass correction. Electric-dipole transitions obey \(\Delta l=\pm1\) and \(\Delta m_l=0,\pm1\), so the labels also decide which lines appear.

Spin-orbit coupling adds a term proportional to \(\mathbf L\cdot\mathbf S=(J^2-L^2-S^2)/2\), which conserves the total angular momentum \(\mathbf J=\mathbf L+\mathbf S\) but not \(m_l\) or \(m_s\) separately. The good quantum numbers become \(n\), \(l\), \(j=l\pm1/2\) and \(m_j\), and the resulting fine-structure splitting is worked out under the Dirac equation.

A magnetic field \(B\) whose Zeeman energies are small compared with that splitting shifts each level according to \(m_j\). The interaction \(\mu_BB(L_z+g_sS_z)/\hbar\) is evaluated within a level of fixed \(j\). The spin moment carries \(g_s\approx2\) and the orbital moment carries one, so projecting both onto \(\mathbf J\) gives the Landé factor:

\begin{centeredalign}
\Delta E=g_J\mu_BBm_j,\qquad g_J=1+\frac{j(j+1)+s(s+1)-l(l+1)}{2j(j+1)},\\
g_J(s_{1/2})=2,\qquad g_J(p_{1/2})=\tfrac23,\qquad g_J(p_{3/2})=\tfrac43.
\end{centeredalign}

An electric-dipole photon carries one unit of angular momentum, so only transitions with \(\Delta m_j=0,\pm1\) appear. Their photon energies are shifted by \(\mu_BB(g'm_j'-gm_j)\), where primes mark the upper level. For the sodium D lines, from \(3p_{1/2}\) and \(3p_{3/2}\) to \(3s_{1/2}\), the shifts in units of \(\mu_BB\) are

\begin{centeredalign}
\mathrm{D}_1:\ \pm\tfrac23,\ \pm\tfrac43\qquad\qquad \mathrm{D}_2:\ \pm\tfrac13,\ \pm1,\ \pm\tfrac53.
\end{centeredalign}

D1 therefore splits into four components and D2 into six. In a field of 1 T the outer D2 components lie 23 GHz from the line centre, far beyond the Doppler width of about 1.5 GHz in a sodium lamp. The field is still weak in the required sense, because the fine-structure splitting of sodium is 515 GHz. Hyperfine shifts from the nuclear spin are smaller and are neglected here. These patterns, the anomalous Zeeman effect, could be explained only after electron spin with \(g_s\approx2\) was introduced. When both levels have zero total spin, as for the red cadmium line at 643.8 nm, \(g_J=1\) in both levels. The components then merge into the normal triplet at \(0\) and \(\pm\mu_BB\). Viewed along the field, the \(\Delta m_j=0\) component is absent, a direct test of the selection rule. A strong field decouples \(\mathbf L\) from \(\mathbf S\) and restores \(m_l\) and \(m_s\) as approximate labels.

Each quantum number thus records a symmetry of the Hamiltonian at a stated level of approximation. Rotational invariance conserves \(j\) and \(m_j\), and reflection symmetry gives the parity \((-1)^l\). The extra degeneracy in \(l\) of the pure Coulomb problem comes from a further conserved quantity, the Laplace--Runge--Lenz vector. An interaction that breaks a symmetry splits a level into components carrying the new labels, and spectroscopy measures those splittings.

\subsection*{References For The Physical Derivation}
\begin{itemize}
\item \href{https://arxiv.org/abs/hep-ph/0002158}{M. I. Eides, H. Grotch and V. A. Shelyuto, Theory of Light Hydrogenlike Atoms; Dirac fine structure and the Lamb shift in hydrogen.}
\end{itemize}





\clearpage

\section{Spin (physics)}

\label{topic:spin-physics}

In 1922 Stern and Gerlach sent a beam of silver atoms through a strongly inhomogeneous magnetic field and collected them on a glass plate. Magnetic moments with random classical orientations would have spread the atoms into a continuous band. The deposit instead showed two separate lines. A silver atom has one unpaired 5s electron with no orbital angular momentum. In modern terms, the two lines reveal an angular momentum intrinsic to the electron, its spin.

Spin obeys the same commutation relations as orbital angular momentum, given under the angular momentum operator, but its quantum number can be half-integer. For the electron \(s=1/2\), so its spin component along any axis takes the two values \(\pm\hbar/2\). The spin produces a magnetic moment and a Zeeman energy:

\begin{centeredalign}
\boldsymbol\mu=-g_s\mu_B\frac{\mathbf S}{\hbar},\qquad U=-\boldsymbol\mu\cdot\mathbf B,
\qquad \mu_B=\frac{e\hbar}{2m_e}.
\end{centeredalign}

Here \(g_s\approx2.0023\) is the electron g-factor and \(\mu_B\) the Bohr magneton. In a field whose strength changes along its own direction \(z\), the force on the atom is \(F_z=-g_sm_s\mu_B\,\partial B_z/\partial z\) with \(m_s=\pm1/2\). The two signs push the atom up or down by equal amounts. For an atom of mass \(M\) and speed \(v\) crossing a magnet of length \(L\), each beam is displaced by

\begin{centeredalign}
\Delta z=\frac{g_s\mu_B}{2}\left|\frac{\partial B_z}{\partial z}\right|\frac{L^2}{2Mv^2}.
\end{centeredalign}

A gradient of \(10^3\) T/m over 3.5 cm acting on silver atoms moving at 550 m/s gives 0.10 mm per beam. The two beams then lie about 0.2 mm apart, comparable to the separation recorded in 1922. The calculation keeps only the gradient along the strong field. Because \(\nabla\cdot\mathbf B=0\), transverse gradients are also present, but the moment precesses rapidly about the strong field and their effect averages out.

The same Zeeman energy makes the spin precess about a uniform field at the Larmor frequency. An oscillating field at that frequency drives transitions, which is magnetic resonance:

\begin{centeredalign}
\omega_L=\frac{g_s\mu_BB}{\hbar},\qquad \frac{\omega_L}{2\pi B}\approx28.0\ \mathrm{GHz/T}.
\end{centeredalign}

Electron spin resonance spectrometers working at 9.5 GHz therefore operate near 0.34 T. The rotation of the spin direction follows from the Pauli matrices, and a rotation by \(2\pi\) changes the sign of the spinor, which neutron interferometers have measured. For protons, the nuclear magneton and the measured proton g-factor replace \(\mu_B\) and \(g_s\). Nuclear magnetic resonance accordingly works at hundreds of megahertz in fields of several tesla.

The value of \(g_s\) and its small departure from 2 are treated under electron. Spin also fixes how identical particles behave under exchange. The spin-statistics theorem of relativistic quantum field theory, stated here without proof, makes states of identical particles with half-integer spin antisymmetric and those with integer spin symmetric. Antisymmetry gives the exclusion principle, which produces the shell structure counted under quantum number.

A composite particle follows the total spin of its constituents: an even number of fermions makes a boson and an odd number a fermion. A helium-4 atom contains six fermions, two protons, two neutrons and two electrons, and is a boson. Helium-3 lacks one neutron and is a fermion. The two liquids have nearly the same interactions but opposite statistics. Liquid helium-4 becomes superfluid near the temperature at which an ideal Bose gas of the same density would condense, as worked out under Bose--Einstein statistics. Liquid helium-3 stays normal down to about 1 mK at low pressure. A superfluid appears only after its atoms bind into Cooper pairs with total spin one.

\subsection*{References For The Physical Derivation}
\begin{itemize}
\item \href{https://arxiv.org/abs/2209.13084}{X. Fan, T. G. Myers, B. A. D. Sukra and G. Gabrielse, Measurement of the Electron Magnetic Moment; Penning-trap value of the electron g-factor.}
\item \href{https://arxiv.org/abs/0801.1134}{D. Hanneke, S. Fogwell and G. Gabrielse, New Measurement of the Electron Magnetic Moment and the Fine Structure Constant; comparison with the QED series.}
\end{itemize}





\clearpage

\section{Schrödinger's cat}

\label{topic:schr-dinger-s-cat}

In 1935 Schrödinger imagined a cat sealed in a chamber with a radioactive atom, a Geiger counter and a flask of poison. If the atom decays, the counter releases a hammer that breaks the flask. After the time at which the decay probability is one half, linear evolution of the whole system gives

\begin{centeredalign}
\ket\Psi=\frac{1}{\sqrt2}\left(\ket{u}\ket{\rm alive}+\ket{d}\ket{\rm dead}\right),
\end{centeredalign}

where \(\ket u\) and \(\ket d\) are the undecayed and decayed atom. Schrödinger presented the state as an absurd consequence of applying the wave function to everyday objects. Two physical questions remain: which observation would distinguish this state from a mixture of living and dead cats, and why no such observation succeeds.

Only interference can make the distinction, because it responds to the element \(\ket{\rm alive}\bra{\rm dead}\) of the density matrix, and that element is fragile. A laboratory version replaces the cat by a microwave cavity field in a superposition of two coherent states of opposite phase, \(\ket\beta+\ket{-\beta}\) up to normalization. The mean photon number is \(\bar n=|\beta|^2\). Photons leak out at the energy decay rate \(\kappa\). Leakage acts as a weak beam splitter between the cavity and the outside field, and a coherent state divides into two coherent states:

\begin{centeredalign}
\ket{\pm\beta}_{\rm cav}\ket0_{\rm out}\;\longrightarrow\;
\ket{\pm\beta e^{-\kappa t/2}}_{\rm cav}\,\ket{\pm\beta\sqrt{1-e^{-\kappa t}}}_{\rm out}.
\end{centeredalign}

Each component leaves its own coherent state in the outside field. Their overlap, \(\langle-\gamma|\gamma\rangle=e^{-2|\gamma|^2}\), multiplies the coherence between the components, as for any record left in an environment. With \(|\gamma|^2=\bar n(1-e^{-\kappa t})\) and \(\beta_t=\beta e^{-\kappa t/2}\), the weight \(C\) of the coherence terms \(\ket{\pm\beta_t}\bra{\mp\beta_t}\) relative to the population terms becomes

\begin{centeredalign}
C(t)=\exp\!\left[-2\bar n\left(1-e^{-\kappa t}\right)\right]
\approx e^{-2\bar n\kappa t}\quad(\kappa t\ll1).
\end{centeredalign}

Each component loses energy at rate \(\kappa\), but the coherence between them disappears \(2\bar n\) times faster. The rate grows with the squared separation \(|2\beta|^2\) of the components in phase space. Coherence falls to \(e^{-2}\) once a single photon, on average, has escaped, whatever the size of \(\bar n\). Brune, Haroche and colleagues in Paris observed this progressive decoherence in 1996, using fields of a few photons probed by Rydberg atoms. Components further apart lost their coherence faster.

For a massive object the environment is usually scattered gas molecules or thermal radiation. In the quantum Brownian motion model, with relaxation rate \(\gamma\) and temperature \(T\), a superposition of two positions a distance \(\Delta x\) apart loses coherence at the rate

\begin{centeredalign}
\tau_D^{-1}=\gamma\left(\frac{\Delta x}{\lambda_{\rm th}}\right)^2,\qquad
\lambda_{\rm th}=\frac{\hbar}{\sqrt{2mk_BT}}.
\end{centeredalign}

For a mass of one gram at 300 K, \(\lambda_{\rm th}\approx4\times10^{-23}\) m, and a separation of 1 cm gives \(\tau_D\approx10^{-41}/\gamma\). The formula assumes high temperature and a separation shorter than the wavelength of the environmental particles. For larger separations each collision resolves the position, and the rate saturates at the collision rate. Either estimate puts decoherence far below any time on which the motion of the object could be observed.

Decoherence explains why the alive and dead components cannot be made to interfere, and why the record forms in macroscopic configurations. It does not state which component is found. The reduced density matrix of the cat becomes indistinguishable from a mixture, while the global state remains a superposition under unitary evolution. Whether a further physical process selects one outcome, as objective-collapse models propose, is the measurement problem. Interference with ever larger molecules and mechanical oscillators continues to constrain such models.

\subsection*{References For The Physical Derivation}
\begin{itemize}
\item \href{https://arxiv.org/abs/quant-ph/0306072}{W. H. Zurek, Decoherence and the Transition from Quantum to Classical, Revisited; decoherence time of a macroscopic superposition in quantum Brownian motion.}
\item \href{https://arxiv.org/abs/0809.1064}{S. Deléglise and colleagues, Reconstruction of Non-Classical Cavity Field States with Snapshots of Their Decoherence; cat states of a microwave cavity field and their decay.}
\end{itemize}





\clearpage

\section{Quantum entanglement}

\label{topic:quantum-entanglement}

Polarization-entangled photon pairs from parametric down-conversion each appear unpolarized when examined alone: a polarizer at any angle passes half of them. When both photons of a pair meet linear polarizers set to a common angle, whichever angle is chosen, their results agree. Pairs with definite and correlated polarizations can reproduce the agreement along one axis, but not along two axes such as \(0^\circ\) and \(45^\circ\). A joint state that does not factor into states of the two photons accounts for both observations.

For subsystems \(A\) and \(B\) the joint state space is the tensor product \(\mathcal H_A\otimes\mathcal H_B\). Expand a pure state as \(\ket\Psi=\sum_{ij}c_{ij}\ket{i}\ket{j}\) in product bases. The singular-value decomposition of the coefficient matrix \(c\) rewrites it as a single sum, the Schmidt form:

\begin{centeredalign}
\ket\Psi=\sum_k\sqrt{\lambda_k}\,\ket{k_A}\ket{k_B},\qquad \sum_k\lambda_k=1,\qquad
\rho_A=\operatorname{Tr}_B\ket\Psi\bra\Psi=\sum_k\lambda_k\ket{k_A}\bra{k_A}.
\end{centeredalign}

The state is a product exactly when only one Schmidt coefficient is nonzero. Otherwise each subsystem is in a mixed state although the pair is pure. Both reduced states have the eigenvalues \(\lambda_k\), so the entanglement entropy is the same on either side:

\begin{centeredalign}
S_A=-\operatorname{Tr}\rho_A\log_2\rho_A=-\sum_k\lambda_k\log_2\lambda_k=S_B.
\end{centeredalign}

Local unitary operations change the Schmidt vectors but not the coefficients. Local operations assisted by classical communication cannot increase this entropy on average. For \(\cos\theta\ket{HH}+\sin\theta\ket{VV}\) the coefficients are \(\cos^2\theta\) and \(\sin^2\theta\). The choice \(\theta=\pi/4\) gives one bit, the largest value for two-level systems, and \(\theta=\pi/8\) gives 0.60 bit.

Noise turns a pure pair into a mixed state, and then the failure of one product decomposition proves nothing. A mixed state is separable when it can be written as \(\rho=\sum_kp_k\rho_k^A\otimes\rho_k^B\) with probabilities \(p_k\), and entangled otherwise. Transposing only the \(B\) factor maps every separable state to a positive operator. The Werner state mixes the singlet \(\ket{\Psi^-}\) with white noise, and its partial transpose has the eigenvalues

\begin{centeredalign}
\rho_W=p\ket{\Psi^-}\bra{\Psi^-}+(1-p)\frac{I}{4},\qquad
\operatorname{spec}\rho_W^{T_B}=\left\{\tfrac{1+p}{4},\tfrac{1+p}{4},\tfrac{1+p}{4},\tfrac{1-3p}{4}\right\}.
\end{centeredalign}

A negative eigenvalue appears for \(p>1/3\), and for two qubits the criterion of Peres and the Horodeckis is also sufficient. The largest Clauser--Horne--Shimony--Holt combination for this state is \(2\sqrt2\,p\), so the inequality is violated only for \(p>1/\sqrt2\approx0.71\). For \(1/3<p\leq1/\sqrt2\) the pairs are entangled, yet their correlations violate no such inequality. Werner constructed a local hidden-variable model for all projective measurements when \(p\leq1/2\). Entanglement and Bell nonlocality are therefore different properties; Bell's theorem concerns the second.

Entanglement is defined relative to a stated division into subsystems. A single photon sent through a balanced beam splitter is entangled between the two output modes, while in the input modes it is a product of one photon and vacuum. The relevant division is the one matched to the available operations, such as two separate laboratories. With that division fixed, a shared Bell pair is the resource consumed in quantum teleportation.

\subsection*{References For The Physical Derivation}
\begin{itemize}
\item \href{https://arxiv.org/abs/quant-ph/0702225}{R. Horodecki, P. Horodecki, M. Horodecki and K. Horodecki, Quantum Entanglement; Schmidt decomposition, entropy, separability, Werner states and Bell inequalities.}
\item \href{https://arxiv.org/abs/quant-ph/9604005}{A. Peres, Separability Criterion for Density Matrices; the partial-transpose test and its application to Werner states.}
\end{itemize}



\subsection*{Equation Sources}
\begin{itemize}
\item \href{https://arxiv.org/abs/math-ph/0211065}{arXiv:math-ph/0211065}
\end{itemize}

\clearpage

\section{Quantum biology}

\label{topic:quantum-biology}

Migratory birds orient in the geomagnetic field, whose strength of about 50 microtesla changes the energy of an electron spin by roughly \(2\times10^{-7}\) of the thermal energy at body temperature. A thermal population effect of that size is negligible. A chemical reaction controlled by the spin dynamics of a short-lived radical pair can still respond to such a field. This radical-pair mechanism is the most developed proposal for a magnetic sense, and its central calculation is a singlet yield.

Light absorption moves an electron between two molecules and creates two radicals, each with one unpaired electron. The pair starts in the singlet state, with total electron spin zero. In the simplest kinetic model the pair reacts at rate \(k\) whatever its spin state, singlet pairs to one product and triplet pairs to another. Between creation and reaction the electron spins evolve under the Zeeman interaction and hyperfine couplings to nuclear spins. With one spin-one-half nucleus \(\mathbf I\) on radical 1, isotropic coupling \(a\) and equal g-factors,

\begin{centeredalign}
\frac{H}{\hbar}=\omega(S_{1z}+S_{2z})+a\,\mathbf S_1\cdot\mathbf I,\qquad \omega=\frac{g\mu_BB}{\hbar},
\qquad \Phi_S=k\int_0^\infty e^{-kt}P_S(t)\,dt.
\end{centeredalign}

Here \(P_S(t)\) is the singlet probability and \(\Phi_S\) the singlet yield; the nuclear spin starts unpolarized. In zero field, the hyperfine term splits electron 1 and the nucleus into total spin one and zero, separated by \(a\), and the singlet probability oscillates at that frequency. In a strong field, \(\omega\gg a\), the triplet states with \(m=\pm1\) are shifted away. Only the triplet with \(m=0\) then mixes with the singlet, at frequency \(a/2\). Integrating against the survival factor gives the two yields:

\begin{centeredalign}
P_S^{(0)}(t)=\tfrac58+\tfrac38\cos at,\qquad P_S^{(\infty)}(t)=\tfrac12+\tfrac12\cos\tfrac{at}{2},\\
\Phi_S^{(0)}=\tfrac58+\tfrac38\,\frac{k^2}{k^2+a^2},\qquad
\Phi_S^{(\infty)}=\tfrac12+\tfrac12\,\frac{k^2}{k^2+a^2/4}.
\end{centeredalign}

For fast reaction, \(k\gg a\), both yields approach one and the field has no effect, because the spins have no time to evolve. For slow reaction the yield falls from 5/8 to 1/2. This 20 percent change is produced by a field whose energy is negligible compared with \(k_BT\). A field weaker than \(a\) acts through the Zeeman splitting of nearly degenerate levels. In this model fields of \(0.05a\) to \(0.2a\) lower the slow-reaction yield to about 0.4. That response requires the pair to live, and its spins to stay coherent, longer than the inverse Larmor frequency, 0.11 microseconds at 50 microtesla.

Sensing the direction of the field requires anisotropic hyperfine couplings, so that the yield depends on the angle between the field and the molecule. Cryptochrome proteins in the retina form flavin-tryptophan radical pairs and are the leading candidates. Magnetic field effects on radical-pair reactions are established in spin chemistry. Effects on isolated cryptochromes have been measured in fields of a few millitesla. Whether the avian compass works in this way at 50 microtesla in vivo, and how spin coherence survives for a microsecond in a protein, remain open.

Other parts of quantum biology show the same division between an established microscopic step and an open biological role. Proton and hydride transfer in some enzymes shows kinetic isotope effects far larger than transfer over the barrier allows, which quantum tunnelling explains. Long-lived oscillations in two-dimensional spectra of light-harvesting complexes were initially attributed to electronic coherence. They are now largely assigned to vibrational or vibronic motion, and their functional role is not established. In each case quantum decoherence sets how long a coherent step can last, and the biological question is whether it changes a rate or yield that matters to the organism.

\subsection*{References For The Physical Derivation}
\begin{itemize}
\item \href{https://arxiv.org/abs/1205.0883}{N. Lambert and colleagues, Functional Quantum Biology in Photosynthesis and Magnetoreception; radical-pair compass model and coherence in light harvesting.}
\item \href{https://arxiv.org/abs/1307.3530}{S. F. Huelga and M. B. Plenio, Vibrations, Quanta and Biology; vibrational and vibronic contributions to energy transfer in light-harvesting complexes.}
\item \href{https://arxiv.org/abs/0906.3725}{E. Gauger and colleagues, Sustained Quantum Coherence and Entanglement in the Avian Compass; spin coherence times required by a radical-pair compass.}
\end{itemize}





\clearpage

\section{Macroscopic quantum phenomena}

\label{topic:macroscopic-quantum-phenomena}

A superconducting ring cooled through its transition in a weak magnetic field traps the flux threading it. The trapped flux takes only discrete values, separated by about \(2\times10^{-15}\) Wb, a step fixed by Planck's constant and the electron charge. The ring can be millimetres across, yet its flux, a macroscopic electrical variable, is quantized.

Below the transition temperature, electrons bound in Cooper pairs share one wave function \(\psi=\sqrt{n_s}\,e^{i\theta}\), where \(n_s\) is the pair density and \(\theta\) a phase common to all pairs. A pair has charge \(q=-2e\) and mass \(m^*=2m_e\). Its current density in a vector potential \(\mathbf A\) depends on the phase gradient. Deep inside a ring thicker than the magnetic penetration depth the current vanishes. Integrating around such a path, and requiring \(\psi\) to be single-valued, gives

\begin{centeredalign}
\mathbf j_s=\frac{qn_s}{m^*}\left(\hbar\nabla\theta-q\mathbf A\right)=0
\;\Rightarrow\;2\pi\hbar n=q\oint\mathbf A\cdot d\mathbf l,\\
\Phi=n\,\frac{h}{2e}=n\Phi_0,\qquad \Phi_0=2.07\times10^{-15}\ \mathrm{Wb}.
\end{centeredalign}

Here \(n\) is an integer and the sign of the charge is absorbed into it. Deaver and Fairbank, and independently Doll and Näbauer, measured this quantization in 1961 in small superconducting cylinders. The step \(h/2e\), not \(h/e\), showed that the condensate carries charge in pairs.

Two superconductors joined through a thin insulating barrier have a phase difference \(\varphi\) across it. Pairs tunnel through the barrier, and the Josephson relations connect the current \(I\) and the voltage \(V\) to that phase:

\begin{centeredalign}
I=I_c\sin\varphi,\qquad \frac{d\varphi}{dt}=\frac{2eV}{\hbar}.
\end{centeredalign}

A constant voltage makes the current oscillate at \(2eV/h\), which is 483.6 GHz per millivolt. Microwaves of frequency \(f\) produce current steps at the voltages \(nhf/2e\). Voltage standards use arrays of such junctions, because the step voltages depend only on \(f\) and fundamental constants.

The phase difference is itself a dynamical variable. For a junction of capacitance \(C\) biased by a current \(I\), it moves like a particle of mass \(C(\hbar/2e)^2\) in a tilted washboard potential:

\begin{centeredalign}
U(\varphi)=-E_J\cos\varphi-\frac{\hbar I}{2e}\varphi,\qquad E_J=\frac{\hbar I_c}{2e},\\
\omega_p=\sqrt{\frac{2eI_c}{\hbar C}}\left(1-\frac{I^2}{I_c^2}\right)^{1/4}.
\end{centeredalign}

Here \(\omega_p\) is the frequency of small oscillations in a well. Near \(I_c\) each well becomes shallow, and escape from it switches the junction to a finite voltage. At high temperature escape is thermally activated. Below a crossover temperature near \(\hbar\omega_p/2\pi k_B\), about 76 mK for \(\omega_p/2\pi=10\) GHz, the escape rate stops depending on temperature because the phase tunnels through the barrier. Dissipation reduces the tunnelling rate, as Caldeira and Leggett showed, so the junction must be weakly damped.

Devoret, Martinis and Clarke observed this macroscopic tunnelling in 1985, and microwave spectroscopy resolved quantized energy levels of the phase in the well. This work, recognized by the 2025 Nobel Prize in Physics, underlies superconducting qubits. These effects are established, and their macroscopic character lies in one collective coordinate of very many pairs. Superpositions of macroscopically distinct states, such as opposite currents circulating in a ring, raise a further question. Several measures of macroscopicity have been proposed, and they rank such experiments differently, which leaves the comparison with Schrödinger's cat open.

\subsection*{References For The Physical Derivation}
\begin{itemize}
\item \href{https://arxiv.org/abs/cond-mat/0411174}{M. H. Devoret, A. Wallraff and J. M. Martinis, Superconducting Qubits: A Short Review; Josephson junction dynamics, macroscopic tunnelling and quantized levels.}
\item \href{https://arxiv.org/abs/1706.06173}{F. Fröwis and colleagues, Macroscopic Quantum States: Measures, Fragility and Implementations; competing measures of macroscopicity.}
\end{itemize}





\clearpage

\section{Quantum fluctuation}

\label{topic:quantum-fluctuation}

Two uncharged, parallel metal plates a micrometre apart in vacuum attract each other with a pressure of about one millipascal. Casimir predicted this force in 1948. It has been measured since 1997 between a sphere and a plate and between parallel plates. The force follows from the zero-point fluctuations of the electromagnetic field between the plates.

Each mode of the field is a harmonic oscillator of angular frequency \(\omega\). As for the quantum harmonic oscillator, its ground state has energy \(\hbar\omega/2\) and nonzero variances of both field quadratures: the mean field vanishes, but its square does not. Conducting plates change the allowed modes and therefore the sum of zero-point energies. The logic is clearest for a scalar field in one dimension between mirrors a distance \(d\) apart, where the mode frequencies are \(\omega_n=n\pi c/d\):

\begin{centeredalign}
E_0(d)=\sum_{n=1}^\infty\frac{\hbar\omega_n}{2}=\frac{\pi\hbar c}{2d}\sum_{n=1}^\infty n.
\end{centeredalign}

The sum diverges, but real mirrors become transparent at high frequency. A factor \(e^{-\epsilon\omega_n}\), with \(\epsilon\) of order the inverse plasma frequency of the metal, represents that transparency. The sum is then elementary, \(\sum_nne^{-\lambda n}=\lambda^{-2}-1/12+O(\lambda^2)\) with \(\lambda=\epsilon\pi c/d\), and gives

\begin{centeredalign}
E_0(d)=\frac{\hbar d}{2\pi c\,\epsilon^2}-\frac{\pi\hbar c}{24d}+O(\epsilon^2).
\end{centeredalign}

The first term is proportional to \(d\), with the same energy per length as in free space. In a long cavity divided by a movable mirror, the linear terms from the two sides add to a constant and exert no net force. The remaining term does not depend on the cutoff and gives an attraction of magnitude \(\pi\hbar c/24d^2\). For the electromagnetic field between perfectly conducting plates of area \(A\), the same steps with two polarizations and all transverse wavevectors give

\begin{centeredalign}
\frac{E}{A}=-\frac{\pi^2\hbar c}{720\,d^3},\qquad \frac{F}{A}=-\frac{\pi^2\hbar c}{240\,d^4}.
\end{centeredalign}

At \(d=1\) micrometre the pressure is 1.3 mPa, and halving the separation multiplies it by 16. This steep rise makes the force relevant to the sticking of nearby parts in micromachined devices. The ideal-conductor formula requires the separation to exceed the plasma wavelength of the metal, about 0.14 micrometre for gold. At smaller separations finite conductivity weakens the force, and Lifshitz theory with measured dielectric functions replaces the formula. At room temperature thermal photons add a correction once \(d\) approaches \(\hbar c/k_BT\approx7.6\) micrometres.

The calculation does not show that empty space has a measurable absolute energy density. The force depends on how the plates respond to the field, and it vanishes as their coupling to the field goes to zero. It can be derived equally as the retarded van der Waals attraction between fluctuating charges in the two plates, without reference to vacuum energy. What is established is how the zero-point energy changes with the boundary conditions. The coupling of bound electrons to the same field fluctuations contributes to the Lamb shift treated by quantum electrodynamics. The absolute vacuum energy and its relation to the cosmological constant remain open.

\subsection*{References For The Physical Derivation}
\begin{itemize}
\item \href{https://arxiv.org/abs/quant-ph/0106045}{M. Bordag, U. Mohideen and V. M. Mostepanenko, New Developments in the Casimir Effect; derivation, finite-conductivity and thermal corrections, measurements.}
\item \href{https://arxiv.org/abs/hep-th/0503158}{R. L. Jaffe, The Casimir Effect and the Quantum Vacuum; the force as a van der Waals interaction independent of absolute vacuum energy.}
\end{itemize}





\chapter{Dynamics and transformations}

A Hamiltonian, a unitary operator or a sum over paths tells how a prepared state changes before it is measured.

Equivalent formulations of dynamics give the same amplitudes and expectation values; perturbation theory, chaos and stochastic descriptions follow particular features of the same evolution.

\section*{Topics}

\begin{longtable}{p{0.92\linewidth}}
\toprule
Page \\
\midrule
\endhead
\phantomsection\label{page:unitary-operator}Unitary operator \\
\phantomsection\label{page:perturbation-theory}Perturbation theory \\
\phantomsection\label{page:quantum-dynamics}Quantum dynamics \\
\phantomsection\label{page:path-integral-formulation}Path integral formulation \emph{(see Path Integral)} \\
\phantomsection\label{page:hamiltonian-mechanics}Hamiltonian mechanics \\
\phantomsection\label{page:path-integral}Path integral \\
\phantomsection\label{page:hamiltonian-quantum-mechanics}Hamiltonian (quantum mechanics) \emph{(see Hamiltonian Mechanics)} \\
\phantomsection\label{page:perturbation-theory-quantum-mechanics}Perturbation theory (quantum mechanics) \emph{(see Perturbation Theory)} \\
\phantomsection\label{page:quantum-chaos}Quantum chaos \\
\phantomsection\label{page:schr-dinger-equation}Schrödinger equation \\
\phantomsection\label{page:symmetry-in-quantum-mechanics}Symmetry in quantum mechanics \\
\phantomsection\label{page:heisenberg-picture}Heisenberg picture \\
\phantomsection\label{page:quantum-stochastic-calculus}Quantum stochastic calculus \\
\phantomsection\label{page:relativistic-quantum-mechanics}Relativistic quantum mechanics \\
\phantomsection\label{page:schr-dinger-picture}Schrödinger picture \\
\phantomsection\label{page:creation-and-annihilation-operators}Creation and annihilation operators \\
\bottomrule
\end{longtable}

\clearpage

\section{Unitary operator}

\label{topic:unitary-operator}

A lossless beam splitter divides a light beam between two output ports. Whatever phase is then added in one arm of an interferometer, the detection probabilities at the outputs add to one. Transformations that preserve probabilities in this way are tightly restricted, and Wigner's theorem states the restriction.

A pure state is a ray, a unit vector defined up to a phase. Consider a bijective map of rays that preserves every transition probability \(|\langle\phi|\psi\rangle|^2\). Wigner's theorem states that the map is implemented on vectors by an operator that is either unitary or antiunitary, unique up to an overall phase:

\begin{centeredalign}
\langle U\phi|U\psi\rangle=\langle\phi|\psi\rangle,\qquad
\langle A\phi|A\psi\rangle=\langle\phi|\psi\rangle^*,\qquad
A(a\phi+b\psi)=a^*A\phi+b^*A\psi.
\end{centeredalign}

A unitary operator is linear and satisfies \(U^\dagger U=UU^\dagger=I\); an antiunitary operator is antilinear. Transformations that can be reached continuously from the identity, such as time evolution, rotations and translations, are unitary. Each is the square of a transformation half as large, and the square of either kind of operator is unitary. Their self-adjoint generators are discussed under symmetry in quantum mechanics.

A single photon in two spatial modes shows the unitary case. A lossless beam splitter with transmission amplitude \(t\) and reflection amplitude \(r\) acts on the two-component amplitude vector. Conservation of total probability for every input is equivalent to \(B^\dagger B=I\):

\begin{centeredalign}
B=\begin{pmatrix}t&r\\r&t\end{pmatrix},\qquad |t|^2+|r|^2=1,\qquad t^*r+r^*t=0.
\end{centeredalign}

The second condition places the reflected amplitude at a phase of \(\pm\pi/2\) relative to the transmitted one. Take the balanced case \(t=1/\sqrt2\) and \(r=i/\sqrt2\). A Mach--Zehnder interferometer applies the splitter, a phase \(\varphi\) in one arm and a second splitter. For a photon entering port 1, the output probabilities are

\begin{centeredalign}
U(\varphi)=B\begin{pmatrix}e^{i\varphi}&0\\0&1\end{pmatrix}B,\qquad
p_1=\sin^2\frac\varphi2,\qquad p_2=\cos^2\frac\varphi2.
\end{centeredalign}

Scanning \(\varphi\) moves probability between the detectors and keeps the sum at one, which is unitarity in measured form. Absorption in one arm, with amplitude transmission \(\eta<1\), makes the \(2\times2\) matrix non-unitary. The missing probability goes into other modes, and the fringe visibility falls to \(2\eta/(1+\eta^2)\).

Time reversal is the antiunitary case. It must leave positions unchanged, reverse momenta and turn \(e^{-iHt/\hbar}\) into \(e^{iHt/\hbar}\). A unitary operator could do this only by anticommuting with \(H\), which would map each energy \(E\) onto \(-E\). That is impossible for an energy bounded below, so time reversal is antiunitary. For one spin one-half it is \(T=-i\sigma_yK\), with \(K\) complex conjugation in the \(\sigma_z\) basis, and

\begin{centeredalign}
T^2=-I,\qquad T^2=(-1)^N\ \text{for }N\text{ electrons},\\
[H,T]=0,\quad H\psi=E\psi\ \Longrightarrow\ H(T\psi)=E\,T\psi,\quad \langle T\psi|\psi\rangle=0.
\end{centeredalign}

The orthogonality follows from antilinearity. The rule \(\langle A\phi|A\chi\rangle=\langle\chi|\phi\rangle\) with \(\phi=\psi\) and \(\chi=T\psi\) gives \(\langle T\psi|T^2\psi\rangle=\langle T\psi|\psi\rangle\). With \(T^2=-I\) the left side equals \(-\langle T\psi|\psi\rangle\), so both vanish. Every energy level of a system with an odd number of electrons is therefore at least two-fold degenerate when no magnetic field acts. This is Kramers' theorem.

The ion \(\mathrm{Er}^{3+}\) has eleven 4f electrons and a ground multiplet with \(J=15/2\), sixteen states. An electric crystal field of any symmetry, however low, is even under time reversal. It splits the multiplet into at most eight levels, each a Kramers doublet. The ion \(\mathrm{Pr}^{3+}\) has two 4f electrons and \(J=4\), and a low-symmetry field can split its nine states into nine singlets. A magnetic field is odd under time reversal and splits each Kramers doublet linearly in the field, which produces the electron paramagnetic resonance lines of odd-electron rare-earth ions.

The degeneracy survives any electric field, spin-orbit coupling and lattice distortion. It is lifted only by perturbations odd under time reversal, such as an applied magnetic field or the internal field of magnetic order. Time evolution generated by the Hamiltonian is the most important unitary map; the Fourier transform and the change between Schrödinger and Heisenberg pictures are others.

\subsection*{References For The Physical Derivation}
\begin{itemize}
\item \href{https://arxiv.org/abs/0808.0779}{R. Simon, N. Mukunda, S. Chaturvedi and V. Srinivasan, Two elementary proofs of the Wigner theorem on symmetry in quantum mechanics; symmetries of rays as unitary or antiunitary operators.}
\item \href{https://arxiv.org/abs/1508.06951}{V. Moretti, Mathematical Foundations of Quantum Mechanics: An Advanced Short Course; unitary and antiunitary operators and the representation of symmetries.}
\end{itemize}





\clearpage

\section{Perturbation theory}

\label{topic:perturbation-theory}

An electric field of 100 kV/cm splits the \(n=2\) level of hydrogen into components hundreds of gigahertz apart, while the ground level moves by only a few megahertz. Both shifts come from the same weak addition to the atomic Hamiltonian. Perturbation theory explains why one shift is linear in the field and the other quadratic.

Write \(H=H_0+\lambda V\), where the eigenstates \(|n\rangle\) and energies \(E_n^{(0)}\) of \(H_0\) are known and \(\lambda\) counts powers of the perturbation. Expand the exact eigenvalue and eigenvector in powers of \(\lambda\) and insert them into \(H|\psi_n\rangle=E_n|\psi_n\rangle\). Projecting the equation at each order onto \(\langle n|\) and onto \(\langle m|\) with \(m\neq n\) gives the energy and the admixed states:

\begin{centeredalign}
E_n=E_n^{(0)}+\lambda V_{nn}
+\lambda^2\sum_{m\neq n}\frac{|V_{mn}|^2}{E_n^{(0)}-E_m^{(0)}}+O(\lambda^3),\\
|\psi_n\rangle=|n\rangle+\lambda\sum_{m\neq n}\frac{V_{mn}}{E_n^{(0)}-E_m^{(0)}}|m\rangle+O(\lambda^2),
\qquad V_{mn}=\langle m|V|n\rangle.
\end{centeredalign}

The expansion is useful when every coupling is small compared with the corresponding gap, \(|V_{mn}|\ll|E_n^{(0)}-E_m^{(0)}|\). For the ground state every denominator is negative, so the second-order shift always lowers the energy. The sum runs over the complete set of unperturbed states, including any continuum.

For an electron in a uniform field \(F\) along \(z\), the perturbation is \(V=eFz\). The hydrogen ground state has even parity and \(z\) is odd, so \(V_{nn}=0\). The leading shift is quadratic, \(\Delta E=-\alpha_{\mathrm{pol}}F^2/2\), with static polarizability \(\alpha_{\mathrm{pol}}=\tfrac92(4\pi\varepsilon_0)a_0^3\). At 100 kV/cm this gives 5.6 MHz. The bound states supply only 81 percent of the polarizability; the ionization continuum supplies the rest.

The \(n=2\) level behaves differently because the \(2s\) and \(2p\) states have the same unperturbed energy. The denominators vanish, and the expansion must start by diagonalizing \(V\) within the degenerate subspace. Only \(2s\) and \(2p_0\) are coupled, and

\begin{centeredalign}
\langle2s|z|2p_0\rangle=-3a_0,\qquad E=E_2\pm3ea_0F,\qquad
|\pm\rangle=\frac{|2s\rangle\mp|2p_0\rangle}{\sqrt2}.
\end{centeredalign}

The mixed states carry permanent electric dipoles of \(\mp3ea_0\), so their shift is linear in \(F\). At 100 kV/cm the outer components lie \(\pm384\) GHz from the unperturbed level, while \(2p_{\pm1}\) remains unshifted. Stark observed this linear splitting of the Balmer lines in 1913.

The degenerate treatment assumes that the Stark coupling exceeds the small splittings omitted from the nonrelativistic \(H_0\). The fine-structure splitting of \(n=2\) is 10.9 GHz, and the Lamb shift between \(2s_{1/2}\) and \(2p_{1/2}\) is 1.06 GHz. Below a few hundred volts per centimetre the coupling is smaller than the Lamb shift, and the shift of \(2s_{1/2}\) becomes quadratic again. In any static field the electron can also tunnel out of the atom, at a rate that vanishes faster than any power of \(F\). The levels are therefore resonances, and the series in \(F\) is asymptotic rather than convergent. A perturbation that oscillates in time produces transitions instead of shifts. The same expansion then gives Fermi's golden rule for transition rates, used for spectral lines and scattering.

\subsection*{References For The Physical Derivation}
\begin{itemize}
\item \href{https://arxiv.org/abs/physics/0010038}{U. D. Jentschura, Resummation of the Divergent Perturbation Series for a Hydrogen Atom in an Electric Field; the Stark perturbation series and its divergence.}
\end{itemize}





\clearpage

\section{Quantum dynamics}

\label{topic:quantum-dynamics}

A picosecond laser pulse excites an atom into a superposition of Rydberg levels near \(n=50\). The electron then moves as a localized packet around a Kepler orbit. After a few revolutions it spreads around the orbit, and later it reassembles. Each stage follows from the spacing of the energy levels.

At short times the centre of a narrow packet follows Newton's law. Differentiating expectation values with the Schrödinger equation gives Ehrenfest's theorem:

\begin{centeredalign}
\frac{d\langle X\rangle}{dt}=\frac{\langle P\rangle}{m},\qquad
\frac{d\langle P\rangle}{dt}=-\langle V'(X)\rangle\approx-V'(\langle X\rangle).
\end{centeredalign}

The approximation holds while the force varies little across the packet. For a harmonic potential it is exact at all times, because \(V'\) is linear. In other potentials the packet spreads, and the centre no longer obeys a closed classical equation.

The energy spectrum shows how long the classical stage lasts. Expand the packet in energy eigenstates with weights \(|c_n|^2\) concentrated around \(\bar n\), and expand the energies to second order in \(n-\bar n\). The amplitude to return to the initial state and the two resulting times are

\begin{centeredalign}
A(t)=\sum_n|c_n|^2e^{-iE_nt/\hbar},\qquad
E_n\approx E_{\bar n}+E'(n-\bar n)+\frac{E''}{2}(n-\bar n)^2,\\
T_{\mathrm{cl}}=\frac{2\pi\hbar}{|E'|},\qquad T_{\mathrm{rev}}=\frac{4\pi\hbar}{|E''|}.
\end{centeredalign}

The linear term brings every component back in phase after the classical period \(T_{\mathrm{cl}}\). The quadratic term slowly dephases the components. With a spread \(\Delta n\) in \(n\), the packet covers the whole orbit after a time of order \(T_{\mathrm{rev}}/(2\pi\Delta n^2)\). At \(T_{\mathrm{rev}}\) every quadratic phase is a multiple of \(2\pi\), because \(n-\bar n\) is an integer, and the packet reforms. At rational fractions of \(T_{\mathrm{rev}}\) it appears as several smaller packets, the fractional revivals.

For hydrogen \(E_n=-R/n^2\) with \(R=13.6\) eV, so \(E'=2R/n^3\) and \(E''=-6R/n^4\). Then \(T_{\mathrm{cl}}=\pi\hbar n^3/R\) equals the Kepler period of the classical orbit with the same energy, and \(T_{\mathrm{rev}}=(2\bar n/3)T_{\mathrm{cl}}\). For \(\bar n=50\) this gives \(T_{\mathrm{cl}}=19\) ps and \(T_{\mathrm{rev}}=0.63\) ns. Yeazell, Mallalieu and Stroud observed the collapse and revival of such Rydberg packets in 1990.

Higher terms limit the picture. The cubic term adds a phase of about \(8(\Delta n)^3/\bar n\) at \(T_{\mathrm{rev}}\), so a clean revival requires \((\Delta n)^3\ll\bar n/8\). A direct sum for \(\bar n=50\) gives a return probability of 0.91 at the revival for \(\Delta n=1\), but only about 0.5 for \(\Delta n=2.5\). An evenly spaced spectrum, as in the quantum harmonic oscillator, never dephases. In a classically chaotic system, agreement with a single trajectory lasts only about the Ehrenfest time, of order \(\lambda^{-1}\ln(S/\hbar)\), where \(\lambda\) is the Lyapunov exponent and \(S\) a typical action. This short time scale is one signature of quantum chaos.

\subsection*{References For The Physical Derivation}
\begin{itemize}
\item \href{https://arxiv.org/abs/quant-ph/0401031}{R. W. Robinett, Quantum wave packet revivals; classical period, revival time, fractional revivals and Rydberg wave packets.}
\end{itemize}





\clearpage

\section{Path integral formulation}

\label{topic:path-integral-formulation}

A neutron interferometer is tilted so that one beam travels higher in the Earth's gravitational field than the other. The interference signal oscillates as the tilt changes, although the two beams follow nearly identical classical trajectories. The path-integral formulation computes such phases directly from the classical action.

In this formulation the amplitude to go from \(x_i\) to \(x_f\) is a sum over every history \(x(t)\) connecting them, each weighted by \(e^{iS[x]/\hbar}\). Here \(S=\int L\,dt\) is the classical action and \(L=m\dot x^2/2-V(x)\) the Lagrangian. Consider a small deformation \(x+\delta x\) with fixed endpoints. The change of action is

\begin{centeredalign}
\delta S=\int_{t_i}^{t_f}\big[-m\ddot x-V'(x)\big]\,\delta x\,dt,\qquad
\delta S=0\ \Longleftrightarrow\ m\ddot x=-V'(x).
\end{centeredalign}

Where \(\delta S\neq0\), neighbouring paths have phases that differ by many times \(2\pi\) once \(S\) is large compared with \(\hbar\), and their contributions cancel. Paths near a stationary path add with nearly equal phases. The amplitude is therefore dominated by thin tubes around classical trajectories, within which the action varies by about \(\hbar\). Expanding the action to second order about each classical path gives the semiclassical propagator

\begin{centeredalign}
K\simeq\sum_{\mathrm{cl}}\frac{1}{\sqrt{2\pi i\hbar}}
\Big|\frac{\partial^2S_{\mathrm{cl}}}{\partial x_i\,\partial x_f}\Big|^{1/2}
\exp\Big(\frac{i}{\hbar}S_{\mathrm{cl}}-\frac{i\pi\mu}{2}\Big).
\end{centeredalign}

Here \(S_{\mathrm{cl}}\) is the action of a classical path between the endpoints, and \(\mu\) counts the focal points the path passes. For a free particle, \(S_{\mathrm{cl}}=m(x_f-x_i)^2/2T\), and the prefactor reproduces \(\sqrt{m/2\pi i\hbar T}\). For Lagrangians at most quadratic in position and velocity the expression is exact. The approximation fails near caustics, where classical paths merge and the prefactor diverges. As \(S/\hbar\) grows, interference fringes become finer than any detector resolves, and only the classical motion remains visible.

A uniform gravitational field gives a linear potential, so the semiclassical form is exact and only classical paths are needed. A neutron beam has a fixed energy rather than a fixed travel time. Evaluating the time integral of the propagator by stationary phase then replaces \(S\) by the reduced action, and the phase along a path becomes \(\hbar^{-1}\int p\,dx\). Let the two beams enclose a parallelogram of area \(A=H\ell\) in a vertical plane. Their horizontal segments of length \(\ell\) differ in height by \(H\), and the sloping segments contribute equally to both. Energy conservation lowers the momentum on the upper segment from \(p=h/\lambda\) to about \(p-m^2gH/p\), so that

\begin{centeredalign}
\Delta\varphi=\frac{\ell}{\hbar}\,\frac{m^2gH}{p}=\frac{mgH\ell}{\hbar v}
=\frac{2\pi m^2g\lambda A}{h^2}.
\end{centeredalign}

For neutrons of wavelength \(0.1445\) nm and an area of \(10\,\mathrm{cm}^2\) in a vertical plane, \(\Delta\varphi\approx57\) rad, about nine fringes. Colella, Overhauser and Werner observed this gravitationally induced phase in 1975 by rotating a silicon-crystal neutron interferometer about the incident beam. The classical trajectories do not depend on mass; the mass enters the phase through the ratio \(m/\hbar\). At fixed wavelength the phase grows as \(m^2\).

Light-pulse atom interferometers use the same calculation with laser phases added at each pulse. Their phase \(k_{\mathrm{eff}}gT^2\) exceeds \(10^6\) rad for rubidium atoms with \(T=0.1\) s, and such gravimeters reach parts-per-billion precision in \(g\). The derivation of the propagator by time slicing, and its free-particle Gaussian form, are given under path integral.

\subsection*{References For The Physical Derivation}
\begin{itemize}
\item \href{https://arxiv.org/abs/quant-ph/0004090}{R. MacKenzie, Path Integral Methods and Applications; stationary phase, the classical limit and the semiclassical propagator.}
\item \href{https://arxiv.org/abs/0712.3703}{A. D. Cronin, J. Schmiedmayer and D. E. Pritchard, Atom Interferometers; gravitationally induced phases in neutron and atom interferometers and light-pulse gravimetry.}
\end{itemize}





\clearpage

\section{Hamiltonian mechanics}

\label{topic:hamiltonian-mechanics}

A pendulum, a planet or a charged particle in a magnet is fully described at one instant by its coordinates and momenta. As time passes, each point of this phase space moves, and a small patch of initial conditions deforms while keeping its area. Hamiltonian mechanics generates this motion from one function, the energy \(H(q,p)\).

Start from a Lagrangian \(L(q,\dot q)\) and define the canonical momentum \(p=\partial L/\partial\dot q\). The Legendre transform \(H=p\dot q-L\), written as a function of \(q\) and \(p\), yields Hamilton's equations. For any phase-space function \(f(q,p)\) they combine into one rule through the Poisson bracket:

\begin{centeredalign}
\dot q_i=\frac{\partial H}{\partial p_i},\qquad \dot p_i=-\frac{\partial H}{\partial q_i},
\qquad \frac{df}{dt}=\{f,H\},\\
\{f,g\}=\sum_i\Big(\frac{\partial f}{\partial q_i}\frac{\partial g}{\partial p_i}
-\frac{\partial f}{\partial p_i}\frac{\partial g}{\partial q_i}\Big),\qquad \{q_i,p_j\}=\delta_{ij}.
\end{centeredalign}

Three consequences follow. A function \(G\) with \(\{G,H\}=0\) is conserved, and the flow that \(G\) generates leaves \(H\) unchanged, which links symmetry to conservation as in Noether's theorem. The flow has zero divergence, so phase-space volume is conserved, which is Liouville's theorem. The bracket of a coordinate with its momentum is the same constant for every system.

Quantum mechanics keeps this structure and changes the bracket. Dirac's rule replaces functions by operators and the Poisson bracket by a commutator divided by \(i\hbar\):

\begin{centeredalign}
\{f,g\}\;\longrightarrow\;\frac{1}{i\hbar}\big[\hat f,\hat g\big],\qquad
[\hat q_i,\hat p_j]=i\hbar\,\delta_{ij},\qquad
\frac{d\hat f}{dt}=\frac{1}{i\hbar}\big[\hat f,\hat H\big].
\end{centeredalign}

The equation of motion keeps its form and becomes the Heisenberg equation. Conserved quantities and symmetry generators carry over, and the conservation of phase-space volume is replaced by the unitarity of time evolution.

A charge \(q\) of mass \(m\) moving in the plane perpendicular to a uniform field \(B\), with \(qB>0\), shows the rule at work. In the symmetric gauge the canonical momenta are \(p_x\) and \(p_y\), but the velocity is fixed by the kinetic momenta \(\pi_x=p_x+qBy/2\) and \(\pi_y=p_y-qBx/2\). Their Poisson bracket does not vanish:

\begin{centeredalign}
\{\pi_x,\pi_y\}=qB\;\longrightarrow\;[\hat\pi_x,\hat\pi_y]=i\hbar qB,\qquad
H=\frac{\hat\pi_x^2+\hat\pi_y^2}{2m},\\
E_n=\hbar\omega_c\Big(n+\frac12\Big),\qquad \omega_c=\frac{qB}{m}.
\end{centeredalign}

Classically, the bracket produces circular motion at the cyclotron frequency \(\omega_c\). Quantum mechanically, \(\hat\pi_x/\sqrt{qB}\) and \(\hat\pi_y/\sqrt{qB}\) form a canonical pair, and \(H\) becomes a harmonic oscillator of frequency \(\omega_c\). These Landau levels are spaced by 0.116 meV for a free electron in 1 T, equivalent to a temperature of 1.3 K. The orbit centre forms a second canonical pair that commutes with \(H\), so each level holds \(qB/h\) states per unit area. Landau quantization underlies the Shubnikov--de Haas oscillations and the quantum Hall effect, seen when the level spacing exceeds the thermal energy.

Dirac's rule is exact for the canonical pairs and for functions at most quadratic in them. It cannot hold for all higher polynomials at once: the Groenewold--van Hove theorem shows that no consistent map sends every Poisson bracket to the corresponding commutator. Operator ordering, such as \(qp\to(\hat q\hat p+\hat p\hat q)/2\), then matters, and the brackets differ at order \(\hbar^2\). The classical Hamiltonian guides the choice of quantum Hamiltonian but does not fix it uniquely. The operator side of this correspondence is developed under the canonical commutation relation and the Heisenberg picture.

\subsection*{References For The Physical Derivation}
\begin{itemize}
\item \href{https://arxiv.org/abs/math-ph/9809011}{M. J. Gotay, Obstructions to Quantization; the Groenewold-van Hove theorem and the limits of the Poisson-bracket to commutator rule.}
\item \href{https://arxiv.org/abs/1606.06687}{D. Tong, Lectures on the Quantum Hall Effect; Landau levels from the kinetic momenta, their degeneracy and the quantum Hall effect.}
\end{itemize}





\clearpage

\section{Path integral}

\label{topic:path-integral}

A particle released near one point at time \(t_i\) is later detected near another point at time \(t_f\). The amplitude for this transition is the propagator \(K(x_f,t_f;x_i,t_i)=\langle x_f|U(t_f-t_i)|x_i\rangle\). It evolves any initial wave function by \(\psi(x_f,t_f)=\int K(x_f,t_f;x_i,t_i)\psi(x_i,t_i)\,dx_i\). The path integral computes \(K\) by cutting the time interval into many short steps.

Divide \(T=t_f-t_i\) into \(N\) steps of length \(\epsilon=T/N\), and insert a complete set of position states at each intermediate time. For \(H=P^2/2m+V(X)\), the Trotter product formula factors each short step into a kinetic and a potential part, with an error of order \(\epsilon^2\) per step. The kinetic factor is Gaussian in momentum, and integrating over momentum gives

\begin{centeredalign}
K=\lim_{N\to\infty}\Big(\frac{m}{2\pi i\hbar\epsilon}\Big)^{N/2}
\int dx_1\cdots dx_{N-1}\,e^{iS_N/\hbar},\\
S_N=\sum_{k=0}^{N-1}\epsilon\Big[\frac m2\Big(\frac{x_{k+1}-x_k}{\epsilon}\Big)^2-V(x_k)\Big],
\qquad x_0=x_i,\quad x_N=x_f.
\end{centeredalign}

Each set of intermediate positions defines a broken-line path, and \(S_N\) is the classical action along it. The limit is written \(\int\mathcal Dx\,e^{iS[x]/\hbar}\). The measure has no meaning apart from this limit: the factors \((m/2\pi i\hbar\epsilon)^{1/2}\) diverge as \(\epsilon\to0\) and are compensated by the oscillating integrals. In imaginary time the weights become Gaussian probabilities, the limit exists as a Wiener integral, and typical paths are continuous but nowhere differentiable.

For a free particle every integral is Gaussian. Integrating over one intermediate point joins two free kernels of durations \(t_1\) and \(t_2\) into one of duration \(t_1+t_2\), because completing the square leaves an exponent quadratic in the endpoints. Repeating the step \(N-1\) times gives the exact free propagator:

\begin{centeredalign}
\int K_0(x_f,x;t_2)\,K_0(x,x_i;t_1)\,dx=K_0(x_f,x_i;t_1+t_2),\\
K_0(x_f,x_i;T)=\sqrt{\frac{m}{2\pi i\hbar T}}\,
\exp\Big[\frac{im(x_f-x_i)^2}{2\hbar T}\Big].
\end{centeredalign}

This kernel solves the free Schrödinger equation and tends to \(\delta(x_f-x_i)\) as \(T\to0\). Its modulus is the same for every endpoint. A particle released from a point therefore reaches all positions with equal amplitude magnitude, while the phase records the velocity \((x_f-x_i)/T\) needed to arrive.

Applied to a packet of width \(\sigma\), the kernel gives the time-of-flight relation. For \(T\gg m\sigma^2/\hbar\) the phase \(mx_i^2/2\hbar T\) is negligible across the packet. The remaining integral is a Fourier transform evaluated at the momentum \(mx/T\):

\begin{centeredalign}
|\psi(x,T)|^2\;\to\;\frac mT\,\Big|\widetilde\psi\Big(\frac{mx}{T}\Big)\Big|^2,
\qquad T\gg\frac{m\sigma^2}{\hbar}.
\end{centeredalign}

Absorption images of ultracold atoms taken after release from a trap therefore record the momentum distribution. Two condensates released a distance \(d\) apart overlap with fringes of spacing \(hT/(md)\), as observed by Andrews and colleagues in 1997. For sodium with \(d=40\,\mu\mathrm m\) and \(T=40\) ms the spacing is about \(17\,\mu\mathrm m\).

The relation assumes free flight. Gravity adds a linear potential, for which the slicing is again exactly Gaussian, whereas interactions during the early expansion must be treated separately. For a charge in a magnetic field the potential term depends on velocity, and the midpoint of each step must be used to obtain the correct Hamiltonian. How the integral is dominated by classical paths when the action is large compared with \(\hbar\) is treated under path integral formulation.

\subsection*{References For The Physical Derivation}
\begin{itemize}
\item \href{https://arxiv.org/abs/quant-ph/0004090}{R. MacKenzie, Path Integral Methods and Applications; time slicing, the free-particle Gaussian integrals and the propagator.}
\item \href{https://arxiv.org/abs/hep-th/9302097}{C. Grosche, An Introduction into the Feynman Path Integral; the lattice definition of the path integral and its ordering prescriptions.}
\end{itemize}





\clearpage

\section{Hamiltonian (quantum mechanics)}

\label{topic:hamiltonian-quantum-mechanics}

Ammonia absorbs microwaves near 24 GHz, and its line at 23.87 GHz was used in the ammonia maser of 1954. The same frequency describes a motion: a nitrogen atom placed on one side of the plane of the three hydrogens tunnels to the other side and back. One operator, the Hamiltonian \(H\), fixes both the energy outcomes and this time evolution.

For a closed system with time-independent \(H\), the energy eigenstates \(|E_n\rangle\) form a basis, and evolution multiplies each component by its own phase. Write the prepared state as \(|\psi\rangle=\sum_nc_n|E_n\rangle\). For any observable \(O\) with matrix elements \(O_{mn}=\langle E_m|O|E_n\rangle\),

\begin{centeredalign}
H|E_n\rangle=E_n|E_n\rangle,\qquad
U(t)=e^{-iHt/\hbar}=\sum_ne^{-iE_nt/\hbar}|E_n\rangle\langle E_n|,\\
\langle O\rangle(t)=\sum_{m,n}c_m^*c_n\,O_{mn}\,e^{i(E_m-E_n)t/\hbar}.
\end{centeredalign}

Only energy differences enter, so a constant added to \(H\) changes no prediction. The weights \(|c_n|^2\) never change, which is energy conservation. Every observable oscillates at the Bohr frequencies \((E_m-E_n)/\hbar\), the same frequencies a spectrometer records as absorption or emission lines. An eigenstate acquires only an overall phase, so all its expectation values are constant. Motion requires a superposition of different energies.

In ammonia the umbrella-inversion coordinate has two equivalent minima. Let \(|L\rangle\) and \(|R\rangle\) be states localized in each well with equal energy \(E_0\), and let \(A>0\) be the tunnelling amplitude through the barrier between them. Within this pair of states,

\begin{centeredalign}
H=\begin{pmatrix}E_0&-A\\-A&E_0\end{pmatrix},\qquad E_\pm=E_0\pm A,\qquad
P_L(t)=\cos^2\frac{\Delta E\,t}{2\hbar},\qquad \Delta E=2A.
\end{centeredalign}

The symmetric combination of \(|L\rangle\) and \(|R\rangle\) is the lower level, and the antisymmetric combination lies higher by \(\Delta E\). A molecule prepared in \(|L\rangle\) is an equal superposition of the two. Its probability \(P_L\) of remaining on the left therefore oscillates at the line frequency \(\Delta E/h\). A line at 23.87 GHz gives \(\Delta E=9.9\times10^{-5}\) eV, or \(0.80\,\mathrm{cm}^{-1}\), and a crossing time \(\pi\hbar/\Delta E\approx21\) ps.

The maser used the same two levels. An inhomogeneous electric field shifts them in opposite directions, so it separates molecules in the upper level from those in the lower level. Molecules in the upper level were focused into a microwave cavity, where they emitted at \(\Delta E/h\).

The two-level description holds because the next vibrational level of the umbrella motion lies near \(950\,\mathrm{cm}^{-1}\), far above the splitting. Rotation changes the splitting slightly, which spreads the inversion lines over a band near 24 GHz. More generally, \(H\) must be self-adjoint for \(U(t)\) to be unitary, and bounded below for a ground state to exist. For a particle in a potential, \(H=P^2/2m+V(X)\) follows from the classical energy, as described under Hamiltonian mechanics. When \(H\) depends on time, energy is not conserved and \(U\) becomes a time-ordered product, treated under the Schrödinger picture.

\subsection*{References For The Physical Derivation}
\begin{itemize}
\item \href{https://arxiv.org/abs/1508.06951}{V. Moretti, Mathematical Foundations of Quantum Mechanics: An Advanced Short Course; self-adjoint Hamiltonians, spectral decomposition and the unitary evolution they generate.}
\end{itemize}





\clearpage
\section{Perturbation theory (quantum mechanics)}
\label{topic:perturbation-theory-quantum-mechanics}
Perturbation theory (quantum mechanics) is an alternative name for Perturbation Theory.
The physical derivation is given in \hyperref[topic:perturbation-theory]{Perturbation Theory}.

\clearpage

\section{Quantum chaos}

\label{topic:quantum-chaos}

Slow neutrons striking heavy nuclei excite many narrow resonances, and neighbouring resonances are rarely found close together. The levels of a particle in a rectangular box with incommensurate sides behave differently, because near coincidences are common there. The difference reflects the classical motion, which is regular in the box and chaotic or highly complex in the nucleus.

Consider two neighbouring levels that may be coupled. When the system is invariant under time reversal, the Hamiltonian can be chosen real and symmetric. The \(2\times2\) block and its level spacing are

\begin{centeredalign}
H=\begin{pmatrix}h_{11}&h_{12}\\h_{12}&h_{22}\end{pmatrix},\qquad
s=E_+-E_-=\sqrt{(h_{11}-h_{22})^2+4h_{12}^2}.
\end{centeredalign}

A degeneracy requires two independent conditions, \(h_{11}=h_{22}\) and \(h_{12}=0\). If the matrix elements are spread smoothly, the probability of a spacing below \(s\) scales as \(s^2\), and the spacing density vanishes linearly. This is level repulsion, described by von Neumann and Wigner. Independent Gaussian matrix elements define the Gaussian orthogonal ensemble. Scaling the spacing to unit mean gives the Wigner surmise. Without time-reversal invariance \(h_{12}\) is complex, three conditions are needed, and the density vanishes as \(s^2\):

\begin{centeredalign}
P_{\mathrm{Poisson}}(s)=e^{-s},\qquad
P_{\mathrm{GOE}}(s)=\frac{\pi s}{2}\,e^{-\pi s^2/4},\\
P_{\mathrm{GUE}}(s)=\frac{32s^2}{\pi^2}\,e^{-4s^2/\pi}.
\end{centeredalign}

In an integrable system, levels carry independent quantum numbers, and levels with different labels do not repel. Their sequence is uncorrelated, which gives the Poisson law, as Berry and Tabor argued for generic integrable systems. The laws differ most at small spacing. The probability of a spacing below one tenth of the mean is 0.095 for Poisson statistics, 0.0078 for the Wigner surmise and 0.0011 for the unitary ensemble.

Bohigas, Giannoni and Schmit conjectured in 1984 that quantum systems with a chaotic classical limit follow these random-matrix laws. The conjecture is supported by stadium and Sinai billiards, by microwave cavities of the same shapes, by hydrogen in strong magnetic fields and by the nuclear data ensemble of neutron resonances. Semiclassical sums over periodic orbits explain the spectral correlations. A general proof remains open.

These statistics require care. Before any comparison the spectrum is unfolded, by rescaling energies with the smooth level density so that the mean spacing is one everywhere. Levels must also be sorted by exact symmetry. Independent sequences of different parity or angular momentum superpose without repulsion and look Poissonian. The random-matrix laws hold only over energy ranges smaller than \(h\) divided by the period of the shortest periodic orbit. Eigenstates concentrated along unstable periodic orbits, called scars, show that eigenfunctions also deviate from random behaviour.

In the time domain, chaos limits how long a wave packet follows a single classical trajectory, a theme of quantum dynamics. The random-matrix character of chaotic eigenstates also underlies thermalization of isolated systems in quantum statistical mechanics.

\subsection*{References For The Physical Derivation}
\begin{itemize}
\item \href{https://arxiv.org/abs/cond-mat/9707301}{T. Guhr, A. Mueller-Groeling and H. A. Weidenmueller, Random Matrix Theories in Quantum Physics: Common Concepts; level spacing statistics, the Wigner surmise, nuclear data and chaotic billiards.}
\item \href{https://arxiv.org/abs/1509.06411}{L. D'Alessio, Y. Kafri, A. Polkovnikov and M. Rigol, From Quantum Chaos and Eigenstate Thermalization to Statistical Mechanics and Thermodynamics; random-matrix statistics and thermalization of isolated systems.}
\end{itemize}





\clearpage

\section{Schrödinger equation}

\label{topic:schr-dinger-equation}

Atomic hydrogen absorbs ultraviolet light only at sharp wavelengths, and its Lyman series ends at 91.2 nm, the ionization threshold of 13.6 eV. A wave equation for the electron's amplitude must reproduce these numbers. Its form follows from the de Broglie relations and the nonrelativistic energy.

A plane wave \(e^{i(\mathbf k\cdot\mathbf x-\omega t)}\) with \(\mathbf p=\hbar\mathbf k\) and \(E=\hbar\omega\) satisfies \(i\hbar\partial_t\psi=E\psi\) and \(-i\hbar\nabla\psi=\mathbf p\psi\). Requiring \(E=p^2/2m+V\) for every component, and linearity so that superpositions of solutions remain solutions, gives

\begin{centeredalign}
i\hbar\frac{\partial\psi}{\partial t}=-\frac{\hbar^2}{2m}\nabla^2\psi+V(\mathbf x)\,\psi,
\qquad i\hbar\frac{d}{dt}|\psi(t)\rangle=H|\psi(t)\rangle.
\end{centeredalign}

For a constant potential this argument is complete. For a varying \(V\) the equation is postulated, and its predictions test it. It contains a single time derivative, so the wave function at one instant fixes it at all later times. For real \(V\), multiplying by \(\psi^*\) and subtracting the complex conjugate equation gives a local conservation law:

\begin{centeredalign}
\frac{\partial|\psi|^2}{\partial t}+\nabla\cdot\mathbf j=0,\qquad
\mathbf j=\frac{\hbar}{m}\,\mathrm{Im}\big(\psi^*\nabla\psi\big).
\end{centeredalign}

Probability flows from one region to another, but it is not created or destroyed. An imaginary part of \(V\) would act as a source or sink, which is how absorption into unobserved channels is sometimes modelled. If \(V\) does not depend on time, separable solutions \(\psi=\phi(\mathbf x)e^{-iEt/\hbar}\) reduce the equation to the eigenvalue problem \(H\phi=E\phi\). Normalizable solutions exist only for particular energies, which is the origin of discrete levels. A general preparation is a superposition of such stationary states, and its probability density changes in time only through interference between components of different energy.

For hydrogen, \(V=-k/r\) with \(k=e^2/4\pi\varepsilon_0\). Try the spherically symmetric function \(\phi=e^{-r/a}\), for which \(\nabla^2\phi=(1/a^2-2/ar)\phi\). Substitution gives

\begin{centeredalign}
H\phi=\Big[-\frac{\hbar^2}{2ma^2}+\Big(\frac{\hbar^2}{ma}-k\Big)\frac1r\Big]\phi,\\
a_0=\frac{\hbar^2}{mk}=0.0529\ \mathrm{nm},\qquad
E_1=-\frac{mk^2}{2\hbar^2}=-13.6\ \mathrm{eV}.
\end{centeredalign}

The \(1/r\) term must vanish for \(\phi\) to be an eigenfunction, which fixes the Bohr radius \(a_0\). The remaining constant is the ground-state energy. The full solution gives \(E_n=E_1/n^2\), so the Lyman limit lies at \(hc/|E_1|=91.1\) nm. Replacing \(m\) by the reduced mass of electron and proton moves it to the observed 91.2 nm.

The equation treats the electron as nonrelativistic. In hydrogen its typical speed is \(\alpha c\), with \(\alpha\approx1/137\), so corrections of relative size \(\alpha^2\approx5\times10^{-5}\) appear. They split levels of equal \(n\) into the fine structure described under relativistic quantum mechanics. Spin, nuclear size and the coupling to the radiation field add smaller corrections, such as the Lamb shift. The abstract form \(i\hbar\,d|\psi\rangle/dt=H|\psi\rangle\) applies to any closed system, and the organization of its solutions in time is treated under the Schrödinger picture.

\subsection*{References For The Physical Derivation}
\begin{itemize}
\item \href{https://arxiv.org/abs/2007.14798}{A. Galler, J. Canfield and J. K. Freericks, Schrödinger's original quantum-mechanical solution for hydrogen; the 1926 wave-mechanical derivation of the hydrogen spectrum.}
\end{itemize}





\clearpage

\section{Symmetry in quantum mechanics}

\label{topic:symmetry-in-quantum-mechanics}

A rotation of an isolated atom changes the description of its orientation while leaving transition probabilities unchanged. Such transformations act on quantum rays and are represented, under the hypotheses of Wigner's theorem, by unitary or antiunitary maps. Continuous transformations connected to the identity are represented unitarily; their generators turn an invariance of probabilities into a relation between operators.

Let a one-parameter symmetry be \(U(\epsilon)=\exp(-i\epsilon G/\hbar)\), with self-adjoint generator \(G\) and a parameter whose units make the exponent dimensionless. For a time-independent Hamiltonian, invariance under this group implies

\begin{centeredalign}
U(\epsilon)HU(\epsilon)^\dagger=H,
\qquad [G,H]=0,\qquad \frac{d}{dt}\langle G\rangle=0.
\end{centeredalign}

The last equality follows from the Heisenberg equation when \(G\) has no explicit time dependence and the operator domains permit the calculation. Spatial translations give momentum conservation; rotations give angular momentum conservation. A boundary or an external field can break the symmetry even if a local bulk term remains invariant. Conservation therefore belongs to the full physical problem.

Symmetry also controls which transitions can occur. Suppose parity \(\Pi\) commutes with \(H\) and two nondegenerate states have definite parities \(\pi_i,\pi_f=\pm1\). Position is odd under parity. Inserting \(\Pi^\dagger\Pi\) into its matrix element yields

\begin{centeredalign}
\langle f|X|i\rangle=-\pi_f\pi_i\langle f|X|i\rangle.
\end{centeredalign}

The electric-dipole matrix element vanishes for states of the same parity. This is a selection rule for the pair consisting of the dynamics and the coupling to the probe. It can suppress a spectral line without creating an additional conserved quantity. Conversely, an accidental equality of two energy differences can merge lines without imposing any selection rule. These mechanisms must be distinguished when a spectral pattern simplifies.

For a proposed transfer between physical systems, a shared symmetry is useful because it restricts the admissible target interactions and observables. It does not determine their coupling constants or all matrix elements. The construction becomes predictive only after the target representation, Hamiltonian and measurement are fixed. Symmetry then supplies exact relations, such as forbidden transitions, that can test the proposed realization.

\subsection*{References For The Physical Derivation}
\begin{itemize}
\item \href{https://arxiv.org/abs/1508.06951}{V. Moretti, Mathematical Foundations of Quantum Mechanics: An Advanced Short Course; states, operators and symmetry.}
\end{itemize}





\clearpage

\section{Heisenberg picture}

\label{topic:heisenberg-picture}

The Schrödinger and Heisenberg pictures assign time dependence to different parts of the same prediction. With unitary evolution \(U(t)\), an initial density operator \(\rho_0\) and an observable \(O\), the expectation can be evaluated either by evolving the state or by evolving the operator:

\begin{centeredalign}
\rho_S(t)=U(t)\rho_0U(t)^\dagger,\qquad
O_H(t)=U(t)^\dagger OU(t),\qquad
\operatorname{Tr}[\rho_S(t)O]=\operatorname{Tr}[\rho_0O_H(t)].
\end{centeredalign}

Differentiating \(O_H\) for a time-independent Hamiltonian \(H\) gives the Heisenberg equation:

\begin{centeredalign}
\dot O_H=\frac{i}{\hbar}[H,O_H].
\end{centeredalign}

The resulting commutators identify which other observables enter the prediction. This is particularly useful when an experiment observes only a small part of a many-body system.

Consider two spins with \(H=\hbar g\,Z_1Z_2\), where \(X_j,Y_j,Z_j\) denote dimensionless Pauli operators and \(g\) is a frequency. If the measured quantity is \(x=\langle X_1\rangle\), its derivative contains the correlation \(c=\langle Y_1Z_2\rangle\). A second commutator closes the pair:

\begin{centeredalign}
\dot x=-2gc,\qquad \dot c=2gx,\qquad
x(t)=x(0)\cos(2gt)-c(0)\sin(2gt).
\end{centeredalign}

The correlation is physically necessary: preparations with the same initial transverse magnetization but different \(c(0)\) give different later magnetizations. Its appearance is fixed by the interaction algebra. Retaining the two expectations supplies a complete closed prediction for this observable even though it does not reconstruct every entry of the two-spin density operator.

Eliminating the correlation replaces the pair of first-order equations by a history-dependent equation for \(x\). The initial correlation survives as a separate term:

\begin{centeredalign}
\dot x(t)=-2g c(0)-4g^2\int_0^t x(s)\,ds.
\end{centeredalign}

Both descriptions give the same signal when they use the same preparation. Discarding the integral or the initial-correlation term gives a different physical prediction. Repeated commutators thus provide a concrete way to construct the set of observables needed by a measurement, and to identify when a smaller description requires memory. For more complicated interactions the sequence may generate a much larger space rather than closing after two steps.

\subsection*{References For The Physical Derivation}
\begin{itemize}
\item \href{https://arxiv.org/abs/1508.06951}{V. Moretti, Mathematical Foundations of Quantum Mechanics: An Advanced Short Course; states, operators and symmetry.}
\end{itemize}

\subsection*{Relations In The Original Papers}

\href{https://arxiv.org/html/1508.06951#S3.Ex177}{arXiv:1508.06951, S3.Ex177}. State evolution and observable evolution give identical spectral probabilities. The source chooses its unitary convention explicitly. The two-spin calculation uses U = exp(-iHt/hbar).



\clearpage

\section{Quantum stochastic calculus}

\label{topic:quantum-stochastic-calculus}

Light leaks out of an optical cavity through a partly transmitting mirror. The emerging beam carries the cavity field, and the same mirror admits vacuum fluctuations from outside. Quantum stochastic calculus describes this exchange with the external field in infinitesimal time steps, in the way Itô calculus describes classical noise.

Treat the external field as a sequence of independent time bins of length \(dt\) that pass the mirror one after another. Let \(b_{\mathrm{in}}(t)\) be the incoming field operator, with \([b_{\mathrm{in}}(t),b_{\mathrm{in}}^\dagger(t')]=\delta(t-t')\). Define the increment \(dA_t=\int_t^{t+dt}b_{\mathrm{in}}(s)\,ds\). For an incoming vacuum, products of increments obey the quantum Itô table:

\begin{centeredalign}
dA_t\,dA_t^\dagger=dt,\qquad dA_t^\dagger\,dA_t=dA_t\,dA_t=dA_t^\dagger\,dA_t^\dagger=0,
\qquad dt\,dA_t=0.
\end{centeredalign}

The asymmetry between the two orderings is the vacuum noise. Creating and then removing a quantum in an empty bin leaves a term proportional to \(dt\), whereas the reverse order gives zero. A classical Wiener increment gives \(dW^2=dt\) for either order. For a thermal input with mean occupation \(N\), the table becomes \(dA\,dA^\dagger=(N+1)\,dt\) and \(dA^\dagger dA=N\,dt\).

Let the system couple to the field through an operator \(L\), for example \(L=\sqrt\kappa\,a\) for a cavity mode \(a\) with energy decay rate \(\kappa\). Hudson and Parthasarathy showed that the joint unitary evolution of system and field obeys

\begin{centeredalign}
dU_t=\Big[-\Big(\frac{i}{\hbar}H+\frac12L^\dagger L\Big)dt
+L\,dA_t^\dagger-L^\dagger dA_t\Big]U_t.
\end{centeredalign}

The term \(-\frac12L^\dagger L\,dt\) is required by unitarity. In \(d(U^\dagger U)=dU^\dagger U+U^\dagger dU+dU^\dagger dU\), the last product contributes \(L^\dagger L\,dA\,dA^\dagger=L^\dagger L\,dt\), which cancels it exactly. Averaging over the vacuum field gives the Lindblad master equation written under quantum jump. The calculus also yields the outgoing field: it equals the incoming field plus the emission of the system,

\begin{centeredalign}
dA^{\mathrm{out}}_t=L\,dt+dA_t,\qquad
b_{\mathrm{out}}=b_{\mathrm{in}}+\sqrt\kappa\,a.
\end{centeredalign}

For a cavity driven through the same mirror by a coherent beam of amplitude \(\beta\), detuned by \(\Delta\) from resonance, the mean cavity field obeys \(\dot a=-(i\Delta+\kappa/2)a-\sqrt\kappa\,\beta\). In the steady state the reflection coefficient is

\begin{centeredalign}
r(\Delta)=\frac{b_{\mathrm{out}}}{\beta}=\frac{2i\Delta-\kappa}{2i\Delta+\kappa},\qquad
|r|=1,\qquad r(0)=-1.
\end{centeredalign}

All light returns, as expected for a lossless single-sided cavity, but its phase \(\pi-2\arctan(2\Delta/\kappa)\) turns by \(2\pi\) across a window of width about \(\kappa\). A superconducting qubit that shifts the cavity frequency by \(\pm\chi\) therefore imprints reflected phases that differ by \(4\arctan(2\chi/\kappa)\). The difference reaches \(\pi\) at \(\chi=\kappa/2\), where the two reflected coherent states are most distinguishable for a given input flux. Pound--Drever--Hall laser locking uses the same phase near resonance.

The calculus rests on a Markov approximation. The coupling to the external field must be flat in frequency over a band much wider than \(\kappa\), and the rotating-wave approximation requires \(\kappa\) far below the cavity frequency. Optical cavities and microwave resonators satisfy both conditions well. Structured reservoirs, such as photonic band edges, require a memory kernel instead. Photon counting on the same output field gives the quantum-jump trajectories of quantum optics.

\subsection*{References For The Physical Derivation}
\begin{itemize}
\item \href{https://arxiv.org/abs/1408.5686}{K. R. Parthasarathy, Quantum Stochastic Calculus and Quantum Gaussian Processes; the Hudson-Parthasarathy calculus on boson Fock space.}
\item \href{https://arxiv.org/abs/1611.00375}{J. Combes, J. Kerckhoff and M. Sarovar, The SLH framework for modeling quantum input-output networks; quantum stochastic differential equations and input-output relations.}
\item \href{https://arxiv.org/abs/0810.4729}{A. A. Clerk, M. H. Devoret, S. M. Girvin, F. Marquardt and R. J. Schoelkopf, Introduction to Quantum Noise, Measurement and Amplification; cavity input-output theory and dispersive qubit measurement.}
\end{itemize}





\clearpage

\section{Relativistic quantum mechanics}

\label{topic:relativistic-quantum-mechanics}

The \(n=2\) level of hydrogen is not a single level. Spectroscopy resolves components 10.9 GHz apart, about \(1.3\times10^{-5}\) of the binding energy, while the Schrödinger equation predicts exact degeneracy. The electron moves at about \(\alpha c\), so relativistic kinematics is the natural place to look for the splitting.

Replace the nonrelativistic energy by \(E^2=p^2c^2+m^2c^4\) and substitute \(E\to i\hbar\partial_t\) and \(\mathbf p\to-i\hbar\nabla\). For a spin-zero particle in an electrostatic potential energy \(V\), this gives the Klein--Gordon equation and its conserved density:

\begin{centeredalign}
\big(i\hbar\partial_t-V\big)^2\phi=\big(-\hbar^2c^2\nabla^2+m^2c^4\big)\phi,\\
\rho=\frac{i\hbar}{2mc^2}\big(\phi^*\partial_t\phi-\phi\,\partial_t\phi^*\big)-\frac{V}{mc^2}|\phi|^2.
\end{centeredalign}

The equation is of second order in time, so \(\phi\) and \(\partial_t\phi\) are independent initial data, and \(\rho\) is not positive. A stationary solution with energy \(E\) has \(\rho=(E-V)|\phi|^2/mc^2\), which is negative for the negative-energy solutions that the quadratic relation allows. The density is therefore a charge density, not a probability density. A single-particle reading is consistent only when the potential is weak compared with \(2mc^2\) and varies slowly over the Compton wavelength \(\hbar/mc\), which is 386 fm for an electron. Stronger or steeper potentials create particle-antiparticle pairs, as the Klein paradox shows, and require quantum field theory.

For an attractive Coulomb potential \(V=-Z\alpha\hbar c/r\), the radial equation keeps the hydrogen form with \(l(l+1)\) replaced by \(l(l+1)-(Z\alpha)^2\). The exact levels and their expansion in \(Z\alpha\) are

\begin{centeredalign}
E_{nl}=mc^2\Big[1+\frac{(Z\alpha)^2}
{\big(n-l-\frac12+\sqrt{(l+\frac12)^2-(Z\alpha)^2}\big)^2}\Big]^{-1/2},\\
E_{nl}\approx mc^2\Big[1-\frac{(Z\alpha)^2}{2n^2}
-\frac{(Z\alpha)^4}{2n^4}\Big(\frac{n}{l+\frac12}-\frac34\Big)\Big].
\end{centeredalign}

The second term is the Bohr energy. The third depends on \(l\), so relativistic kinematics does lift the degeneracy. For hydrogen it is smaller than the Bohr energy by a factor of order \(\alpha^2\), about \(10^{-5}\), which matches the observed scale of the splitting. At \(n=2\) it places \(2s\) below \(2p\) by \(mc^2\alpha^4/12\), which is 29 GHz. Observation disagrees: \(2s_{1/2}\) and \(2p_{1/2}\) nearly coincide, and the 10.9 GHz splitting separates \(2p_{1/2}\) from \(2p_{3/2}\). Schrödinger met this mismatch before publishing his nonrelativistic equation. The Dirac equation, which includes spin, replaces \(l+\frac12\) by \(j+\frac12\) and reproduces the observed pattern, apart from the 1.06 GHz Lamb shift.

For spin-zero particles the Klein--Gordon equation is correct. In pionic atoms a negative pion replaces the electron. X-ray energies between orbits of high \(l\), where the strong interaction is negligible, are calculated with this equation and small QED corrections. Such measurements determine the charged-pion mass. The exact formula also marks a limit. For \(s\) states the square root becomes imaginary when \(Z\alpha>\frac12\), that is for \(Z\geq69\). A point nucleus is then unphysical, and the finite nuclear size must be included.

\subsection*{References For The Physical Derivation}
\begin{itemize}
\item \href{https://arxiv.org/abs/quant-ph/9905076}{A. Calogeracos and N. Dombey, History and Physics of the Klein Paradox; single-particle relativistic wave equations and pair creation in strong potentials.}
\item \href{https://arxiv.org/abs/1605.03300}{M. Trassinelli, D. F. Anagnostopoulos, G. Borchert and colleagues, Measurement of the charged pion mass using X-ray spectroscopy of exotic atoms; Klein-Gordon energies of pionic atoms.}
\end{itemize}





\clearpage

\section{Schrödinger picture}

\label{topic:schr-dinger-picture}

Two atomic levels are swept through an avoided crossing by a ramped electric field. If the sweep is slow, the atom ends in the other diabatic level; if it is fast, the atom stays where it started. A prediction for a Hamiltonian that changes in time is most directly written for a time-dependent state, which is the Schrödinger picture.

In this picture observables, such as a level projector, stay fixed, and the state carries all the time dependence. For a Hamiltonian \(H(t)\), the evolution operator obeys \(i\hbar\,\partial_tU(t,t_0)=H(t)U(t,t_0)\). When Hamiltonians at different times do not commute, the solution is a time-ordered exponential:

\begin{centeredalign}
|\psi(t)\rangle=U(t,t_0)|\psi(t_0)\rangle,\qquad
U(t,t_0)=\mathcal T\exp\Big[-\frac{i}{\hbar}\int_{t_0}^{t}H(s)\,ds\Big].
\end{centeredalign}

The ordering symbol \(\mathcal T\) places operators at later times to the left. Expanding it gives the Dyson series, the starting point of time-dependent perturbation theory. Without the ordering, the exponential gives wrong predictions whenever \([H(s),H(s')]\neq0\). A measurement at an intermediate time is applied to the state at that time, so a sequence of operations is followed in the order in which it occurs.

For the swept atom, let \(|1\rangle\) and \(|2\rangle\) be the diabatic states, whose energies cross linearly in time at rate \(\alpha\), and let \(\Delta\) be the minimum gap set by their coupling. Then

\begin{centeredalign}
H(t)=\frac12\big(\alpha t\,\sigma_z+\Delta\,\sigma_x\big),\qquad
P_{1\to1}=\exp\Big(-\frac{\pi\Delta^2}{2\hbar\alpha}\Big).
\end{centeredalign}

Here \(\sigma_z\) and \(\sigma_x\) are Pauli matrices in the diabatic basis, and \(P_{1\to1}\) is the probability of remaining in \(|1\rangle\) after the crossing. Landau, Zener, Stückelberg and Majorana obtained this result in 1932 by solving the time-dependent equation exactly. For a sweep of duration \(t\) centred on the crossing, the naive exponential of the integrated Hamiltonian would give a rotation about \(x\) by the angle \(\Delta t/\hbar\). That prediction oscillates with the duration and does not depend on the sweep rate.

For a slow sweep, \(\hbar\alpha\ll\Delta^2\), the state follows the instantaneous eigenstate, whose character changes from \(|1\rangle\) to \(|2\rangle\); this adiabatic passage transfers the population completely. For a fast sweep the state cannot follow and remains in \(|1\rangle\). In frequency units the exponent is \(\pi^2f_\Delta^2/\dot\nu\), where \(f_\Delta=\Delta/h\) and \(\dot\nu\) is the sweep rate. A gap of 10 MHz crossed at \(1\,\mathrm{GHz}/\mu\mathrm s\) leaves 37 percent in the initial diabatic state.

Rubbmark, Kash, Littman and Kleppner measured this dependence with Rydberg atoms in 1981, and superconducting qubits reproduce it with swept flux. To lowest order in \(\Delta\) the transition probability is \(\pi\Delta^2/2\hbar\alpha\), which is what time-dependent perturbation theory gives. In the slow limit the diabatic probability vanishes faster than any power of \(\alpha\), so no finite order of an expansion in the sweep rate captures it.

The formula assumes an isolated pair of levels and a sweep extending far beyond the crossing time, the larger of \(\Delta/\alpha\) and \(\sqrt{\hbar/\alpha}\). Finite sweeps add small oscillations, and dephasing during the passage reduces the adiabatic transfer. The Heisenberg picture yields the same probabilities, but its equations contain \(U^\dagger H(t)U\) and therefore require the same time-ordered evolution. The position-space form of the same evolution law is the Schrödinger equation.

\subsection*{References For The Physical Derivation}
\begin{itemize}
\item \href{https://arxiv.org/abs/0911.1917}{S. N. Shevchenko, S. Ashhab and F. Nori, Landau-Zener-Stuckelberg interferometry; the Landau-Zener transition probability, adiabatic and fast limits, and realizations with qubits.}
\end{itemize}





\clearpage

\section{Creation and annihilation operators}

\label{topic:creation-and-annihilation-operators}

A short density pulse made at the centre of an elongated Bose--Einstein condensate travels along it at a few millimetres per second. That speed, and the change from sound to free-particle motion at short wavelengths, follow from the algebra of the operators that add or remove one atom from a mode.

Let \(i\) label orthonormal single-particle modes, such as plane waves in a box or the vibrational modes of a trap. A state with \(n_i\) quanta in mode \(i\) is written \(|\dots,n_i,\dots\rangle\). For bosons the creation operator \(a_i^\dagger\) and the annihilation operator \(a_i\) act as follows, while fermion operators \(c_i\) obey anticommutators:

\begin{centeredalign}
a_i^\dagger|\dots,n_i,\dots\rangle=\sqrt{n_i+1}\,|\dots,n_i+1,\dots\rangle,\qquad
a_i|\dots,n_i,\dots\rangle=\sqrt{n_i}\,|\dots,n_i-1,\dots\rangle,\\
N_i=a_i^\dagger a_i,\qquad [a_i,a_j^\dagger]=\delta_{ij},\qquad
\{c_i,c_j^\dagger\}=\delta_{ij},\qquad \{c_i,c_j\}=0.
\end{centeredalign}

The square-root factors make \(N_i\) count the quanta in mode \(i\), and they follow from the commutator alone. Setting \(i=j\) in \(\{c_i,c_j\}=0\) gives \((c_i^\dagger)^2=0\), so no mode can hold two fermions. This is the Pauli exclusion principle, obtained without writing an antisymmetric wave function.

One-body observables are sums of single products: if \(o_{ij}\) is the matrix element of a single-particle operator between modes, the many-particle operator is \(O=\sum_{ij}o_{ij}a_i^\dagger a_j\). A two-body interaction needs two creation and two annihilation operators. Consider bosons of mass \(m\) in a box of volume \(V\), with plane-wave modes \(\mathbf k\) and a contact interaction of strength \(g=4\pi\hbar^2a_s/m\) fixed by the s-wave scattering length \(a_s\):

\begin{centeredalign}
H=\sum_{\mathbf k}\epsilon_k\,a_{\mathbf k}^\dagger a_{\mathbf k}
+\frac{g}{2V}\sum_{\mathbf k,\mathbf k',\mathbf q}
a_{\mathbf k+\mathbf q}^\dagger a_{\mathbf k'-\mathbf q}^\dagger a_{\mathbf k'}a_{\mathbf k},
\qquad \epsilon_k=\frac{\hbar^2k^2}{2m}.
\end{centeredalign}

Each interaction term removes atoms from modes \(\mathbf k\) and \(\mathbf k'\) and returns them with momentum \(\hbar\mathbf q\) exchanged, so total momentum is conserved. In the condensate nearly all \(N\) atoms occupy \(\mathbf k=0\), where \(a_0\) and \(a_0^\dagger\) multiply by about \(\sqrt N\). Bogoliubov replaced both by \(\sqrt N\) and kept the terms quadratic in the other operators. With density \(n=N/V\),

\begin{centeredalign}
H\approx\frac{gnN}{2}+\sum_{\mathbf k\neq0}\Big[(\epsilon_k+gn)\,a_{\mathbf k}^\dagger a_{\mathbf k}
+\frac{gn}{2}\big(a_{\mathbf k}^\dagger a_{-\mathbf k}^\dagger+a_{\mathbf k}a_{-\mathbf k}\big)\Big].
\end{centeredalign}

The pair terms move two atoms with opposite momenta out of the condensate or back into it. They are removed by mixing \(a_{\mathbf k}\) with \(a_{-\mathbf k}^\dagger\):

\begin{centeredalign}
b_{\mathbf k}=u_ka_{\mathbf k}+v_ka_{-\mathbf k}^\dagger,\qquad
u_k^2=\frac12\Big(\frac{\epsilon_k+gn}{E_k}+1\Big),\qquad v_k^2=u_k^2-1,\\
H=E_0+\sum_{\mathbf k\neq0}E_k\,b_{\mathbf k}^\dagger b_{\mathbf k},\qquad
E_k=\sqrt{\epsilon_k(\epsilon_k+2gn)}.
\end{centeredalign}

The condition \(u_k^2-v_k^2=1\) keeps \([b_{\mathbf k},b_{\mathbf k}^\dagger]=1\), so \(b_{\mathbf k}\) removes one quasiparticle with bosonic statistics. At long wavelength \(E_k\approx\hbar ck\), a sound wave with speed \(c=\sqrt{gn/m}\). At short wavelength \(E_k\approx\epsilon_k+gn\), a free atom raised by the mean interaction energy. The crossover lies at the healing length \(\xi=\hbar/\sqrt{2mgn}\), where \(\epsilon_k=gn\) for \(k=1/\xi\).

For sodium, \(a_s=2.75\) nm. At a density of \(10^{14}\,\mathrm{cm}^{-3}\) this gives \(c=5.1\) mm/s and \(\xi=0.38\,\mu\mathrm m\). Andrews and colleagues excited density pulses in elongated sodium condensates in 1997 and imaged them as they moved. The measured speeds were several millimetres per second and grew with density as the theory predicts. In a cigar-shaped cloud the pulse samples the whole cross-section. The speed is then set by the density averaged over that section, half the central value for a Thomas--Fermi profile, which lowers \(c\) by a factor \(\sqrt2\).

The treatment requires a dilute gas, \(na_s^3\ll1\). Even at zero temperature the pair terms keep a fraction \((8/3)\sqrt{na_s^3/\pi}\) of the atoms outside the condensate, 0.2 percent for these numbers. Because \(E_k/\hbar k\) never falls below \(c\), an impurity moving slower than \(c\) cannot create excitations, which is Landau's criterion for superfluidity. A single mode reproduces the ladder of the quantum harmonic oscillator, and the same pair terms describe squeezed light in quantum optics.

\subsection*{References For The Physical Derivation}
\begin{itemize}
\item \href{https://arxiv.org/abs/cond-mat/9806038}{F. Dalfovo, S. Giorgini, L. P. Pitaevskii and S. Stringari, Theory of Bose-Einstein condensation in trapped gases; Bogoliubov theory, the excitation spectrum, sound, healing length and depletion.}
\item \href{https://arxiv.org/abs/cond-mat/9708237}{E. Zaremba, Sound propagation in a cylindrical Bose-condensed gas; pulses in elongated condensates travel at a speed set by the density averaged over the cross-section.}
\end{itemize}





\part{Observables, Compatibility, And Probability}

\chapter{Observables and spectra}

An observable is a measurable quantity whose operator fixes the possible outcomes; the state then fixes their probabilities.

Spectral theory, self-adjointness and operator algebras make this correspondence precise.

\section*{Topics}

\begin{longtable}{p{0.92\linewidth}}
\toprule
Page \\
\midrule
\endhead
\phantomsection\label{page:spectral-theory}Spectral theory \\
\phantomsection\label{page:pauli-matrices}Pauli matrices \\
\phantomsection\label{page:angular-momentum-operator}Angular momentum operator \\
\phantomsection\label{page:observable}Observable \\
\phantomsection\label{page:self-adjoint-operator}Self-adjoint operator \\
\phantomsection\label{page:operator-theory}Operator theory \\
\phantomsection\label{page:operator-physics}Operator (physics) \emph{(see Operator Theory)} \\
\phantomsection\label{page:eigenvalues-and-eigenvectors}Eigenvalues and eigenvectors \\
\phantomsection\label{page:heisenberg-group}Heisenberg group \\
\phantomsection\label{page:von-neumann-algebra}Von Neumann algebra \\
\phantomsection\label{page:quantum-logic}Quantum logic \\
\bottomrule
\end{longtable}

\clearpage

\section{Spectral theory}

\label{topic:spectral-theory}

An observable specifies both possible outcomes and how a prepared state distributes probability among them. For a self-adjoint operator \(A\), the spectral measure \(E_A\) assigns an orthogonal projector to each measurable set of real outcomes. The operator and its measurement probabilities are reconstructed from this measure:

\begin{centeredalign}
A=\int_{\mathbb R}\lambda\,dE_A(\lambda),\qquad
p(\lambda\in B)=\operatorname{Tr}[\rho E_A(B)].
\end{centeredalign}

Here \(B\) is an interval or other measurable outcome set and \(\rho\) the density operator. For a discrete spectrum the integral becomes a sum over eigenspace projectors. For position or free-particle momentum it is genuinely continuous; their ideal eigenvectors are generalized distributions rather than normalizable prepared states. A detector with finite resolution measures an interval probability, not a normalizable state at one exact continuum value.

The same spectral measure determines functions of the operator. If \(H\) is the time-independent Hamiltonian, the function \(\exp(-iEt/\hbar)\) gives its unitary evolution. An energy decomposition therefore connects a spectroscopic question to dynamics:

\begin{centeredalign}
U(t)=\int e^{-iEt/\hbar}\,dE_H(E),\qquad
\langle\psi|U(t)|\psi\rangle=\int e^{-iEt/\hbar}\,d\mu_\psi(E),
\end{centeredalign}

where \(\mu_\psi(B)=\langle\psi|E_H(B)|\psi\rangle\) is the energy distribution of the preparation. The survival amplitude is its Fourier transform. A narrow energy distribution changes phase slowly relative to itself; a broad distribution can dephase rapidly. This connects spectral width to temporal evolution without assuming irreversible decay.

The spectrum belongs to the operator with its domain. On an interval of length \(L\), the same differential expression \(-\hbar^2\partial_x^2/(2m)\) has a zero-energy constant mode with Neumann conditions, while Dirichlet conditions exclude it and begin at \(\pi^2\hbar^2/(2mL^2)\). A change of boundary therefore changes both the energy outcomes and the time evolution, even though the bulk equation is unchanged.

Two Hamiltonians with identical energy values can still have different observable matrix elements. Transporting a mechanism requires a correspondence between states and observables, not just a matched list of eigenvalues. Spectral theory supplies the predictions once that correspondence and the operator domains are specified; it also identifies precisely what an isospectral comparison leaves undecided.

\subsection*{References For The Physical Derivation}
\begin{itemize}
\item \href{https://arxiv.org/abs/1508.06951}{V. Moretti, Mathematical Foundations of Quantum Mechanics: An Advanced Short Course; states, operators and symmetry.}
\end{itemize}

\subsection*{Relations In The Original Papers}

\href{https://arxiv.org/html/1508.06951#S2.EGx25}{arXiv:1508.06951, S2.EGx25}. Functions of a self-adjoint observable are defined by integration against its spectral measure. The domain of an unbounded function of the operator must be retained.



\clearpage

\section{Pauli matrices}

\label{topic:pauli-matrices}

A two-level system has three independent traceless Hermitian observables. Choosing an orthonormal basis of states identifies them with the Pauli matrices. Their multiplication law determines rotations, transition amplitudes and the closure of equations for spin expectations:

\begin{centeredalign}
\sigma_x=\begin{pmatrix}0&1\\1&0\end{pmatrix},\quad
\sigma_y=\begin{pmatrix}0&-i\\i&0\end{pmatrix},\quad
\sigma_z=\begin{pmatrix}1&0\\0&-1\end{pmatrix},
\end{centeredalign}

\begin{centeredalign}
\sigma_i\sigma_j=\delta_{ij}I+i\sum_k\epsilon_{ijk}\sigma_k.
\end{centeredalign}

For a unit vector \(\mathbf n\), this identity gives \((\mathbf n\cdot\boldsymbol\sigma)^2=I\). Exponentiating a Hamiltonian proportional to that operator therefore separates into a cosine and a sine. The rotation of a spinor through angle \(\theta\) is

\begin{centeredalign}
U(\theta)=e^{-i\theta\mathbf n\cdot\boldsymbol\sigma/2}
=\cos(\theta/2)I-i\sin(\theta/2)\mathbf n\cdot\boldsymbol\sigma.
\end{centeredalign}

A full \(2\pi\) rotation gives \(-I\). This phase leaves an isolated spin probability unchanged, but becomes observable relative to a coherent reference branch that was not rotated. The Bloch vector itself completes one ordinary spatial rotation. The distinction between spinor amplitude and observable vector is the reason the half-angle appears.

Writing a two-level Hamiltonian as \(H=E_0I+(\hbar/2)\boldsymbol\omega\cdot\boldsymbol\sigma\) and a state as \(\rho=(I+\mathbf r\cdot\boldsymbol\sigma)/2\) gives \(\dot{\mathbf r}=\boldsymbol\omega\times\mathbf r\). The identity term contributes only an overall phase, while the traceless part produces precession. The vector's length is conserved by this isolated unitary dynamics; relaxation requires an additional physical coupling.

For multiple spins, tensor products of Pauli matrices form an operator basis. Commuting a Hamiltonian with a one-spin observable can produce a two-spin product, showing how a local signal becomes coupled to a correlation. Counting the independent operators generated by those commutators measures the size of that particular observable dynamics. A reduction in this size can result from symmetry, selection rules or equal transition frequencies; the Pauli algebra permits these possibilities to be tested separately.

\subsection*{References For The Physical Derivation}
\begin{itemize}
\item \href{https://arxiv.org/abs/1508.06951}{V. Moretti, Mathematical Foundations of Quantum Mechanics: An Advanced Short Course; states, operators and symmetry.}
\end{itemize}

\subsection*{Relations In The Original Papers}

\href{https://arxiv.org/html/quant-ph/9705052#Ch1.E3}{arXiv:quant-ph/9705052, Ch1.E3}. The Pauli matrices represent the three spin components used to calculate rotations and measurement probabilities. One spin-1/2 degree of freedom; these finite matrices do not realize a canonical position-momentum pair.



\clearpage

\section{Angular momentum operator}

\label{topic:angular-momentum-operator}

Rotations about different axes change orientation in an order-dependent way. In quantum mechanics their generators \(J_x,J_y,J_z\) obey an algebra whose representations determine the allowed angular-momentum values. With \(\hbar\) the reduced Planck constant,

\begin{centeredalign}
[J_i,J_j]=i\hbar\sum_k\epsilon_{ijk}J_k,\qquad
[J^2,J_i]=0,\qquad J^2=J_x^2+J_y^2+J_z^2.
\end{centeredalign}

A simultaneous eigenstate of \(J^2\) and \(J_z\) can be labeled \(|j,m\rangle\). Define \(J_\pm=J_x\pm iJ_y\). Their commutators with \(J_z\) show that they raise or lower its eigenvalue by \(\hbar\), while preserving \(j\). Positivity of the squared norms fixes the endpoints of the ladder:

\begin{centeredalign}
J^2|j,m\rangle=\hbar^2j(j+1)|j,m\rangle,\qquad
J_z|j,m\rangle=\hbar m|j,m\rangle,
\end{centeredalign}

\begin{centeredalign}
J_\pm|j,m\rangle=\hbar\sqrt{j(j+1)-m(m\pm1)}\,|j,m\pm1\rangle.
\end{centeredalign}

The ladder terminates at \(m=\pm j\), giving \(2j+1\) states. Single-valued orbital wave functions on ordinary three-dimensional space carry integer orbital angular momentum. Intrinsic spin uses representations of the covering group SU(2), allowing half-integer \(j\). The same local commutator algebra is therefore realized on physical systems with different global transformation properties.

For two spins, total angular momentum is \(\mathbf J=\mathbf J_1+\mathbf J_2\). An isotropic exchange interaction can be rewritten using the total Casimir:

\begin{centeredalign}
H=\frac{K}{\hbar^2}\mathbf J_1\cdot\mathbf J_2
=\frac{K}{2\hbar^2}(J^2-J_1^2-J_2^2).
\end{centeredalign}

Here \(K\) is an energy. Two spin-one-half particles then have a singlet energy \(-3K/4\) and a triplet energy \(K/4\). The gap is fixed by the coupling and the representation, without diagonalizing a general four-dimensional matrix. A magnetic field separates magnetic sublevels; an anisotropic interaction can also mix sectors that the isotropic model kept separate.

This calculation connects symmetry to a measurable excitation energy. Transferring the exchange formula to another object requires identifying which angular-momentum representation its degrees of freedom carry and whether anisotropy, orbital coupling or the environment changes the generator. A common algebra supplies the ladder relations, while the Hamiltonian selects which of those states a physical system occupies.

\subsection*{References For The Physical Derivation}
\begin{itemize}
\item \href{https://arxiv.org/abs/1508.06951}{V. Moretti, Mathematical Foundations of Quantum Mechanics: An Advanced Short Course; states, operators and symmetry.}
\end{itemize}





\clearpage

\section{Observable}

\label{topic:observable}

A spin-1 atom passing a Stern--Gerlach magnet leaves three spots, one for each value of \(S_z\). An apparatus that recombines the two outer beams without recording which one each atom took asks a coarser question. Comparing the two apparatuses shows what an observable specifies.

An observable is a list of outcomes together with orthogonal projectors that sum to the identity. The operator collects them in its spectral decomposition, and the projectors give the probabilities:

\begin{centeredalign}
A=\sum_ia_iP_i,\qquad P_iP_j=\delta_{ij}P_i,\qquad \sum_iP_i=I,\qquad
p(a_i)=\operatorname{Tr}(\rho P_i).
\end{centeredalign}

The eigenvalues \(a_i\) label detector outputs; the projectors carry the physics. The mean \(\operatorname{Tr}(\rho A)\) and the moments \(\operatorname{Tr}(\rho A^k)\) follow from the same probabilities. A function \(f(A)=\sum_if(a_i)P_i\) relabels the outputs and merges those that receive the same label. After a recorded outcome with minimal disturbance, the state becomes \(P_i\rho P_i\) divided by its trace.

For spin 1, \(S_z=\hbar(P_{+1}-P_{-1})\), while merging the outer beams measures \(S_z^2=\hbar^2(P_{+1}+P_{-1})\). Prepare \(|\psi\rangle=(|{+}1\rangle+|{-}1\rangle)/\sqrt2\). A measurement of \(S_z\) gives \(\pm\hbar\) with probability one half each and leaves \(|{+}1\rangle\) or \(|{-}1\rangle\). A measurement of \(S_z^2\) gives \(\hbar^2\) with certainty and leaves \(|\psi\rangle\) unchanged, including the relative phase of its components. A later experiment sensitive to that phase distinguishes the two apparatuses, even though one of them gave a certain result.

For spin 1 the three squared components commute with each other. They obey a sum rule:

\begin{centeredalign}
[S_x^2,S_y^2]=[S_y^2,S_z^2]=[S_z^2,S_x^2]=0,\qquad
S_x^2+S_y^2+S_z^2=2\hbar^2.
\end{centeredalign}

One apparatus can therefore measure all three. Each has eigenvalues \(0\) and \(\hbar^2\), and the sum rule forces exactly one zero in every run. The state \(|\psi\rangle\) above gives \(S_x^2=\hbar^2\), \(S_y^2=0\) and \(S_z^2=\hbar^2\) with certainty. By contrast, \(S_x\) and \(S_z\) do not commute, and no apparatus measures both sharply. Commuting triples of this kind underlie the Kochen--Specker theorem discussed under quantum logic.

A real detector can assign an outcome to the wrong beam. Suppose a spin-1/2 analyzer reports the wrong spot with probability \(\varepsilon\), for instance because the spots overlap. The report up then corresponds to the positive operator \(E_+=(1-\varepsilon)P_\uparrow+\varepsilon P_\downarrow\), and

\begin{centeredalign}
E_\pm=\frac{I\pm\eta\,\sigma_z}{2},\qquad \eta=1-2\varepsilon,\qquad
p_\pm=\frac{1\pm\eta\,r_z}{2},
\end{centeredalign}

where \(r_z\) is the \(z\) component of the Bloch vector. These operators are not projectors, and together they form a positive-operator-valued measure. The mean outcome still determines \(r_z\) after division by \(\eta\), but no single run gives certainty when \(\varepsilon>0\). For continuous spectra the projectors are assigned to intervals of outcomes, as described in spectral theory. Two devices measure the same observable when they assign the same operators to their outcomes, whatever their construction.

\subsection*{References For The Physical Derivation}
\begin{itemize}
\item \href{https://arxiv.org/abs/1508.06951}{V. Moretti, Mathematical Foundations of Quantum Mechanics: An Advanced Short Course; observables, projection-valued and positive-operator-valued measures.}
\end{itemize}





\clearpage

\section{Self-adjoint operator}

\label{topic:self-adjoint-operator}

A charged particle circulates in a thin ring threaded by magnetic flux. Its momentum values shift with the flux, although the field vanishes where the particle moves. The shift enters through a boundary condition. That boundary condition is exactly what makes the momentum operator self-adjoint.

Let \(P=-i\hbar\,d/dx\) act on wave functions on an interval \([0,L]\). Integration by parts leaves a boundary term:

\begin{centeredalign}
\langle\phi|P\psi\rangle-\langle P\phi|\psi\rangle
=-i\hbar\bigl[\phi^*(L)\psi(L)-\phi^*(0)\psi(0)\bigr].
\end{centeredalign}

On functions that vanish at both ends the term is zero, so \(P\) is symmetric and its expectation values are real. It is not self-adjoint, however. Its adjoint acts on all differentiable functions without any boundary condition, and on that larger domain every \(e^{ipx/\hbar}\) with complex \(p\) is an eigenfunction. The spectrum of such an operator is not a set of real outcomes.

The deficiency indices measure the mismatch. Count the normalizable solutions of \(P^\dagger\psi=\pm i\kappa\psi\) for a fixed \(\kappa>0\); they are \(e^{\mp\kappa x/\hbar}\). A symmetric operator has self-adjoint extensions only when the two counts \(n_\pm\) are equal. The extensions are then labelled by unitary maps between the two solution spaces:

\begin{centeredalign}
[0,L]:\ (n_+,n_-)=(1,1),\qquad \psi(L)=e^{i\theta}\psi(0),\qquad
p_k=\frac{\hbar(2\pi k+\theta)}{L},\\
[0,\infty):\ (1,0),\ \text{no extension};\qquad
\mathbb R:\ (0,0),\ \text{unique}.
\end{centeredalign}

On the interval each phase \(\theta\) defines a different self-adjoint momentum with its own real spectrum, \(k\) an integer. On a ring threaded by flux \(\Phi\), a gauge transformation moves the vector potential into this boundary condition, with \(\theta=2\pi\Phi/\Phi_0\) and \(\Phi_0=h/e=4.14\times10^{-15}\) weber. The energies \(p_k^2/2m\) are periodic in the flux with period \(\Phi_0\); this is the Aharonov--Bohm effect. Persistent currents in mesoscopic metal rings show this flux period. In a single ring, disorder, temperature and the many electrons of a metal reduce the current amplitude, while the period \(h/e\) remains.

On the half line the counts differ, and no self-adjoint momentum exists. A translation toward the wall would push probability through the wall, so no unitary translation group exists there. Stone's theorem then provides no self-adjoint generator. The momentum of a particle confined to \(x>0\) can be measured only as a generalized observable.

The kinetic energy on the half line behaves differently, with indices \((1,1)\). Its extensions impose \(\psi'(0)=-\psi(0)/b\) for real \(b\), with the Dirichlet condition as the limit \(b\to0\). For \(b>0\) the function \(e^{-x/b}\) is a bound state with energy \(-\hbar^2/(2mb^2)\), describing a short-range attraction at the wall. On a finite box the kinetic energy has indices \((2,2)\), a family of conditions that includes the Dirichlet, Neumann and periodic cases compared in spectral theory.

Self-adjointness is required for two physical reasons. Only self-adjoint operators have spectral decompositions with real outcomes, and only they generate unitary evolution. Symmetry alone, real expectation values on a restricted domain, provides neither. The choice of domain is part of the physical description of the system, as in the treatments of observables and of the mathematical formulation.

\subsection*{References For The Physical Derivation}
\begin{itemize}
\item \href{https://arxiv.org/abs/quant-ph/0103153}{G. Bonneau, J. Faraut and G. Valent, Self-adjoint extensions of operators and the teaching of quantum mechanics; momentum and Hamiltonian extensions on an interval and on the half line.}
\item \href{https://arxiv.org/abs/0906.4780}{A. C. Bleszynski-Jayich and colleagues, Persistent currents in normal metal rings: comparing high-precision experiment with theory; flux-periodic currents in single rings and arrays.}
\end{itemize}





\clearpage

\section{Operator theory}

\label{topic:operator-theory}

The energy of an electron in an atom can take arbitrarily large values, while a spin-1/2 component never exceeds \(\hbar/2\) in magnitude. Operator theory distinguishes these cases as unbounded and bounded operators. The distinction decides which mathematical statements hold, and it explains why a hydrogen atom does not collapse.

The norm \(\|A\|\) is the largest factor by which \(A\) can stretch a unit vector. For an observable it is the largest magnitude of an outcome. A bounded operator acts on every state. An unbounded symmetric operator, such as the kinetic energy, cannot act on every state. By the Hellinger--Toeplitz theorem, a symmetric operator defined everywhere is bounded.

Position and momentum cannot both be bounded. Suppose \([A,B]=cI\) with \(c\neq0\). Induction gives \([A,B^n]=ncB^{n-1}\). Taking norms of both sides yields

\begin{centeredalign}
n|c|\,\|B^{n-1}\|\leq2\|A\|\,\|B\|\,\|B^{n-1}\|.
\end{centeredalign}

No power of \(B\) vanishes, because \(B^{n-1}=0\) would force \(B^{n-2}=0\) and so on down to \(I=0\). Dividing by \(\|B^{n-1}\|\) gives \(n\leq2\|A\|\|B\|/|c|\) for every \(n\), which is impossible. This argument of Wielandt holds in infinite dimensions; the finite-dimensional trace argument appears under the canonical commutation relation. Canonical pairs therefore require unbounded operators, and with them stated domains.

In hydrogen, \(H=p^2/2m-k/r\) with \(k=e^2/(4\pi\varepsilon_0)\). The Coulomb term is unbounded below, and a classical electron would radiate and spiral into the nucleus. The kinetic energy controls the singularity through Hardy's inequality, valid for wave functions in three dimensions:

\begin{centeredalign}
\int|\nabla\psi|^2d^3x\geq\frac14\int\frac{|\psi|^2}{r^2}\,d^3x
\quad\Longrightarrow\quad
H\geq\frac{\hbar^2}{8mr^2}-\frac{k}{r}\geq-\frac{2mk^2}{\hbar^2}.
\end{centeredalign}

The last step minimizes the function of \(r\) at \(r=\hbar^2/(4mk)\). The bound equals four rydbergs below zero, \(-54.4\) eV, compared with the exact ground energy of \(-13.6\) eV. It is not sharp, but it proves that the energy is bounded below. Compressing the wave function raises its kinetic energy faster than the Coulomb attraction lowers its potential energy.

The same comparison shows that the Coulomb term is small relative to the kinetic energy. The Kato--Rellich theorem then makes \(H\) self-adjoint on the domain of the kinetic energy, as Kato proved for atoms in 1951. The evolution \(e^{-iHt/\hbar}\) is therefore uniquely defined.

The argument fails for stronger singularities. For \(V=-g/r^2\), Hardy's inequality gives \(H\geq0\) when \(g\leq\hbar^2/(8m)\). For larger \(g\) the energy is unbounded below, and the Hamiltonian needs an extra condition at the origin to define the dynamics. A point electric dipole produces an inverse-square potential with angular dependence. It can bind an electron only above a critical moment of about 1.63 debye. The relativistic kinetic energy \(\sqrt{p^2c^2+m^2c^4}\) grows only linearly with momentum. A Coulomb potential of nuclear charge \(Ze\) then keeps the energy bounded below only for \(Z\alpha\leq2/\pi\). Self-adjointness and spectral theory build on these estimates.

\subsection*{References For The Physical Derivation}
\begin{itemize}
\item \href{https://arxiv.org/abs/quant-ph/9907069}{F. Gieres, Mathematical surprises and Dirac's formalism in quantum mechanics; unbounded operators, domains and the canonical commutation relation.}
\item \href{https://arxiv.org/abs/math-ph/0209034}{E. H. Lieb, The Stability of Matter and Quantum Electrodynamics; stability of the hydrogen atom and of matter, with and without relativistic kinetic energy.}
\end{itemize}





\clearpage
\section{Operator (physics)}
\label{topic:operator-physics}
Operator (physics) is an alternative name for Operator Theory.
The physical derivation is given in \hyperref[topic:operator-theory]{Operator Theory}.

\clearpage

\section{Eigenvalues and eigenvectors}

\label{topic:eigenvalues-and-eigenvectors}

Sodium vapour illuminated by light with two optical frequencies fluoresces brightly, except where the difference of the two frequencies matches the splitting of two ground levels. There the fluorescence drops sharply. Alzetta and coworkers observed such dark lines in 1976. They are explained by an eigenvector of the atom-light Hamiltonian that has no excited-state component.

An eigenvector of the Hamiltonian, \(H|n\rangle=E_n|n\rangle\), changes in time only by a phase, so all of its probabilities stay constant. A general state is a sum of eigenvectors, each with its own phase. Expectation values therefore oscillate at the Bohr frequencies \((E_n-E_m)/\hbar\):

\begin{centeredalign}
|\psi(t)\rangle=\sum_nc_ne^{-iE_nt/\hbar}|n\rangle,\qquad
\langle A\rangle(t)=\sum_{m,n}c_m^*c_nA_{mn}\,e^{-i(E_n-E_m)t/\hbar}.
\end{centeredalign}

Spectral lines appear at differences of eigenvalues for this reason. For a Hermitian \(H\), evaluate \(\langle m|H|n\rangle\) by acting to the right and to the left. The two results give \((E_n-E_m^*)\langle m|n\rangle=0\). Setting \(m=n\) shows that eigenvalues are real, and distinct eigenvalues have orthogonal eigenvectors. In finite dimensions the eigenvalues are the roots of \(\det(H-EI)=0\).

In the dark-resonance experiment, two ground levels \(|1\rangle\) and \(|2\rangle\) couple to one excited level \(|e\rangle\). Two light fields drive the transitions with Rabi frequencies \(\Omega_p\) and \(\Omega_c\). In a frame rotating with the fields, within the rotating-wave approximation and with both fields on resonance, the Hamiltonian in the basis \((|1\rangle,|2\rangle,|e\rangle)\) is

\begin{centeredalign}
H=\frac\hbar2\begin{pmatrix}0&0&\Omega_p\\0&0&\Omega_c\\ \Omega_p&\Omega_c&0\end{pmatrix},\qquad
\det(H-EI)=-E\left(E^2-\frac{\hbar^2\Omega^2}{4}\right),
\end{centeredalign}

with \(\Omega^2=\Omega_p^2+\Omega_c^2\). The eigenvalues are \(0\) and \(\pm\hbar\Omega/2\). The eigenvector with eigenvalue zero is

\begin{centeredalign}
|D\rangle=\frac{\Omega_c|1\rangle-\Omega_p|2\rangle}{\Omega}.
\end{centeredalign}

Its excited-state amplitude vanishes because the two excitation paths interfere destructively. An atom in \(|D\rangle\) neither absorbs nor fluoresces. Spontaneous decay from \(|e\rangle\) returns atoms to the ground levels at random, and atoms accumulate in \(|D\rangle\). This is coherent population trapping. The state remains an exact eigenvector for any common detuning of the two fields, as long as their frequency difference matches the ground splitting. A mismatch \(\delta\) couples \(|D\rangle\) to the other states in proportion to \(\delta\), and fluorescence returns.

The width of the dark resonance is set by the decay of the coherence between the ground levels and by power broadening. It can be far narrower than the optical linewidth. Compact atomic clocks use this resonance to lock a laser frequency difference to the cesium ground-state splitting of 9.192631770 gigahertz.

Eigenvalues alone would not explain the effect: the spectrum \(0,\pm\hbar\Omega/2\) says nothing about fluorescence. The prediction depends on the composition of the eigenvector. When eigenvalues are degenerate, the eigenvectors are not unique, and a perturbation selects the physically relevant combinations, as in perturbation theory. The two-state quantum system gives the simplest case of the same calculation.

\subsection*{References For The Physical Derivation}
\begin{itemize}
\item \href{https://arxiv.org/abs/0810.2071}{N. Belcher, E. E. Mikhailov and I. Novikova, Atomic Clocks and Coherent Population Trapping: Experiments for Undergraduate Laboratories; dark states and the measurement of a ground-state splitting.}
\end{itemize}





\clearpage

\section{Heisenberg group}

\label{topic:heisenberg-group}

A displacement of a particle's position followed by a momentum kick traces a rectangle in phase space when both operations are undone. Although the final position and momentum shifts cancel, the quantum amplitude acquires a phase. The Heisenberg group includes this phase as part of the composition law.

For canonical operators satisfying \([X,P]=i\hbar I\), define the displacement with position increment \(a\) and momentum increment \(b\) by \(D(a,b)=\exp[i(bX-aP)/\hbar]\). The commutator of the exponents is a scalar, so the Baker--Campbell--Hausdorff series stops after that commutator:

\begin{centeredalign}
D(a,b)D(a',b')=
e^{i(ba'-ab')/(2\hbar)}D(a+a',b+b').
\end{centeredalign}

The antisymmetric combination \(ba'-ab'\) is the oriented symplectic area. Its appearance explains why a description retaining only the two displacement coordinates has a projective composition law, whereas adding the phase gives an ordinary group representation. The central coordinate records how operations compose, rather than another position of the particle.

\begin{centeredalign}
D(a,0)D(0,b)D(-a,0)D(0,-b)=e^{-iab/\hbar}I.
\end{centeredalign}

A common overall phase has no effect on the probability of a single prepared state. To measure the loop phase, the displaced sequence must interfere with a reference sequence, for example through two coherent branches of a controlled operation. The observable is then a relative phase between the branches. A phase generated by a chosen path is compatible with fully specified dynamics.

The construction also relates particle motion to oscillator modes. Appropriate dimensionless quadratures express the same displacement algebra through creation and annihilation operators. Acting on the oscillator vacuum produces a coherent state; displacement changes its mean quadratures while preserving the vacuum covariance. The physical system and its units differ, but the central commutation relation is the same mathematical relation being realized.

The uniqueness of the usual irreducible representation requires finitely many canonical pairs, regularity and a fixed nonzero central parameter. Quantum fields contain infinitely many modes and can admit inequivalent representations. Consequently the finite-dimensional symplectic picture cannot be promoted without qualification to an equivalence of arbitrary field-theory vacua. The group provides a precise reusable mechanism of composition within its stated representation assumptions.

\subsection*{References For The Physical Derivation}
\begin{itemize}
\item \href{https://arxiv.org/abs/1508.06951}{V. Moretti, Mathematical Foundations of Quantum Mechanics: An Advanced Short Course; states, operators and symmetry.}
\end{itemize}

\subsection*{Relations In The Original Papers}

\href{https://arxiv.org/html/1508.06951#S1.E21}{arXiv:1508.06951, S1.E21}. The central commutator supplies the infinitesimal relation used in the displacement composition law. This display establishes the infinitesimal relation, not by itself the global Weyl representation; the chapter states that representation explicitly.



\clearpage

\section{Von Neumann algebra}

\label{topic:von-neumann-algebra}

An experimenter who controls one of two entangled spins, or the field inside one region of space, can apply only some operations. The set of those operations is an algebra closed under products, adjoints and limits. Von Neumann algebras describe such sets. Their structure decides whether the subsystem has a density matrix and how thermal equilibrium can be defined.

The commutant \(M'\) of a set \(M\) of operators contains every operator that commutes with all of \(M\). Von Neumann's bicommutant theorem states that an algebra containing the identity, closed under adjoints and under weak limits, equals its double commutant. For two spins the operations on \(A\) and on \(B\) are each other's commutants:

\begin{centeredalign}
M'=\{B:\ [A,B]=0\ \text{for all}\ A\in M\},\qquad M''=M,\\
\bigl(\mathcal B(\mathcal H_A)\otimes I\bigr)'=I\otimes\mathcal B(\mathcal H_B).
\end{centeredalign}

The centre \(M\cap M'\) contains operators that commute with everything in the algebra. If it holds more than multiples of the identity, the states split into sectors with no coherent superpositions between them. Total electric charge separates sectors in this way.

Algebras with trivial centre are called factors, and they come in types. A subsystem of finitely many degrees of freedom gives type I, where every state is a density matrix \(\rho_A\) with \(\langle A\rangle=\operatorname{Tr}(\rho_AA)\). An infinite chain of Bell pairs gives, on one side, an algebra with a normalized trace but no minimal projector, of type \(\mathrm{II}_1\). The algebra of a region in relativistic quantum field theory has type III, with no trace at all. There the entanglement entropy of a region diverges unless a short-distance cutoff is imposed.

Thermal equilibrium can be stated without a density matrix. For a finite system, the Gibbs state \(\rho=e^{-\beta H}/Z\) with \(\beta=1/k_BT\) satisfies the Kubo--Martin--Schwinger condition. It follows from cyclicity of the trace, with \(e^{-\beta H}\) acting as an evolution in imaginary time:

\begin{centeredalign}
\langle A(t)B\rangle_\beta=\langle B\,A(t+i\hbar\beta)\rangle_\beta,\qquad
A(t)=e^{iHt/\hbar}Ae^{-iHt/\hbar}.
\end{centeredalign}

In an infinite system \(e^{-\beta H}\) has no trace, but the condition still makes sense on the algebra and defines equilibrium. Several solutions at the same temperature describe coexisting phases, such as the two magnetizations of a ferromagnet. Tomita--Takesaki theory shows that every faithful normal state of a von Neumann algebra satisfies this condition for a flow of its own, the modular flow.

The condition has a measured consequence. Fourier transforming it for the correlation function of an observable gives detailed balance, with \(\omega>0\) the energy taken up by the system:

\begin{centeredalign}
S(\omega)=\int dt\,e^{i\omega t}\langle A(t)A(0)\rangle_\beta,\qquad
S(-\omega)=e^{-\beta\hbar\omega}S(\omega).
\end{centeredalign}

In Raman scattering from silicon at 300 kelvin, the optical phonon at 520 inverse centimetres gives an anti-Stokes to Stokes ratio of \(e^{-2.49}\approx0.083\). This value holds after correction for the frequency dependence of scattering. The ratio serves as a contactless thermometer. The Bisognano--Wichmann theorem applies the same condition to the vacuum restricted to a wedge of spacetime, with boosts as the flow. It predicts the Unruh temperature \(\hbar a/(2\pi ck_B)\), about 0.41 kelvin for an acceleration of \(10^{20}\) metres per second squared. The theorem is established; a direct observation of the Unruh effect is still lacking. Quantum statistical mechanics uses the finite-system form of the same condition.

\subsection*{References For The Physical Derivation}
\begin{itemize}
\item \href{https://arxiv.org/abs/1803.04993}{E. Witten, Notes on Some Entanglement Properties of Quantum Field Theory; local algebras, modular theory and the Bisognano--Wichmann theorem.}
\item \href{https://arxiv.org/abs/math-ph/0411058}{J. Yngvason, The Role of Type III Factors in Quantum Field Theory; type I and type III algebras and their physical consequences.}
\end{itemize}





\clearpage

\section{Quantum logic}

\label{topic:quantum-logic}

Two crossed polarizers block all light. A third polarizer inserted between them at 45 degrees lets one eighth of unpolarized light through. The propositions tested by polarizers do not combine like the propositions of classical logic, and quantum logic describes how they do combine.

The proposition that a photon passes a polarizer at angle \(\theta\) corresponds to the projector \(P_\theta\) onto that polarization. Its negation is the orthogonal polarization. The conjunction of two propositions is certain exactly when both are certain; it corresponds to the intersection of their subspaces. The disjunction corresponds to the smallest closed subspace containing both. Birkhoff and von Neumann showed in 1936 that these operations form a lattice:

\begin{centeredalign}
a\wedge b\leftrightarrow \operatorname{ran}P_a\cap\operatorname{ran}P_b,\qquad
a\vee b\leftrightarrow \overline{\operatorname{ran}P_a+\operatorname{ran}P_b},\qquad
\neg a\leftrightarrow I-P_a.
\end{centeredalign}

No polarization passes the \(0^\circ\) and \(45^\circ\) polarizers with certainty, so their conjunction is the zero subspace. Yet light does pass the sequence of three polarizers. A sequence of filters is not a conjunction. Von Neumann showed that the conjunction is the limit of infinitely many alternations of the two projectors. For \(0^\circ\) and \(45^\circ\), each further polarizer of an alternating stack transmits half the light:

\begin{centeredalign}
(P_aP_b)^n\longrightarrow P_{a\wedge b}\quad(n\to\infty),\qquad
I_m=I_1\cos^{2m}45^\circ=2^{-m}I_1.
\end{centeredalign}

The intensity after \(m\) further polarizers tends to zero, the empty conjunction. For commuting projectors the product \(P_aP_b\) is already the conjunction.

The distributive law fails. Let \(a\) be horizontal polarization, \(b\) the diagonal one and \(\neg b\) the antidiagonal one. The disjunction \(b\vee\neg b\) is the whole space, while neither \(a\wedge b\) nor \(a\wedge\neg b\) contains a nonzero vector:

\begin{centeredalign}
a\wedge(b\vee\neg b)=a\neq0=(a\wedge b)\vee(a\wedge\neg b).
\end{centeredalign}

Physically, a diagonal analyzer always gives one of its two outcomes for a horizontal photon, but neither outcome is certain in advance. No polarization makes both horizontal and diagonal certain. A description that assigned every photon a definite answer to both questions would satisfy the distributive law.

A weaker law remains, the orthomodular law: if \(a\leq b\) then \(b=a\vee(b\wedge\neg a)\). Within any set of mutually commuting projectors the lattice is Boolean, so classical logic holds for compatible propositions. The failures appear only across incompatible measurements. The Kochen--Specker theorem sharpens this. In dimension three or more, no assignment of true or false to all projectors makes exactly one member of every complete orthogonal set true. Kirchmair and coworkers tested a state-independent form with trapped ions in 2009, measuring triples of commuting two-qubit observables. The measured correlations exceeded the bound obeyed by any such assignment.

The lattice structure, Gleason's theorem and the Kochen--Specker theorem are established mathematics with tested consequences. The proposal that this lattice should replace classical logic in reasoning, argued by Putnam in 1968, is not generally accepted. Reasoning about experiments uses ordinary logic, while the lattice describes the propositions that experiments test. Projectors also generate the operator algebras discussed under von Neumann algebras.

\subsection*{References For The Physical Derivation}
\begin{itemize}
\item \href{https://arxiv.org/abs/quant-ph/0101028}{M. L. Dalla Chiara and R. Giuntini, Quantum Logic; lattices of sharp and unsharp quantum propositions.}
\item \href{https://arxiv.org/abs/0904.1655}{G. Kirchmair and colleagues, State-independent experimental test of quantum contextuality; trapped-ion test of noncontextual value assignments.}
\end{itemize}





\chapter{Noncommuting observables}

Noncommuting observables generally require different measurement bases. Their commutator determines how their dispersions are constrained in a given preparation.

The commutator and the uncertainty relations measure this incompatibility, which is a property of the observable algebra and state. Noncommuting operators may share particular eigenvectors; their noncommutation excludes a common complete eigenbasis, not every state with two sharp values.

\section*{Topics}

\begin{longtable}{p{0.92\linewidth}}
\toprule
Page \\
\midrule
\endhead
\phantomsection\label{page:commutator}Commutator \\
\phantomsection\label{page:uncertainty-principle}Uncertainty principle \\
\phantomsection\label{page:canonical-commutation-relation}Canonical commutation relation \\
\bottomrule
\end{longtable}

\clearpage

\section{Commutator}

\label{topic:commutator}

Rotating a spin by a small angle about \(x\), then about \(y\), and then reversing both rotations in the same order leaves it rotated about \(z\). The net rotation is of second order in the angle. Its size and axis are fixed by the commutator of the two generators, not by either rotation alone.

The commutator \([A,B]=AB-BA\) compares the two orders of applying two operators. For generators \(X\) and \(Y\) of transformations, the Baker--Campbell--Hausdorff formula gives \(e^Xe^Y=\exp(X+Y+\tfrac12[X,Y]+\cdots)\). The four-step loop therefore has the commutator as its leading term. For spin rotations \(R_{\mathbf n}(\varepsilon)=e^{-i\varepsilon\,\mathbf n\cdot\boldsymbol\sigma/2}\), with \([\sigma_x,\sigma_y]=2i\sigma_z\), this gives

\begin{centeredalign}
e^Xe^Ye^{-X}e^{-Y}=e^{[X,Y]+O(3)},\\
R_x(\varepsilon)R_y(\varepsilon)R_x(-\varepsilon)R_y(-\varepsilon)=R_z(\varepsilon^2)+O(\varepsilon^3).
\end{centeredalign}

The operator product acts from right to left. For \(\varepsilon=0.1\) radian the exact product is a rotation by 0.0100 radian about an axis tilted about four degrees from \(z\). Reversing the orientation of the loop reverses the sense of the rotation. Pulse sequences in magnetic resonance generate rotations about \(z\) from \(x\) and \(y\) pulses alone. More generally, repeated commutators of the available Hamiltonians determine which operations a set of controls can reach.

The same term measures the error of splitting an evolution. Approximating \(e^{-i(A+B)t/\hbar}\) by alternating \(A\) and \(B\) gives

\begin{centeredalign}
e^{-iAt/\hbar}e^{-iBt/\hbar}
=\exp\!\left(-\frac{i(A+B)t}{\hbar}-\frac{t^2}{2\hbar^2}[A,B]+\cdots\right).
\end{centeredalign}

Divided into \(n\) steps, the error is bounded to leading order by \(t^2\|[A,B]\|/(2n\hbar^2)\), which sets the cost of digital quantum simulation. Commuting terms can be split without error.

Two identities make commutators computable. The Leibniz rule \([A,BC]=[A,B]C+B[A,C]\) reduces composite expressions to simple ones, and the Jacobi identity relates commutators of commutators. For example, \([L_z,x]=i\hbar y\) and \([L_z,y]=-i\hbar x\) give \([L_z,x^2+y^2]=0\). The generator of rotations about \(z\) therefore commutes with any potential that depends on \(x\) and \(y\) only through \(x^2+y^2\).

Commuting Hermitian operators share an eigenbasis. If \([A,B]=0\) and \(A|a\rangle=a|a\rangle\), then \(A(B|a\rangle)=aB|a\rangle\), so \(B\) maps each eigenspace of \(A\) into itself and can be diagonalized there. The labels \(n,\ell,m\) of hydrogen come from the commuting set \(H\), \(L^2\), \(L_z\). A measurement of \(L_x\) on such an eigenstate preserves \(n\) and \(\ell\), because \(L_x\) also commutes with \(H\) and \(L^2\), but it generally leaves no definite value of \(m\). A nonzero commutator excludes a complete common eigenbasis but not every common eigenvector: all three components of \(L\) annihilate an \(\ell=0\) state.

For unbounded operators, commuting on a common dense domain does not guarantee that the spectral projectors commute. Nelson constructed operators with that property whose unitary groups fail to commute. Compatibility is therefore defined through the projectors or the unitary groups. The uncertainty principle turns a nonzero commutator into a bound on spreads, and the Heisenberg picture uses \([H,A]\) for the time dependence of \(A\).

\subsection*{References For The Physical Derivation}
\begin{itemize}
\item \href{https://arxiv.org/abs/1912.08854}{A. M. Childs, Y. Su, M. C. Tran, N. Wiebe and S. Zhu, A Theory of Trotter Error; product-formula errors controlled by commutators of the Hamiltonian terms.}
\end{itemize}



\subsection*{Equation Sources}
\begin{itemize}
\item \href{https://arxiv.org/abs/quant-ph/0005019}{arXiv:quant-ph/0005019}
\item \href{https://arxiv.org/abs/quant-ph/0001056}{arXiv:quant-ph/0001056}
\item \href{https://arxiv.org/abs/hep-ph/0104068}{arXiv:hep-ph/0104068}
\end{itemize}

\clearpage

\section{Uncertainty principle}

\label{topic:uncertainty-principle}

A beam of molecules passing a slit spreads as the slit is narrowed. Nairz, Arndt and Zeilinger measured this broadening for C70 fullerene molecules at a temperature of 900 kelvin. The slit width times the momentum transfer to the first diffraction minimum equalled Planck's constant \(h\). A narrower position distribution forces a wider momentum distribution, a statement that holds for every preparation, not only for a disturbed one.

For observables \(A\) and \(B\) and a normalized state \(\psi\), define \(f=(A-\langle A\rangle)\psi\) and \(g=(B-\langle B\rangle)\psi\). The variances are \(\sigma_A^2=\langle f|f\rangle\) and \(\sigma_B^2=\langle g|g\rangle\). The overlap \(\langle f|g\rangle\) has a real part equal to the covariance and an imaginary part equal to \(\langle[A,B]\rangle/2i\). The Cauchy--Schwarz inequality \(|\langle f|g\rangle|^2\leq\langle f|f\rangle\langle g|g\rangle\) then gives the Robertson--Schrödinger relation:

\begin{centeredalign}
\sigma_A^2\sigma_B^2\geq
\left(\tfrac12\langle\{A,B\}\rangle-\langle A\rangle\langle B\rangle\right)^2
+\left|\frac{\langle[A,B]\rangle}{2i}\right|^2.
\end{centeredalign}

Dropping the covariance gives Robertson's form \(\sigma_A\sigma_B\geq\tfrac12|\langle[A,B]\rangle|\). For position and momentum it becomes \(\hbar/2\), as derived under the canonical commutation relation.

The bound depends on the state. For a spin 1/2 with Bloch vector \(\mathbf r\), the spin components give \(\sigma_{S_x}^2=\tfrac{\hbar^2}{4}(1-r_x^2)\), \(\sigma_{S_y}^2=\tfrac{\hbar^2}{4}(1-r_y^2)\) and \(\langle[S_x,S_y]\rangle=i\hbar^2r_z/2\). For a spin along \(x\) the Robertson bound is zero, although \(S_y\) is maximally spread. Pure states satisfy the full relation with equality, because in two dimensions \(f\) and \(g\) are both orthogonal to \(\psi\) and hence parallel:

\begin{centeredalign}
(1-r_x^2)(1-r_y^2)=r_z^2+r_x^2r_y^2\qquad (|\mathbf r|=1).
\end{centeredalign}

An entropic form removes the state dependence. Maassen and Uffink proved that the Shannon entropies of two measurements obey \(H(A)+H(B)\geq-2\log_2c\), where \(c\) is the largest overlap between their eigenvectors. For \(S_z\) and \(S_x\) of a spin 1/2, \(c=1/\sqrt2\):

\begin{centeredalign}
H(Z)+H(X)\geq1\ \text{bit}.
\end{centeredalign}

A \(z\) eigenstate attains the bound with \(H(Z)=0\) and \(H(X)=1\). This form enters the security proofs of quantum key distribution.

Standard deviations can also fail to describe a measurement. A slit of width \(w\) with uniform amplitude has position spread \(w/\sqrt{12}\). Its momentum amplitude is proportional to \(\sin(pw/2\hbar)/p\), whose variance is infinite, so the variance relation holds trivially and says nothing. The central diffraction peak, between zeros at \(p=\pm h/w\), contains 90.3 percent of the probability. The product of slit width and peak half-width is \(h\), the relation found in the fullerene experiment. A width that contains a fixed fraction of the probability is therefore the meaningful measure of spread here.

These relations concern spreads in identically prepared ensembles. Relations between the error of one measurement and the disturbance it causes to another are separate inequalities with their own definitions. One such relation, proposed by Ozawa, has been tested with successive measurements of neutron spins. The commutator and the Fourier transform give the operator and wave pictures of the same incompatibility.

\subsection*{References For The Physical Derivation}
\begin{itemize}
\item \href{https://arxiv.org/abs/quant-ph/0105061}{O. Nairz, M. Arndt and A. Zeilinger, Experimental verification of the Heisenberg uncertainty principle for hot fullerene molecules; slit width times momentum transfer to the first minimum equals h.}
\item \href{https://arxiv.org/abs/1511.04857}{P. J. Coles, M. Berta, M. Tomamichel and S. Wehner, Entropic Uncertainty Relations and their Applications; the Maassen--Uffink bound and its cryptographic uses.}
\item \href{https://arxiv.org/abs/1201.1833}{J. Erhart and colleagues, Experimental demonstration of a universally valid error-disturbance uncertainty relation in spin-measurements; neutron test of Ozawa's relation.}
\end{itemize}





\clearpage

\section{Canonical commutation relation}

\label{topic:canonical-commutation-relation}

Translations in position and changes in momentum act differently when their order is reversed. For a particle on the real line, let \(X\) multiply a wave function by \(x\) and let \(P=-i\hbar\partial_x\). On smooth rapidly decreasing wave functions, direct differentiation gives

\begin{centeredalign}
[X,P]\psi=X(-i\hbar\partial_x\psi)+i\hbar\partial_x(x\psi)
=i\hbar\psi.
\end{centeredalign}

The additional term comes from differentiating the coordinate itself. It fixes the relation between the two operations, independently of the Hamiltonian subsequently chosen. With standard deviations defined in the prepared state and finite relevant moments, the Cauchy--Schwarz inequality applied to \((X-\langle X\rangle)\psi\) and \((P-\langle P\rangle)\psi\) then gives

\begin{centeredalign}
\Delta X\,\Delta P\geq\frac{\hbar}{2}.
\end{centeredalign}

This bound concerns the spread of outcomes in identically prepared ensembles. A measurement-disturbance experiment introduces an apparatus and a sequence of measurements and requires its own dynamical model. The preparation bound should not be substituted for that separate calculation.

The exponentiated relation makes domain issues more transparent. Let \(T(a)=\exp(-iaP/\hbar)\) translate a wave function by length \(a\), and \(B(b)=\exp(ibX/\hbar)\) shift its momentum by \(b\). Acting on a wave function in either order gives

\begin{centeredalign}
[T(a)\psi](x)=\psi(x-a),\qquad
T(a)B(b)=e^{-iab/\hbar}B(b)T(a).
\end{centeredalign}

The phase is proportional to the phase-space area enclosed by the two moves. It is common to the state for a single path, but can become a relative phase when the two operation sequences form coherent alternatives. This is the physical content of the central extension represented by the Heisenberg group.

No finite-dimensional matrices obey \([X,P]=i\hbar I\) exactly: taking a trace makes the left side zero and the right side nonzero. A truncated oscillator basis or a finite numerical grid must therefore depart from the relation somewhere. Such departures can be approximation effects rather than additional interactions. An exact construction records the representation and tests the states on which the commutator identity is being used.

\subsection*{References For The Physical Derivation}
\begin{itemize}
\item \href{https://arxiv.org/abs/1508.06951}{V. Moretti, Mathematical Foundations of Quantum Mechanics: An Advanced Short Course; states, operators and symmetry.}
\end{itemize}

\subsection*{Relations In The Original Papers}

\href{https://arxiv.org/html/1508.06951#S1.E21}{arXiv:1508.06951, S1.E21}. Position and momentum obey the canonical commutation relations on a common invariant domain. The surrounding derivation specifies Schwartz functions; the finite-matrix obstruction is calculated separately in the chapter.



\chapter{Measurement, instruments, and probabilities}

A measurement couples a system to an apparatus, produces a recorded outcome with a definite probability and changes the state of the system.

Projective and generalized measurements, the Born rule, collapse, and the experiments of Bell, Einstein, Podolsky and Rosen and of the quantum eraser examine this coupling from different sides.

\section*{Topics}

\begin{longtable}{p{0.92\linewidth}}
\toprule
Page \\
\midrule
\endhead
\phantomsection\label{page:povm}POVM \\
\phantomsection\label{page:measurement-in-quantum-mechanics}Measurement in quantum mechanics \\
\phantomsection\label{page:quantum-jump}Quantum jump \\
\phantomsection\label{page:wave-function-collapse}Wave function collapse \\
\phantomsection\label{page:born-rule}Born rule \\
\phantomsection\label{page:measurement-problem}Measurement problem \\
\phantomsection\label{page:quantum-eraser-experiment}Quantum eraser experiment \\
\phantomsection\label{page:projection-valued-measure}Projection-valued measure \\
\phantomsection\label{page:quantum-amplifier}Quantum amplifier \\
\phantomsection\label{page:bell-s-theorem}Bell's theorem \\
\phantomsection\label{page:quantum-nonlocality}Quantum nonlocality \\
\phantomsection\label{page:delayed-choice-quantum-eraser}Delayed-choice quantum eraser \\
\phantomsection\label{page:quantum-imaging}Quantum imaging \\
\phantomsection\label{page:wave-particle-duality}Wave–particle duality \\
\phantomsection\label{page:einstein-podolsky-rosen-paradox}Einstein–Podolsky–Rosen paradox \\
\phantomsection\label{page:electron-microscope}Electron microscope \\
\bottomrule
\end{longtable}

\clearpage

\section{POVM}

\label{topic:povm}

A photon arrives polarized either horizontally or diagonally, with equal probability. The two states are not orthogonal, so no measurement identifies them every time without error. A measurement can nevertheless be built that never errs and sometimes declares the result inconclusive.

Write the states as \(|\psi_0\rangle\) and \(|\psi_1\rangle\) with phases chosen so that \(s=\langle\psi_0|\psi_1\rangle\) is real and positive. Let \(|\psi_j^\perp\rangle\) be the state orthogonal to \(|\psi_j\rangle\). A measurement with three outcomes is defined by the positive operators

\begin{centeredalign}
E_0=\frac{|\psi_1^\perp\rangle\langle\psi_1^\perp|}{1+s},\qquad
E_1=\frac{|\psi_0^\perp\rangle\langle\psi_0^\perp|}{1+s},\qquad
E_?=I-E_0-E_1.
\end{centeredalign}

Outcome 0 never occurs for \(|\psi_1\rangle\), and outcome 1 never occurs for \(|\psi_0\rangle\). The prefactor \(1/(1+s)\) is the largest value that keeps \(E_?\) positive; its smaller eigenvalue then vanishes. Each state is identified with probability \(|\langle\psi_0|\psi_1^\perp\rangle|^2/(1+s)=1-s\). This is the Ivanovic--Dieks--Peres limit, optimal for equal prior probabilities.

For horizontal and diagonal polarization, \(s=1/\sqrt2\) and the success probability is \(0.293\). A projective measurement in the basis containing \(|\psi_0^\perp\rangle\) also never errs, but identifies only \(|\psi_1\rangle\), giving an average of \((1-s^2)/2=0.25\). Three nonzero orthogonal projectors cannot fit in a two-dimensional space. This is the general structure: a positive operator-valued measure (POVM) has effects \(E_k\geq0\) that sum to the identity, with \(p(k)=\operatorname{Tr}(\rho E_k)\).

Naimark's theorem shows how such a measurement is performed. The photon is coupled to an additional mode, the ancilla, prepared in \(|0\rangle_a\). A unitary interaction acts as

\begin{centeredalign}
U|\psi_j\rangle|0\rangle_a=\sqrt{1-s}\,|u_j\rangle|0\rangle_a+\sqrt{s}\,|\chi\rangle|1\rangle_a,
\qquad\langle u_0|u_1\rangle=0.
\end{centeredalign}

Inner products are preserved, since \((1-s)\langle u_0|u_1\rangle+s=s\), so such a unitary exists. When the ancilla is found in \(|0\rangle_a\), the photon is in one of two orthogonal states and a projective measurement identifies it without error. When the ancilla is found in \(|1\rangle_a\), both inputs leave the same state \(|\chi\rangle\), and the result is inconclusive. Restricting this larger projective measurement to the photon reproduces \(E_0\), \(E_1\) and \(E_?\) exactly.

In general, a POVM fixes only the outcome probabilities. The state after outcome \(k\) depends on how the measurement is carried out, through operators \(K_{k\alpha}\) with \(E_k=\sum_\alpha K_{k\alpha}^\dagger K_{k\alpha}\). Different implementations of the same effects can leave the photon in different states, or absorb it altogether.

Huttner and colleagues realized this scheme in 1996 using polarization-dependent loss in optical fibre. Photons removed by the loss element play the role of the inconclusive outcome. The same measurement underlies the B92 version of quantum key distribution. Missed detections add inconclusive events without creating errors. The sharp special case, with orthogonal projectors, is the projection-valued measure.

\subsection*{References For The Physical Derivation}
\begin{itemize}
\item \href{https://arxiv.org/abs/0810.1970}{S. M. Barnett and S. Croke, Quantum state discrimination; unambiguous discrimination, the Ivanovic-Dieks-Peres limit and its optical realizations.}
\item \href{https://arxiv.org/abs/0810.3536}{T. Heinosaari and M. Ziman, Guide to Mathematical Concepts of Quantum Theory; POVMs and Naimark dilation.}
\end{itemize}





\clearpage

\section{Measurement in quantum mechanics}

\label{topic:measurement-in-quantum-mechanics}

A Stern--Gerlach magnet sends silver atoms up or down according to their spin, and a screen records where each atom lands. The spin is never observed directly. It is inferred from a position that the magnet has correlated with it.

Von Neumann's model makes this correlation explicit. Let the measured observable be \(A=\sum_aa|a\rangle\langle a|\), and let a pointer have position \(q\), momentum \(p_q\) and initial wave function \(\phi(q)\) of width \(\sigma\). A brief interaction \(H=g(t)A\otimes p_q\) with total strength \(G=\int g\,dt\) translates the pointer by an amount proportional to the eigenvalue:

\begin{centeredalign}
e^{-iGA\otimes p_q/\hbar}\sum_ac_a|a\rangle\phi(q)
=\sum_ac_a|a\rangle\phi(q-Ga).
\end{centeredalign}

Finding the pointer at \(q\) then has probability density \(\sum_a|c_a|^2|\phi(q-Ga)|^2\), a sum of displaced pointer profiles. When \(G\) times the eigenvalue spacing greatly exceeds \(\sigma\), the peaks separate and outcome \(a\) occurs with probability \(|c_a|^2\). When the peaks overlap, the record yields only partial information about \(A\). The corresponding effects \(\sum_a|\phi(q-Ga)|^2|a\rangle\langle a|\) form a POVM rather than a set of projectors.

The same coupling disturbs the system. Tracing over the pointer multiplies each coherence \(\rho_{aa'}\) of the system by the overlap of the two displaced pointer states. For a Gaussian pointer,

\begin{centeredalign}
\rho_{aa'}\;\longmapsto\;\rho_{aa'}\exp\!\left[-\frac{G^2(a-a')^2}{8\sigma^2}\right].
\end{centeredalign}

Resolution and disturbance are therefore controlled by the same ratio \(G|a-a'|/\sigma\). A pointer that separates two eigenvalues cleanly removes the coherence between them, and with it the interference that observables not commuting with \(A\) would show. The diagonal elements are untouched because the interaction commutes with \(A\). This is the quantum nondemolition condition: an immediate repetition returns the same eigenvalue.

In the Stern--Gerlach magnet, the field gradient \(B'\) couples the spin component \(\sigma_z\) to the transverse position \(z\) through the energy \(\mu_BB'z\sigma_z\), with \(\mu_B\) the Bohr magneton. The pointer is the transverse momentum, shifted by \(\pm\mu_BB'\tau\) during the transit time \(\tau\). For silver atoms at 500 m/s in a 3.5 cm magnet with \(B'=10^3\) T/m, the two components separate by about 0.25 mm at the exit. This is comparable to the splitting seen in 1922.

The momentum kick far exceeds the momentum spread of a collimated beam. The interference between the spin components is therefore lost, and restoring it would require undoing the kick to within that spread. The update that follows a recorded outcome is discussed under wave function collapse. Continuous monitoring of an emitting atom, in which the pointer is the radiation field, leads to the quantum jump description.

\subsection*{References For The Physical Derivation}
\begin{itemize}
\item \href{https://arxiv.org/abs/0810.4729}{A. A. Clerk, M. H. Devoret, S. M. Girvin, F. Marquardt and R. J. Schoelkopf, Introduction to Quantum Noise, Measurement and Amplification; von Neumann pointer model, back-action and quantum limits.}
\end{itemize}





\clearpage

\section{Quantum jump}

\label{topic:quantum-jump}

A fluorescence detector records discrete photons, whereas an unobserved excited atom is often described by a smoothly decaying density operator. These descriptions refer to different information about the same atom and its radiation field. Conditioning on the detection record gives quantum trajectories; averaging over records recovers the master equation.

For Hamiltonian \(H\) and monitored decay operators \(L_k\), assume a Markovian reservoir and time-independent detection channels. In a short interval \(dt\), the no-count operator \(M_0\) and the count operators \(M_k\) are

\begin{centeredalign}
M_0=I-\frac{iH\,dt}{\hbar}-\frac{dt}{2}\sum_kL_k^\dagger L_k,
\qquad M_k=\sqrt{dt}\,L_k.
\end{centeredalign}

They satisfy \(\sum_{k\geq0}M_k^\dagger M_k=I+O(dt^2)\). The probability of a count in channel \(k\) is \(dt\,\operatorname{Tr}(L_k^\dagger L_k\rho)\); the state after that count is \(L_k\rho L_k^\dagger\) divided by its trace. Even the absence of a count changes the conditioned state, because it supplies information about its excited-state population.

\begin{centeredalign}
\dot\rho=-\frac{i}{\hbar}[H,\rho]+\sum_k
\left(L_k\rho L_k^\dagger-\frac12\{L_k^\dagger L_k,\rho\}\right).
\end{centeredalign}

This equation follows by summing the unnormalized states for all outcomes, expanding to first order in \(dt\) and subtracting the previous density operator. It is the ensemble prediction, not one particular observed sequence. For a two-level atom with \(L=\sqrt{\gamma}|g\rangle\langle e|\), no drive and an initially excited state, the survival probability is \(e^{-\gamma t}\) and the first-photon waiting-time density is \(\gamma e^{-\gamma t}\).

Changing the measurement of the outgoing field can change the conditioned trajectories while leaving the ensemble master equation unchanged. Direct photon counting gives jumps; quadrature detection gives a continuous noisy record. Consequently a trajectory belongs to the atom-field dynamics together with the measurement, not to the reduced Hamiltonian alone.

Finite efficiency separates the decay operator into observed and unobserved channels. A missed photon still affects the atom, but does not produce an observed jump. Feedback adds a further dependence: the detector record controls a subsequent Hamiltonian or operation. A mechanism assembled from decay, observation and feedback must retain these distinctions to predict both the mean population and the distribution of individual records.

\subsection*{References For The Physical Derivation}
\begin{itemize}
\item \href{https://arxiv.org/abs/quant-ph/9702007}{M. B. Plenio and P. L. Knight, The Quantum Jump Approach to Dissipative Dynamics in Quantum Optics; conditional and ensemble evolution.}
\end{itemize}

\subsection*{Relations In The Original Papers}

\href{https://arxiv.org/html/quant-ph/9702007#S4.E66}{arXiv:quant-ph/9702007, S4.E66}. Evolution conditioned on no detected emission is generated by an effective non-Hermitian Hamiltonian. The neighbouring equation gives the no-emission probability as the squared norm; detector assumptions determine the conditional evolution.



\clearpage

\section{Wave function collapse}

\label{topic:wave-function-collapse}

A trapped ion is driven from one hyperfine level to another by a resonant radio-frequency field. If the level is checked repeatedly during the drive, the transfer slows down. Frequent enough checks nearly stop it.

The calculation uses the update rule for a recorded outcome. For a measurement with orthogonal projectors \(P_k\), the Lüders rule replaces the density operator \(\rho\) after outcome \(k\) by the normalized projection. If the outcome is not recorded, the state becomes the probability-weighted average:

\begin{centeredalign}
\rho\;\longmapsto\;\frac{P_k\rho P_k}{\operatorname{Tr}(\rho P_k)},\qquad
\rho\;\longmapsto\;\sum_kP_k\rho P_k.
\end{centeredalign}

Because \(P_k^2=P_k\), an immediate repetition returns outcome \(k\) with certainty. The second map removes the coherences between different outcomes and keeps their populations.

Let the drive rotate the Bloch vector of the two-level ion about the \(x\) axis, completing a \(\pi\) rotation in time \(T\). Without interruption it transfers the ion completely. Divide \(T\) into \(N\) intervals, each rotating the vector by \(\pi/N\) and followed by a measurement of the level. Each measurement removes the transverse components, so the \(z\) component is multiplied by \(\cos(\pi/N)\) at every step:

\begin{centeredalign}
P_N=\frac12\left[1-\cos^N\frac{\pi}{N}\right]\approx\frac{\pi^2}{4N}\qquad(N\gg1).
\end{centeredalign}

One interruption gives full transfer, two give 0.5, and 64 give 0.037. Itano, Heinzen, Bollinger and Wineland observed this suppression in 1990. They drove a hyperfine transition of beryllium ions with a 256 ms radio-frequency pulse, interrupted by up to 64 short laser pulses. Ions in one of the levels scattered light from each pulse, and the measured transfer followed the prediction. This slowing is the quantum Zeno effect.

The suppression relies on the quadratic start of the transfer: with Rabi frequency \(\Omega\), the probability of leaving the initial level grows as \((\Omega t/2)^2\) at short times. A decay into a broad continuum is quadratic only for times shorter than the correlation time of the continuum. Ordinary radioactive decay is therefore not slowed by any realistic observation. Each laser pulse must also dephase the two levels completely without driving the radio-frequency transition.

The same suppression follows when each laser pulse is treated as a coupling of the ion to the scattered light, without invoking the update rule. The experiment therefore confirms that a strong probe removes the coherence between levels. It does not show that collapse is a physical process distinct from that interaction. Whether the update describes a physical change or a revision of information is the measurement problem. The quantum jump description gives the same conditioning for continuous photon counting.

\subsection*{References For The Physical Derivation}
\begin{itemize}
\item \href{https://arxiv.org/abs/quant-ph/0612187}{W. M. Itano, Perspectives on the quantum Zeno paradox; the 1990 trapped-ion experiment and its interpretation.}
\end{itemize}





\clearpage

\section{Born rule}

\label{topic:born-rule}

A spin prepared up along \(z\) enters a Stern--Gerlach magnet tilted by an angle \(\theta\). Each atom lands in one of two spots. Over many atoms, the fraction in the upper spot approaches \(\cos^2(\theta/2)\), although no single run reveals that number.

The spin state up along the tilted axis \(\mathbf n\) is \(\cos(\theta/2)|{\uparrow}\rangle+\sin(\theta/2)|{\downarrow}\rangle\). Its inner product with the prepared state is the amplitude \(\cos(\theta/2)\), and the Born rule assigns probability equal to its squared modulus. For a density operator \(\rho\) and an outcome projector \(P_k\) the same rule reads

\begin{centeredalign}
p_+=|\langle{+_{\mathbf n}}|{\uparrow_z}\rangle|^2=\cos^2\frac{\theta}{2},\qquad
p(k)=\operatorname{Tr}(\rho P_k).
\end{centeredalign}

The rule refers to relative frequencies in repeated, identically prepared runs. For \(N\) independent atoms the observed fraction scatters about \(p\) with standard deviation \(\sqrt{p(1-p)/N}\). At \(\theta=120^\circ\), \(p=1/4\), and \(10^4\) atoms fix the fraction to about \(\pm0.004\). Comparing such fractions at many angles tests the rule, provided detector efficiency and background are calibrated independently.

The squared modulus is not an arbitrary choice. Gleason's theorem states that, in a space of dimension three or more, any assignment of probabilities to projectors must take the form \(\operatorname{Tr}(\rho P)\). The assumption is that probabilities add over mutually orthogonal projectors and do not depend on which other projectors are measured alongside. The rule thus follows from the projector structure and this additivity, not from the Schrödinger equation.

A sharp consequence concerns interference among more than two paths. Let \(a\), \(b\) and \(c\) be the amplitudes at one screen point from three slits \(A\), \(B\) and \(C\). Write \(P_{AB}\) for the probability with slits \(A\) and \(B\) open, and so on. Expanding the squared moduli gives

\begin{centeredalign}
|a+b+c|^2=|a|^2+|b|^2+|c|^2+2\operatorname{Re}(a^*b+b^*c+a^*c),\\
I_3=P_{ABC}-P_{AB}-P_{BC}-P_{AC}+P_A+P_B+P_C=0.
\end{centeredalign}

Every cross term appears once in \(P_{ABC}\) and once in a two-slit pattern, so the combination \(I_3\) cancels identically. A rule in which probability were a higher power of the amplitude would leave a third-order interference term. Sinha and colleagues measured the patterns for every combination of open slits, using single photons and three slits. They bounded the three-path term below one percent of the two-path interference.

The cancellation assumes that opening a slit only adds an amplitude. Near-field paths that loop between neighbouring open slits violate this assumption slightly. Within ordinary quantum mechanics they produce a small nonzero \(I_3\), which precise tests must subtract.

For unsharp detectors the projector is replaced by a positive effect, as described under POVM. The sharp case, with orthogonal projectors that sum to the identity, is described under projection-valued measure. Wave interference supplies the amplitudes; the Born rule converts them into counting rates.

\subsection*{References For The Physical Derivation}
\begin{itemize}
\item \href{https://arxiv.org/abs/1007.4193}{U. Sinha, C. Couteau, T. Jennewein, R. Laflamme and G. Weihs, Ruling Out Multi-Order Interference in Quantum Mechanics; three-slit bound on third-order interference.}
\item \href{https://arxiv.org/abs/gr-qc/9401003}{R. D. Sorkin, Quantum Mechanics as Quantum Measure Theory; hierarchy of sum rules in which quantum probabilities have only pairwise interference.}
\end{itemize}





\clearpage

\section{Measurement problem}

\label{topic:measurement-problem}

A detector whose pointer can swing left or right interacts with a spin prepared in a superposition of up and down. Linear quantum dynamics applied to spin and detector predicts a superposition of both pointer positions. Each actual run shows one.

Suppose the detector works correctly on the basis states. Starting in a ready state \(|A_0\rangle\), the interaction \(U\) produces \(|{\uparrow}\rangle|A_\uparrow\rangle\) from \(|{\uparrow}\rangle|A_0\rangle\) and \(|{\downarrow}\rangle|A_\downarrow\rangle\) from \(|{\downarrow}\rangle|A_0\rangle\). Linearity then fixes its action on any superposition:

\begin{centeredalign}
U(\alpha|{\uparrow}\rangle+\beta|{\downarrow}\rangle)|A_0\rangle
=\alpha|{\uparrow}\rangle|A_\uparrow\rangle+\beta|{\downarrow}\rangle|A_\downarrow\rangle.
\end{centeredalign}

Adding an observer, a camera or a printer extends this correlation to more systems without ending it. The Born rule assigns the probabilities \(|\alpha|^2\) and \(|\beta|^2\), but the evolution contains no step at which one term is selected.

Decoherence explains part of what is seen. Air molecules and thermal photons scatter from the pointer and acquire states that depend on its position. With environment states \(|E_\uparrow\rangle\) and \(|E_\downarrow\rangle\), the reduced state of spin and pointer becomes

\begin{centeredalign}
\rho_{SA}=|\alpha|^2|{\uparrow}A_\uparrow\rangle\langle{\uparrow}A_\uparrow|
+|\beta|^2|{\downarrow}A_\downarrow\rangle\langle{\downarrow}A_\downarrow|\\
+\left(\alpha\beta^*\Gamma|{\uparrow}A_\uparrow\rangle\langle{\downarrow}A_\downarrow|+{\rm h.c.}\right),
\qquad \Gamma=\langle E_\downarrow|E_\uparrow\rangle.
\end{centeredalign}

If \(N\) independent scatterers each leave states with overlap \(c\), then \(\Gamma=c^N\). With \(c=\cos0.1\approx0.995\) and \(N=10^4\), \(|\Gamma|\approx2\times10^{-22}\), so interference between pointer positions cannot be detected on spin and pointer alone. The pointer basis is selected by the interaction, here position, because scattering depends on where the pointer is.

Decoherence does not select an outcome. The global state still contains both terms, and \(\rho_{SA}\) is a reduced description of an entangled state, not a list of which outcome occurred. The measurement problem is that unitary evolution alone yields neither single records nor their probabilities; both enter as separate assumptions. Interpretations divide on how to supply the missing step: by keeping all branches, by adding particle positions, or by changing the dynamics.

The last option makes predictions. In the Ghirardi--Rimini--Weber model each nucleon is spontaneously localized to a length \(r_C\approx10^{-7}\) m at a rate \(\lambda\approx10^{-16}\) s\(^{-1}\). For a pointer whose \(N\) nucleons are displaced by more than \(r_C\), one localization of any nucleon collapses the whole superposition:

\begin{centeredalign}
\Gamma_{\rm collapse}=N\lambda\approx10^{23}\times10^{-16}\,{\rm s^{-1}}=10^{7}\,{\rm s^{-1}}.
\end{centeredalign}

A single nucleon is localized about once in \(3\times10^8\) years, while a macroscopic pointer resolves within 0.1 microsecond. The modified dynamics also predicts slight heating and spontaneous radiation. Interferometry with massive molecules and searches for this radiation bound \(\lambda\) at each \(r_C\). Current upper bounds on \(\lambda\) at \(r_C=10^{-7}\) m remain above the original value, so the model is constrained but not excluded. Schrödinger's cat poses the same problem for a larger system.

\subsection*{References For The Physical Derivation}
\begin{itemize}
\item \href{https://arxiv.org/abs/quant-ph/0312059}{M. Schlosshauer, Decoherence, the measurement problem, and interpretations of quantum mechanics; pointer states and what decoherence does not resolve.}
\item \href{https://arxiv.org/abs/1204.4325}{A. Bassi, K. Lochan, S. Satin, T. P. Singh and H. Ulbricht, Models of Wave-function Collapse, Underlying Theories, and Experimental Tests; GRW amplification and experimental bounds.}
\end{itemize}





\clearpage

\section{Quantum eraser experiment}

\label{topic:quantum-eraser-experiment}

Light passing two slits forms fringes on a screen. If each path leaves a mark in another degree of freedom, the fringes fade. A later measurement of that mark in a suitable basis can recover fringes in sorted subsets of the records.

Let the screen amplitudes from the two slits be \(\psi_1(x)\) and \(\psi_2(x)\), and let the paths leave marker states \(|m_1\rangle\) and \(|m_2\rangle\). The joint state is proportional to \(\psi_1(x)|m_1\rangle+\psi_2(x)|m_2\rangle\). Ignoring the marker, the screen intensity is

\begin{centeredalign}
P(x)\propto|\psi_1|^2+|\psi_2|^2+2\operatorname{Re}\big[\psi_1^*\psi_2\langle m_1|m_2\rangle\big].
\end{centeredalign}

The fringe term is multiplied by the marker overlap. For equal intensities the visibility equals \(|\langle m_1|m_2\rangle|\), and orthogonal markers remove the fringes entirely. The marking need not deliver any momentum kick. The interference term disappears because the screen amplitude is now correlated with an orthogonal marker state.

Now take orthogonal markers and measure them in a rotated basis, with an angle \(\beta\) and a phase \(\gamma\) that the experimenter sets:

\begin{centeredalign}
|e_+\rangle=\cos\tfrac{\beta}{2}|m_1\rangle+e^{i\gamma}\sin\tfrac{\beta}{2}|m_2\rangle,\\
|e_-\rangle=\sin\tfrac{\beta}{2}|m_1\rangle-e^{i\gamma}\cos\tfrac{\beta}{2}|m_2\rangle.
\end{centeredalign}

Conditioning on each outcome projects the marker and leaves a definite screen amplitude. The two conditional patterns have the same visibility:

\begin{centeredalign}
\psi_+=\cos\tfrac{\beta}{2}\,\psi_1+e^{-i\gamma}\sin\tfrac{\beta}{2}\,\psi_2,\qquad
\psi_-=\sin\tfrac{\beta}{2}\,\psi_1-e^{-i\gamma}\cos\tfrac{\beta}{2}\,\psi_2,\\
V_{\rm cond}=|\sin\beta|.
\end{centeredalign}

At \(\beta=0\) the marker measurement reveals the path and the subsets show no fringes. At \(\beta=\pi/2\) the marker basis is unbiased with respect to the paths; each subset shows full-contrast fringes, one shifted by half a period relative to the other. The two subsets always add to the unmarked, fringe-free pattern. Erasure therefore selects subsets, and the total distribution on the screen is unchanged.

Walborn, Terra Cunha, Pádua and Monken carried this out with polarization-entangled photon pairs from parametric down-conversion. Quarter-wave plates in front of the two slits marked the path in the circular polarization of the signal photon, and the coincidence fringes vanished. A polarizer on the distant idler photon then projected the pair onto a basis unbiased with respect to the paths. Fringes and antifringes reappeared in the coincidence counts.

The entangled partner matters. If a single photon carries its own polarization mark, a polarizer at \(45^\circ\) before the screen recovers fringes. That version is reproduced by classical polarization optics, through the Fresnel--Arago laws. With an entangled partner, the basis is chosen on a photon that never passes the slits, and the fringes exist only in coincidence data.

The trade-off between partial marking and partial visibility is the subject of wave-particle duality. The delayed-choice quantum eraser applies the same conditioning to an idler recorded after its partner.

\subsection*{References For The Physical Derivation}
\begin{itemize}
\item \href{https://arxiv.org/abs/quant-ph/0106078}{S. P. Walborn, M. O. Terra Cunha, S. Padua and C. H. Monken, A double-slit quantum eraser; polarization marking and erasure with entangled photons.}
\item \href{https://arxiv.org/abs/1407.2930}{X.-S. Ma, J. Kofler and A. Zeilinger, Delayed-choice gedanken experiments and their realizations; quantum erasers and complementarity.}
\end{itemize}





\clearpage

\section{Projection-valued measure}

\label{topic:projection-valued-measure}

A row of detectors along a line registers where a particle lands. Each region of the line corresponds to a yes-or-no question about the particle. Answers for disjoint regions exclude one another, and some region always answers yes.

For a particle on the real line, the question for a region \(\Delta\) is represented by an operator \(E(\Delta)\). It multiplies the wave function by the indicator function \(\chi_\Delta\), which equals 1 inside \(\Delta\) and 0 outside. These operators are orthogonal projectors, and they combine like the regions themselves:

\begin{centeredalign}
[E(\Delta)\psi](x)=\chi_\Delta(x)\psi(x),\qquad
E(\Delta_1)E(\Delta_2)=E(\Delta_1\cap\Delta_2),\qquad E(\mathbb R)=I.
\end{centeredalign}

For disjoint regions the projectors add, including for countable unions. This family is a projection-valued measure (PVM). The probability of a yes answer in a state \(\psi\) is \(\langle\psi|E(\Delta)|\psi\rangle=\int_\Delta|\psi|^2dx\). By the spectral theorem every self-adjoint operator has exactly one such measure on the real line, so a PVM is the complete description of a sharp observable. For a discrete observable such as a spin component, the measure is concentrated on the eigenvalues. Then \(E(\{m\})\) is the projector onto the eigenspace with eigenvalue \(m\), and a Stern--Gerlach magnet sends each eigenspace to its own spot.

Consider the ground state \(\sqrt{2/L}\sin(\pi x/L)\) of a particle in a box of length \(L\). The probability that a detector covering the middle third responds is

\begin{centeredalign}
p\left(\tfrac{L}{3}<x<\tfrac{2L}{3}\right)=\frac{2}{L}\int_{L/3}^{2L/3}\sin^2\frac{\pi x}{L}\,dx
=\frac13+\frac{\sqrt3}{2\pi}\approx0.609.
\end{centeredalign}

The \(n=2\) level, with a node at the centre, gives \(1/3-\sqrt3/(4\pi)\approx0.196\) instead. Because \(E(\Delta)^2=E(\Delta)\), a state that has just answered yes lies in the range of \(E(\Delta)\). An immediate repetition then answers yes with certainty, which is the repeatability that defines a sharp measurement.

Commuting observables share a joint measure. The \(x\) and \(y\) coordinates of a particle give rectangle projectors \(E_X(\Delta_x)E_Y(\Delta_y)\), and a camera measures both at once. Position and momentum have no such joint measure. A nonzero wave function cannot vanish outside a bounded interval while its Fourier transform also vanishes outside a bounded interval, so

\begin{centeredalign}
\operatorname{ran}E_X(\Delta)\cap\operatorname{ran}E_P(\Gamma)=\{0\}
\qquad\text{for bounded }\Delta,\Gamma.
\end{centeredalign}

No state answers yes with certainty to both questions. This does not follow from the variance bound of the uncertainty principle. It forbids exact confinement to bounded regions of any size, including regions whose product of widths greatly exceeds \(\hbar\).

Real pixels have blurred edges. A detector whose response function \(\mu_\Delta(x)\) takes values strictly between 0 and 1 near the edges is represented by \(F(\Delta)=\int\mu_\Delta(x)\,|x\rangle\langle x|\,dx\), which satisfies \(F^2\neq F\). It is a POVM rather than a PVM, and its records need not repeat. The PVM is the ideal limit in which the response is exactly 0 or 1.

\subsection*{References For The Physical Derivation}
\begin{itemize}
\item \href{https://arxiv.org/abs/0810.3536}{T. Heinosaari and M. Ziman, Guide to Mathematical Concepts of Quantum Theory; projection-valued measures, joint measurability and smeared observables.}
\item \href{https://arxiv.org/abs/1508.06951}{V. Moretti, Mathematical Foundations of Quantum Mechanics: An Advanced Short Course; spectral theorem and projection-valued measures of self-adjoint operators.}
\end{itemize}





\clearpage

\section{Quantum amplifier}

\label{topic:quantum-amplifier}

An amplifier that increases both quadratures of a bosonic signal must also introduce fluctuations. Let \(a\) be the input annihilation operator, with \([a,a^\dagger]=1\), and let \(G>1\) be the intensity gain. Multiplying \(a\) by \(\sqrt G\) alone would make the output commutator equal to \(G\), so it cannot describe an output mode with the original normalization.

A second, independent mode \(b\) supplies the required additional degree of freedom. The phase-insensitive amplification relation is

\begin{centeredalign}
a_{\rm out}=\sqrt G\,a_{\rm in}+\sqrt{G-1}\,b_{\rm in}^\dagger,
\qquad [a_{\rm out},a_{\rm out}^\dagger]=G-(G-1)=1.
\end{centeredalign}

The creation operator has the opposite commutator sign, exactly compensating the excess from the amplified signal. A parametric interaction realizes this transformation by producing correlated excitations in a signal and an idler. Energy is supplied by a pump; the idler is a physical channel rather than an adjustable numerical correction.

Define \(X=(a+a^\dagger)/\sqrt2\) and \(P=(a-a^\dagger)/(i\sqrt2)\), so vacuum has variance \(1/2\) in each quadrature. For uncorrelated inputs and a vacuum idler,

\begin{centeredalign}
\operatorname{Var}X_{\rm out}=G\operatorname{Var}X_{\rm in}+\frac{G-1}{2},
\qquad \langle n_{\rm out}\rangle=G\langle n_{\rm in}\rangle+G-1.
\end{centeredalign}

The second expression includes spontaneous output even for vacuum signal input. Referring the quadrature noise back to the input gives an added variance \((G-1)/(2G)\), approaching half a quantum at large gain. A thermally occupied idler adds still more noise. The commutator determines the necessary channel, while its state determines how much fluctuation that channel contributes.

A phase-sensitive amplifier follows a different relation: one quadrature is stretched and the conjugate one is compressed, preserving their commutator without amplifying both equally. It can amplify a known quadrature without the same added-noise bound. Thus the bound is attached to a specific physical task, not to every process called amplification.

This is an explicit construction from a failed relation. The proposed gain-only map has a nonzero commutator defect; adding an independent idler with the appropriate coefficient restores the algebra and predicts the noise. The correction is already established quantum optics. Its value here is that every additional ingredient changes a calculable experimental consequence.

\subsection*{References For The Physical Derivation}
\begin{itemize}
\item \href{https://arxiv.org/abs/1110.3234}{C. Weedbrook and colleagues, Gaussian Quantum Information; Gaussian channels and amplification.}
\end{itemize}

\subsection*{Relations In The Original Papers}

\href{https://arxiv.org/html/1110.3234v1#S5.E110}{arXiv:1110.3234v1, S5.E110}. Complete positivity ties the noise covariance to the gain of a one-mode Gaussian channel. The source uses vacuum quadrature variance one; the chapter uses one half. Its minimum noise follows after this normalization change.



\clearpage

\section{Bell's theorem}

\label{topic:bell-s-theorem}

Two spin-1/2 particles leave a common source in opposite directions. Each meets a Stern--Gerlach magnet whose orientation is chosen at random just before arrival, and each detector records \(+1\) or \(-1\). Bell asked whether the correlations of these records can be produced by properties fixed at the source together with the local settings alone.

Let each pair carry a variable \(\lambda\) with distribution \(\rho(\lambda)\). A local model gives outcomes \(A_a(\lambda)=\pm1\) and \(B_b(\lambda)=\pm1\), where \(A\) does not depend on the distant setting \(b\) and \(\rho\) depends on neither setting. With two settings per side, one bracket below vanishes and the other equals \(\pm2\) for every pair:

\begin{centeredalign}
S_\lambda=A_a(B_b+B_{b'})+A_{a'}(B_b-B_{b'})=\pm2\\
|E(a,b)+E(a,b')+E(a',b)-E(a',b')|\leq2
\end{centeredalign}

Here \(E(a,b)\) is the average of \(A_aB_b\) over \(\lambda\). Random local responses do not help, because they are mixtures of deterministic ones and the bound is linear in the probabilities. This is the Clauser--Horne--Shimony--Holt (CHSH) inequality.

For the spin singlet, quantum mechanics gives \(E(a,b)=-\cos(\theta_a-\theta_b)\), with \(\theta\) the magnet angle in a common plane. The angles \(0^\circ\) and \(90^\circ\) on one side, with \(45^\circ\) and \(-45^\circ\) on the other, give \(|S|=4/\sqrt2=2\sqrt2\approx2.83\).

No quantum state and no choice of \(\pm1\) observables does better. For the Bell operator \(\mathcal B=A\otimes(B+B')+A'\otimes(B-B')\), squaring and using \(A^2=B^2=I\) gives

\begin{centeredalign}
\mathcal B^2=4I-[A,A']\otimes[B,B'],\qquad
|\langle\mathcal B\rangle|\leq2\sqrt2.
\end{centeredalign}

Each commutator has norm at most 2, so no eigenvalue of \(\mathcal B^2\) exceeds 8. This is the Tsirelson bound. It also shows that a violation requires noncommuting settings on both sides; commuting pairs return the classical value 2. If the singlet is mixed with white noise so that correlations shrink by a factor \(v\), then \(S=2\sqrt2\,v\) and a violation requires \(v>1/\sqrt2\).

An experiment must also close two gaps in the derivation. If detectors miss many particles, the detected subset may be unrepresentative; for a maximally entangled state the CHSH test needs efficiency above \(2(\sqrt2-1)\approx0.83\). If one setting can influence the distant detector before its record is final, a local signal could fix the correlations. The settings must also be statistically independent of \(\lambda\).

In 2015 Hensen and colleagues entangled electron spins in two diamond defects 1.3 km apart. Settings were chosen randomly and each spin was measured within the light travel time between the sites. In 245 heralded trials, each giving an outcome at both sites, they found \(S=2.42\pm0.20\). Photon experiments in Vienna and Boulder closed both gaps in the same year.

The theorem constrains local models of the records; it does not modify quantum mechanics. The Einstein--Podolsky--Rosen argument, which motivated it, used predictions with certainty; quantum nonlocality extends the comparison to three particles without an inequality. Device-independent quantum key distribution uses the observed value of \(S\) to bound an eavesdropper's information.

\subsection*{References For The Physical Derivation}
\begin{itemize}
\item \href{https://arxiv.org/abs/1303.2849}{N. Brunner, D. Cavalcanti, S. Pironio, V. Scarani and S. Wehner, Bell nonlocality; CHSH inequality, Tsirelson bound and loopholes.}
\item \href{https://arxiv.org/abs/1508.05949}{B. Hensen and colleagues, Experimental loophole-free violation of a Bell inequality using entangled electron spins separated by 1.3 km; S = 2.42 in 245 trials.}
\end{itemize}





\clearpage

\section{Quantum nonlocality}

\label{topic:quantum-nonlocality}

Three photons are emitted together and sent to three separated laboratories. Each laboratory measures one of two polarization observables. For a Greenberger--Horne--Zeilinger (GHZ) state, certain products of the three records are predicted with certainty. These certain predictions alone contradict any local assignment of values, without averaging an inequality.

Write the state as \((|000\rangle+|111\rangle)/\sqrt2\), with \(|0\rangle\) and \(|1\rangle\) the eigenstates of the Pauli operator \(Z\). Applying the Pauli operators \(X\) and \(Y\) directly gives the eigenvalue relations

\begin{centeredalign}
X_1X_2X_3|{\rm GHZ}\rangle=+|{\rm GHZ}\rangle,\\
X_1Y_2Y_3|{\rm GHZ}\rangle=Y_1X_2Y_3|{\rm GHZ}\rangle=Y_1Y_2X_3|{\rm GHZ}\rangle=-|{\rm GHZ}\rangle.
\end{centeredalign}

Suppose each photon carries predetermined values \(x_j=\pm1\) and \(y_j=\pm1\), the records it would give if \(X\) or \(Y\) were measured. Locality requires that these values do not depend on which observables the other two laboratories choose. The three \(XYY\)-type relations then fix the product of all \(x_j\), because each \(y_j\) appears twice:

\begin{centeredalign}
(x_1y_2y_3)(y_1x_2y_3)(y_1y_2x_3)=x_1x_2x_3=(-1)^3=-1.
\end{centeredalign}

Quantum mechanics predicts \(x_1x_2x_3=+1\) with certainty. No local assignment satisfies all four relations. For two spin-1/2 particles, perfect correlations alone do not suffice: the perfect anticorrelations of a singlet along equal axes have a local model. Bell's inequality needs the imperfect correlations at unequal angles. Each run measures only one combination of observables, so the contradiction is established across runs with different settings. Pan, Bouwmeester, Daniell, Weinfurter and Zeilinger performed this test with three-photon GHZ states in 2000. The \(XXX\) products came out predominantly \(+1\), as quantum mechanics predicts and opposite to the local inference from the \(XYY\)-type data.

Imperfect states require a bound that tolerates noise. Mermin's combination \(M=\langle XXX\rangle-\langle XYY\rangle-\langle YXY\rangle-\langle YYX\rangle\) is at most 2 for local values, while the GHZ state gives 4. If the state is mixed with white noise so that correlations shrink by a factor \(v\), then \(M=4v\), and a violation requires only \(v>1/2\). The two-particle CHSH test of Bell's theorem requires \(v>1/\sqrt2\).

These correlations cannot carry a message. Let laboratory 1 measure with effects \(E_a\) and laboratory 2 with effects \(F_b^y\) for setting \(y\). Because the effects of each setting sum to the identity, the record statistics in laboratory 1 do not depend on the distant choice:

\begin{centeredalign}
\sum_b\operatorname{Tr}[\rho(E_a\otimes F^y_b)]=\operatorname{Tr}(\rho_1E_a),
\qquad \rho_1=\operatorname{Tr}_2\rho.
\end{centeredalign}

Nonlocality appears only when records from separated laboratories are compared. Entanglement is necessary for it but not sufficient. Werner showed that a singlet mixed with enough white noise remains entangled yet admits a local model for all projective measurements. Device-independent quantum key distribution relies on observed nonlocal correlations rather than on a model of the devices. The Einstein--Podolsky--Rosen paradox is the two-particle argument that preceded these tests.

\subsection*{References For The Physical Derivation}
\begin{itemize}
\item \href{https://arxiv.org/abs/0712.0921}{D. M. Greenberger, M. A. Horne and A. Zeilinger, Going Beyond Bell's Theorem; the three-particle argument without inequalities.}
\item \href{https://arxiv.org/abs/1303.2849}{N. Brunner, D. Cavalcanti, S. Pironio, V. Scarani and S. Wehner, Bell nonlocality; Mermin inequality, no-signalling and local models of entangled states.}
\end{itemize}





\clearpage

\section{Delayed-choice quantum eraser}

\label{topic:delayed-choice-quantum-eraser}

A photon reaches a detector that scans across a screen, and its position is recorded. Several nanoseconds later its partner photon is sent at random along one of two routes. One leads to a detector that reveals where the pair was emitted; the other leads to a beam splitter that hides this information. Sorting the screen records by the partner's later result reveals interference fringes in some subsets.

In the experiment of Kim, Yu, Kulik, Shih and Scully, a pump laser illuminates a nonlinear crystal in two regions \(A\) and \(B\). Either region may emit a signal and idler pair, so the pair is in the state \((|A\rangle_s|A\rangle_i+|B\rangle_s|B\rangle_i)/\sqrt2\). A lens maps the two signal paths to amplitudes \(\psi_A(x)\) and \(\psi_B(x)\) at the screen detector \(D_0\).

The idler meets a 50:50 beam splitter on each path. One output goes to a detector that receives only idlers from \(A\) (\(D_3\)) or only from \(B\) (\(D_4\)). The other outputs are combined on a central beam splitter whose two ports, \(D_1\) and \(D_2\), cannot tell \(A\) from \(B\). Projecting both photons onto the detected states gives the joint probabilities

\begin{centeredalign}
P(x,D_{1,2})=\tfrac18|\psi_A(x)\pm\psi_B(x)|^2,\qquad
P(x,D_3)=\tfrac14|\psi_A(x)|^2,\qquad P(x,D_4)=\tfrac14|\psi_B(x)|^2.
\end{centeredalign}

Coincidences with \(D_1\) show fringes and those with \(D_2\) show the complementary antifringes, shifted by half a period. Coincidences with \(D_3\) or \(D_4\) show the smooth single-path profile. Each idler detector receives a quarter of the pairs.

The time order of the two detections plays no part in this calculation. The projector for the signal position and the projector for an idler detector act on different photons and commute. The joint probability is therefore the same whichever record is made earlier, and in any reference frame:

\begin{centeredalign}
P(x,k)=\big\|(\langle x|_s\otimes\langle k|_i)\Psi\big\|^2,\qquad
\sum_kP(x,k)=\tfrac12\big(|\psi_A(x)|^2+|\psi_B(x)|^2\big).
\end{centeredalign}

The screen pattern by itself has no fringes, whatever happens to the idler and whenever it happens. The idler path did not affect where the signal landed. It determines which subset of the existing records is selected, and the fringes appear only when the subsets are compared using the idler records. Those records travel by ordinary means, so no information is sent to the earlier detection.

In the 2000 experiment the idler paths were longer, so the idler was registered several nanoseconds after its partner. Ma and colleagues later made the choice with a random number generator spacelike separated from the interference, over a 144 km free-space link between two Canary Islands. Interference again appeared only in the sorted subsets.

The calculation is the conditional-fringe analysis of the quantum eraser experiment, applied to an entangled pair whose members are recorded at different times. The commuting projectors are also why the correlations in Bell experiments cannot be used to send messages.

\subsection*{References For The Physical Derivation}
\begin{itemize}
\item \href{https://arxiv.org/abs/quant-ph/9903047}{Y.-H. Kim, R. Yu, S. P. Kulik, Y. H. Shih and M. O. Scully, A Delayed Choice Quantum Eraser; the two-region down-conversion experiment analysed here.}
\item \href{https://arxiv.org/abs/1206.6578}{X.-S. Ma and colleagues, Quantum erasure with causally disconnected choice; spacelike-separated choice over 144 km.}
\item \href{https://arxiv.org/abs/1407.2930}{X.-S. Ma, J. Kofler and A. Zeilinger, Delayed-choice gedanken experiments and their realizations; review of the time-order analysis.}
\end{itemize}





\clearpage

\section{Quantum imaging}

\label{topic:quantum-imaging}

Imaging a faintly absorbing sample with few photons, such as a cell that light damages, is limited by the random arrival of photons. Light from parametric down-conversion comes in pairs. Counting the partner photons measures that randomness, which can then be subtracted.

Consider one pixel and an object with transmission \(t=1-\alpha\). With coherent light, the transmitted count \(N_t\) is Poisson distributed with mean \(t\bar N\). The estimate \(\hat t=N_t/\bar N\) then has variance \(t/\bar N\), the shot-noise limit. In twin beams, signal and idler pixel pairs contain equal photon numbers. The idler count \(N_i\) therefore tells how many photons entered the sample in the matching signal pixel.

Each photon then survives the object independently with probability \(t\). The signal count \(N_s\) is binomially distributed given \(N_i\), and the ratio \(\hat t=N_s/N_i\) has a smaller variance:

\begin{centeredalign}
\operatorname{Var}\hat t_{\rm coh}=\frac{t}{\bar N},\qquad
\operatorname{Var}\hat t_{\rm twin}=\frac{t(1-t)}{\bar N},\qquad
\frac{\operatorname{Var}\hat t_{\rm twin}}{\operatorname{Var}\hat t_{\rm coh}}=\alpha.
\end{centeredalign}

Fluctuations in the number of pairs cancel, because they appear equally in both beams. Only the randomness of absorption itself remains. For a sample absorbing \(\alpha=1\%\), the variance falls a hundredfold. The same precision is then reached with a hundredth of the photons, an important saving for light-sensitive samples.

Real detectors miss photons. If each beam is detected with efficiency \(\eta\), a missed idler is mistaken for an absent photon. The statistic used in experiments is the noise reduction factor. For a transparent region of the sample it becomes

\begin{centeredalign}
\frac{\operatorname{Var}(N_s-N_i)}{\langle N_s+N_i\rangle}=1-\eta\qquad(t=1),
\end{centeredalign}

compared with 1 for two independent beams at the shot-noise level. The advantage thus requires high efficiency. Each pixel must also collect both members of most pairs, so pixels must be larger than the coherence area of the down-converted light. That area sets the spatial resolution.

Brida, Genovese and Ruo Berchera demonstrated sub-shot-noise imaging in 2010. They recorded twin beams from parametric down-conversion on a CCD camera, placed a weakly absorbing object in one beam, and subtracted the matching pixels of the other. The object appeared with a larger signal-to-noise ratio than with classical illumination at the same photon number.

Ghost imaging also uses pair correlations. A detector without spatial resolution collects light behind the object, while a camera in the other beam forms the image in coincidence counts. Ghost images can also be formed with classical thermal light at lower visibility, so a ghost image alone is not a signature of quantum correlations. The photon-number difference above is quantum, because classical fields cannot push the noise reduction factor below 1. Quantum metrology treats the same comparison of precision per photon in general form.

\subsection*{References For The Physical Derivation}
\begin{itemize}
\item \href{https://arxiv.org/abs/1004.1274}{G. Brida, M. Genovese and I. Ruo Berchera, Experimental realization of sub-shot-noise quantum imaging; twin-beam absorption imaging.}
\item \href{https://arxiv.org/abs/1601.06066}{M. Genovese, Real applications of quantum imaging; sub-shot-noise and ghost imaging.}
\end{itemize}





\clearpage

\section{Wave–particle duality}

\label{topic:wave-particle-duality}

Electrons sent one at a time through a two-path interferometer arrive at single points, yet their accumulated positions form fringes. A detector that records which path each electron took reduces the fringe contrast. Englert's inequality states exactly how much path information can coexist with how much contrast.

The wave behaviour follows from the de Broglie relation \(\lambda=h/p\), which the diffraction of electrons by a nickel crystal confirmed, as described under the wave function.

Now let a beam splitter send the particle into path 1 or 2 with probabilities \(w_1\) and \(w_2\). An internal degree of freedom, the marker, is left in state \(|m_1\rangle\) or \(|m_2\rangle\). After recombination with relative phase \(\phi\), the two output ports have probabilities \(\frac12[1\pm2\sqrt{w_1w_2}\operatorname{Re}(e^{i\phi}\langle m_1|m_2\rangle)]\). The visibility and the a priori predictability of the path are

\begin{centeredalign}
V=2\sqrt{w_1w_2}\,|\langle m_1|m_2\rangle|,\qquad P=|w_1-w_2|.
\end{centeredalign}

Path information stored in the marker is quantified by the distinguishability \(D\). A measurement on the marker identifies the path correctly with probability at most \((1+D)/2\). Helstrom's result for discriminating two states gives \(D\) as a trace norm, which for pure markers can be evaluated in closed form:

\begin{centeredalign}
D=\big\|w_1|m_1\rangle\langle m_1|-w_2|m_2\rangle\langle m_2|\big\|_1
=\sqrt{1-4w_1w_2|\langle m_1|m_2\rangle|^2},\\
V^2+D^2=1\ \ \text{(pure markers)},\qquad V^2+D^2\leq1\ \ \text{(general)}.
\end{centeredalign}

If the marker is mixed, or part of it escapes to an environment that is not examined, the equality becomes an inequality. Without a marker, \(D\) reduces to \(P\) and the relation becomes \(V^2+P^2\leq1\). For equal paths and a marker overlap of \(\cos30^\circ\), the fringes keep 87 percent contrast while the path can be guessed correctly 75 percent of the time. Intermediate cases have partial fringes and partial path knowledge, related by this single inequality.

The bound refers to the information available in the marker, not to what an experimenter actually measures. A marker measurement in a basis unrelated to the paths extracts no path information. Fringes then reappear in subsets of the records, as in the quantum eraser experiment.

Dürr, Nonn and Rempe tested the mechanism in an atom interferometer in 1998. A microwave field stored the path in the internal hyperfine state of each atom, and the fringes vanished. The momentum transferred by the microwave photons was far too small to wash out the fringes by random kicks. The loss of contrast came from the correlation between path and marker, as the formula states.

The relation applies to any two-path interferometer whose marker can be modelled as a quantum state. With a thermal environment as the unread marker, the same overlap describes decoherence. The electron microscope uses the wave side of this relation, forming images from interference between scattered and unscattered electrons.

\subsection*{References For The Physical Derivation}
\begin{itemize}
\item \href{https://arxiv.org/abs/1403.4687}{P. J. Coles, J. Kaniewski and S. Wehner, Equivalence of wave-particle duality to entropic uncertainty; visibility-distinguishability relations.}
\item \href{https://arxiv.org/abs/2507.19801}{H. Lin, Y.-K. Lu, V. Fedoseev, Y. K. Lee, J. Lyu and W. Ketterle, Fringe visibility and which-way information in Young's double slit experiments with light scattered from single atoms; motional states of trapped atoms as which-path markers.}
\end{itemize}





\clearpage

\section{Einstein–Podolsky–Rosen paradox}

\label{topic:einstein-podolsky-rosen-paradox}

Two particles interact and fly apart so that both their separation and their total momentum are sharply fixed. Measuring the position of one particle then predicts the position of the other, and measuring its momentum predicts the other's momentum. In 1935 Einstein, Podolsky and Rosen argued that both predicted values must correspond to physical properties of the distant particle.

Such a state exists because the relative position and the total momentum commute. With \(x_j\) and \(p_j\) the position and momentum of particle \(j\),

\begin{centeredalign}
[x_1-x_2,\;p_1+p_2]=[x_1,p_1]-[x_2,p_2]=i\hbar-i\hbar=0,\\
\psi(x_1,x_2)\propto\delta(x_1-x_2-L)=\frac{1}{2\pi\hbar}\int e^{ip(x_1-x_2-L)/\hbar}dp.
\end{centeredalign}

The integral shows that the same state contains every momentum pair \((p,-p)\). A position measurement on particle 1 that gives \(\xi\) predicts \(x_2=\xi-L\), and a momentum measurement giving \(\kappa\) predicts \(p_2=-\kappa\). Neither measurement acts on particle 2. EPR concluded that particle 2 has definite position and momentum, which no quantum state of particle 2 alone can describe. Hence, they argued, the quantum description is incomplete.

The ideal state is not normalizable, but an optical version is. A nondegenerate parametric amplifier produces two beams in a two-mode squeezed state. Use dimensionless field quadratures with \([x_j,p_j]=i\) and vacuum variance 1/2, and squeezing parameter \(r\). Then \(\operatorname{Var}(x_1-x_2)=\operatorname{Var}(p_1+p_2)=e^{-2r}\), while each beam alone shows enlarged noise \(\cosh(2r)/2\).

Reid turned the argument into a measurable test. Let \(\Delta^2_{\rm inf}x_2\) be the smallest mean-square error with which \(x_2\) can be predicted from a record of \(x_1\), and similarly for \(p_2\). If beam 2 had a local quantum state, the uncertainty principle would require the product of these inferred variances to be at least 1/4. The best linear prediction gives

\begin{centeredalign}
\Delta^2_{\rm inf}x_2=\Delta^2_{\rm inf}p_2=\frac{1}{2\cosh2r},\qquad
\Delta^2_{\rm inf}x_2\,\Delta^2_{\rm inf}p_2=\frac{1}{4\cosh^22r}<\frac14.
\end{centeredalign}

At \(r=0.5\) the product is \(0.105\), well below the bound. Ou, Pereira, Kimble and Peng observed an inferred-variance product below the bound in 1992. Loss on either beam adds vacuum noise and raises the inferred variances, so the demonstration requires efficient transmission and detection.

The EPR argument does not by itself exclude local hidden variables. The Wigner function of the two-mode squeezed state is a positive Gaussian, which can serve as a local probability distribution for all quadrature records. Bell showed that other measurements, such as spin components along several axes, separate the options. Bell's theorem decides between local models and quantum predictions, whereas the EPR test shows that local states of one particle cannot explain the predictions made from the other. Quantum nonlocality treats the stronger correlations in multipartite states.

\subsection*{References For The Physical Derivation}
\begin{itemize}
\item \href{https://arxiv.org/abs/0806.0270}{M. D. Reid and colleagues, The Einstein-Podolsky-Rosen paradox: from concepts to applications; inferred-variance criterion and two-mode squeezed states.}
\end{itemize}





\clearpage

\section{Electron microscope}

\label{topic:electron-microscope}

A transmission electron microscope sends electrons, accelerated by about 200 kV, through a specimen a few tens of nanometres thick. Its images resolve columns of atoms 0.1 to 0.2 nm apart. The electron wavelength is a few picometres, so lens aberrations rather than diffraction set the resolution.

An electron of rest mass \(m_0\) and charge \(e\) accelerated through a potential difference \(U\) gains kinetic energy \(eU\). The relativistic relation \((pc)^2=(eU+m_0c^2)^2-(m_0c^2)^2\) fixes its momentum \(p\), and the de Broglie relation gives the wavelength:

\begin{centeredalign}
\lambda=\frac{h}{\sqrt{2m_0eU\left(1+\dfrac{eU}{2m_0c^2}\right)}},\qquad
\lambda(200\,{\rm kV})=2.51\,{\rm pm}.
\end{centeredalign}

At 200 kV the electron moves at 0.70 \(c\), and the nonrelativistic formula would give 2.74 pm, an error of 9 percent. At 300 kV the wavelength is 1.97 pm. Scattering angles from atomic spacings are then of order 10 mrad, so the electron optics works close to the axis.

A thin specimen barely absorbs the beam. It delays the electron wave by an amount proportional to the projected electrostatic potential \(V_p(\mathbf x)\). The transmitted wave is \(e^{i\sigma V_p}\approx1+i\sigma V_p\), where \(\sigma=2\pi me\lambda/h^2\) with the relativistic mass \(m\). At 200 kV, \(\sigma=7.29\times10^{-3}\) rad per volt-nanometre. An ideal lens would reproduce \(|1+i\sigma V_p|^2\approx1\), and the atoms would be invisible.

A real objective lens adds a phase \(\chi(k)\) to the wave scattered at spatial frequency \(k\), the scattering angle divided by \(\lambda\). Defocus \(\Delta f\) and spherical aberration \(C_s\) give the leading terms. Interference with the unscattered beam then turns the phase delay into intensity, with image Fourier transform

\begin{centeredalign}
\chi(k)=\pi\Delta f\,\lambda k^2+\frac{\pi}{2}C_s\lambda^3k^4,\qquad
\widetilde I(\mathbf k)=\delta(\mathbf k)+2\sigma\widetilde V_p(\mathbf k)\sin\chi(k).
\end{centeredalign}

The factor \(\sin\chi\) is the contrast transfer function in this sign convention. Scherzer chose the defocus \(\Delta f=-\sqrt{4C_s\lambda/3}\), so that \(\sin\chi\) stays close to \(-1\) over a broad band while \(\chi\) never falls below \(-2\pi/3\). The lowest zero of \(\sin\chi\) beyond that band gives the point resolution \(d\approx0.66\,C_s^{1/4}\lambda^{3/4}\). For \(C_s=1\) mm at 200 kV this gives \(\Delta f\approx-58\) nm and \(d\approx0.23\) nm. The diffraction limit for a 10 mrad aperture would be \(0.61\lambda/\alpha\approx0.15\) nm.

Aberration correctors made from multipole lenses reduce \(C_s\) and bring the resolution below 0.1 nm. A phase plate that shifts the unscattered beam by \(\pi/2\) converts \(\sin\chi\) into \(\cos\chi\), the electron analogue of Zernike phase contrast in light microscopy. The linear relation holds only for weak phase objects with \(\sigma V_p\ll1\). Thicker crystals scatter electrons several times and require dynamical calculations, and partial coherence of the beam damps the highest frequencies. For biological specimens radiation damage limits the dose, so the shot noise of the electron count limits contrast.

The instrument is a direct application of wave-particle duality: the same electrons that strike the detector one at a time form an interference pattern. The specimen potential enters as in the Born approximation of scattering theory, with single scattering from the projected potential.

\subsection*{References For The Physical Derivation}
\begin{itemize}
\item \href{https://arxiv.org/abs/1009.4038}{B. Gamm, K. Schultheiss, R. R. Schroder and D. Gerthsen, Effect of a Physical Phase Plate on Contrast Transfer in an Aberration-Corrected Transmission Electron Microscope; weak-phase-object contrast transfer, spherical aberration, defocus and phase plates.}
\end{itemize}





\part{Controlled And Spatial Realizations}

\chapter{Control sequences and quantum channels}

Sequences of gates, measurements and conditional operations turn quantum dynamics into computation, communication and sensing.

The order of operations is part of the physics: a circuit, a channel or an error-correcting cycle is defined by what acts and when.

\section*{Topics}

\begin{longtable}{p{0.92\linewidth}}
\toprule
Page \\
\midrule
\endhead
\phantomsection\label{page:quantum-computing}Quantum computing \\
\phantomsection\label{page:quantum-information-science}Quantum information science \emph{(see Quantum Information)} \\
\phantomsection\label{page:quantum-network}Quantum network \\
\phantomsection\label{page:quantum-algorithm}Quantum algorithm \\
\phantomsection\label{page:quantum-error-correction}Quantum error correction \\
\phantomsection\label{page:quantum-channel}Quantum channel \\
\phantomsection\label{page:quantum-logic-gate}Quantum logic gate \\
\phantomsection\label{page:quantum-neural-network}Quantum neural network \\
\phantomsection\label{page:quantum-circuit}Quantum circuit \\
\phantomsection\label{page:quantum-information}Quantum information \\
\phantomsection\label{page:quantum-key-distribution}Quantum key distribution \\
\phantomsection\label{page:quantum-sensor}Quantum sensor \\
\phantomsection\label{page:quantum-machine}Quantum machine \\
\phantomsection\label{page:quantum-teleportation}Quantum teleportation \\
\phantomsection\label{page:quantum-finite-automaton}Quantum finite automaton \\
\phantomsection\label{page:quantum-machine-learning}Quantum machine learning \\
\phantomsection\label{page:quantum-programming}Quantum programming \\
\phantomsection\label{page:quantum-simulator}Quantum simulator \\
\phantomsection\label{page:quantum-cryptography}Quantum cryptography \\
\phantomsection\label{page:quantum-metrology}Quantum metrology \\
\phantomsection\label{page:quantum-bus}Quantum bus \\
\phantomsection\label{page:quantum-image-processing}Quantum image processing \\
\phantomsection\label{page:quantum-engineering}Quantum engineering \\
\phantomsection\label{page:quantum-cellular-automaton}Quantum cellular automaton \\
\phantomsection\label{page:quantum-complexity-theory}Quantum complexity theory \\
\bottomrule
\end{longtable}

\clearpage

\section{Quantum computing}

\label{topic:quantum-computing}

A function \(f\) on \(n\)-bit strings returns \(0\) or \(1\) and is promised to be either constant or balanced, giving \(0\) on exactly half of its inputs. A deterministic classical procedure may need \(2^{n-1}+1\) evaluations to decide which. A quantum computer decides it with one evaluation, through interference among all \(2^n\) inputs. This Deutsch--Jozsa calculation shows where the advantage of quantum computing comes from, and where it does not.

The register holds \(n\) qubits and one ancilla. The function is accessed through the unitary \(U_f|x\rangle|y\rangle=|x\rangle|y\oplus f(x)\rangle\), where \(\oplus\) is addition modulo two. With the ancilla in \(|-\rangle\), phase kickback turns the call into the sign \((-1)^{f(x)}\). Hadamard gates on all \(n\) qubits before and after the call give

\begin{centeredalign}
|0\rangle^{\otimes n}\ \xrightarrow{H^{\otimes n}}\ \frac{1}{2^{n/2}}\sum_x|x\rangle
\ \xrightarrow{U_f}\ \frac{1}{2^{n/2}}\sum_x(-1)^{f(x)}|x\rangle
\ \xrightarrow{H^{\otimes n}}\ \sum_yc_y|y\rangle,
\end{centeredalign}

\begin{centeredalign}
c_y=\frac{1}{2^n}\sum_x(-1)^{f(x)+x\cdot y},\qquad c_{0}=\frac{1}{2^n}\sum_x(-1)^{f(x)}.
\end{centeredalign}

Here \(x\cdot y\) is the bitwise inner product modulo two. For a constant function \(c_0=\pm1\), so measuring all qubits returns \(00\ldots0\) with certainty. For a balanced function the terms cancel exactly, \(c_0=0\), and the all-zero outcome never occurs. The decisive step is the second layer of Hadamards, which makes the \(2^n\) phase-carrying paths interfere at a single output.

The state after the call contains all \(2^n\) function values as signs, but a measurement returns only \(n\) bits. Holevo's bound forbids extracting more than \(n\) classical bits from \(n\) qubits. A quantum computer therefore does not evaluate every input and report the results. It gains only when a global property of \(f\) can be turned into an interference pattern. The comparison also has a limit. A randomized classical procedure that checks \(k\) random inputs errs with probability at most \(2^{1-k}\), independent of \(n\). The exponential separation holds for exact computation only.

The circuit has been run with nuclear spins in liquid-state NMR. Gulde and colleagues (2003) ran it with a single trapped calcium ion, using an electronic transition as one qubit and a vibrational mode as the other. The same pattern, a superposition, a phase oracle and an interfering transform, underlies Simon's algorithm and the Fourier transform in Shor's factoring algorithm.

Several facts are established. Universal gate sets exist, as described under quantum circuit. Quantum error correction can protect stored information, and the threshold theorem shows that below a constant error rate per gate, arbitrarily long computations succeed with polylogarithmic overhead. Processors with hundreds of physical qubits operate, and logical qubits whose errors decrease with code size have been demonstrated. Open questions remain: building machines large enough to factor numbers of cryptographic size, and finding practical problems with proven speedups beyond the simulation of quantum systems.

\subsection*{References For The Physical Derivation}
\begin{itemize}
\item \href{https://arxiv.org/abs/quant-ph/9708016}{R. Cleve, A. Ekert, C. Macchiavello and M. Mosca, Quantum Algorithms Revisited; the Deutsch-Jozsa circuit with phase kickback and interference of computational paths.}
\item \href{https://arxiv.org/abs/1801.00862}{J. Preskill, Quantum Computing in the NISQ era and beyond; noisy processors, error correction and open questions about practical speedups.}
\end{itemize}





\clearpage
\section{Quantum information science}
\label{topic:quantum-information-science}
Quantum information science is an alternative name for Quantum Information.
The physical derivation is given in \hyperref[topic:quantum-information]{Quantum Information}.

\clearpage

\section{Quantum network}

\label{topic:quantum-network}

Telecom fiber at a wavelength of 1550 nm attenuates light by about 0.2 decibels per kilometre. One photon in a hundred survives 100 km, and one in \(10^{20}\) survives 1000 km. Entanglement between distant laboratories cannot simply be amplified on the way, because amplifying an unknown quantum state adds noise, as the quantum amplifier shows. Quantum networks instead combine heralded short links, quantum memories and entanglement swapping.

For attenuation \(\alpha\) in decibels per kilometre, the transmission over length \(L\) is \(\eta=10^{-\alpha L/10}=e^{-L/L_{\rm att}}\), with \(L_{\rm att}\approx22\) km for \(\alpha=0.2\). A source emitting pairs at rate \(R_0\) delivers \(\eta R_0\) of them. Pirandola and colleagues proved a stronger statement for a pure-loss channel. Without intermediate stations, no protocol distributes entanglement faster than

\begin{centeredalign}
R\leq-\log_2(1-\eta)\approx1.44\,\eta\quad\text{ebits per channel use}\qquad(\eta\ll1).
\end{centeredalign}

For a 10 GHz source and 1000 km this allows about \(10^{-10}\) pairs per second, roughly one pair every two centuries. Loss lowers the rate but need not lower the quality of delivered pairs. If the arrival of a photon is detected and announced, failed attempts are discarded. Imperfect sources, memory dephasing and dark counts lower instead the fidelity \(F=\langle\Phi^+|\rho|\Phi^+\rangle\) with the Bell state \(|\Phi^+\rangle=(|00\rangle+|11\rangle)/\sqrt2\). A common model is the Werner state \(\rho=F|\Phi^+\rangle\langle\Phi^+|+(1-F)(I-|\Phi^+\rangle\langle\Phi^+|)/3\).

Entanglement swapping joins two links. Node B holds one qubit entangled with A and another entangled with C. A Bell measurement on the two qubits at B projects A and C onto a Bell state labelled by the outcome, by the same expansion as in quantum teleportation. A Pauli correction then gives \(|\Phi^+\rangle\). For Werner links of fidelities \(F_1\) and \(F_2\), direct evaluation gives

\begin{centeredalign}
F_{AC}=F_1F_2+\frac{(1-F_1)(1-F_2)}{3},\qquad w_{AC}=w_1w_2,\qquad w=\frac{4F-1}{3}.
\end{centeredalign}

The parameter \(w\) multiplies along a chain. With \(F=0.95\) per link, eight links give \(F=0.68\) and sixteen give \(0.499\). A Werner state with \(F\leq1/2\) is separable, because its partial transpose is then positive. Long chains therefore need entanglement purification, which consumes several noisy pairs to produce fewer pairs of higher fidelity, much as quantum error correction uses redundancy.

Swapping changes the rate scaling when the links are stored. Divide \(L\) into \(N\) segments with quantum memories at the nodes. Each segment retries until its photon arrives, with success probability \(\eta^{1/N}\) per attempt, and a memory holds each success while the other segments retry. Swaps then join the segments, and the rate scales as

\begin{centeredalign}
R\propto\eta^{1/N},\qquad \eta^{1/8}\approx3\times10^{-3}\ \text{for 1000 km, against}\ \eta=10^{-20}.
\end{centeredalign}

Factors from waiting for the slowest segment and from the swap success probability multiply this scaling. The memories must stay coherent over many attempts, each lasting at least the signalling time across a segment, about 0.6 ms for 125 km. Bell measurements with linear optics and no auxiliary photons succeed at most half the time, which lowers the rate further.

Hensen and colleagues used this scheme with two electron spins in diamond, 1.3 km apart. Each spin emitted a photon entangled with it, and a Bell measurement on the two photons at a middle station heralded entanglement of the spins. In 245 heralded trials the spins violated the CHSH form of Bell's inequality, with \(S=2.42\pm0.20\) against the local bound of 2.

\subsection*{References For The Physical Derivation}
\begin{itemize}
\item \href{https://arxiv.org/abs/1510.08863}{S. Pirandola, R. Laurenza, C. Ottaviani and L. Banchi, Fundamental Limits of Repeaterless Quantum Communications; the rate bound for a pure-loss channel.}
\item \href{https://arxiv.org/abs/0906.2699}{N. Sangouard, C. Simon, H. de Riedmatten and N. Gisin, Quantum repeaters based on atomic ensembles and linear optics; heralding, memories, swapping and rate scaling.}
\item \href{https://arxiv.org/abs/1508.05949}{B. Hensen and colleagues, Experimental loophole-free violation of a Bell inequality using entangled electron spins separated by 1.3 km; event-ready entanglement of distant spins.}
\end{itemize}





\clearpage

\section{Quantum algorithm}

\label{topic:quantum-algorithm}

A list of \(N\) entries contains one marked entry, and the only access is a test that answers whether a given entry is marked. A classical search needs about \(N/2\) tests on average. Grover's algorithm finds the entry with about \((\pi/4)\sqrt N\) tests by rotating a quantum state, step by step, toward the answer.

Hadamard gates prepare a register of \(n=\log_2N\) qubits in the uniform superposition \(|s\rangle=N^{-1/2}\sum_x|x\rangle\). Let \(|w\rangle\) be the marked state and \(|r\rangle\) the normalized superposition of the other \(N-1\) states. Then \(|s\rangle=\sin\theta\,|w\rangle+\cos\theta\,|r\rangle\) with \(\sin\theta=1/\sqrt N\). The test acts as a phase oracle \(O\) that flips the sign of the marked amplitude, and a second operation \(D\) reflects about \(|s\rangle\):

\begin{centeredalign}
O=I-2|w\rangle\langle w|,\qquad D=2|s\rangle\langle s|-I,\qquad G=DO.
\end{centeredalign}

Both reflections map the plane spanned by \(|w\rangle\) and \(|r\rangle\) into itself. The oracle reflects about \(|r\rangle\) and \(D\) reflects about \(|s\rangle\). Two reflections about lines separated by the angle \(\theta\) combine into a rotation by \(2\theta\). After \(k\) iterations the state and the probability of finding the marked entry are

\begin{centeredalign}
G^k|s\rangle=\sin\big((2k+1)\theta\big)|w\rangle+\cos\big((2k+1)\theta\big)|r\rangle,\qquad
P_k=\sin^2\big((2k+1)\theta\big).
\end{centeredalign}

The probability is nearest one when \((2k+1)\theta\approx\pi/2\), after \(k\approx\pi/(4\theta)-1/2\approx(\pi/4)\sqrt N\) iterations. The oracle sign flip is realized with an ancilla in the state \(|-\rangle\), by the phase kickback described under quantum circuit. The reflection \(D\) equals \(H^{\otimes n}(2|0\rangle\langle0|-I)H^{\otimes n}\), a sign flip of the all-zero state between two layers of Hadamards.

For \(N=4\) the angle is \(\theta=30^\circ\). One iteration gives \((2k+1)\theta=90^\circ\), so the marked entry is found with certainty after a single query. A classical strategy with one query followed by a guess succeeds with probability one half. Chuang, Gershenfeld and Kubinec ran this two-qubit search with nuclear spins in 1998. For \(N=2^{20}\), 804 iterations give a success probability above \(0.9999997\), against about half a million classical tests on average.

The rotation does not stop at the answer. Continuing beyond the optimum turns the state past \(|w\rangle\); for \(N=16\), six iterations give a success probability of only \(0.02\), while three give \(0.96\). The number of marked entries must therefore be known, or estimated from the eigenphases \(\pm2\theta\) of \(G\) by phase estimation.

The same geometry applies to any preparation circuit \(\mathcal A\) that succeeds with probability \(a=\sin^2\theta\). Alternating a reflection that flags success with a reflection about \(\mathcal A|0\rangle\) raises the success probability close to one:

\begin{centeredalign}
k\approx\frac{\pi}{4\sqrt a}\ \text{rounds}\qquad\text{instead of}\qquad\sim\frac1a\ \text{classical repetitions}.
\end{centeredalign}

This amplitude amplification is a quadratic gain, and for unstructured search it cannot be improved, as shown under quantum complexity theory. In practice noise limits it. Every iteration needs an oracle circuit and an \(n\)-qubit reflection, and errors accumulated over about \(\sqrt N\) iterations dephase the coherent rotation.

\subsection*{References For The Physical Derivation}
\begin{itemize}
\item \href{https://arxiv.org/abs/quant-ph/9605043}{L. K. Grover, A fast quantum mechanical algorithm for database search; the search algorithm with about the square root of N queries.}
\item \href{https://arxiv.org/abs/quant-ph/0005055}{G. Brassard, P. Hoyer, M. Mosca and A. Tapp, Quantum Amplitude Amplification and Estimation; the rotation angle and its estimation for general preparation circuits.}
\end{itemize}





\clearpage

\section{Quantum error correction}

\label{topic:quantum-error-correction}

A quantum memory must distinguish errors without measuring the unknown logical amplitudes it is intended to preserve. The state is encoded into a subspace in which different correctable errors leave distinguishable information in auxiliary degrees of freedom, while their probabilities reveal no logical-state information.

Let \(P\) project onto the code subspace and let \(E_a\) be the error operators in a specified noise model. Exact correction of their linear span is possible when there is a matrix \(c\) such that

\begin{centeredalign}
P E_a^\dagger E_b P=c_{ab}P.
\end{centeredalign}

The right side is proportional to the identity within the code. Thus overlaps of the error-affected states depend on the error labels, but not on which logical superposition was stored. Diagonalizing \(c\) organizes the error spaces so that a measurement can identify the required correction without resolving the encoded state.

For a simple illustration, encode \(\alpha|0\rangle+\beta|1\rangle\) as \(\alpha|000\rangle+\beta|111\rangle\) and consider at most one bit flip. The commuting observables \(Z_1Z_2\) and \(Z_2Z_3\) give four distinct syndromes:

\begin{centeredalign}
\begin{array}{c|rrrr}
\text{error}&I&X_1&X_2&X_3\\\hline
(Z_1Z_2,Z_2Z_3)&(+,+)&(-,+)&(-,-)&(+,-)
\end{array}
\end{centeredalign}

Both basis codewords have the same syndrome before an error. The syndrome measurement therefore preserves their coherence. Applying the indicated bit flip restores the encoded state for any \(\alpha\) and \(\beta\). Measuring the individual \(Z_j\) instead would reveal which codeword was present and destroy an unknown superposition.

The noise specification matters. A phase flip acts within this code as a logical error and is not identified by those two syndromes. Protecting arbitrary single-qubit errors requires additional structure, such as a code that corrects both bit and phase errors. Increasing the number of spins without changing the encoded subspace and syndrome observables does not automatically supply that protection.

Error correction connects interaction, preparation, measurement and conditional dynamics in one mechanism. The environment determines the error operators; encoding determines the code subspace; syndrome extraction determines what information becomes observable; and the conditional operation restores the logical state. The algebraic correction condition is an exact test of this assembly. Fault tolerance adds another problem: the physical operations used to extract a syndrome can themselves propagate errors, so those operations must be analyzed as part of the noise model.

\subsection*{References For The Physical Derivation}
\begin{itemize}
\item \href{https://arxiv.org/abs/quant-ph/9705052}{D. Gottesman, Stabilizer Codes and Quantum Error Correction; error correction conditions and syndrome measurements.}
\end{itemize}

\subsection*{Relations In The Original Papers}

\href{https://arxiv.org/html/quant-ph/9705052#Ch2.E10}{arXiv:quant-ph/9705052, Ch2.E10}. Correctable errors have overlaps independent of the logical state inside the code space. A specified code and error set; the three-qubit example corrects single bit flips, not arbitrary single-qubit errors.

\href{https://arxiv.org/html/quant-ph/0302006#S3.E22}{arXiv:quant-ph/0302006, S3.E22}. A detected error can be reversed on the code when its norm does not distinguish logical codewords. Detected-emission feedback with a specified error record; it is narrower than correction of an arbitrary unobserved error set.



\clearpage

\section{Quantum channel}

\label{topic:quantum-channel}

An excited atom in free space emits a photon that flies away undetected. The state of the atom afterwards follows from the joint unitary evolution of atom and radiation field, after which the field is ignored. Let \(|g\rangle\) and \(|e\rangle\) be the ground and excited states, let the field start in its vacuum \(|0\rangle_E\), and let \(\Gamma\) be the spontaneous emission rate. During a time \(t\) the excitation leaves the atom with probability \(\gamma=1-e^{-\Gamma t}\):

\begin{centeredalign}
U|g\rangle|0\rangle_E=|g\rangle|0\rangle_E,\qquad
U|e\rangle|0\rangle_E=\sqrt{1-\gamma}\,|e\rangle|0\rangle_E+\sqrt{\gamma}\,|g\rangle|1\rangle_E.
\end{centeredalign}

Tracing out the field gives the density operator of the atom alone. Each orthogonal field state \(|k\rangle_E\) contributes one term, with the Kraus operator \(K_k={}_E\langle k|U|0\rangle_E\). Unitarity of \(U\) makes the Kraus operators complete, which is conservation of probability. Any later measurement with operators \(M_y\) then has outcome probabilities \(p(y)=\operatorname{Tr}[M_y\mathcal E(\rho)]\):

\begin{centeredalign}
\mathcal E(\rho)=\sum_k K_k\rho K_k^\dagger,\qquad \sum_k K_k^\dagger K_k=I,\\
K_0=|g\rangle\langle g|+\sqrt{1-\gamma}\,|e\rangle\langle e|,\qquad K_1=\sqrt{\gamma}\,|g\rangle\langle e|.
\end{centeredalign}

In the basis \((|g\rangle,|e\rangle)\) this amplitude-damping channel moves population downward and shrinks the coherence by a smaller factor:

\begin{centeredalign}
\mathcal E\begin{pmatrix}\rho_{gg}&\rho_{ge}\\ \rho_{eg}&\rho_{ee}\end{pmatrix}
=\begin{pmatrix}\rho_{gg}+\gamma\rho_{ee}&\sqrt{1-\gamma}\,\rho_{ge}\\
\sqrt{1-\gamma}\,\rho_{eg}&(1-\gamma)\rho_{ee}\end{pmatrix}.
\end{centeredalign}

The excited population decays as \(e^{-\Gamma t}\), while the coherence decays as \(\sqrt{1-\gamma}=e^{-\Gamma t/2}\). Energy relaxation alone therefore gives a coherence time \(T_2=2T_1\), with \(T_1=1/\Gamma\). Measured atomic and superconducting qubits obey \(T_2\leq2T_1\), and any shortfall comes from additional dephasing, such as slow fluctuations of the transition frequency. Two successive damping intervals compose as \(1-\gamma=(1-\gamma_1)(1-\gamma_2)\), consistent with the exponential law.

The Kraus operators are not unique. A random phase flip \(Z\) applied with probability \(p\) and an unrecorded measurement in the \(Z\) basis made with probability \(2p\) produce the same channel. Both multiply the coherence by \(1-2p\). Two Kraus sets of one channel are always related by a unitary mixing of their operators. Physically, this corresponds to measuring the discarded environment in a different basis.

A channel must remain physical when its input is entangled with an untouched partner, and positivity of \(\mathcal E(\rho)\) for single inputs does not ensure this. Matrix transposition \(T\) maps every density operator to a density operator. Applied to one half of the Bell state \(|\Phi^+\rangle=(|00\rangle+|11\rangle)/\sqrt2\), however, it gives

\begin{centeredalign}
(\mathrm{id}\otimes T)\big(|\Phi^+\rangle\langle\Phi^+|\big)=\tfrac12\,\mathrm{SWAP},
\qquad \text{eigenvalues } \tfrac12,\tfrac12,\tfrac12,-\tfrac12.
\end{centeredalign}

The negative eigenvalue would be a negative probability, so no physical process acts as a transposition. The test generalizes: a linear map is completely positive exactly when its Choi matrix \(\sum_{ij}|i\rangle\langle j|\otimes\mathcal E(|i\rangle\langle j|)\) is positive. For amplitude damping the eigenvalues are \(0\), \(0\), \(\gamma\) and \(2-\gamma\).

The description assumes that the atom and field start uncorrelated and that emitted light never returns. A mirror that reflects the photon back to the atom violates both conditions and makes the evolution depend on its history. Over a short interval, \(\gamma=\Gamma\,dt\) turns \(K_1\) into the jump operator of the quantum jump description. Quantum error correction asks which channels of this kind can be undone on an encoded subspace.

\subsection*{References For The Physical Derivation}
\begin{itemize}
\item \href{https://arxiv.org/abs/1106.1445}{M. M. Wilde, From Classical to Quantum Shannon Theory; Kraus and Stinespring representations, complete positivity, the Choi criterion and the amplitude-damping channel.}
\end{itemize}





\clearpage

\section{Quantum logic gate}

\label{topic:quantum-logic-gate}

In a trapped-ion processor, laser pulses rotate the internal two-level state of each ion, and a further pulse couples two ions through their shared motion. Every computation on the register must be assembled from these few physical operations. Two facts make this possible. Three rotations produce any single-qubit unitary, and one entangling gate between pairs completes the set.

A resonant pulse rotates the Bloch vector about an axis fixed by the drive phase, as in the two-state system. Let \(X,Y,Z\) be the Pauli matrices, \(R_y(\gamma)=e^{-i\gamma Y/2}\) and \(R_z(\beta)=e^{-i\beta Z/2}\). Every single-qubit unitary has the form

\begin{centeredalign}
U=e^{i\alpha}R_z(\beta)R_y(\gamma)R_z(\delta),\qquad H=e^{i\pi/2}R_y(\pi/2)R_z(\pi).
\end{centeredalign}

After a global phase is removed, \(U\) has determinant one. It is then the matrix with rows \((a,-b^*)\) and \((b,a^*)\), where \(|a|^2+|b|^2=1\). Setting \(|a|=\cos(\gamma/2)\) fixes \(\gamma\), and the phases of \(a\) and \(b\) fix \(\beta+\delta\) and \(\beta-\delta\). The Hadamard gate \(H\), which exchanges the \(Z\) and \(X\) bases, is shown as an example.

Single-qubit gates cannot create entanglement, so an interaction is required. An Ising coupling proportional to \(Z_1Z_2\) occurs naturally in spin and ion systems. The controlled-\(Z\) gate changes the sign of \(|11\rangle\) only. Since \(|11\rangle\langle11|=(I-Z_1)(I-Z_2)/4\) and commuting terms exponentiate separately,

\begin{centeredalign}
CZ=e^{i\pi/4}e^{-i\pi Z_1/4}e^{-i\pi Z_2/4}e^{i\pi Z_1Z_2/4},\qquad
\mathrm{CNOT}=(I\otimes H)\,CZ\,(I\otimes H).
\end{centeredalign}

The coupling, applied until its phase reaches \(\pi/4\), supplies the last factor, and local \(z\) rotations remove the rest. Hadamard gates on the target exchange \(Z\) and \(X\), converting the sign change into the bit flip of the CNOT gate.

Barenco and colleagues showed that CNOT together with all single-qubit gates is exactly universal. The key step builds any controlled single-qubit gate. Choose \(A=R_z(\beta)R_y(\gamma/2)\), \(B=R_y(-\gamma/2)R_z(-(\beta+\delta)/2)\) and \(C=R_z((\delta-\beta)/2)\). Then \(ABC=I\). Conjugation by \(X\) reverses both rotation angles, so \(XBX=R_y(\gamma/2)R_z((\beta+\delta)/2)\) and \(AXBXC=e^{-i\alpha}U\). With \(P_\alpha=\mathrm{diag}(1,e^{i\alpha})\) acting on the control,

\begin{centeredalign}
C\text{-}U=(P_\alpha\otimes A)\,\mathrm{CNOT}\,(I\otimes B)\,\mathrm{CNOT}\,(I\otimes C).
\end{centeredalign}

Operators act from right to left. If the control is 0, the target receives \(ABC=I\). If it is 1, the target receives \(e^{-i\alpha}U\), and the phase gate restores \(e^{i\alpha}\). For the controlled Hadamard, \(\alpha=\pi/2\), \(\beta=0\), \(\gamma=\pi/2\) and \(\delta=\pi\). Then \(A=R_y(\pi/4)\), \(B=R_y(-\pi/4)R_z(-\pi/2)\), \(C=R_z(\pi/2)\) and \(P_\alpha=\mathrm{diag}(1,i)\). Multiplying the five matrices reproduces the controlled Hadamard exactly. Toffoli gates and general multi-qubit unitaries follow by repeating such constructions.

Universality says nothing about efficiency. Each single-qubit gate has at most three angles and a CNOT has none. A generic \(n\)-qubit unitary has \(4^n-1\) real parameters and therefore needs at least \((4^n-1)/3\) single-qubit gates. Fault-tolerant hardware offers only a finite set, such as \(H\), \(T=\mathrm{diag}(1,e^{i\pi/4})\) and CNOT. The Solovay--Kitaev theorem guarantees that any single-qubit gate is approximated to accuracy \(\epsilon\) by a sequence from this set. The Dawson--Nielsen algorithm uses \(O(\log^{3.97}(1/\epsilon))\) gates. Errors of individual approximated gates add at most linearly:

\begin{centeredalign}
\|U_m\cdots U_1-V_m\cdots V_1\|\leq\sum_{k=1}^{m}\|U_k-V_k\|.
\end{centeredalign}

The bound follows from the triangle inequality and unitarity by replacing one factor at a time. A circuit of \(m\) gates therefore needs accuracy of about \(\epsilon/m\) per gate for a total error \(\epsilon\). The translation of whole programs into such gate sequences is treated under quantum programming.

\subsection*{References For The Physical Derivation}
\begin{itemize}
\item \href{https://arxiv.org/abs/quant-ph/9503016}{A. Barenco and colleagues, Elementary gates for quantum computation; universality of CNOT with single-qubit gates and controlled-U constructions.}
\item \href{https://arxiv.org/abs/quant-ph/0505030}{C. M. Dawson and M. A. Nielsen, The Solovay-Kitaev algorithm; approximation of single-qubit gates from a finite set with O(log\textasciicircum{}3.97(1/epsilon)) gates.}
\end{itemize}





\clearpage

\section{Quantum neural network}

\label{topic:quantum-neural-network}

A parametrized circuit used as a neural network takes an input number \(x\) as a rotation angle, applies trainable gates, and reports the average of a measured spin. Repeating the experiment for many inputs traces a curve \(f(x)\). The frequencies present in such curves are fixed by the eigenvalues of the generators that encode \(x\), much as level differences fix the lines of a spectrum.

Let the circuit alternate trainable unitaries \(W_\ell\) with \(L\) encoding gates \(S(x)=e^{-ixG/2}\), where the Hermitian generator \(G\) has eigenvalues \(\lambda_j\). With measured observable \(M\) and initial state \(|0\rangle\),

\begin{centeredalign}
f_\theta(x)=\langle0|U^\dagger MU|0\rangle,\qquad
U=W_L\,S(x)\,W_{L-1}\cdots W_1\,S(x)\,W_0.
\end{centeredalign}

In the eigenbasis of \(G\), each encoding multiplies an amplitude by \(e^{-ix\lambda_j/2}\). After \(L\) encodings each amplitude is a sum of terms carrying \(e^{-ix\Lambda/2}\), where \(\Lambda\) is a sum of \(L\) eigenvalues, one from each encoding. The average multiplies one amplitude by the complex conjugate of another, so only differences of such sums appear:

\begin{centeredalign}
f_\theta(x)=\sum_{\omega\in\Omega}c_\omega(\theta)\,e^{i\omega x},\qquad
\Omega=\Big\{\frac{\Lambda_k-\Lambda_j}{2}\Big\},\qquad c_{-\omega}=c_\omega^*.
\end{centeredalign}

For a Pauli generator, \(\lambda=\pm1\) and \(\Omega=\{-L,\ldots,L\}\). The model is then a trigonometric polynomial of degree \(L\) in \(x\). The trainable gates set only the coefficients \(c_\omega\), and adding trainable layers adds no frequencies (Schuld, Sweke and Meyer). Encoding \(x\) once on each of \(n\) qubits in parallel also gives degree \(n\).

A one-qubit example makes this concrete. Apply \(R_y(\theta)\) to \(|0\rangle\), encode with \(R_x(x)\), apply \(R_y(\phi)\) and measure \(Z\). Following the Bloch vector through the three rotations gives the first expression below. The circuit \(R_x(x)R_y(\phi)R_x(x)\) acting on \(|0\rangle\), with the input encoded twice, gives the second:

\begin{centeredalign}
L=1:\ f=\cos\theta\cos\phi\cos x-\sin\theta\sin\phi,\\
L=2:\ f=\frac{\cos\phi-1}{2}+\frac{\cos\phi+1}{2}\cos2x.
\end{centeredalign}

The single-encoding circuit cannot represent \(\cos2x\) for any angles. Its best mean-square approximation within degree one is zero, because \(\cos2x\) is orthogonal to \(1\), \(\cos x\) and \(\sin x\). The two-encoding circuit produces it exactly at \(\phi=0\). Repeated encoding of the data is how such models gain frequencies. Their coefficients are trained by the parameter-shift rule of quantum machine learning, and their gradients face the same barren plateaus.

For a single encoding layer and a fixed initial layer \(W_0\), the model is \(f_\theta(x)=\operatorname{Tr}[\rho(x)M_\theta]\). Here \(\rho(x)=S(x)W_0|0\rangle\langle0|W_0^\dagger S(x)^\dagger\) is the encoded state and \(M_\theta=W_1^\dagger MW_1\). The model is linear in \(\rho(x)\), and the dependence on \(x\) enters only through the encoding. Schuld showed that supervised models of this kind are kernel methods, in which data enter through overlaps of encoded states. The name neural network refers to the layered trainable structure, not to neuron-like nonlinear units.

Established are the Fourier structure, the role of the encoding and the kernel form. Liu, Arunachalam and Temme proved a speed-up for one classification problem constructed from the discrete logarithm, assuming that this logarithm is classically hard to compute. That result requires a fault-tolerant computer and specially built data. Whether such models outperform classical networks on practical data is open.

\subsection*{References For The Physical Derivation}
\begin{itemize}
\item \href{https://arxiv.org/abs/2008.08605}{M. Schuld, R. Sweke and J. J. Meyer, The effect of data encoding on the expressive power of variational quantum machine learning models; Fourier spectrum fixed by the encoding gates.}
\item \href{https://arxiv.org/abs/2101.11020}{M. Schuld, Supervised quantum machine learning models are kernel methods; linearity in the encoded state and the kernel form.}
\item \href{https://arxiv.org/abs/2010.02174}{Y. Liu, S. Arunachalam and K. Temme, A rigorous and robust quantum speed-up in supervised machine learning; separation for a classification problem built on the discrete logarithm.}
\end{itemize}





\clearpage

\section{Quantum circuit}

\label{topic:quantum-circuit}

Two laboratory operations on a pair of qubits, a Hadamard pulse on one and a controlled-NOT between them, give different results in the two possible orders. In one order the qubits end entangled, with perfectly correlated outcomes; in the other they remain uncorrelated. A quantum circuit specifies such orders. Each gate is a unitary \(U_j\) acting on a few qubits, and the register evolves by the time-ordered product:

\begin{centeredalign}
\rho_{\rm out}=U_m\cdots U_2U_1\,\rho_{\rm in}\,U_1^\dagger U_2^\dagger\cdots U_m^\dagger,
\qquad p(y)=\operatorname{Tr}(M_y\rho_{\rm out}).
\end{centeredalign}

Here \(M_y\) is the measurement operator for outcome \(y\), usually a projector on a computational basis state. The Hadamard gate is \(H=(X+Z)/\sqrt2\), and the controlled-NOT flips the target qubit when the control qubit is in \(|1\rangle\). Starting from \(|00\rangle\), the two orders give

\begin{centeredalign}
\mathrm{CNOT}\,(H\otimes I)|00\rangle=\frac{|00\rangle+|11\rangle}{\sqrt2},\qquad
(H\otimes I)\,\mathrm{CNOT}|00\rangle=\frac{|00\rangle+|10\rangle}{\sqrt2}.
\end{centeredalign}

In the left state, \(Z\) measurements on the two qubits always agree. In the right state, the target qubit is always found in \(0\) and the control qubit is random. The circuit diagram is therefore a record of the time order of physical couplings, not a list of available operations.

The roles of control and target depend on the basis in which a gate is viewed. Let \(|-\rangle=(|0\rangle-|1\rangle)/\sqrt2\), an eigenstate of \(X\) with eigenvalue \(-1\). Flipping the target then leaves it unchanged except for a sign, which appears as a phase of the control:

\begin{centeredalign}
\mathrm{CNOT}\,|x\rangle|-\rangle=(-1)^x|x\rangle|-\rangle,\qquad
(H\otimes H)\,\mathrm{CNOT}_{12}\,(H\otimes H)=\mathrm{CNOT}_{21}.
\end{centeredalign}

The reversal of control and target follows from the kickback, because Hadamards exchange the \(X\) and \(Z\) bases on both qubits. This phase kickback is the step by which an oracle writes a function value into a phase, as in the Deutsch--Jozsa circuit described under quantum computing.

Single-qubit rotations together with the controlled-NOT suffice to compose any unitary on \(n\) qubits (Barenco and colleagues, 1995). Universality says nothing about cost. A general unitary on \(n\) qubits has \(4^n-1\) real parameters, while each gate supplies a bounded number. Most unitaries therefore require a number of gates that grows exponentially with \(n\). Efficient algorithms are the special unitaries with circuits of polynomial size.

The circuit description assumes that each gate acts only on its intended qubits and that the qubits do not leak into other levels. Noise is then represented by a quantum channel after each gate. If errors are independent with probability \(\epsilon\) per gate, a circuit of \(n_g\) gates runs without any error with probability

\begin{centeredalign}
P_{\rm ok}=(1-\epsilon)^{n_g}\approx e^{-n_g\epsilon}.
\end{centeredalign}

With \(\epsilon=0.005\) per two-qubit gate, a circuit of 200 such gates runs without error with probability \(0.37\). This fixes the depth that can be reached without quantum error correction. Error correction replaces the exponential decay by repeated syndrome measurements and corrections, which are themselves further gates of the circuit.

\subsection*{References For The Physical Derivation}
\begin{itemize}
\item \href{https://arxiv.org/abs/quant-ph/9503016}{A. Barenco, C. H. Bennett, R. Cleve, D. P. DiVincenzo, N. Margolus, P. Shor and colleagues, Elementary gates for quantum computation; universality of one-qubit gates with the controlled-NOT and gate counts.}
\item \href{https://arxiv.org/abs/quant-ph/9708016}{R. Cleve, A. Ekert, C. Macchiavello and M. Mosca, Quantum Algorithms Revisited; phase kickback and circuits as multi-path interference.}
\end{itemize}





\clearpage

\section{Quantum information}

\label{topic:quantum-information}

A sender encodes bits in single photons, horizontal polarization for 0 and diagonal polarization for 1. The two states overlap, so no detector distinguishes them with certainty. The physical question is how many bits per photon a receiver can extract, measuring the photons one at a time or jointly.

The uncertainty of a preparation is measured by the von Neumann entropy of its density operator \(\rho\). For a qubit with Bloch vector \(\mathbf r\), as in the treatment of the quantum state, the eigenvalues are \((1\pm|\mathbf r|)/2\). Hence

\begin{centeredalign}
S(\rho)=-\operatorname{Tr}\rho\log_2\rho,\qquad S=h_2\Big(\frac{1+|\mathbf r|}{2}\Big),\\
h_2(q)=-q\log_2q-(1-q)\log_2(1-q).
\end{centeredalign}

A pure state has zero entropy and the maximally mixed qubit has one bit. Schumacher compression gives the entropy an operational meaning. Long sequences from a source of pure states with average \(\rho\) can be stored in \(S(\rho)\) qubits per signal.

Let letter \(x\) occur with probability \(p_x\) and be sent as the state \(\rho_x\). The receiver's system is then described by \(\rho=\sum_xp_x\rho_x\), and a measurement yields an outcome \(Y\) correlated with \(X\). Holevo's bound limits their classical mutual information:

\begin{centeredalign}
I(X{:}Y)\leq\chi=S(\rho)-\sum_xp_xS(\rho_x)\leq\log_2d.
\end{centeredalign}

The argument has two steps. Record the letter in a classical register, giving the joint state \(\sum_xp_x|x\rangle\langle x|\otimes\rho_x\). Evaluating the entropies of this block-diagonal state shows that its quantum mutual information \(S(X)+S(Q)-S(XQ)\) equals \(\chi\). A measurement is a physical operation on the system alone, and mutual information cannot increase under local operations. This data-processing inequality follows from strong subadditivity of the entropy. After the measurement the quantity becomes the classical \(I(X{:}Y)\). The second inequality uses \(S(\rho)\leq\log_2d\) for a \(d\)-dimensional system and \(S(\rho_x)\geq0\). A qubit therefore conveys at most one bit, although its pure states form a continuous sphere.

For the two polarizations sent with equal probability, the overlap is \(s=|\langle0|+\rangle|=1/\sqrt2\). The average state has eigenvalues \((1\pm s)/2\), and the pure signals have zero entropy. For one photon at a time, the best measurement is the symmetric projective one that minimizes the error, with error probability \(P_e=[1-(1-s^2)^{1/2}]/2\). It turns each photon into a binary symmetric channel:

\begin{centeredalign}
\chi=h_2\Big(\frac{1+s}{2}\Big)=0.601,\qquad
I_{\rm single}=1-h_2(P_e)=0.399\ \text{bits}.
\end{centeredalign}

A numerical search over all projective measurements and random three-outcome measurements finds no single-photon value above 0.399 bits. The gap to 0.601 bits is physical. Long code words built from the two states, measured jointly, reach any rate below \(\chi\) with vanishing error; this is the Holevo--Schumacher--Westmoreland theorem. The additional 0.2 bits per photon therefore require measurements that entangle many photons, not better single-photon detectors.

The bound concerns the system that actually reaches the receiver. In superdense coding one transmitted qubit carries two bits. The receiver, however, already holds the other half of an entangled pair, and the bound for the two-qubit system is \(\log_24=2\). Quantum teleportation performs the reverse trade, two classical bits and a shared pair for one qubit. In quantum key distribution the same quantity \(\chi\), evaluated for an eavesdropper's probe, bounds what that probe reveals about the key.

\subsection*{References For The Physical Derivation}
\begin{itemize}
\item \href{https://arxiv.org/abs/1604.07450}{J. Preskill, Quantum Shannon Theory; von Neumann entropy, compression, accessible information and the Holevo bound.}
\item \href{https://arxiv.org/abs/1106.1445}{M. M. Wilde, From Classical to Quantum Shannon Theory; data processing, the Holevo bound and the Holevo--Schumacher--Westmoreland theorem.}
\end{itemize}





\clearpage

\section{Quantum key distribution}

\label{topic:quantum-key-distribution}

Quantum key distribution creates correlated classical data whose secrecy is inferred from a specified quantum communication model. Its physical basis is that information about nonorthogonal signal states cannot be acquired perfectly while leaving those states unchanged. The task is key generation, not the transmission of an encrypted message or the authentication of an unknown partner.

In ideal BB84 the sender chooses between the eigenstates of \(Z\) and \(X\). Write \(|\pm\rangle=(|0\rangle\pm|1\rangle)/\sqrt2\). Suppose an eavesdropper's unitary interaction preserves both \(|0\rangle\) and \(|1\rangle\) while recording them in probe states \(|e_0\rangle\) and \(|e_1\rangle\). Linearity then requires

\begin{centeredalign}
|+\rangle|e\rangle\longmapsto
\frac{|0\rangle|e_0\rangle+|1\rangle|e_1\rangle}{\sqrt2}.
\end{centeredalign}

The signal remains a pure \(|+\rangle\) only when the two probe states coincide up to the phase compatible with the signal. Distinguishable probe records suppress its off-diagonal coherence. A check in the complementary basis can therefore detect the disturbance associated with that information gain.

After transmission, the legitimate parties disclose their basis choices, retain the compatible outcomes and estimate an error rate from a random sample. Classical error correction reconciles their strings and privacy amplification reduces the information that may remain with an adversary. For the ideal asymptotic single-photon BB84 setting with the relevant bit and phase error rates both bounded by \(Q\), the one-way secret fraction per sifted bit has the familiar form

\begin{centeredalign}
r\geq1-2h_2(Q),\qquad
h_2(Q)=-Q\log_2Q-(1-Q)\log_2(1-Q).
\end{centeredalign}

This expression is tied to its security model. Finite data require statistical confidence terms; multiphoton emission, detector imperfections and information disclosed during reconciliation require explicit treatment. An observed low bit-error rate alone does not establish secrecy for an arbitrary device. Loss and source statistics can carry information that a two-level idealization omits.

The classical channel must be authenticated. Otherwise an adversary can impersonate each party in a separate quantum exchange. The quantum mechanism supplies a relation between information and disturbance under stated source and measurement assumptions; authentication supplies the identity of the parties using that relation. Keeping these contributions distinct is necessary to transfer a security argument from an ideal qubit model to an optical implementation.

\subsection*{References For The Physical Derivation}
\begin{itemize}
\item \href{https://arxiv.org/abs/0802.4155}{V. Scarani and colleagues, The Security of Practical Quantum Key Distribution; ideal protocols, device assumptions and security bounds.}
\end{itemize}

\subsection*{Relations In The Original Papers}

\href{https://arxiv.org/html/0802.4155#S4.EGx43}{arXiv:0802.4155, S4.EGx43}. For the specified entanglement-based protocol, the asymptotic key rate subtracts phase-error information and error-correction leakage. Ideal single-photon assumptions and asymptotic rates; setting leakage to binary entropy gives the bound discussed in the chapter.



\clearpage

\section{Quantum sensor}

\label{topic:quantum-sensor}

A nitrogen-vacancy centre in diamond, a trapped ion or an atom in a vapour cell detects a magnetic field through the precession of its spin. A quantum sensor is a probe with well-separated levels whose energy difference depends on the quantity to be measured. The signal accumulates as a relative phase between two states, which is then converted into a population that a detector counts.

In a Ramsey sequence the probe starts in \(|0\rangle\), and a \(\pi/2\) pulse creates an equal superposition of \(|0\rangle\) and \(|1\rangle\). During a free evolution time \(\tau\), a detuning \(\delta\) between the transition and a reference oscillator adds the relative phase \(\delta\tau\). A second \(\pi/2\) pulse converts this phase into the probability of finding the excited state. For \(\nu\) independent probes or repetitions, binomial counting adds projection noise:

\begin{centeredalign}
P_1=\frac{1+\cos\delta\tau}{2},\qquad \Delta P_1=\sqrt{\frac{P_1(1-P_1)}{\nu}},\qquad
\Delta\delta=\frac{\Delta P_1}{|\partial P_1/\partial\delta|}=\frac{1}{\tau\sqrt\nu}.
\end{centeredalign}

The frequency uncertainty is independent of the operating point and defines the standard quantum limit for uncorrelated probes. A longer phase accumulation improves it until the coherence time \(T_2\) limits \(\tau\). Correlated probe states can exceed this limit, as described under quantum metrology.

A different sensor records a weak field in the choice between two macroscopic states. A Kerr parametric oscillator, such as a Josephson microwave resonator pumped at twice its frequency, has two stable oscillation states of opposite phase above a threshold. Swept through the threshold, it amplifies vacuum fluctuations into one of the two states with equal probability. A resonant bias \(h\) far below one photon tilts that choice. In the frame rotating at half the pump frequency, with \(\hbar=1\),

\begin{centeredalign}
H=-Ka^{\dagger2}a^2+\frac{i\epsilon_2(t)}{2}\big(a^{\dagger2}-a^2\big)+ih\big(a^\dagger-a\big),\\
\dot\rho=-i[H,\rho]+\kappa(\bar n+1)\mathcal D[a]\rho+\kappa\bar n\,\mathcal D[a^\dagger]\rho.
\end{centeredalign}

Here \(K\) is the Kerr coefficient, \(\epsilon_2=rt\) the two-photon drive swept at rate \(r\), \(\kappa\) the loss rate, \(\bar n\) the thermal occupation of the bath, and \(\mathcal D[a]\rho=a\rho a^\dagger-\{a^\dagger a,\rho\}/2\). Above threshold, \(\epsilon_2>\kappa/2\), each stored state holds \(\sqrt{\epsilon_2^2-\kappa^2/4}/(2K)\) photons. While the amplitude is small the Kerr term is negligible. A Gaussian Wigner function evolving under a quadratic Hamiltonian with linear loss obeys exactly the classical Fokker--Planck equation. The amplified quadrature \(x=\mathrm{Re}\,\alpha\) therefore follows a classical noisy equation:

\begin{centeredalign}
\dot x=\big(\epsilon_2-\tfrac{\kappa}{2}\big)x+h+\sqrt{2D}\,\xi(t),\qquad 2D=\frac{\kappa(2\bar n+1)}{4},\\
P=\Phi\!\left(\frac{hI_1}{\sqrt{\sigma_0^2+2DI_2}}\right),\qquad I_k=\int_0^\infty e^{-k\phi(s)}ds,\qquad
\phi(s)=\frac{rs^2-\kappa s}{2}.
\end{centeredalign}

Here \(\xi\) is unit white noise, \(\sigma_0^2=(2\bar n_0+1)/4\) the initial variance of \(x\), and \(\Phi\) the normal distribution function. The bias shifts the mean of the amplified seed, the noise sets its width, and \(P\) is the probability of the favoured state. The noise strength \(D\) is not a free parameter; vacuum and thermal fluctuations fix it. Evaluating the integrals gives two limits:

\begin{centeredalign}
\kappa\gg\sqrt r:\ P=\Phi\!\left(\frac{\pi^{1/4}h}{D^{1/2}r^{1/4}}\right),\qquad
\kappa\to0:\ P=\Phi\!\left(h\sqrt{2\pi/r}\right).
\end{centeredalign}

The strong-loss limit is the selection law of a classical swept pitchfork, derived by Kondepudi and Nelson in 1985, here with the noise fixed by the vacuum. Loss renews the seed during the sweep, and the rate enters with the exponent \(-1/4\). Without loss the seed is the initial vacuum, and the exponent becomes \(-1/2\). A higher temperature lowers the argument of \(\Phi\) as \((2\bar n+1)^{-1/2}\). Gu and colleagues derived the corresponding relation for a sudden step of the gain. Roques-Carmes and colleagues measured it in an optical parametric oscillator, where bias pulses below one photon set the two probabilities.

In units where \(\kappa=1\), with \(r=0.01\) and \(h=0.108\), the law predicts \(P=0.90\). Exact density-operator evolution at 27 stored photons gives \(0.900\), within \(0.001\) of the law. The law requires the choice to be complete before the Kerr term acts, which fails at few photons. With \(K=\kappa\) and 1.5 stored photons the same sweep gives \(0.68\), because the oscillator switches between its shallow wells while the sweep proceeds. The probability then relaxes toward the selection \(P_{\rm eq}\) of the biased steady state, at the slowest decay rate \(\Gamma\) of the Lindblad equation:

\begin{centeredalign}
\frac{dp}{dt}=\Gamma(\epsilon_2)\big[P_{\rm eq}(\epsilon_2)-p\big].
\end{centeredalign}

This one-mode relaxation reproduces the exact few-photon probability within \(0.01\) over two decades of sweep rate. The rate dependence separates the regimes: a probability that rises for slower sweeps comes from the Gaussian stage, and a rate-independent probability from equilibrating wells. Yamaji and colleagues (2025) measured the probability of the favoured state against a one-photon bias in a Josephson oscillator holding about twelve photons. A related selection under a residual bias, studied against the crossing rate, occurs at the Dicke transition of a Bose--Einstein condensate in an optical cavity.

\subsection*{References For The Physical Derivation}
\begin{itemize}
\item \href{https://arxiv.org/abs/1611.02427}{C. L. Degen, F. Reinhard and P. Cappellaro, Quantum sensing; Ramsey sequences, projection noise and the standard quantum limit.}
\item \href{https://arxiv.org/abs/2412.02887}{A. Gu, J. Sloan, C. Roques-Carmes and colleagues, Quantum sensitivity of parametric oscillators; the bias-probability relation for a step of the gain with the initial quantum state and the gain noise.}
\item \href{https://arxiv.org/abs/2303.03455}{C. Roques-Carmes, Y. Salamin, J. Sloan and colleagues, Biasing the quantum vacuum to control macroscopic probability distributions; bias pulses below one photon set the probabilities of the two states of an optical parametric oscillator.}
\item \href{https://arxiv.org/abs/2506.23539}{T. Yamaji, S. Masuda, Y. Kano and colleagues, Quantum annealing in capacitively coupled Kerr parametric oscillators using frequency-chirped drives; Josephson Kerr parametric oscillators raised through threshold with one-photon drives, probability of the selected state.}
\item \href{https://arxiv.org/abs/1105.0426}{K. Baumann, R. Mottl, F. Brennecke and T. Esslinger, Exploring Symmetry Breaking at the Dicke Quantum Phase Transition; selection between two ordered phases under a small symmetry-breaking field, against the rate of crossing.}
\end{itemize}





\clearpage

\section{Quantum machine}

\label{topic:quantum-machine}

A single spin in a magnetic field can run a heat engine. Its two levels are separated by \(\hbar\omega\), and the field sets \(\omega\). Contact with a hot bath fixes one level population, contact with a cold bath another. Changing the field between the contacts exchanges energy with the field. The efficiency of this quantum Otto cycle follows from the populations alone.

Let \(H(\omega)=\hbar\omega\sigma_z/2\) and let \(p\) be the probability of the upper level, so the mean energy is \(\hbar\omega(p-1/2)\). A change of \(p\) at fixed \(\omega\) is heat, and a change of \(\omega\) at fixed \(p\) is work. Thermal contact at temperature \(T\) gives the two-level occupation of quantum statistical mechanics:

\begin{centeredalign}
p(\omega,T)=\frac{1}{e^{\hbar\omega/k_BT}+1}.
\end{centeredalign}

The cycle has four strokes. Contact with the hot bath at splitting \(\omega_h\) raises the population from \(p_c\) to \(p_h=p(\omega_h,T_h)\). The spin is then isolated and the field lowered to \(\omega_c\). Contact with the cold bath returns the population to \(p_c=p(\omega_c,T_c)\). Finally the isolated spin is returned to \(\omega_h\). Adding the energy changes of the four strokes gives

\begin{centeredalign}
Q_h=\hbar\omega_h(p_h-p_c),\qquad W=\hbar(\omega_h-\omega_c)(p_h-p_c),\\
\eta=\frac{W}{Q_h}=1-\frac{\omega_c}{\omega_h}.
\end{centeredalign}

The engine delivers work only if \(p_h>p_c\), which means \(\hbar\omega_h/k_BT_h<\hbar\omega_c/k_BT_c\). This is the condition \(\omega_c/\omega_h>T_c/T_h\), so the efficiency stays below the Carnot value \(1-T_c/T_h\). It reaches that value only where \(p_h=p_c\) and the work per cycle vanishes. Run backwards with \(\omega_c/\omega_h<T_c/T_h\), the cycle pumps heat out of the cold bath. Its coefficient of performance \(\omega_c/(\omega_h-\omega_c)\) again stays below the Carnot bound \(T_c/(T_h-T_c)\).

\begin{centeredalign}
\eta=1-\frac{\omega_c}{\omega_h}<1-\frac{T_c}{T_h},\qquad
\frac{\omega_c}{\omega_h-\omega_c}<\frac{T_c}{T_h-T_c}.
\end{centeredalign}

As an example take \(\hbar\omega_h=k_BT_h\), \(\omega_c=\omega_h/2\) and \(T_c=T_h/4\). Then \(p_h=1/(e+1)=0.269\) and \(p_c=1/(e^2+1)=0.119\). Each cycle delivers \(W=0.075\,k_BT_h\) at efficiency 0.50, while the Carnot value for these baths is 0.75. Peterson and colleagues ran such a cycle on a nuclear spin with magnetic-resonance techniques. At maximum power they measured an efficiency near 42 percent, close to the ideal value of 44 percent for their cycle.

The derivation assumes unchanged populations during the isolated strokes. For \(H\propto\sigma_z\) the Hamiltonians at different times commute, so the populations are conserved at any speed. If the field also turns, the eigenbasis rotates, and a fast stroke creates coherence between the instantaneous levels. That coherence carries extra energy, which the next bath contact dissipates. This loss, called quantum friction, lowers the work, and slow strokes avoid it at the cost of power. Incomplete thermalization in a contact reduces the population difference. It lowers the work per cycle but leaves the ratio \(1-\omega_c/\omega_h\) unchanged.

A harmonic oscillator whose frequency is switched between \(\omega_h\) and \(\omega_c\) has the same Otto efficiency, because every level scales in proportion to \(\omega\). A sudden frequency change, however, does alter its occupations. The strokes must then be slow compared with the oscillation period to avoid friction. This oscillator version has been realized with the motion of a single trapped ion.

\subsection*{References For The Physical Derivation}
\begin{itemize}
\item \href{https://arxiv.org/abs/quant-ph/0611275}{H. T. Quan, Y.-x. Liu, C. P. Sun and F. Nori, Quantum Thermodynamic Cycles and quantum heat engines; heat and work in quantum Otto and Carnot cycles.}
\item \href{https://arxiv.org/abs/1612.03582}{R. Kosloff and Y. Rezek, The quantum harmonic Otto cycle; oscillator Otto engine, friction in finite-time strokes and the single-ion realization.}
\item \href{https://arxiv.org/abs/1803.06021}{J. P. S. Peterson and colleagues, Experimental characterization of a spin quantum heat engine; NMR Otto cycle with efficiency near 42 percent against a limit of 44 percent.}
\end{itemize}





\clearpage

\section{Quantum teleportation}

\label{topic:quantum-teleportation}

An unknown qubit can be transferred using a shared entangled pair and two classical bits. The entangled resource alone does not transmit a usable signal: a joint measurement at the sender and a correction conditioned on its outcome are essential parts of the mechanism.

Let the input be \(|\psi\rangle=a|0\rangle+b|1\rangle\) and let the shared pair be \(|\Phi^+\rangle=(|00\rangle+|11\rangle)/\sqrt2\), with the second member held by the receiver. Define the Bell states \(|B_{mn}\rangle=(I\otimes X^nZ^m)|\Phi^+\rangle\), where \(m,n\in\{0,1\}\). Expansion in this joint basis gives

\begin{centeredalign}
|\psi\rangle_1|\Phi^+\rangle_{23}
=\frac12\sum_{m,n=0}^{1}|B_{mn}\rangle_{12}\,X^nZ^m|\psi\rangle_3.
\end{centeredalign}

Each Bell outcome has probability one quarter, independent of the input amplitudes. Once the sender communicates \(m,n\), the receiver applies \(Z^mX^n\) and obtains the original state. The measurement determines a known transformation of the state, rather than its unknown amplitudes. This is why transmitting two classical bits suffices when the entangled pair has already been supplied.

Without the outcome record the receiver averages the four possible transformed states. For any input density operator \(\rho\), that average is

\begin{centeredalign}
\frac14\sum_{m,n}X^nZ^m\rho Z^mX^n=\frac I2.
\end{centeredalign}

The receiver's local statistics then contain no dependence on the input. The same calculation establishes that the protocol cannot be used for faster-than-light signalling. The original input has participated in a destructive joint measurement, so the transfer does not create two independently available copies.

A nonideal entangled resource changes the resulting channel. For a Bell-diagonal resource, its Bell weights become probabilities of Pauli errors in the corrected output. Resource quality, Bell-measurement fidelity and the conditional correction can therefore be distinguished experimentally; all three can reduce the final state fidelity, but they enter the calculation at different steps.

Teleportation illustrates how a mechanism can be transferred between physical systems. Photonic polarization, internal atomic levels and superconducting circuits may realize the same qubit relation, provided their entangling resource, joint measurement and correction implement the required maps. Changing only the name of the physical system leaves those requirements unresolved; demonstrating the maps establishes the actual correspondence.

\subsection*{References For The Physical Derivation}
\begin{itemize}
\item \href{https://arxiv.org/abs/quant-ph/9705052}{D. Gottesman, Stabilizer Codes and Quantum Error Correction; entanglement, measurements and teleportation.}
\item \href{https://arxiv.org/abs/1110.3234}{C. Weedbrook and colleagues, Gaussian Quantum Information; continuous-variable realizations and finite-resource effects.}
\end{itemize}

\subsection*{Relations In The Original Papers}

\href{https://arxiv.org/html/1110.3234v1#S4.E101}{arXiv:1110.3234v1, S4.E101}. Teleportation of coherent states with a finite two-mode squeezed resource has fidelity below one. A continuous-variable comparison to the exact ideal qubit protocol; it is not the derivation of the qubit Bell-state identity.



\clearpage

\section{Quantum finite automaton}

\label{topic:quantum-finite-automaton}

A photon passes through a row of identical polarization rotators, each turning its linear polarization by the same angle. A polarizer at the end transmits or absorbs it. The transmission probability depends only on the number of rotators. A quantum finite automaton uses this arrangement to decide a property of an input word. Each letter applies a fixed unitary to a small quantum system, and one measurement after the last letter gives accept or reject.

In the measure-once model of Moore and Crutchfield the system has a finite orthonormal basis. Letter \(\sigma\) acts by a unitary \(U_\sigma\), the initial state is \(|\psi_0\rangle\), and \(P_{\rm acc}\) projects onto the accepting states. A word \(w=\sigma_1\sigma_2\cdots\sigma_k\) is accepted with probability

\begin{centeredalign}
P_{\rm acc}(w)=\big\|P_{\rm acc}\,U_{\sigma_k}\cdots U_{\sigma_2}U_{\sigma_1}|\psi_0\rangle\big\|^2.
\end{centeredalign}

The measure-many model of Kondacs and Watrous instead measures after every letter. Consider words \(a^k\) over a single letter and the language \(L_p\) of words whose length is divisible by \(p\). Two basis states suffice. Let \(U_a\) rotate the real plane spanned by \(|0\rangle\) and \(|1\rangle\) through the angle \(2\pi/p\), start in \(|0\rangle\) and accept on \(|0\rangle\):

\begin{centeredalign}
U_a^k|0\rangle=\cos\frac{2\pi k}{p}|0\rangle+\sin\frac{2\pi k}{p}|1\rangle,\qquad
P_{\rm acc}(a^k)=\cos^2\frac{2\pi k}{p}.
\end{centeredalign}

Every word in \(L_p\) returns the state to its start and is accepted with certainty. Other words are accepted with probability below one. For odd \(p\), however, the largest such value is \(\cos^2(\pi/p)\), reached at \(k=(p\pm1)/2\). For \(p=7\) it equals 0.81, so a single run separates the two classes poorly.

The separation improves with a superposition of \(d\) rotation planes. Plane \(j\) is rotated through \(2\pi k_j/p\) per letter, with integer multipliers \(k_j\). The initial state is the equal superposition of the starting vectors of all planes, and acceptance projects back onto it. The acceptance probability is then a squared return amplitude:

\begin{centeredalign}
P_{\rm acc}(a^n)=\Big[\frac1d\sum_{j=1}^{d}\cos\frac{2\pi k_jn}{p}\Big]^2.
\end{centeredalign}

For odd \(p\) choose \(k_j=1,2,\ldots,(p-1)/2\). If \(p\) does not divide \(n\), the phases \(e^{2\pi ijn/p}\) with \(j=1,\ldots,p-1\) sum to \(-1\). Pairing \(j\) with \(p-j\) shows that the cosines in the bracket sum to \(-1/2\). Every non-member is therefore accepted with probability exactly \(1/(p-1)^2\). For \(p=7\), three planes with six basis states accept members with certainty and every other length with probability \(1/36\approx0.028\).

This choice needs almost as many states as a classical counter modulo \(p\). Ambainis and Freivalds showed that a number of planes of order \(\log p\), with suitable multipliers, keeps every non-member below a fixed acceptance bound. For \(p=10007\), a search over random multipliers finds 37 planes, or 74 basis states, with no non-member accepted above probability 0.15. Any deterministic automaton for \(L_p\) needs \(p\) states. For prime \(p\) the same holds for a probabilistic automaton with bounded error.

Mereghetti and colleagues built this automaton with single-photon polarization and linear optics, reducing the error by combining independent runs. The eigenphases \(\pm2\pi k_j/p\) of the letter unitary play the role of energies. The acceptance amplitude is thus the survival amplitude of spectral theory, evaluated at integer times.

The limits of the model follow from unitarity. Each letter acts reversibly, so the automaton cannot erase information about earlier letters. With bounded error, measure-once automata recognize exactly the group languages, whose letters permute the states of the minimal classical automaton (Brodsky and Pippenger). The regular language of words ending in \(a\) is excluded, because recognizing it requires forgetting every earlier letter. Intermediate measurements, as in the measure-many model, enlarge the class but still miss some regular languages. General dissipative operations recover all of them, since a quantum operation can map each basis state to the next state of any classical automaton.

\subsection*{References For The Physical Derivation}
\begin{itemize}
\item \href{https://arxiv.org/abs/quant-ph/9802062}{A. Ambainis and R. Freivalds, 1-way quantum finite automata: strengths, weaknesses and generalizations; the automaton for lengths divisible by p with O(log p) states and the classical state count.}
\item \href{https://arxiv.org/abs/quant-ph/9903014}{A. Brodsky and N. Pippenger, Characterizations of 1-Way Quantum Finite Automata; measure-once automata with bounded error and the measure-many model.}
\item \href{https://arxiv.org/abs/2001.10749}{C. Mereghetti and colleagues, Photonic Realization of a Quantum Finite Automaton; polarization implementation of the periodic-language automaton with error reduction.}
\end{itemize}





\clearpage

\section{Quantum machine learning}

\label{topic:quantum-machine-learning}

A quantum processor runs a fixed sequence of gates in which some rotation angles can be set. Each run ends with a measurement of a spin, and the average over many runs estimates an expectation value. Training adjusts the angles to lower a cost built from such averages, and it needs derivatives with respect to the angles. These derivatives can be measured on the same device by shifting each angle by a quarter turn.

Let one gate be \(U(\theta)=e^{-i\theta G/2}\), where the generator \(G\) is a product of Pauli matrices, so \(G^2=I\). Collect all earlier gates into the input state \(|\psi\rangle\). Collect all later gates, together with the measured observable, into a Hermitian operator \(A\). The measured average is \(f(\theta)=\langle\psi|U(\theta)^\dagger AU(\theta)|\psi\rangle\). Because \(G^2=I\), the gate equals \(\cos(\theta/2)I-i\sin(\theta/2)G\). The average contains only products of \(\cos(\theta/2)\) and \(\sin(\theta/2)\), and hence

\begin{centeredalign}
f(\theta)=a+b\cos\theta+c\sin\theta,\qquad
\frac{df}{d\theta}=\frac{f(\theta+\pi/2)-f(\theta-\pi/2)}{2}.
\end{centeredalign}

The rule follows from \(\cos(\theta+\pi/2)-\cos(\theta-\pi/2)=-2\sin\theta\) and the corresponding identity for the sine. It is exact rather than a finite-difference estimate. Because the shifts are large, the difference is not divided by a small step, and shot noise is not amplified. For a qubit prepared in \(|0\rangle\), rotated by \(R_y(\theta)\) and measured along \(z\), \(f=\cos\theta\) and the rule returns \(-\sin\theta\). If \(A\) has eigenvalues \(\pm1\) and each average uses \(M\) runs, the gradient estimate has variance at most \(1/(2M)\). Generators with more than two eigenvalues need rules with more shifts.

Whether the gradients are large enough to measure is a separate question. McClean and colleagues considered circuits in which the gates on each side of the parameter reproduce the first two moments of uniformly random unitaries. Such an ensemble is called a 2-design. For \(n\) qubits the gradient then averages to zero, and its variance decreases exponentially:

\begin{centeredalign}
\langle\partial_kf\rangle=0,\qquad \operatorname{Var}(\partial_kf)\propto2^{-n}.
\end{centeredalign}

For Pauli generator and observable, a direct sampling of random unitaries on both sides gives a variance close to \(2^{-(n+1)}\). Resolving the sign of such a gradient requires a shot-noise variance \(1/(2M)\) below it. For 50 qubits this means of order \(10^{15}\) runs per gradient component.

The condition matters. Cerezo and colleagues proved that observables acting on few qubits, with circuit depth of order \(\log n\), give at worst polynomially small gradients. Global observables give exponentially small gradients even in shallow circuits. A simulation of layered \(y\) rotations and controlled-\(Z\) gates shows both regimes. For the product of \(Z\) over all qubits, the gradient variance decreased roughly as \(2^{-n}\) from \(n=2\) to 12 at depth two. For \(Z\) on one qubit at the same depth it stayed between 0.35 and 0.39.

Established are the shift rule, the trigonometric dependence on each angle, and the exponential suppression of gradients under the stated randomness. Open is whether such circuits, trained on classical data, outperform classical learning on tasks of practical interest. Proven separations exist for specially constructed problems, discussed under quantum neural networks. The same derivative of an average with respect to a phase sets the sensitivity in quantum metrology.

\subsection*{References For The Physical Derivation}
\begin{itemize}
\item \href{https://arxiv.org/abs/1811.11184}{M. Schuld, V. Bergholm, C. Gogolin, J. Izaac and N. Killoran, Evaluating analytic gradients on quantum hardware; parameter-shift rules for gate derivatives.}
\item \href{https://arxiv.org/abs/1803.11173}{J. R. McClean and colleagues, Barren plateaus in quantum neural network training landscapes; exponentially small gradients for circuits forming 2-designs.}
\item \href{https://arxiv.org/abs/2001.00550}{M. Cerezo and colleagues, Cost Function Dependent Barren Plateaus in Shallow Parametrized Quantum Circuits; global versus local observables at logarithmic depth.}
\end{itemize}





\clearpage

\section{Quantum programming}

\label{topic:quantum-programming}

A quantum program written in a language such as OpenQASM states operations convenient to its author, for example a phase \(e^{i\phi}\) applied when two qubits are both 1. A superconducting or trapped-ion processor executes only its native rotations and one entangling gate, between qubits with fixed connections. A compiler translates one description into the other. That translation sets the gate count and hence the accumulated error.

The controlled phase \(CP(\phi)=\mathrm{diag}(1,1,1,e^{i\phi})\) illustrates the translation. It multiplies only \(|11\rangle\), and \(|11\rangle\langle11|=(I-Z_1)(I-Z_2)/4\). The four terms commute, so the exponential factorizes into single-qubit rotations \(R_z(\beta)=e^{-i\beta Z/2}\) and one two-qubit factor:

\begin{centeredalign}
CP(\phi)=e^{i\phi/4}\,[R_z(\phi/2)\otimes R_z(\phi/2)]\,e^{i\phi Z_1Z_2/4}.
\end{centeredalign}

Conjugating \(Z_2\) by a CNOT with control on qubit 1 gives \(Z_1Z_2\). The two-qubit factor therefore needs two CNOT gates:

\begin{centeredalign}
e^{i\phi Z_1Z_2/4}=\mathrm{CNOT}\,[I\otimes R_z(-\phi/2)]\,\mathrm{CNOT}.
\end{centeredalign}

The global phase \(e^{i\phi/4}\) has no observable effect unless the whole gate is later made conditional on another qubit. A compiler applies such rules throughout a program. The quantum Fourier transform on \(n\) qubits uses \(n\) Hadamard gates and \(n(n-1)/2\) controlled phases with angles \(\pi/2^m\), hence \(n(n-1)\) CNOTs. For three qubits this gives six CNOTs and three Hadamards. The reversal of qubit order at the end costs nothing, because the compiler relabels the outputs. If the qubits couple only along a line, the gate between the outer qubits needs a SWAP, three more CNOTs. With independent CNOT failures of probability 0.01, the success estimate drops from \(0.99^6\approx0.94\) to \(0.99^9\approx0.91\).

Programs can also measure a qubit during the run and use the result to choose later gates. The corrections in quantum teleportation are of this kind. Measuring qubit 1 in the \(Z\) basis and then applying \(X^m\) to qubit 2 for outcome \(m\) has exactly the effect of a CNOT followed by that measurement:

\begin{centeredalign}
\sum_{m}(\Pi_m\otimes X^m)\,\rho\,(\Pi_m\otimes X^m)\\
=\sum_{m}(\Pi_m\otimes I)\,\mathrm{CNOT}\,\rho\,\mathrm{CNOT}\,(\Pi_m\otimes I).
\end{centeredalign}

Here \(\Pi_m=|m\rangle\langle m|\) for \(m=0,1\). This deferred-measurement identity holds because the CNOT is block diagonal in the basis of its control. Read from right to left, it removes a two-qubit gate whose control is measured immediately afterward. Griffiths and Niu applied it to the Fourier transform followed by measurement, as in phase estimation. In that circuit each qubit, after its Hadamard gate, takes part only in controlled phases that are diagonal in its computational basis. Its final measurement can therefore be moved forward to directly after the Hadamard. Each controlled phase then becomes a \(z\) rotation on a later qubit, applied only if the earlier outcome was 1.

The three-qubit transform then needs three Hadamard gates, three measurements and classically conditioned rotations, but no two-qubit gate. A simulation with random input states gives output distributions identical to those of the coherent circuit. The cost moves into classical control, since each outcome must reach the next rotation within the coherence time of the remaining qubits. OpenQASM 3 includes mid-circuit measurement and conditional gates for this purpose. Arbitrary rotation angles are afterwards approximated in a fault-tolerant gate set, as described under quantum logic gates.

\subsection*{References For The Physical Derivation}
\begin{itemize}
\item \href{https://arxiv.org/abs/quant-ph/9503016}{A. Barenco and colleagues, Elementary gates for quantum computation; decomposition of controlled gates into CNOTs and single-qubit rotations.}
\item \href{https://arxiv.org/abs/quant-ph/9511007}{R. B. Griffiths and C.-S. Niu, Semiclassical Fourier Transform for Quantum Computation; replacing two-qubit gates by measurement and classically controlled rotations.}
\item \href{https://arxiv.org/abs/2104.14722}{A. W. Cross and colleagues, OpenQASM 3: A broader and deeper quantum assembly language; mid-circuit measurement and classical control in a program language.}
\end{itemize}





\clearpage

\section{Quantum simulator}

\label{topic:quantum-simulator}

A chain of fifty interacting spins has \(2^{50}\approx10^{15}\) amplitudes, about 18 petabytes in double precision. A controllable quantum device can instead be made to follow the same dynamics, for example trapped ions with laser-driven spin couplings or ultracold atoms in optical lattices. In an analog simulator the device Hamiltonian is engineered to match the target. In a digital simulator the target evolution is decomposed into gates.

Let the target Hamiltonian be a sum \(H=A+B\) of two parts that the device can apply separately, such as spin couplings \(A\) and a transverse field \(B\). Set \(\hbar=1\), divide the time \(t\) into \(n\) steps of length \(\delta=t/n\), and alternate the two evolutions. The Baker--Campbell--Hausdorff expansion of one step gives

\begin{centeredalign}
e^{-iA\delta}e^{-iB\delta}=\exp\!\Big[-i(A+B)\delta-\frac{\delta^2}{2}[A,B]+O(\delta^3)\Big].
\end{centeredalign}

Each step therefore carries an error of size \(\delta^2\|[A,B]\|/2\), where \(\|\cdot\|\) is the operator norm. Unitary evolution preserves norms, so the errors of the \(n\) steps add at most linearly:

\begin{centeredalign}
\Big\|e^{-i(A+B)t}-\big(e^{-iAt/n}e^{-iBt/n}\big)^n\Big\|\leq\frac{t^2}{2n}\,\big\|[A,B]\big\|.
\end{centeredalign}

The error falls as \(1/n\) and vanishes when \(A\) and \(B\) commute. A symmetric step places a half step of \(A\) on each side of a full step of \(B\). It cancels the commutator term and leaves an error of order \(t^3/n^2\), set by double commutators.

For two spins with \(A=JZ_1Z_2\) and \(B=h(X_1+X_2)\), the commutator is \([A,B]=2iJh(Y_1Z_2+Z_1Y_2)\), with norm \(4Jh\). For \(J=h\), \(Jt=1\) and ten steps the bound is \(0.2\). Direct computation gives an actual error of \(0.071\), which halves when the number of steps doubles. The symmetric step gives \(0.006\). For a ring of \(N\) spins the triangle inequality bounds the commutator by \(4NJh\), so the required number of steps grows in proportion to \(NJht^2/\epsilon\) for a target error \(\epsilon\).

Real gates are imperfect. If every step adds an error of at most \(\epsilon_s\), the total error is bounded by \(at^2/n+n\epsilon_s\) with \(a=\|[A,B]\|/2\). The bound is smallest at \(n^*=t\sqrt{a/\epsilon_s}\), where it equals \(2t\sqrt{a\epsilon_s}\). Beyond this number of steps, finer time steps make the simulation worse.

A simulation claim also requires an encoding \(V\) of target states into device states, with \(V^\dagger V=I\). If the device evolution differs from the encoded target evolution by at most \(\varepsilon(t)\), every bounded observable \(O\) inherits a bounded error:

\begin{centeredalign}
\big\|U_{\rm dev}(t)V-VU_{\rm tar}(t)\big\|\leq\varepsilon(t)\ \Longrightarrow\
\big|\langle O\rangle_{\rm tar}-\langle VOV^\dagger\rangle_{\rm dev}\big|\leq2\|O\|\,\varepsilon(t).
\end{centeredalign}

Analog simulators are validated differently. In the experiment of Greiner and colleagues (2002), atoms in an optical lattice realized the Bose--Hubbard model. Raising the lattice depth drove the transition from a superfluid to a Mott insulator, seen as the loss of interference peaks after release from the lattice. No Trotter bound applies there. The error budget comes instead from calibrating tunnelling, interaction and residual trapping against independent measurements. The commutator that sets the Trotter error is the same one that makes gate order matter in a quantum circuit.

\subsection*{References For The Physical Derivation}
\begin{itemize}
\item \href{https://arxiv.org/abs/1912.08854}{A. M. Childs, Y. Su, M. C. Tran, N. Wiebe and S. Zhu, A Theory of Trotter Error; commutator bounds for product formulas.}
\item \href{https://arxiv.org/abs/1308.6253}{I. M. Georgescu, S. Ashhab and F. Nori, Quantum Simulation; analog and digital simulators and their platforms.}
\end{itemize}





\clearpage

\section{Quantum cryptography}

\label{topic:quantum-cryptography}

A banknote that cannot be copied needs a physical feature that no measurement reproduces. Classical features can in principle always be measured precisely and copied. Around 1970 Stephen Wiesner proposed storing, together with a serial number, qubits prepared in randomly chosen states. The proposal, published in 1983, later led to quantum key distribution. Its security rests on the no-cloning theorem.

Suppose a unitary \(U\) copied two states onto a blank qubit, \(U|\psi\rangle|0\rangle=|\psi\rangle|\psi\rangle\) and \(U|\phi\rangle|0\rangle=|\phi\rangle|\phi\rangle\). Unitary maps preserve inner products, so

\begin{centeredalign}
\langle\psi|\phi\rangle=\langle\psi|\phi\rangle^2\quad\Longrightarrow\quad
\langle\psi|\phi\rangle\in\{0,1\}.
\end{centeredalign}

Copying is possible only for orthogonal or identical states. For \(|0\rangle\) and \(|+\rangle=(|0\rangle+|1\rangle)/\sqrt2\) the overlap \(1/\sqrt2\) differs from its square \(1/2\), so no device copies both. An ancilla or an intermediate measurement does not help. Any physical operation is a unitary on a larger system, and the same inner-product argument applies there.

In Wiesner's scheme the bank prepares each of \(n\) qubits in one of the four states \(|0\rangle\), \(|1\rangle\), \(|+\rangle\) and \(|-\rangle\), chosen at random. It keeps the list of choices secret under the serial number of the note. To verify a note, the bank measures each qubit in its recorded basis, and a genuine note passes with certainty.

A counterfeiter with one note tries to produce two notes that both pass. The simplest attack measures each qubit in the \(Z\) basis and prepares two copies of the outcome. With the right basis both copies pass; with the wrong basis each copy passes with probability one half. The success probabilities per qubit, and for a whole note, are

\begin{centeredalign}
p_{\rm meas}=\frac12\cdot1+\frac12\cdot\frac14=\frac58,\qquad p_{\rm opt}=\frac34,\qquad
P_{\rm forge}=\Big(\frac34\Big)^n.
\end{centeredalign}

The optimum over all physical operations, \(3/4\) per qubit, was found by Molina, Vidick and Watrous (2012) using semidefinite programming over channels from one qubit to two. A dual bound and an explicit attack both give \(3/4\). They also showed that entangled attacks on several qubits do no better than independent ones. A note of \(n=50\) qubits is therefore forged with probability \((3/4)^{50}\approx6\times10^{-7}\).

The scheme assumes that only the bank verifies notes and that the bases stay secret. If the bank hands back notes that fail verification, repeated queries reveal the states one qubit at a time (Lutomirski, 2010). Even returning only valid notes permits an attack based on interaction-free measurement. Issuing a fresh note after every verification removes both attacks. A practical bank must also tolerate a small fraction of failed qubits, which raises the forging probability by an amount fixed by the tolerated error.

Notes that anyone could verify without the secret list, public-key quantum money, cannot be secure against unbounded adversaries. Proposed constructions rely on computational assumptions whose security against quantum attacks is still being tested. The four Wiesner states and their trade-off between information and disturbance reappear in BB84 quantum key distribution. The counterfeiter's operation is a quantum channel from one qubit to two, constrained by complete positivity.

\subsection*{References For The Physical Derivation}
\begin{itemize}
\item \href{https://arxiv.org/abs/1202.4010}{A. Molina, T. Vidick and J. Watrous, Optimal counterfeiting attacks and generalizations for Wiesner's quantum money; the optimal success probability (3/4)\textasciicircum{}n.}
\item \href{https://arxiv.org/abs/1404.1507}{A. Brodutch, D. Nagaj, O. Sattath and D. Unruh, An adaptive attack on Wiesner's quantum money; attacks when valid notes are returned, based on interaction-free measurement.}
\item \href{https://arxiv.org/abs/1510.06120}{A. Broadbent and C. Schaffner, Quantum Cryptography Beyond Quantum Key Distribution; quantum money and other cryptographic tasks.}
\end{itemize}





\clearpage

\section{Quantum metrology}

\label{topic:quantum-metrology}

A sensor estimates a physical parameter from how that parameter changes a prepared state. Sensitivity depends on the interaction that encodes the parameter, on the fluctuations of its generator in the preparation, and on the measurement used to distinguish nearby states. Entanglement is useful only when it improves this complete estimation problem.

Let a dimensionless parameter \(\theta\) be encoded by \(|\psi_\theta\rangle=e^{-i\theta G}|\psi_0\rangle\), with Hermitian dimensionless generator \(G\). For this pure-state unitary family, the quantum Fisher information is

\begin{centeredalign}
F_Q=4(\langle G^2\rangle-\langle G\rangle^2),\qquad
\operatorname{Var}\widehat\theta\geq\frac{1}{\nu F_Q}.
\end{centeredalign}

Here \(\nu\) counts independent repetitions and the bound applies to locally unbiased estimation, with attainability requiring a suitable measurement and statistical regime. Large generator variance makes neighboring parameter-dependent states more distinguishable. A state that is an eigenstate of \(G\) acquires only an overall phase and has no sensitivity to \(\theta\) in this task.

For \(N\) spins exposed to the same phase, choose \(G=\tfrac12\sum_j Z_j\). Independent spins prepared along \(x\) have \(F_Q=N\). The coherent superposition \((|0\rangle^{\otimes N}+|1\rangle^{\otimes N})/\sqrt2\) has \(F_Q=N^2\), because the two components acquire phases separated by \(N\theta\).

\begin{centeredalign}
\Delta\theta_{\rm product}\geq\frac{1}{\sqrt{\nu N}},\qquad
\Delta\theta_{\rm GHZ}\geq\frac{1}{N\sqrt\nu}.
\end{centeredalign}

These expressions compare the same single-spin coupling with a counted number of uses. They are not universal bounds for arbitrary many-body Hamiltonians or uncounted preparation resources. The enhanced oscillation also creates phase ambiguities over a broad prior interval, so an estimation scheme must establish which fringe contains the parameter.

Independent dephasing exposes the cost of the collective coherence. If a single-spin off-diagonal element decays as \(e^{-\gamma t}\), the coherence between the two GHZ branches decays as \(e^{-N\gamma t}\). For phase encoding over a fixed duration this gives \(F_Q=N^2e^{-2N\gamma t}\). Increasing \(N\) can then reduce rather than improve the usable information. Optimizing frequency estimation further requires counting the interrogation time and available repetitions.

A proposed sensing mechanism must therefore specify what is being estimated and which resources are fixed. The generator identifies the useful preparation, while environmental coupling and measurement determine how much of its distinguishability is available. This makes sensitivity a calculable property of the assembled experiment rather than a consequence of entanglement alone.

\subsection*{References For The Physical Derivation}
\begin{itemize}
\item \href{https://arxiv.org/abs/quant-ph/0509179}{V. Giovannetti, S. Lloyd and L. Maccone, Quantum metrology; generator fluctuations and counted resources.}
\end{itemize}

\subsection*{Relations In The Original Papers}

\href{https://arxiv.org/html/quant-ph/0509179#S0.EGx4}{arXiv:quant-ph/0509179, S0.EGx4}. An entangled probe gives an uncertainty bound scaling inversely with the number of uses of the parameter generator. Unitary encoding, counted resources and the estimator assumptions of the source; noise changes the attainable scaling.



\clearpage

\section{Quantum bus}

\label{topic:quantum-bus}

Two superconducting qubits on a chip, several millimetres apart, have no direct coupling. Both can couple to one microwave resonator that runs between them. When the qubits are detuned from the resonator, it is excited only virtually, yet it transmits an exchange interaction between them. The resonator then acts as a quantum bus, as demonstrated by Majer and colleagues in 2007.

Let qubits \(j=1,2\) have transition frequencies \(\omega_j\) and couplings \(g_j\) to a resonator mode \(a\) of frequency \(\omega_r\), with detunings \(\Delta_j=\omega_j-\omega_r\). With \(\hbar=1\) and the rotating-wave approximation, the Tavis--Cummings Hamiltonian is

\begin{centeredalign}
H=\omega_ra^\dagger a+\sum_{j=1,2}\Big[\frac{\omega_j}{2}\sigma_z^{(j)}
+g_j\big(a\sigma_+^{(j)}+a^\dagger\sigma_-^{(j)}\big)\Big].
\end{centeredalign}

Consider the dispersive regime \(|\Delta_j|\gg g_j\). The state \(|e,g,0\rangle\), with the excitation on qubit 1, couples to \(|g,e,0\rangle\) only through the intermediate state \(|g,g,1\rangle\). In that state the photon sits in the resonator. The amplitude leaves qubit 1 with strength \(g_1\), pays the energy denominator \(\Delta\), and arrives at qubit 2 with strength \(g_2\). Second-order perturbation theory, or a Schrieffer--Wolff transformation, eliminates the resonator. With the resonator in its vacuum,

\begin{centeredalign}
H_{\rm eff}=\sum_j\frac{\omega_j+\chi_j}{2}\sigma_z^{(j)}
+J\big(\sigma_+^{(1)}\sigma_-^{(2)}+\sigma_-^{(1)}\sigma_+^{(2)}\big),\\
\chi_j=\frac{g_j^2}{\Delta_j},\qquad J=\frac{g_1g_2}{2}\Big(\frac{1}{\Delta_1}+\frac{1}{\Delta_2}\Big).
\end{centeredalign}

For equal detunings \(\Delta_1=\Delta_2=\Delta\) the exchange rate is \(J=g_1g_2/\Delta\). The shifts \(\chi_j\) also move each qubit frequency. When the shifted frequencies coincide, an excitation prepared on qubit 1 oscillates to qubit 2:

\begin{centeredalign}
P_{1\to2}(t)=\sin^2(Jt),\qquad t_{\rm swap}=\frac{\pi}{2J}.
\end{centeredalign}

With \(g/2\pi=100\) MHz and \(\Delta/2\pi=1\) GHz, the exchange is \(J/2\pi=10\) MHz, and a full transfer takes \(25\) ns. Exact diagonalization of the three single-excitation states gives a splitting within 2 percent of \(2J\) for these values. At \(\Delta=5g\) the deviation grows to 7 percent, a correction of relative order \((g/\Delta)^2\). Stopping the transfer halfway, at \(t=\pi/(4J)\), produces an entangling gate of the \(\sqrt{i\mathrm{SWAP}}\) type. Detuning the qubits from each other by much more than \(J\) switches the exchange off, because it no longer conserves energy.

The derivation needs \(g_j\ll|\Delta_j|\) and a resonator photon number below the critical value \(n_{\rm crit}=\Delta^2/(4g^2)\). Terms dropped in the rotating-wave approximation contribute with denominators \(\omega_j+\omega_r\) and are small when \(|\Delta_j|\ll\omega_j+\omega_r\). The bus also costs coherence. Resonator loss at rate \(\kappa\) makes each qubit decay at the Purcell rate \(\kappa g^2/\Delta^2\).

Beyond the vacuum, the qubit frequency shift becomes \(\chi(2a^\dagger a+1)\), so the resonator frequency depends on the qubit state. A microwave tone transmitted through the resonator then measures the qubit without exchanging energy with it. The bus is the Jaynes--Cummings coupling of quantum optics applied to two qubits at once, and it supplies two-qubit gates for a quantum circuit.

\subsection*{References For The Physical Derivation}
\begin{itemize}
\item \href{https://arxiv.org/abs/0709.2135}{J. Majer, J. M. Chow, J. M. Gambetta, J. Koch and colleagues, Coupling Superconducting Qubits via a Cavity Bus; resonator-mediated exchange between two transmons.}
\item \href{https://arxiv.org/abs/2005.12667}{A. Blais, A. L. Grimsmo, S. M. Girvin and A. Wallraff, Circuit Quantum Electrodynamics; dispersive regime, exchange coupling, critical photon number and Purcell decay.}
\end{itemize}





\clearpage

\section{Quantum image processing}

\label{topic:quantum-image-processing}

A camera stores a 1024 by 1024 grayscale image as about \(10^6\) intensities. Twenty qubits have a state space of dimension \(2^{20}\), the same number. Writing the intensities into the amplitudes of those qubits is the most compact quantum encoding of the image. The physical questions are how many operations prepare that state, what a gate does to it, and how many runs recover the pixels.

Number the \(N=2^n\) pixels by an index \(i\) written in binary on \(n\) qubits, and let \(x_i\geq0\) be the intensities. Amplitude encoding prepares

\begin{centeredalign}
|\psi_x\rangle=\frac{1}{\|x\|}\sum_{i=0}^{N-1}x_i|i\rangle,\qquad
\|x\|^2=\sum_i x_i^2,\qquad n=\log_2N.
\end{centeredalign}

The total brightness \(\|x\|\) disappears on normalization and must be stored separately. A generic image fixes \(N-1\) real parameters of this state. A single-qubit gate has at most three continuous angles, and a CNOT gate has none. A circuit layout with \(g\) single-qubit gates therefore reaches at most a \(3g\)-dimensional set of states. Preparing every image requires

\begin{centeredalign}
g\geq\frac{N-1}{3}\approx3.5\times10^5\qquad(N=2^{20}).
\end{centeredalign}

Known constructions reach this order, \(O(N)\) gates. Shorter circuits exist for structured inputs, for example intensity profiles whose sums over dyadic intervals can be computed efficiently (Grover and Rudolph). A quantum random-access memory is sometimes assumed instead, but its hardware also grows in proportion to \(N\).

Processing, by contrast, can be cheap. A Hadamard gate on the least significant qubit mixes each pair of neighbouring pixels \(2j\) and \(2j+1\):

\begin{centeredalign}
(x_{2j},x_{2j+1})\mapsto\frac{(x_{2j}+x_{2j+1},\;x_{2j}-x_{2j+1})}{\sqrt2},\qquad
P_1=\frac{\sum_j(x_{2j}-x_{2j+1})^2}{2\|x\|^2}.
\end{centeredalign}

Outcome 1 on that qubit occurs with probability \(P_1\). It leaves a state whose amplitudes are the differences of neighbouring pixels, a discrete gradient. Yao and colleagues used this single-gate operation for edge detection. For a row of eight pixels with intensities \((0,0,0,1,1,1,1,1)\), the step lies inside the pair of pixels 2 and 3, and \(P_1=1/10\). Edges between pairs appear after the pixel index is shifted by one, which is a permutation of basis states.

The cost returns when the gradient image is to be recovered. A computational-basis measurement returns one pixel index per run, with probability \(p_i\) equal to its squared amplitude. In \(M\) runs pixel \(i\) is counted about \(Mp_i\) times, with fluctuation \((Mp_i)^{1/2}\). A relative precision \(\delta\) on a typical pixel with \(p_i\approx1/N\) therefore needs

\begin{centeredalign}
M\approx\frac{N}{\delta^2}\approx10^8\qquad(N=2^{20},\ \delta=0.1).
\end{centeredalign}

Each run consumes a fresh preparation of order \(N\) gates, because the measurement destroys the state. The signs of the differences need further measurement bases, as in state tomography. Loading, one gate and complete measurement therefore cost far more than the roughly \(N\) operations of a classical gradient filter.

This accounting is established. The qubit number is logarithmic in \(N\), generic loading needs order \(N\) gates, and recovering a full image needs order \(N/\delta^2\) runs. A gain remains possible when the input is already a quantum state, such as light from a scene collected by a quantum sensor, or when only a few numbers are wanted. The edge weight \(P_1\) is such a number. Binomial statistics give its estimate to precision \(\epsilon\) from about \(1/(4\epsilon^2)\) runs, independent of \(N\). Which practical image tasks have such inputs and outputs remains open. Quantum imaging treats the related case in which the detected light itself carries the quantum correlations.

\subsection*{References For The Physical Derivation}
\begin{itemize}
\item \href{https://arxiv.org/abs/1801.01465}{X.-W. Yao and colleagues, Quantum Image Processing and Its Application to Edge Detection: Theory and Experiment; amplitude encoding and edge detection with one single-qubit operation.}
\item \href{https://arxiv.org/abs/quant-ph/0406176}{V. V. Shende, S. S. Bullock and I. L. Markov, Synthesis of Quantum Logic Circuits; asymptotically optimal circuits for initializing a register in a generic state.}
\item \href{https://arxiv.org/abs/quant-ph/0208112}{L. Grover and T. Rudolph, Creating superpositions that correspond to efficiently integrable probability distributions; efficient loading of structured amplitude profiles.}
\end{itemize}





\clearpage

\section{Quantum engineering}

\label{topic:quantum-engineering}

Two superconducting pads joined by a thin insulating barrier behave as a single quantum degree of freedom, although they contain a macroscopic number of electrons. Such circuits are designed and fabricated like electrical circuits, and quantum engineering chooses their parameters so that the lowest levels form a usable qubit. The transmon, a widely used design, shows these choices in one calculation.

Let \(n\) be the number of Cooper pairs that have tunnelled across the Josephson junction and \(\varphi\) the superconducting phase difference, with \([\varphi,n]=i\). The charging energy \(E_C=e^2/(2C)\) sets the cost of charge on the capacitance \(C\), and the Josephson energy \(E_J\) sets the strength of pair tunnelling. The Hamiltonian is

\begin{centeredalign}
H=4E_C(n-n_g)^2-E_J\cos\varphi.
\end{centeredalign}

Here \(n_g\) is an offset charge set by the electrostatic environment, which fluctuates in time. For \(E_J\gg E_C\) the phase stays near zero. Expanding the cosine to fourth order gives a harmonic oscillator with a quartic correction. With \(b\) the oscillator lowering operator and \(\varphi=(2E_C/E_J)^{1/4}(b+b^\dagger)\),

\begin{centeredalign}
H\approx\sqrt{8E_JE_C}\Big(b^\dagger b+\frac12\Big)-E_J-\frac{E_C}{12}\big(b+b^\dagger\big)^4.
\end{centeredalign}

In lowest order the quartic term shifts level \(m\) by \(-(E_C/2)(m^2+m+1/2)\). The level spacings therefore decrease with \(m\), and the two lowest transitions differ by the anharmonicity \(\alpha\):

\begin{centeredalign}
\hbar\omega_{01}\approx\sqrt{8E_JE_C}-E_C,\qquad
\hbar\alpha=\hbar(\omega_{12}-\omega_{01})\approx-E_C.
\end{centeredalign}

The anharmonicity turns the oscillator into a qubit. A microwave pulse resonant with \(\omega_{01}\) drives the \(1\to2\) transition with a detuning \(|\alpha|\). Pulses much longer than \(1/|\alpha|\) therefore keep the population in the two lowest levels. For \(E_C/h=250\) MHz and \(E_J/E_C=50\), the formulas give \(\omega_{01}/2\pi=4.75\) GHz and \(\alpha/2\pi=-250\) MHz. Exact diagonalization in the charge basis gives \(4.74\) GHz and \(-287\) MHz. Since \(1/|\alpha|\) is about \(0.6\) ns, pulses of tens of nanoseconds keep leakage small, and shaped pulses suppress it further.

The large ratio \(E_J/E_C\) is the other design choice. The sensitivity of the levels to the fluctuating offset charge falls as \(\exp(-\sqrt{8E_J/E_C})\). In the example, the transition frequency varies with \(n_g\) by about 10 kHz. At \(E_J/E_C=1\), as in the earlier Cooper-pair box, it varies by almost 800 MHz. The price is a weaker relative anharmonicity, which decreases only as \((8E_J/E_C)^{-1/2}\) (Koch and colleagues, 2007).

The device must also be cold. At \(4.74\) GHz the transition energy corresponds to \(0.23\) K. At a refrigerator temperature of 20 mK the thermal excited population \(e^{-h\nu/k_BT}\) is about \(10^{-5}\), although residual excitation measured in devices is often larger. Loss in dielectric interfaces and broken Cooper pairs then limit the coherence times.

The transmon is a quantum harmonic oscillator made weakly anharmonic, and its Josephson term describes tunnelling of Cooper pairs through the barrier. Transmons coupled through a common resonator exchange excitations as described under quantum bus.

\subsection*{References For The Physical Derivation}
\begin{itemize}
\item \href{https://arxiv.org/abs/cond-mat/0703002}{J. Koch, T. M. Yu, J. Gambetta, A. A. Houck, D. I. Schuster, J. Majer and colleagues, Charge insensitive qubit design derived from the Cooper pair box; transmon spectrum, anharmonicity and charge dispersion.}
\item \href{https://arxiv.org/abs/1904.06560}{P. Krantz, M. Kjaergaard, F. Yan, T. P. Orlando, S. Gustavsson and W. D. Oliver, A Quantum Engineer's Guide to Superconducting Qubits; device design, control pulses and loss mechanisms.}
\end{itemize}





\clearpage

\section{Quantum cellular automaton}

\label{topic:quantum-cellular-automaton}

A particle hops on a lattice in discrete time steps, and its internal two-level state is rotated before each hop. It spreads in a way no classical random walk can. The rule is local, identical at every site and unitary. Such rules define quantum cellular automata: lattices of finite-dimensional cells updated by one global unitary \(U\) that is translation invariant and causal. Causality means that for any operator \(A_x\) on cell \(x\), the evolved operator \(U^\dagger A_xU\) acts only on a fixed neighbourhood of \(x\).

The simplest automaton carries one particle with a two-level internal state, the coin. Each step applies a coin rotation \(C(\theta)\) and then moves the particle one site right if the coin points up and one site left if it points down. Translation invariance lets plane waves with quasimomentum \(k\) evolve independently:

\begin{centeredalign}
U(k)=\begin{pmatrix}e^{-ik}&0\\0&e^{ik}\end{pmatrix}C(\theta),\qquad
C(\theta)=\begin{pmatrix}\cos\theta&-\sin\theta\\ \sin\theta&\cos\theta\end{pmatrix}.
\end{centeredalign}

Since \(\det U(k)=1\), its eigenvalues are \(e^{\pm i\omega(k)}\) with \(2\cos\omega=\operatorname{Tr}U(k)\). The dispersion relation and the group velocity follow directly:

\begin{centeredalign}
\cos\omega(k)=\cos\theta\cos k,\qquad v(k)=\frac{d\omega}{dk}=\frac{\cos\theta\sin k}{\sin\omega(k)}.
\end{centeredalign}

The group velocity reaches its largest value \(\cos\theta\) at \(k=\pi/2\). The shift itself imposes a strict light cone of one site per step. A particle started at one site has a flat distribution of \(k\), and both bands share the same \(v^2\). Its spread therefore grows ballistically, with a coefficient set by the average of \(v^2\) over the Brillouin zone:

\begin{centeredalign}
\frac{\langle x^2\rangle}{t^2}\ \xrightarrow{\ t\to\infty\ }\ \frac{1}{2\pi}\int_{-\pi}^{\pi}v(k)^2\,dk=1-|\sin\theta|,
\qquad\text{classical walk: }\langle x^2\rangle=t.
\end{centeredalign}

For the balanced coin \(\theta=\pi/4\), whose spreading equals that of the Hadamard walk, the standard deviation is \(\sigma\approx0.54\,t\). After 100 steps the distribution is about 54 sites wide, against 10 sites for an unbiased classical random walk. A direct simulation of 400 steps gives \(\sigma/t=0.5412\), and no probability appears outside the light cone. Karski and colleagues (2009) implemented such a walk with single neutral atoms in a spin-dependent optical lattice and imaged the final distribution site by site.

Near \(k=0\) and for small \(\theta\), the dispersion becomes \(\omega^2\approx k^2+\theta^2\). With lattice spacing \(a\) and time step \(\tau\) this is the relativistic relation \(E^2=p^2c^2+m^2c^4\), with \(c=a/\tau\) and \(mc^2=\hbar\theta/\tau\). The automaton is a discretization of the one-dimensional Dirac equation that keeps exact unitarity and a strict light cone.

For many interacting cells the same definition applies. Schumacher and Werner showed that every such automaton is obtained from two layers of blockwise unitaries in a generalized Margolus partitioning of the lattice. In one dimension, an index measures the net flow of quantum information through the lattice (Gross, Nesme, Vogts and Werner). A shift of qubits by one cell has index 2, while every finite-depth circuit of local gates on fixed cells has index 1. A translation is therefore local and reversible, yet no such circuit produces it. The classification in three and more dimensions remains open.

The strict light cone is the discrete counterpart of the Lieb--Robinson bound for lattice Hamiltonians. A digital quantum simulator built from Trotter steps is a circuit of the same layered form, and its gates act within the same kind of light cone.

\subsection*{References For The Physical Derivation}
\begin{itemize}
\item \href{https://arxiv.org/abs/1904.12956}{P. Arrighi, An overview of Quantum Cellular Automata; definitions by causality and translation invariance, quantum walks and Dirac limits.}
\item \href{https://arxiv.org/abs/0910.3675}{D. Gross, V. Nesme, H. Vogts and R. F. Werner, Index theory of one dimensional quantum walks and cellular automata; the index that separates shifts from local circuits.}
\item \href{https://arxiv.org/abs/0907.1565}{M. Karski, L. Forster, J.-M. Choi, A. Steffen, W. Alt, D. Meschede and A. Widera, Quantum Walk in Position Space with Single Optically Trapped Atoms; a walk of single atoms in a spin-dependent optical lattice with site-resolved imaging.}
\end{itemize}





\clearpage

\section{Quantum complexity theory}

\label{topic:quantum-complexity-theory}

Grover's algorithm finds one marked entry among \(N\) with about \((\pi/4)\sqrt N\) queries. Whether a cleverer quantum algorithm could use far fewer is a statement about all possible circuits, not about one of them. Quantum complexity theory answers such questions by fixing a model of computation and counting resources.

The standard model is a uniform family of circuits. For each input length \(m\), a classical procedure running in time polynomial in \(m\) writes down a circuit of polynomially many gates from a fixed universal set. The circuit acts on the input and on ancillas prepared in \(|0\rangle\), and one qubit is measured at the end. A decision problem belongs to BQP, bounded-error quantum polynomial time, when the correct answer appears with probability at least \(2/3\) for every input:

\begin{centeredalign}
\Pr[\text{output}=\text{answer}]\geq\frac23,\qquad
\Pr[\text{majority of }k\text{ runs wrong}]\leq e^{-k/18}.
\end{centeredalign}

The constant \(2/3\) is conventional. Repeating the circuit \(k\) times and taking the majority lowers the error exponentially by the Hoeffding inequality, so 373 repetitions bring it below \(10^{-9}\). Classical bounded-error probabilistic computation, BPP, is contained in BQP. BQP is contained in PSPACE, because each outcome probability can be computed as a sum over paths using polynomial memory. Factoring lies in BQP through Shor's algorithm and has no known efficient classical algorithm. Whether BQP is strictly larger than BPP is not proven.

Lower bounds can be proven in the query model, where the input is a black box. For search, the oracle \(O_x\) flips the sign of the query basis state \(|x\rangle\). Any algorithm alternates fixed unitaries with \(T\) oracle calls. Compare a run with the oracle for marked entry \(x\) against a run in which every oracle call is replaced by the identity. Let \(|\psi_t\rangle\) be the state of the second run after \(t\) calls, and \(\alpha_{x,t}\) the norm of its component with query input \(x\), so that \(\sum_x|\alpha_{x,t}|^2=1\). One call can separate the two runs by at most \(2|\alpha_{x,t}|\), and the Cauchy--Schwarz inequality then bounds the total separation:

\begin{centeredalign}
\big\|\psi_T^x-\psi_T\big\|\leq2\sum_{t=0}^{T-1}|\alpha_{x,t}|,\qquad
\sum_x\big\|\psi_T^x-\psi_T\big\|^2\leq4T\sum_{t=0}^{T-1}\sum_x|\alpha_{x,t}|^2=4T^2.
\end{centeredalign}

Success requires that the final state for each \(x\) outputs \(x\) with probability at least \(2/3\). For at least half of the values \(x\), the run without a marked entry outputs \(x\) with probability at most \(2/N\). For each such \(x\) the two final states differ in norm by at least \(2/3-2/N\), which gives

\begin{centeredalign}
\frac N2\Big(\frac23-\frac2N\Big)^2\leq4T^2\quad\Longrightarrow\quad
T\geq\sqrt{\frac N8}\Big(\frac23-\frac2N\Big)\approx0.24\sqrt N.
\end{centeredalign}

For \(N=10^6\) entries any quantum algorithm needs at least 236 queries. Grover's algorithm uses 785, and a classical search about \(5\times10^5\) on average. The square-root gain is therefore optimal up to a constant factor, and Zalka showed that \((\pi/4)\sqrt N\) is exactly optimal for near-certain success. Grover's algorithm almost saturates the inequality: for \(N=1024\) and 25 iterations the sum of squared separations is 2032, against the bound 2500.

This hybrid argument is due to Bennett, Bernstein, Brassard and Vazirani (1997). It shows that unstructured search admits no exponential quantum speedup. It also yields an oracle relative to which NP is not contained in BQP, although the unrelativized question remains open. Problems with algebraic structure escape the argument; period finding in Shor's algorithm is the standard example. The matching upper bound is the amplitude-amplification angle under quantum algorithm, and the polynomial-size condition counts gates of a quantum circuit.

\subsection*{References For The Physical Derivation}
\begin{itemize}
\item \href{https://arxiv.org/abs/quant-ph/9701001}{C. H. Bennett, E. Bernstein, G. Brassard and U. Vazirani, Strengths and Weaknesses of Quantum Computing; the hybrid argument and the square-root lower bound for search.}
\item \href{https://arxiv.org/abs/quant-ph/9711070}{C. Zalka, Grover's quantum searching algorithm is optimal; the exact constant pi/4 for near-certain success.}
\item \href{https://arxiv.org/abs/0804.3401}{J. Watrous, Quantum Computational Complexity; uniform circuit families, BQP and its relations to classical classes.}
\end{itemize}





\chapter{Boundaries and operator domains}

Walls, barriers, interfaces and the conditions far from a scatterer determine which wave functions are allowed, and with them the spectrum and the scattering of a particle.

Confinement quantizes energies, finite barriers allow tunnelling, and open boundaries turn bound states into resonances and scattering channels.

\section*{Topics}

\begin{longtable}{p{0.92\linewidth}}
\toprule
Page \\
\midrule
\endhead
\phantomsection\label{page:potential-well}Potential well \\
\phantomsection\label{page:scattering}Scattering \\
\phantomsection\label{page:particle-in-a-box}Particle in a box \\
\phantomsection\label{page:quantum-optics}Quantum optics \\
\phantomsection\label{page:wave-interference}Wave interference \\
\phantomsection\label{page:spectral-line}Spectral line \\
\phantomsection\label{page:quantum-tunnelling}Quantum tunnelling \\
\phantomsection\label{page:s-matrix}S-matrix \\
\phantomsection\label{page:quantum-harmonic-oscillator}Quantum harmonic oscillator \\
\phantomsection\label{page:quantum-metamaterial}Quantum metamaterial \\
\bottomrule
\end{longtable}

\clearpage

\section{Potential well}

\label{topic:potential-well}

An electron in a thin layer of gallium arsenide between layers of aluminium gallium arsenide is trapped in one direction. The trap has a finite depth, set by the offset between the conduction bands of the two materials. The electron occupies a small number of discrete levels, and its wave function extends into the barriers. Both the energies and the number of levels follow from one matching condition.

Take a well of depth \(V_0\) and width \(2a\), with \(V=-V_0\) for \(|x|<a\) and \(V=0\) outside. A bound state has \(-V_0<E<0\). Inside the well it oscillates with wave number \(k=\sqrt{2m(E+V_0)}/\hbar\); outside it decays as \(e^{-\kappa|x|}\) with \(\kappa=\sqrt{-2mE}/\hbar\). The potential is even, so each state is even or odd. Continuity of \(\psi'/\psi\) at \(x=a\) gives

\begin{centeredalign}
\text{even: }\kappa=k\tan ka,\qquad \text{odd: }\kappa=-k\cot ka,\qquad
k^2+\kappa^2=\frac{2mV_0}{\hbar^2}.
\end{centeredalign}

The last relation holds because \(E\) cancels. In the dimensionless variables \(\xi=ka\) and \(\eta=\kappa a\), every solution lies on a circle in the positive quadrant. The bound states are the intersections of that circle with the curves \(\eta=\xi\tan\xi\) and \(\eta=-\xi\cot\xi\). Even and odd branches alternate, each starting from \(\eta=0\) at a multiple of \(\pi/2\):

\begin{centeredalign}
\xi^2+\eta^2=R^2,\qquad R=\frac{a\sqrt{2mV_0}}{\hbar},\qquad
N=1+\left\lfloor\frac{2R}{\pi}\right\rfloor.
\end{centeredalign}

A branch is intersected only if it starts inside the circle. The number \(N\) of bound states therefore grows by one each time \(R\) passes a multiple of \(\pi/2\). The even branch starting at \(\xi=0\) is always intersected, so a one-dimensional well of any depth binds at least one state. The count has a direct physical reading. \(N-1\) is the integer part of \(\sqrt{V_0/E_1}\), where \(E_1=\pi^2\hbar^2/[2m(2a)^2]\) is the ground energy of an infinite box of the same width.

In a shallow well, \(R\ll1\), the single bound state has \(\eta\approx R^2\), so \(E\approx-2ma^2V_0^2/\hbar^2\). This energy depends only on the product \(2aV_0\), as for a delta-function well of the same strength, and the decay length \(1/\kappa\) far exceeds the well width. In a deep well the roots approach \(\xi=n\pi/2\), the Dirichlet levels of particle in a box. Penetration into the walls lowers every level below the corresponding box value.

Consider a 10 nm GaAs layer with effective mass \(0.067\,m_e\) and a conduction-band offset of 0.24 eV, appropriate for an aluminium fraction near 0.3 in the barriers. Then \(R=3.25\) and the well holds three bound states. The ground level lies 32 meV above the well bottom, compared with 56 meV for an infinite box of the same width. Its wave function decays into the barrier over 1.7 nm. The third level lies only about 2 meV below the barrier top, so its existence depends on the precise offset. The calculation assumes the same effective mass in both materials. Otherwise current conservation requires \(\psi'/m^*\) rather than \(\psi'\) to be continuous, which shifts the levels slightly. Transitions between such levels are used in quantum-well infrared detectors and quantum cascade lasers.

Dimension changes the binding threshold. In a spherical well of radius \(a\) with zero angular momentum, the radial function \(u=r\psi\) must vanish at the origin. Only the odd one-dimensional solutions survive, so a bound state requires \(R>\pi/2\). The deuteron, bound by 2.2 MeV in a well tens of MeV deep, is a weakly bound state of this kind. At positive energy the same matching condition gives the phase shifts and resonances treated under scattering and S-matrix.

\subsection*{References For The Physical Derivation}
\begin{itemize}
\item \href{https://arxiv.org/abs/1505.03376}{K. Razi Naqvi and S. Waldenstrøm, The finite square well: whatever is worth teaching at all is worth teaching well; tutorial review of the matching conditions, bound-state count and quantum-well applications.}
\end{itemize}





\clearpage

\section{Scattering}

\label{topic:scattering}

Geiger and Marsden directed alpha particles at a thin gold foil and counted the scintillations at each angle. The count fell steeply with angle, as \(1/\sin^4(\theta/2)\), yet a few particles were scattered backwards. The angular distribution of scattered particles measures the interaction. When the interaction is weak, it measures the Fourier transform of the potential.

A beam of particles with mass \(m\) and wave vector \(\mathbf k\) meets a potential \(V(\mathbf r)\) of finite extent. Far from the target the stationary wave is the incident plane wave plus an outgoing spherical wave. The amplitude \(f\) of the spherical wave fixes the number of particles per unit solid angle, divided by the incident flux:

\begin{centeredalign}
\psi(\mathbf r)\to e^{i\mathbf k\cdot\mathbf r}+f(\theta,\varphi)\frac{e^{ikr}}{r},\qquad
\frac{d\sigma}{d\Omega}=|f(\theta,\varphi)|^2.
\end{centeredalign}

The Schrödinger equation with this boundary condition can be rewritten as an integral equation. In it, the wave at the target acts as the source of the outgoing wave. If the potential changes the wave only slightly, the incident plane wave can be used inside the integral. This Born approximation gives the amplitude as the Fourier transform of the potential at the momentum transfer \(\mathbf q=\mathbf k'-\mathbf k\):

\begin{centeredalign}
f_B(\mathbf q)=-\frac{m}{2\pi\hbar^2}\int e^{-i\mathbf q\cdot\mathbf r}V(\mathbf r)\,d^3r,
\qquad |\mathbf q|=2k\sin\frac\theta2.
\end{centeredalign}

Each volume element of the target emits a secondary wave, and \(f_B\) adds those waves with phases set by their path differences. Large angles require large \(q\) and therefore probe short distances. The approximation is valid when \(|V|\ll\hbar^2/(mb^2)\) for a potential of range \(b\). At high energy the weaker condition \(|V|b\ll\hbar v\), with \(v\) the particle speed, suffices.

For a screened Coulomb, or Yukawa, potential \(V=V_0e^{-\mu r}/r\), the integral is elementary. With \(E=\hbar^2k^2/(2m)\) the amplitude and the unscreened limit are

\begin{centeredalign}
f_B=-\frac{2mV_0/\hbar^2}{\mu^2+q^2},\qquad
\frac{d\sigma}{d\Omega}\Big|_{\mu\to0}=\frac{V_0^2}{16E^2\sin^4(\theta/2)}.
\end{centeredalign}

With \(V_0=Z_1Z_2e^2/(4\pi\varepsilon_0)\) the limit is the Rutherford cross-section. A finite range \(1/\mu\) removes the divergence at small angles: the total cross-section is \(\sigma=4\pi(2mV_0/\hbar^2)^2/[\mu^2(\mu^2+4k^2)]\), which is finite only for \(\mu>0\). In a neutral atom the electrons screen the nucleus over roughly the atomic size. They set \(\mu\) and cut off the forward peak that the bare Coulomb law would make infinite.

For 5 MeV alpha particles on gold, the Coulomb parameter \(Z_1Z_2e^2/(4\pi\varepsilon_0\hbar v)\) is about 22, so the Born approximation is not justified. The Rutherford formula nevertheless holds, because the exact Coulomb amplitude differs from the Born amplitude only by a phase. Classical mechanics gives the same cross-section, which is why Rutherford's analysis succeeded. Deviations appear only when the closest approach in a head-on collision, \(Z_1Z_2e^2/(4\pi\varepsilon_0E)\), reaches the nuclear radius. For 5 MeV alpha particles on gold this distance is 45 fm, far outside the nucleus.

For a target with an extended charge distribution, the point-charge Born amplitude acquires a factor \(F(q)\). This form factor is the Fourier transform of the charge density. Electron scattering measured \(F(q)\) for nuclei and for the proton. Its slope at small \(q\), \(F\approx1-q^2\langle r^2\rangle/6\), gives the mean square charge radius. The partial-wave description of the same process, with exact unitarity, is developed under S-matrix. The sum over secondary waves is the same superposition that produces wave interference.

\subsection*{References For The Physical Derivation}
\begin{itemize}
\item \href{https://arxiv.org/abs/hep-ph/0612014}{C. F. Perdrisat, V. Punjabi and M. Vanderhaeghen, Nucleon Electromagnetic Form Factors; electron scattering, form factors and charge radii.}
\end{itemize}





\clearpage

\section{Particle in a box}

\label{topic:particle-in-a-box}

A cyanine dye carries electrons along a conjugated chain from one nitrogen atom to the other and absorbs visible light. Lengthening the chain moves the absorption to the red. A semiconductor nanocrystal shifts its emission from red to blue as it is made smaller. In both cases electrons move almost freely through a region of fixed size. The simplest model of such confinement is a particle of mass \(m\) on an interval \(0<x<L\) with no force inside:

\begin{centeredalign}
-\frac{\hbar^2}{2m}\psi''(x)=E\psi(x),\qquad
\psi(x)=A\sin kx+B\cos kx,\qquad E=\frac{\hbar^2k^2}{2m}.
\end{centeredalign}

The differential equation is the same for every box. Physically different boxes differ in the condition at the ends. The Hamiltonian must conserve probability, so the boundary term produced by integrating by parts has to vanish for every pair of allowed wave functions \(\phi\) and \(\psi\):

\begin{centeredalign}
\int_0^L\big[\phi^*(H\psi)-(H\phi)^*\psi\big]dx
=-\frac{\hbar^2}{2m}\Big[\phi^*\psi'-\phi'^*\psi\Big]_0^L=0.
\end{centeredalign}

Three common choices satisfy this requirement and describe different devices. Impenetrable walls give the Dirichlet condition \(\psi(0)=\psi(L)=0\), the limit of a finite well whose depth tends to infinity. A ring of circumference \(L\) gives the periodic condition \(\psi(L)=\psi(0)\) and \(\psi'(L)=\psi'(0)\). The Neumann condition \(\psi'(0)=\psi'(L)=0\) allows a nonzero density at the walls but no current through them. The allowed wave numbers and energies follow directly:

\begin{centeredalign}
\text{Dirichlet: }E_n=\frac{\pi^2\hbar^2n^2}{2mL^2},\ n=1,2,\dots\qquad
\text{Neumann: same }E_n,\ n=0,1,2,\dots\\
\text{periodic: }E_n=\frac{2\pi^2\hbar^2n^2}{mL^2},\ n=0,\pm1,\pm2,\dots
\end{centeredalign}

The Dirichlet ground state has energy \(\pi^2\hbar^2/(2mL^2)\), and its wave function \(\sqrt{2/L}\,\sin(\pi x/L)\) vanishes at both walls. Confinement costs kinetic energy because the wave must curve to meet the walls. The Neumann box has a constant state at zero energy, with the density as large at the ends as in the middle. The periodic box also starts at zero energy, but each higher level is twofold degenerate: \(\pm n\) describe the two senses of circulation around the ring. In every case, doubling \(L\) lowers each level and each spacing by a factor of four.

In the free-electron model of a cyanine dye, \(N\) delocalized electrons fill the Dirichlet levels in pairs up to \(n=N/2\). Absorption promotes one electron to \(n=N/2+1\), at an energy cost \(\pi^2\hbar^2(N+1)/(2mL^2)\). With \(h=2\pi\hbar\) and \(c\) the speed of light, the absorbed wavelength is

\begin{centeredalign}
\lambda=\frac{8mcL^2}{h(N+1)}.
\end{centeredalign}

For \(L=1.0\) nm and \(N=6\) the formula gives 471 nm, in the blue part of the visible spectrum. Each added CH=CH unit lengthens the box by two bonds of 0.139 nm and adds two electrons. The predicted absorption then moves about 125 nm to the red, while symmetric cyanine dyes show shifts close to 100 nm per unit. The model fails for polyenes, whose alternating bond lengths open an additional gap, so that their absorption wavelength approaches a finite limit.

The periodic box accounts for the electron counts of stable ring molecules. The \(n=0\) level holds two electrons and each degenerate pair \(\pm n\) holds four. Shells therefore close at 2, 6, 10 electrons, that is, at \(4j+2\) for integer \(j\geq0\). Benzene, with six \(\pi\) electrons, fills \(n=0\) and \(n=\pm1\) exactly. The free-electron ring thus reproduces the electron count of Hückel's rule for aromatic rings.

Semiconductor nanocrystals, or quantum dots, follow the same \(1/L^2\) scaling in three dimensions, with electron and hole effective masses in place of \(m\). The scaling holds when the crystal is smaller than the exciton Bohr radius of the bulk material, so that confinement energy dominates the electron--hole attraction. A finite barrier lets the wave function leak into the walls and lowers every level, as calculated under potential well. Under quantum tunnelling the same boundary term becomes a condition matching incoming and outgoing flux.

\subsection*{References For The Physical Derivation}
\begin{itemize}
\item \href{https://arxiv.org/abs/quant-ph/0103153}{G. Bonneau, J. Faraut and G. Valent, Self-adjoint extensions of operators and the teaching of quantum mechanics; boundary conditions that make the box Hamiltonian self-adjoint and their spectra.}
\end{itemize}





\clearpage

\section{Quantum optics}

\label{topic:quantum-optics}

A single atom inside a small cavity of two facing mirrors can absorb a photon, re-emit it into the same mode and absorb it again. When this exchange is faster than the loss of the photon or of the atomic excitation, a weak probe beam no longer finds one cavity resonance. It finds two lines, split although no photon was present before the probe arrived. The Jaynes--Cummings model accounts for the splitting and predicts how it grows with photon number.

Let one cavity mode of angular frequency \(\omega_c\) have annihilation operator \(a\), and let a two-level atom with transition frequency \(\omega_a\) have lowering operator \(\sigma_-=\ket{g}\bra{e}\). The electric-dipole coupling, after the rotating-wave approximation, gives

\begin{centeredalign}
H=\hbar\omega_c a^\dagger a+\hbar\omega_a\sigma_+\sigma_-
+\hbar g\,(a\sigma_++a^\dagger\sigma_-),\\
g=\frac{d}{\hbar}\sqrt{\frac{\hbar\omega_c}{2\varepsilon_0V}}.
\end{centeredalign}

Here \(d\) is the transition dipole moment along the cavity field and \(V\) the mode volume. The coupling grows as the cavity shrinks, because the field of one photon is concentrated in a smaller volume. The rotating-wave approximation drops the terms \(a\sigma_-\) and \(a^\dagger\sigma_+\), which change the number of excitations by two. It is valid when \(g\) and the detuning \(\Delta=\omega_a-\omega_c\) are small compared with \(\omega_c\). The remaining Hamiltonian conserves \(a^\dagger a+\sigma_+\sigma_-\). It couples only the pairs \(\ket{e,n-1}\) and \(\ket{g,n}\), each a two-state system with coupling \(\hbar g\sqrt n\):

\begin{centeredalign}
E_{n,\pm}=n\hbar\omega_c+\frac{\hbar\Delta}{2}\pm\frac{\hbar}{2}\sqrt{\Delta^2+4g^2n},
\qquad n=1,2,\dots
\end{centeredalign}

The factor \(\sqrt n\) comes from \(a^\dagger\ket{n-1}=\sqrt n\ket{n}\). On resonance the lowest doublet is split by \(2\hbar g\), the vacuum Rabi splitting, although the cavity holds no photon. An atom prepared in \(\ket{e,0}\) then oscillates with excited-state probability \(\cos^2(gt)\). The higher doublets are split by \(2\hbar g\sqrt n\), so the ladder is not equally spaced. A probe tuned to the lowest doublet cannot add a second photon at the same frequency, an effect called photon blockade.

In circuit quantum electrodynamics the atom is a superconducting qubit and the cavity is a coplanar microwave resonator. Wallraff and coworkers measured a transmission doublet split by \(2g/2\pi\approx11.6\) MHz. The photon loss rate was about 0.8 MHz and the qubit decay rate about 0.7 MHz, so the two peaks were well resolved. The condition \(g\gg\kappa,\gamma\), with \(\kappa\) and \(\gamma\) the photon and atomic loss rates, defines strong coupling. Single atoms in optical cavities of high finesse show the same doublet.

The \(\sqrt n\) ladder distinguishes a quantized field from a classical one. Suppose the atom starts excited and the field starts in a coherent state with mean photon number \(\bar n\). Each photon-number component makes the atom oscillate at its own frequency:

\begin{centeredalign}
P_e(t)=\frac12+\frac12\sum_{n=0}^\infty p_n\cos\big(2g\sqrt{n+1}\,t\big),
\qquad p_n=e^{-\bar n}\frac{\bar n^{\,n}}{n!}.
\end{centeredalign}

The oscillations dephase after a time of order \(1/g\) and revive near \(t=2\pi\sqrt{\bar n}/g\), when neighbouring frequencies return into phase. A classical field of the same mean intensity would give undamped Rabi oscillations instead. Brune and coworkers observed the frequencies in the ratio \(1:\sqrt2:\sqrt3\) with Rydberg atoms crossing a microwave cavity. Each doublet follows the dynamics of the two-state quantum system, while the cavity mode is a quantum harmonic oscillator.

\subsection*{References For The Physical Derivation}
\begin{itemize}
\item \href{https://arxiv.org/abs/cond-mat/0407325}{A. Wallraff and colleagues, Circuit Quantum Electrodynamics: Coherent Coupling of a Single Photon to a Cooper Pair Box; measured vacuum Rabi splitting in a microwave resonator.}
\item \href{https://arxiv.org/abs/2005.12667}{A. Blais, A. L. Grimsmo, S. M. Girvin and A. Wallraff, Circuit Quantum Electrodynamics; Jaynes-Cummings model, strong coupling and circuit realizations.}
\item \href{https://arxiv.org/abs/2202.00330}{J. Larson and Th. K. Mavrogordatos, The Jaynes-Cummings model and its descendants; excitation-number doublets, collapse and revival.}
\end{itemize}





\clearpage

\section{Wave interference}

\label{topic:wave-interference}

Electrons or molecules sent one at a time through two slits arrive at single points on a screen. After many arrivals the points form bright and dark fringes. The spacing of the fringes is fixed by the de Broglie wavelength and the geometry. The fringes disappear when anything in the surroundings records which slit was taken, even if the particle itself is not deflected.

Let \(\psi_1(x)\) and \(\psi_2(x)\) be the amplitudes at screen position \(x\) for passage through slit 1 or slit 2 alone. With both slits open the amplitude is their sum, and the detection probability contains a cross term:

\begin{centeredalign}
P(x)=|\psi_1(x)|^2+|\psi_2(x)|^2+2\operatorname{Re}\big[\psi_1^*(x)\psi_2(x)\big].
\end{centeredalign}

For slits separated by \(d\) and a screen at distance \(D\gg d\), the two paths to \(x\) differ in length by about \(dx/D\). The relative phase is \(2\pi dx/(\lambda D)\), with \(\lambda=h/p\) for a particle of momentum \(p\). The cross term therefore oscillates with period

\begin{centeredalign}
\Delta x=\frac{\lambda D}{d},\qquad \lambda=\frac{h}{p}.
\end{centeredalign}

This far-field form requires \(D\gg d^2/\lambda\); closer to the slits the near-field pattern applies instead. The envelope of the fringes is the diffraction pattern of a single slit. A spread of velocities in the beam smears the higher-order fringes, because each velocity has its own spacing.

Arndt and coworkers diffracted C\(_{60}\) fullerene molecules of mass 720 u, moving at about 220 m/s, through a grating with period 100 nm. The de Broglie wavelength is 2.5 pm, about four hundred times smaller than the molecule. The diffraction angle \(\lambda/d\) is 25 microradians. At 1.25 m behind the grating the first-order peaks therefore lie 31 micrometres from the central peak, and the far-field condition holds by a wide margin.

Suppose passage through slit \(j\) also changes another system, such as an emitted photon or a spin, into the state \(\ket{E_j}\). The joint state after the slits is \(\psi_1(x)\ket{E_1}+\psi_2(x)\ket{E_2}\). A screen that does not detect the marker averages over its states:

\begin{centeredalign}
P(x)=|\psi_1|^2+|\psi_2|^2+2\operatorname{Re}\big[\psi_1^*\psi_2\langle E_1|E_2\rangle\big],
\qquad \mathcal V=|\langle E_1|E_2\rangle|.
\end{centeredalign}

For equal intensities, the fringe visibility \(\mathcal V\) equals the overlap of the two marker states. Orthogonal marker states remove the fringes completely, and the screen shows the sum of the two single-slit patterns. No momentum kick is needed; the loss comes from the correlation between path and marker. The trade-off between visibility and the information available about the path is the quantitative content of wave--particle duality. Measuring the marker in a basis that does not distinguish \(\ket{E_1}\) from \(\ket{E_2}\), and sorting arrivals by its outcome, recovers fringes, as in the quantum eraser experiment.

A hot molecule may emit one photon of wave number \(k\) in a random direction. If the molecule was at one of two positions separated by \(s\), the marker states have overlap

\begin{centeredalign}
\langle E_1|E_2\rangle=\int\frac{d\Omega}{4\pi}\,e^{ik\hat{\mathbf n}\cdot\mathbf s}=\frac{\sin ks}{ks}.
\end{centeredalign}

Photons with wavelengths much longer than \(s\) leave the overlap close to one and the fringes intact. Photons with wavelengths comparable to \(s\) or shorter carry path information and wash the fringes out. Hackermüller and coworkers heated fullerene molecules before they entered a Talbot--Lau interferometer. The visibility fell as thermal emission increased, in agreement with an overlap calculation of this kind summed over the emission spectrum. The experiment is a direct observation of quantum decoherence by the environment.

\subsection*{References For The Physical Derivation}
\begin{itemize}
\item \href{https://arxiv.org/abs/1109.5937}{K. Hornberger, S. Gerlich, P. Haslinger, S. Nimmrichter and M. Arndt, Colloquium: Quantum interference of clusters and molecules; far-field and Talbot-Lau interference of fullerenes.}
\item \href{https://arxiv.org/abs/quant-ph/0402146}{L. Hackermüller, K. Hornberger, B. Brezger, A. Zeilinger and M. Arndt, Decoherence of matter waves by thermal emission of radiation; loss of fringe visibility by photon emission.}
\end{itemize}





\clearpage

\section{Spectral line}

\label{topic:spectral-line}

The yellow light of a sodium lamp, examined with a high-resolution spectrometer, does not have one sharp frequency. Each line has a width and a shape. Part of the width comes from the finite lifetime of the excited atom, and part from the thermal motion of the emitting atoms. The two mechanisms produce different line shapes, with widths that differ by two orders of magnitude.

An atom in its excited state \(\ket e\) decays to \(\ket g\) by emitting into the continuum of electromagnetic modes. When the coupling to each mode is weak and the mode density varies smoothly near the transition frequency \(\omega_0\), the excited amplitude decays exponentially. This is the Wigner--Weisskopf result. The emitted field inherits this time dependence, and its Fourier transform fixes the spectrum:

\begin{centeredalign}
c_e(t)=e^{-i\omega_0t-\gamma t/2},\qquad
I(\omega)\propto\Big|\int_0^\infty e^{i\omega t}c_e(t)\,dt\Big|^2=\frac{1}{(\omega-\omega_0)^2+\gamma^2/4}.
\end{centeredalign}

This Lorentzian has a full width at half maximum \(\gamma\), equal to the inverse lifetime \(1/\tau\). The decay rate follows from the transition dipole moment \(d=|\langle g|e\mathbf r|e\rangle|\) by summing the emission probability over all field modes:

\begin{centeredalign}
\gamma=\frac{\omega_0^3d^2}{3\pi\varepsilon_0\hbar c^3},\qquad
\Delta\nu_{\rm nat}=\frac{\gamma}{2\pi}=\frac{1}{2\pi\tau}.
\end{centeredalign}

For the hydrogen \(2p\to1s\) transition, \(d=0.745\,ea_0\) and \(\hbar\omega_0=10.2\) eV give \(\gamma=6.3\times10^8\ \mathrm{s}^{-1}\), a lifetime of 1.6 ns. For the sodium D2 line at 589 nm the measured lifetime is 16.2 ns, so the natural width is 9.8 MHz. The cubic dependence on frequency makes spontaneous emission far slower for microwave transitions than for optical ones. The broadening is homogeneous: every atom has the same decay rate and the same centre frequency.

Thermal motion shifts the frequency measured in the laboratory. An atom with velocity component \(v_z\) along the line of sight emits at \(\omega_0(1+v_z/c)\), to first order in \(v_z/c\). For atoms of mass \(M\) at temperature \(T\), the Maxwell distribution of \(v_z\) is Gaussian, and so is the line:

\begin{centeredalign}
I_D(\omega)\propto\exp\Big[-\frac{Mc^2(\omega-\omega_0)^2}{2k_BT\omega_0^2}\Big],\qquad
\frac{\Delta\nu_D}{\nu_0}=\sqrt{\frac{8k_BT\ln2}{Mc^2}}.
\end{centeredalign}

For sodium vapour at 400 K the Doppler width is 1.5 GHz, about 150 times the natural width. This broadening is inhomogeneous: each atom keeps a narrow Lorentzian response centred at its own shifted frequency. Saturated-absorption spectroscopy sends two counter-propagating beams through the vapour. Both beams act on the same atoms only when \(v_z\approx0\), which produces a feature much narrower than the Doppler profile at the line centre.

When both mechanisms act, the profile is their convolution, the Voigt profile. Its core is Gaussian when the Doppler width dominates, but its far wings remain Lorentzian. The Gaussian falls off as \(\exp(-\delta^2)\) in the detuning \(\delta\), whereas the Lorentzian falls only as \(1/\delta^2\). Collisions add a further Lorentzian width proportional to the gas pressure, because each collision interrupts the phase of the emitted wave. The relation between width and lifetime is the Fourier relation between an energy distribution and the survival amplitude, discussed under spectral theory. The decay of a single atom, photon by photon, is described under quantum jump.







\clearpage

\section{Quantum tunnelling}

\label{topic:quantum-tunnelling}

A scanning tunnelling microscope holds a metal tip a few tenths of a nanometre above a conducting surface. A bias of a fraction of a volt drives about a nanoampere across the vacuum gap, although the electrons lack the energy to pass over it. Moving the tip by 0.1 nm changes the current by nearly a factor of ten. A rectangular barrier shows where this sensitivity comes from.

Let a particle of mass \(m\) and energy \(E\) meet a barrier of height \(V_0>E\) occupying \(0<x<a\). Outside the barrier the solutions are travelling waves with wave number \(k=\sqrt{2mE}/\hbar\). Inside, the kinetic energy is negative, and the solutions grow or decay with \(\kappa=\sqrt{2m(V_0-E)}/\hbar\):

\begin{centeredalign}
\psi(x)=\begin{cases}e^{ikx}+r\,e^{-ikx}, & x<0,\\
C e^{\kappa x}+D e^{-\kappa x}, & 0<x<a,\\
t\,e^{ikx}, & x>a.\end{cases}
\end{centeredalign}

Continuity of \(\psi\) and \(\psi'\) at both edges gives four linear equations for \(r\), \(C\), \(D\) and \(t\). The probability current \(j=(\hbar/m)\operatorname{Im}(\psi^*\psi')\) is the same on both sides, so \(|r|^2+|t|^2=1\). Solving for the transmitted fraction gives

\begin{centeredalign}
T=|t|^2=\left[1+\frac{V_0^2\sinh^2\kappa a}{4E(V_0-E)}\right]^{-1}
\\ \kappa a\gg1:\qquad T\approx\frac{16E(V_0-E)}{V_0^2}\,e^{-2\kappa a}.
\end{centeredalign}

The approximation holds for a thick barrier, where the decaying exponential inside dominates. The prefactor varies slowly with energy, while the exponent depends on the width and on the square root of the barrier height above \(E\). For \(E>V_0\) the hyperbolic sine becomes an ordinary sine. Transmission is then incomplete, except at the resonances \(a\sqrt{2m(E-V_0)}/\hbar=n\pi\) where reflections from the two edges cancel. The amplitudes \(r\) and \(t\) for waves arriving from either side form the two-channel S-matrix of the barrier.

In the microscope the barrier height is approximately the work function of the metals, 4 to 5 eV above the Fermi level. With \(\phi=4.5\) eV, the decay constant is \(\kappa=1.09\) per ångström. The current at small bias is proportional to \(e^{-2\kappa d}\) for a gap \(d\), so a change of 0.1 nm alters it by a factor of 8.8. A height change of 10 pm, a small fraction of an atomic diameter, changes the current by about 24 per cent. Feedback that holds the current constant therefore follows the surface height with picometre precision. A full treatment adds the image-charge potential, which lowers and rounds the barrier, and the electronic states of tip and sample. These change the prefactor and the effective barrier height, but keep the exponential dependence on distance.

For a smoothly varying barrier, the local decay constant \(\kappa(x)=\sqrt{2m[V(x)-E]}/\hbar\) replaces \(\kappa\), and the exponent becomes an integral between the classical turning points. This WKB estimate requires \(\kappa\) to change little over one decay length, \(|d\kappa/dx|\ll\kappa^2\), and the exponent to be large:

\begin{centeredalign}
T\approx\exp\Big(-2\int_{x_1}^{x_2}\kappa(x)\,dx\Big),\\
2G=\frac{2\sqrt{2\mu E}\,b}{\hbar}\Big[\arccos\sqrt{R/b}-\sqrt{(R/b)(1-R/b)}\Big].
\end{centeredalign}

In alpha decay the barrier is the Coulomb repulsion between the alpha particle and a daughter nucleus of charge \(Ze\). It extends from the nuclear radius \(R\) to the turning point \(b=2Ze^2/(4\pi\varepsilon_0E)\), where \(E\) is the decay energy and \(\mu\) the reduced mass. The integral gives the Gamow exponent \(2G\) above. Take \(R=1.2\,(A^{1/3}+4^{1/3})\) fm, with \(A\) the mass number of the daughter. The exponents for uranium-238 (\(E=4.27\) MeV) and polonium-212 (\(E=8.95\) MeV) then differ by about 56. The exponent alone predicts a half-life ratio near \(10^{24}\), and the measured ratio is \(5\times10^{23}\). A factor of two in energy changes the lifetime by 24 orders of magnitude, which is the Geiger--Nuttall systematics. As \(V_0\) grows without limit, the decay length vanishes and the walls of particle in a box are recovered.

\subsection*{References For The Physical Derivation}
\begin{itemize}
\item \href{https://arxiv.org/abs/1405.5633}{C. Qi, A. N. Andreyev, M. Huyse, R. J. Liotta, P. Van Duppen and R. Wyss, On the Validity of the Geiger-Nuttall Alpha-Decay Law and its Microscopic Basis; barrier penetration and the energy dependence of alpha-decay half-lives.}
\end{itemize}





\clearpage

\section{S-matrix}

\label{topic:s-matrix}

A detector far from a target sees only outgoing waves. If the target cannot absorb, every particle that enters the interaction region must leave it again, and only the phase of each outgoing wave can change. The S-matrix records those changes. Probability conservation then limits how large a cross-section can be and ties the forward amplitude to the total number of scattered particles.

Consider a particle of wave number \(k\) scattered by a spherically symmetric potential of finite range. Angular momentum is conserved, so each partial wave of angular momentum \(\ell\) scatters independently. Outside the potential, the radial function \(u_\ell(r)=rR_\ell(r)\) combines an incoming and an outgoing spherical wave. The interaction can change only the coefficient of the outgoing wave:

\begin{centeredalign}
u_\ell(r)\propto e^{-i(kr-\ell\pi/2)}-S_\ell\,e^{i(kr-\ell\pi/2)},\qquad
|S_\ell|=1\ \Rightarrow\ S_\ell=e^{2i\delta_\ell}.
\end{centeredalign}

Equal incoming and outgoing flux in each partial wave forces \(|S_\ell|=1\). With several channels, such as inelastic excitation of the target, the same conservation becomes unitarity of the full matrix, \(S^\dagger S=I\). The phase shift \(\delta_\ell\) measures how far the outgoing wave is advanced or delayed relative to free motion. Subtracting the free wave, for which \(S_\ell=1\), gives the scattering amplitude and the total cross-section:

\begin{centeredalign}
f(\theta)=\frac1k\sum_{\ell=0}^\infty(2\ell+1)\,e^{i\delta_\ell}\sin\delta_\ell\,P_\ell(\cos\theta),\\
\sigma=\frac{4\pi}{k^2}\sum_{\ell=0}^\infty(2\ell+1)\sin^2\delta_\ell=\frac{4\pi}{k}\operatorname{Im}f(0).
\end{centeredalign}

Each partial wave contributes at most \(4\pi(2\ell+1)/k^2\), the unitarity limit, reached when \(\delta_\ell=\pi/2\). The last equality is the optical theorem: at \(\theta=0\) every Legendre polynomial equals one, and the imaginary part of \(f\) reproduces the sum. Physically, the forward scattered wave interferes destructively with the incident wave. It removes from the beam exactly the flux that appears at other angles, and this deficit is the shadow behind the target.

A hard sphere of radius \(a\) forces the wave to vanish at \(r=a\). For \(\ell=0\) the outside solution is \(\sin(kr+\delta_0)\), so \(\delta_0=-ka\). Higher partial waves have \(\delta_\ell\propto(ka)^{2\ell+1}\), because the centrifugal barrier keeps them away from the sphere, and they are negligible when \(ka\ll1\). The cross-section has two limits:

\begin{centeredalign}
\sigma\to\frac{4\pi}{k^2}\sin^2ka\approx4\pi a^2\quad(ka\ll1),\qquad
\sigma\to2\pi a^2\quad(ka\gg1).
\end{centeredalign}

A wave much longer than the sphere cannot resolve its shape. It is scattered equally in all directions with amplitude \(f=-a\), so the cross-section is four times the geometric area \(\pi a^2\). At high energy, about \(ka\) partial waves contribute and the sum approaches \(2\pi a^2\). Half of it is reflected as a classical particle would be. The other half forms a diffraction peak in the forward direction, of angular width about \(1/(ka)\). That peak is the shadow required by the optical theorem.

For any short-range potential at low energy, \(\delta_0\approx-ka_s\) defines the scattering length \(a_s\), and \(\sigma\approx4\pi a_s^2\). In ultracold gases a magnetic field tunes \(a_s\) through Feshbach resonances, where a bound state of another channel crosses the collision threshold. In one dimension the same unitarity gives \(|r|^2+|t|^2=1\) for the reflection and transmission amplitudes of quantum tunnelling. When every phase shift is small, the amplitude becomes linear in the phase shifts. To first order in the potential it then coincides with the Born approximation treated under scattering.

\subsection*{References For The Physical Derivation}
\begin{itemize}
\item \href{https://arxiv.org/abs/0812.1496}{C. Chin, R. Grimm, P. Julienne and E. Tiesinga, Feshbach Resonances in Ultracold Gases; s-wave phase shift, scattering length and its magnetic tuning.}
\end{itemize}





\clearpage

\section{Quantum harmonic oscillator}

\label{topic:quantum-harmonic-oscillator}

Small oscillations about a stable equilibrium have a quadratic energy. Quantizing that local approximation produces equally spaced excitations and a nonzero ground-state variance. For mass \(m\), angular frequency \(\omega>0\), position \(X\) and momentum \(P\), take

\begin{centeredalign}
H=\frac{P^2}{2m}+\frac12m\omega^2X^2,\qquad
a=\sqrt{\frac{m\omega}{2\hbar}}X+\frac{iP}{\sqrt{2m\hbar\omega}}.
\end{centeredalign}

The canonical commutator gives \([a,a^\dagger]=1\). Substituting \(a\) and \(a^\dagger\) into the energy leaves an additive term from their noncommutation:

\begin{centeredalign}
H=\hbar\omega\left(a^\dagger a+\frac12\right),\qquad
E_n=\hbar\omega\left(n+\frac12\right).
\end{centeredalign}

Positivity of \(a^\dagger a\) bounds the energy below. Repeated lowering reaches a state annihilated by \(a\); raising generates the nonnegative integer occupations. The half-quantum is consequently fixed by the operator algebra. It is not a fitted offset introduced to match a spectrum.

The ground state has position variance \(\hbar/(2m\omega)\) and momentum variance \(m\hbar\omega/2\). Their product saturates the preparation uncertainty relation. Displacing that state changes the mean position and momentum while leaving the variances fixed. Under the harmonic Hamiltonian, those means obey the classical oscillator equation, even though the state retains quantum fluctuations.

\begin{centeredalign}
X(t)=X(0)\cos\omega t+\frac{P(0)}{m\omega}\sin\omega t.
\end{centeredalign}

The equation is an operator identity. Taking expectations gives the classical motion of the center; taking products gives correlations and variances. A thermal mixture has zero mean displacement but fluctuating energy and a larger position variance. A squeezed state instead redistributes variance between the two quadratures, and that anisotropy rotates in phase space.

The same construction describes a normal mode of a vibrating solid or electromagnetic cavity after its generalized coordinate, conjugate momentum and normalization have been identified. The material and boundary conditions determine the frequency and mode shape. Anharmonic interactions couple modes or shift their spacing, while damping adds an environment and its fluctuations. The oscillator is therefore a transferable quadratic mechanism, not an assertion that all its realizations have identical lifetimes, couplings or measurement units.

For a finite matrix truncation, the highest retained state breaks the exact ladder algebra. This can be checked directly rather than interpreted as a physical correction. Predictions for low occupations converge as the truncation is enlarged; strong driving that reaches the cutoff requires a larger state space.

\subsection*{References For The Physical Derivation}
\begin{itemize}
\item \href{https://arxiv.org/abs/1110.3234}{C. Weedbrook and colleagues, Gaussian Quantum Information; canonical quadratures, Gaussian states and transformations.}
\end{itemize}

\subsection*{Relations In The Original Papers}

\href{https://arxiv.org/html/1110.3234v1#S2.E3}{arXiv:1110.3234v1, S2.E3}. The annihilation operator lowers an oscillator number state with amplitude equal to the square root of its occupation. An infinite bosonic Fock space; finite truncations change the commutator at the upper boundary.



\clearpage

\section{Quantum metamaterial}

\label{topic:quantum-metamaterial}

Microwaves in a superconducting transmission line can be sent through a row of artificial atoms. Each atom is a superconducting qubit with a transition frequency of several gigahertz, and the qubits sit far closer together than one wavelength. A microwave then responds to the row as to a continuous medium with a refractive index. That index, and the frequencies at which the medium stops transmitting, follow from the scattering of one qubit.

Let a qubit of transition frequency \(\omega_0\) emit into the line at rate \(\Gamma_{1D}\), shared equally between the two directions. For weak probe power the qubit stays close to its ground state and responds linearly. A wave of frequency \(\omega\) and wave number \(k=\omega/v\), with \(v\) the phase velocity of the line, is then reflected and transmitted with amplitudes

\begin{centeredalign}
r=-\frac{i\Gamma_{1D}/2}{\omega-\omega_0+i\Gamma_{1D}/2},\qquad t=1+r.
\end{centeredalign}

On resonance the reflection is complete, so a single qubit acts as a mirror one atom wide. Now place identical qubits at spacing \(d\). Between neighbours the wave propagates freely, and at each qubit it is scattered with \(r\) and \(t\). The transfer matrix of one period has unit determinant. Bloch's theorem requires its eigenvalues to be \(e^{\pm iqd}\), where \(q\) is the wave number of the collective wave. Half the trace of the matrix gives the dispersion relation

\begin{centeredalign}
\cos qd=\cos kd+\frac{\Gamma_{1D}}{2(\omega-\omega_0)}\sin kd.
\end{centeredalign}

When the right side exceeds one in magnitude, \(q\) is complex and the wave decays along the array. For subwavelength spacing, \(kd\ll1\), expanding both sides to second order gives an effective permittivity of Lorentz form:

\begin{centeredalign}
\varepsilon_{\rm eff}(\omega)=\frac{q^2}{k^2}\approx1+\frac{\Gamma_{1D}}{kd\,(\omega_0-\omega)},
\qquad \varepsilon_{\rm eff}<0\ \text{ for }\ \omega_0<\omega<\omega_0+\frac{\Gamma_{1D}}{k_0d}.
\end{centeredalign}

Below resonance the index exceeds one and the microwaves slow down. Just above resonance the permittivity is negative, and a stop band opens with width about \(\Gamma_{1D}/(k_0d)\), where \(k_0=\omega_0/v\). The band is much wider than the linewidth of one qubit, because many qubits within a wavelength respond together. For \(\Gamma_{1D}/2\pi=10\) MHz and a spacing of about one sixtieth of a wavelength, \(k_0d\approx0.1\), the stop band is about 100 MHz wide. Arrays of transmon qubits on a coplanar waveguide show such a band in transmission, and tuning the qubits with magnetic flux moves it.

Three approximations enter. The probe must be weak, because a qubit that holds one excitation cannot absorb a second; at higher power the reflection saturates and the medium becomes nonlinear. Losses into other channels and dephasing must be small compared with \(\Gamma_{1D}\); they add an imaginary part to \(\omega_0\) and round the band edges. The coupling must also be weak compared with \(\omega_0\), so that the resonant form of \(r\) holds.

The medium differs from a classical metamaterial because each element is a two-level quantum system. Its saturation at the level of single photons produces strong photon--photon interactions, and its state can be prepared in superpositions. The linear band structure and the tunable stop band are established experimentally. Using the quantum state of the whole array to control propagation remains largely a proposal. The reflection by one qubit is the waveguide form of the atom--field coupling treated under quantum optics. The band calculation repeats the Kronig--Penney construction for an electron in a periodic array of potential wells.

\subsection*{References For The Physical Derivation}
\begin{itemize}
\item \href{https://arxiv.org/abs/2006.03330}{J. D. Brehm, A. N. Poddubny, A. Stehli, T. Wolz, H. Rotzinger and A. V. Ustinov, Waveguide Bandgap Engineering with an Array of Superconducting Qubits; polaritonic band gap of eight tunable transmons in a waveguide.}
\item \href{https://arxiv.org/abs/2103.06824}{A. S. Sheremet, M. I. Petrov, I. V. Iorsh, A. V. Poshakinskiy and A. N. Poddubny, Waveguide quantum electrodynamics: collective radiance and photon-photon correlations; single-emitter reflection and band structure of emitter arrays.}
\item \href{https://arxiv.org/abs/0709.1314}{A. L. Rakhmanov, A. M. Zagoskin, S. Savel'ev and F. Nori, Quantum metamaterials: Electromagnetic waves in a Josephson qubit line; proposal in which the qubit quantum state controls propagation.}
\end{itemize}





\part{Fields, Constraints, And Scale}

\chapter{Fields, constraints, and scale}

In quantum field theory particles are excitations of fields, and their number can change.

Exchange symmetry divides particles into bosons and fermions, gauge invariance fixes how they interact, and renormalization describes how the effective interactions change with the scale of observation.

\section*{Topics}

\begin{longtable}{p{0.92\linewidth}}
\toprule
Page \\
\midrule
\endhead
\phantomsection\label{page:gauge-theory}Gauge theory \\
\phantomsection\label{page:photon}Photon \\
\phantomsection\label{page:renormalization}Renormalization \\
\phantomsection\label{page:quantum-field-theory}Quantum field theory \\
\phantomsection\label{page:fermi-dirac-statistics}Fermi–Dirac statistics \\
\phantomsection\label{page:dirac-equation}Dirac equation \\
\phantomsection\label{page:quantum-electrodynamics}Quantum electrodynamics \\
\phantomsection\label{page:boson}Boson \\
\phantomsection\label{page:quantum-chromodynamics}Quantum chromodynamics \\
\phantomsection\label{page:fermion}Fermion \\
\phantomsection\label{page:electron}Electron \\
\phantomsection\label{page:spin-foam}Spin foam \\
\phantomsection\label{page:ads-cft-correspondence}AdS/CFT correspondence \\
\phantomsection\label{page:loop-quantum-gravity}Loop quantum gravity \\
\phantomsection\label{page:spin-network}Spin network \\
\phantomsection\label{page:quantum-chemistry}Quantum chemistry \\
\phantomsection\label{page:standard-model}Standard Model \\
\phantomsection\label{page:klein-gordon-equation}Klein–Gordon equation \\
\phantomsection\label{page:bose-einstein-statistics}Bose–Einstein statistics \\
\phantomsection\label{page:string-theory}String theory \\
\phantomsection\label{page:quantum-statistical-mechanics}Quantum statistical mechanics \\
\phantomsection\label{page:quantum-gravity}Quantum gravity \\
\phantomsection\label{page:quantum-geometry}Quantum geometry \\
\phantomsection\label{page:quantum-spacetime}Quantum spacetime \\
\phantomsection\label{page:fock-space}Fock space \\
\bottomrule
\end{longtable}

\clearpage

\section{Gauge theory}

\label{topic:gauge-theory}

An electron beam split around a long, shielded solenoid recombines with its interference fringes displaced. The electrons travel only where the magnetic field vanishes. The displacement follows from requiring that the phase of a wave function can be chosen independently at every point. Let \(\psi(\mathbf r,t)\) be the wave function of a particle of charge \(q\). Consider the local change \(\psi\to e^{iq\chi/\hbar}\psi\), with \(\chi(\mathbf r,t)\) an arbitrary real function.

The probability density is unchanged, but the gradient acquires an extra term: \(\nabla(e^{iq\chi/\hbar}\psi)=e^{iq\chi/\hbar}(\nabla\psi+iq\psi\nabla\chi/\hbar)\). The Schrödinger equation with the ordinary gradient therefore changes form. The extra term is cancelled by a vector potential that shifts as \(\mathbf A\to\mathbf A+\nabla\chi\), entering through the covariant derivative

\begin{centeredalign}
\mathbf D=\nabla-\frac{iq}{\hbar}\mathbf A,\qquad
\mathbf D'\big(e^{iq\chi/\hbar}\psi\big)=e^{iq\chi/\hbar}\,\mathbf D\psi,\\
[D_x,D_y]=-\frac{iq}{\hbar}\big(\partial_xA_y-\partial_yA_x\big)=-\frac{iq}{\hbar}B_z.
\end{centeredalign}

The same argument for the time derivative requires a scalar potential with \(\varphi\to\varphi-\partial_t\chi\). The Hamiltonian becomes \((\mathbf p-q\mathbf A)^2/2m+q\varphi\), the minimal coupling of electrodynamics. The commutator of covariant derivatives is the magnetic field. It is the part of the potential that no choice of \(\chi\) can remove. Carrying the phase around a closed path \(\gamma\) and applying Stokes' theorem gives

\begin{centeredalign}
\delta_{\rm AB}=\frac{q}{\hbar}\oint_\gamma\mathbf A\cdot d\mathbf l
=\frac{q}{\hbar}\int_\Sigma\mathbf B\cdot d\mathbf S=\frac{q\Phi}{\hbar},
\end{centeredalign}

where \(\Sigma\) is a surface bounded by \(\gamma\) and \(\Phi\) the enclosed flux. Outside an ideal solenoid the potential \(A_\theta=\Phi/2\pi r\) has zero curl, so the field vanishes on both paths. Locally \(\mathbf A=\nabla(\Phi\theta/2\pi)\), but the angle \(\theta\) is not single valued around the solenoid. The potential therefore cannot be removed everywhere. The beams passing on either side differ in phase by \(q\Phi/\hbar\), and only this phase modulo \(2\pi\) is observable. For electrons the fringes move by one period when \(\Phi\) changes by \(h/e=4.14\times10^{-15}\) Wb.

Tonomura and colleagues tested this in 1986 with a toroidal magnet coated by a superconducting layer. The layer kept the electron wave out of the field region and quantized the trapped flux in units of \(h/2e\). The measured phase shift was \(0\) or \(\pi\), as the relation predicts for an even or odd number of flux quanta.

When the internal state has several components, as for the three colours of a quark, the phase becomes a position-dependent unitary matrix \(U(x)\). The potential becomes a matrix \(A_\mu=A_\mu^aT^a\) built from the generators \(T^a\) of the group, and the field strength acquires a commutator:

\begin{centeredalign}
D_\mu=\partial_\mu-igA_\mu,\qquad
F_{\mu\nu}=\partial_\mu A_\nu-\partial_\nu A_\mu-ig[A_\mu,A_\nu],\\
F_{\mu\nu}\to UF_{\mu\nu}U^{-1},\qquad
W(\gamma)=\operatorname{Tr}\mathcal P\exp\Big(ig\oint_\gamma A_\mu dx^\mu\Big).
\end{centeredalign}

Here \(g\) is the coupling and \(\mathcal P\) orders the matrices along the path. The field strength is no longer invariant; it rotates in the internal space. Observables are traces, such as \(\operatorname{Tr}F_{\mu\nu}F^{\mu\nu}\) and the Wilson loop \(W(\gamma)\) of a test charge carried around \(\gamma\). In the Hamiltonian description Gauss's law generates the time-independent gauge transformations, and physical states are those it annihilates. The commutator term makes the gauge field respond to its own charge, so gluons interact with one another; quantum chromodynamics develops this case. The Aharonov--Bohm phase is the Abelian Wilson loop, and the same loop variables underlie the spin networks of loop quantum gravity.

\subsection*{References For The Physical Derivation}
\begin{itemize}
\item \href{https://arxiv.org/abs/1912.12596}{E. Cohen, H. Larocque, F. Bouchard, F. Nejadsattari, Y. Gefen and E. Karimi, Geometric phase from Aharonov-Bohm to Pancharatnam-Berry and beyond; the Aharonov-Bohm phase and its measurements.}
\item \href{https://arxiv.org/abs/hep-ph/0001312}{G. S. Bali, QCD forces and heavy quark bound states; non-Abelian covariant derivative, field strength and Wilson loops.}
\end{itemize}



\subsection*{Equation Sources}
\begin{itemize}
\item \href{https://arxiv.org/abs/hep-th/0106242}{arXiv:hep-th/0106242}
\item \href{https://arxiv.org/abs/hep-th/0010277}{arXiv:hep-th/0010277}
\end{itemize}

\clearpage

\section{Photon}

\label{topic:photon}

Light from an attenuated laser falling on a 50:50 beam splitter occasionally produces simultaneous clicks in detectors behind both outputs. Light emitted by a single atom or quantum dot after one excitation does not: one detector clicks, never both. Grangier, Roger and Aspect observed this anticorrelation in 1986 with photons from an atomic cascade, using the partner photon of each pair to announce the event. The quantity that distinguishes the two cases is the second-order correlation \(g^{(2)}(0)\).

The electromagnetic field in a volume \(V\) is a set of modes labeled by a wave vector \(\mathbf k\) and a transverse polarization vector \(\boldsymbol\epsilon_\lambda\), with \(\mathbf k\cdot\boldsymbol\epsilon_\lambda=0\). Each mode is a harmonic oscillator with angular frequency \(\omega=c|\mathbf k|\) and annihilation operator \(a_{\mathbf k\lambda}\), where \([a,a^\dagger]=1\). A photon is one quantum of such a mode:

\begin{centeredalign}
E=\hbar\omega,\qquad \mathbf p=\hbar\mathbf k,\qquad
\ket{1_{\mathbf k\lambda}}=a_{\mathbf k\lambda}^\dagger\ket{0},\\
\mathbf E^{(+)}(\mathbf r,t)=i\sum_{\mathbf k,\lambda}\sqrt{\frac{\hbar\omega}{2\epsilon_0V}}\,
\boldsymbol\epsilon_\lambda a_{\mathbf k\lambda}e^{i(\mathbf k\cdot\mathbf r-\omega t)}.
\end{centeredalign}

Here \(\mathbf E^{(+)}\) is the positive-frequency part of the electric field operator, and \(\epsilon_0\) is the vacuum permittivity. A one-photon state has zero mean electric field, although it carries energy \(\hbar\omega\) above the vacuum. An absorbing detector registers normally ordered products of field operators. For two detectors viewing the same mode, the ratio of the coincidence rate to the product of the single rates at zero delay is

\begin{centeredalign}
g^{(2)}(0)=\frac{\langle a^\dagger a^\dagger aa\rangle}{\langle a^\dagger a\rangle^2}
=\frac{\langle n(n-1)\rangle}{\langle n\rangle^2}.
\end{centeredalign}

For a number state \(|n\rangle\) this equals \(1-1/n\), which vanishes for one photon. A coherent state, the output of an ideal laser, has a Poisson number distribution and \(g^{(2)}(0)=1\). Thermal light has \(g^{(2)}(0)=2\), the bunching measured by Hanbury Brown and Twiss. A classical fluctuating intensity \(I\) gives \(\langle I^2\rangle/\langle I\rangle^2\), which the Cauchy--Schwarz inequality keeps at or above one. A value below one therefore has no description by a classical wave.

A two-level emitter, such as a semiconductor quantum dot, must be excited again after each emission before it can emit another photon. Under pulsed excitation the coincidence histogram consists of peaks separated by the pulse period. The value of \(g^{(2)}(0)\) is the area of the peak at zero delay divided by the area of a distant peak. If a pulse yields one photon with probability \(p_1\) and two photons with a small probability \(p_2\), then \(g^{(2)}(0)\approx2p_2/p_1^2\). Quantum dots in optical microcavities reach values of order 0.01 and below.

An attenuated laser behaves differently. With mean photon number \(\bar n\) per pulse, the ratio of two-photon to one-photon pulses is \(\bar n/2\), and \(g^{(2)}(0)\) remains 1. Multiphoton events can then be suppressed only by reducing the rate of single photons as well. This difference matters in quantum key distribution, where a second photon in a pulse can leak information.

Linear loss removes photons independently and leaves \(g^{(2)}(0)\) unchanged. Background light, detector dark counts and timing jitter raise the measured value above that of the emitted state. A small \(g^{(2)}(0)\) certifies few multiphoton events but not that successive photons are identical. That property is tested by the Hong--Ou--Mandel interference described under Boson. The temporal shape of each photon is set by the spontaneous decay treated under Quantum jump.

\subsection*{References For The Physical Derivation}
\begin{itemize}
\item \href{https://arxiv.org/abs/1210.1234}{S. Buckley, K. Rivoire and J. Vuckovic, Engineered Quantum Dot Single Photon Sources; second-order correlation, antibunching and cavity-enhanced single-photon emission.}
\end{itemize}





\clearpage

\section{Renormalization}

\label{topic:renormalization}

The strength of the electromagnetic interaction depends on the momentum transfer at which it is measured. Atomic experiments give \(\alpha=1/137.036\), while electron-positron scattering at the Z-boson mass of 91 GeV requires \(\alpha\approx1/129\). The difference is calculable. Virtual charged pairs screen a charge, and a probe with larger momentum transfer penetrates more of the screening cloud. Renormalization organizes this dependence on scale.

In quantum electrodynamics, loops of charged fermions modify the photon propagator, an effect called vacuum polarization. The interaction between two charges at momentum transfer \(Q\) keeps its lowest-order form, with a modified coupling. Let fermion \(f\) have mass \(m_f\), charge \(e_fe\) and \(N_c\) colours. For \(Q\gg m_f\), and with \(\alpha\) defined at zero momentum transfer, the one-loop result is

\begin{centeredalign}
\alpha(Q)=\frac{\alpha}{1-\Delta\alpha(Q)},\\
\Delta\alpha(Q)=\frac{\alpha}{3\pi}\sum_fN_ce_f^2\left[\ln\frac{Q^2}{m_f^2}-\frac53\right].
\end{centeredalign}

Differentiating with respect to \(\ln Q\) removes the constants and gives the rate of change of the coupling, the beta function. Integrating it between a reference scale \(\mu\) and \(Q\) gives

\begin{centeredalign}
\mu\frac{d\alpha}{d\mu}=\beta(\alpha)=\frac{2\alpha^2}{3\pi}\sum_fN_ce_f^2,\\
\frac{1}{\alpha(Q)}=\frac{1}{\alpha(\mu)}-\frac{2}{3\pi}\sum_fN_ce_f^2\ln\frac{Q}{\mu}.
\end{centeredalign}

The positive beta function expresses screening: the effective charge grows at short distance. Only fermions lighter than \(Q\) contribute appreciably, while heavier ones decouple, so the sum changes at each mass threshold. A prediction written in terms of \(\alpha(Q)\) does not depend on the reference scale \(\mu\) used to define the coupling. This independence is the content of the renormalization group equation.

The three charged leptons give \(\Delta\alpha_{\rm lep}(M_Z)=0.0314\) at one loop. Quarks are confined at low momentum, so their contribution is obtained from the measured cross section for \(e^+e^-\to\mathrm{hadrons}\) through a dispersion relation, which gives \(\Delta\alpha_{\rm had}\approx0.0277\). Together they yield \(1/\alpha(M_Z)\approx137.036\times(1-0.059)\approx128.9\). Bhabha scattering at LEP observed the predicted increase of the coupling directly. This running value enters the W-mass relation of the Standard Model.

Wilson's formulation explains why the same structure governs very different systems. Integrating out fluctuations between a cutoff \(\Lambda\) and \(\Lambda/b\) changes the couplings \(g_i\) of the remaining long-wavelength theory. Near a fixed point \(g^*\), where all beta functions vanish, small deviations scale with powers of \(b\):

\begin{centeredalign}
\beta_i(g^*)=0,\qquad \delta g_i(b)=b^{y_i}\,\delta g_i(1).
\end{centeredalign}

Directions with \(y_i>0\) grow and are called relevant; those with \(y_i<0\) shrink and are called irrelevant. Long-distance behaviour then depends only on the few relevant couplings. For this reason a liquid near its critical point and a uniaxial magnet share the same critical exponents.

The one-loop QED formula holds while \((\alpha/3\pi)\ln(Q^2/m^2)\) remains small. Taken literally with the electron loop alone, it would make \(\alpha\) diverge near \(10^{280}\) MeV, far beyond the range where the theory is applied. In Quantum chromodynamics the self-interaction of gluons reverses the sign of the beta function, and the strong coupling decreases at short distance.

\subsection*{References For The Physical Derivation}
\begin{itemize}
\item \href{https://arxiv.org/abs/0807.4206}{F. Jegerlehner, The running fine structure constant alpha(E) via the Adler function; leptonic and hadronic contributions to the running of alpha.}
\end{itemize}



\subsection*{Equation Sources}
\begin{itemize}
\item \href{https://arxiv.org/abs/hep-ph/0106305}{arXiv:hep-ph/0106305}
\item \href{https://arxiv.org/abs/hep-th/0104197}{arXiv:hep-th/0104197}
\end{itemize}

\clearpage

\section{Quantum field theory}

\label{topic:quantum-field-theory}

A mirror moving at a speed close to that of light emits photons even when no light was present before. In 2011 Wilson and coworkers observed this dynamical Casimir effect in a superconducting circuit, where microwave photons appeared in correlated pairs. Two features of quantum field theory explain it: particle number is not conserved, and the vacuum depends on the Hamiltonian.

Consider a real scalar field \(\phi(\mathbf x,t)\) in a box of volume \(V\), in units with \(\hbar=c=1\). Its conjugate momentum is \(\pi=\dot\phi\), and its energy contains kinetic, gradient and mass terms. Expanding in plane waves decouples the field into one oscillator for each wave vector \(\mathbf k\):

\begin{centeredalign}
H=\int d^3x\left[\frac{\pi^2}{2}+\frac{(\nabla\phi)^2}{2}+\frac{m^2\phi^2}{2}\right]
=\sum_{\mathbf k}\omega_{\mathbf k}\left(a_{\mathbf k}^\dagger a_{\mathbf k}+\frac12\right),\\
\phi(\mathbf x)=\sum_{\mathbf k}\frac{a_{\mathbf k}e^{i\mathbf k\cdot\mathbf x}
+a_{\mathbf k}^\dagger e^{-i\mathbf k\cdot\mathbf x}}{\sqrt{2\omega_{\mathbf k}V}},\qquad
\omega_{\mathbf k}=\sqrt{\mathbf k^2+m^2}.
\end{centeredalign}

The operators obey \([a_{\mathbf k},a_{\mathbf k'}^\dagger]=\delta_{\mathbf k\mathbf k'}\). Each mode is a Quantum harmonic oscillator, and its occupation number counts particles of momentum \(\mathbf k\) and energy \(\omega_{\mathbf k}\). States with any number of particles form a Fock space. The vacuum is the state annihilated by every \(a_{\mathbf k}\), so it is defined with respect to these particular oscillators. Correlation functions such as \(\langle0|T\{\phi(x)\phi(y)\}|0\rangle\) describe the propagation of a quantum between two spacetime points and are the elements of perturbative scattering amplitudes.

A single mode shows what happens when the Hamiltonian changes. Let a mode of unit mass have angular frequency \(\omega_0\) for \(t<0\) and be prepared in its ground state, which \(a\) annihilates. At \(t=0\) the frequency jumps to \(\omega_1\). An instantaneous change leaves the state unchanged, but the new Hamiltonian is diagonalized by different operators. Expressing the coordinate and momentum through \(a\) and \(a^\dagger\) gives the Bogoliubov transformation and the mean number of new quanta:

\begin{centeredalign}
b=u\,a+v\,a^\dagger,\qquad
u=\frac{\omega_0+\omega_1}{2\sqrt{\omega_0\omega_1}},\qquad
v=\frac{\omega_1-\omega_0}{2\sqrt{\omega_0\omega_1}},\\
\langle n_1\rangle=\langle0|b^\dagger b|0\rangle=v^2=\frac{(\omega_0-\omega_1)^2}{4\omega_0\omega_1}.
\end{centeredalign}

The relation \(u^2-v^2=1\) preserves \([b,b^\dagger]=1\). The old ground state is not the new one: it contains new quanta, always in even numbers, because the transformation mixes \(a\) with \(a^\dagger\). Halving the frequency gives \(\langle n_1\rangle=1/8\), with probability 0.943 of finding no quanta and 0.052 of finding two. The final mean energy is \(E_f=\omega_1(\langle n_1\rangle+\frac12)\) in the units \(\hbar=1\) used here. Its excess above the final ground-state energy is \(E_{\rm exc}=E_f-\omega_1/2\), whereas the mean work done on the mode is \(\langle W\rangle=E_f-\omega_0/2\). For \(\omega_0=1\) and \(\omega_1=1/2\), these are \(E_{\rm exc}=1/16\) and \(\langle W\rangle=-3/16\): the mode is excited relative to its new ground state while releasing energy to the agent that lowers its frequency. If the frequency changes slowly, with \(|\dot\omega|\ll\omega^2\), the mode follows its instantaneous ground state and no quanta appear; the work is then the ground-state energy change \((\omega_1-\omega_0)/2\).

In a field, a change that is uniform in space, such as a sudden change of the mass \(m\), conserves momentum. It therefore mixes each \(a_{\mathbf k}\) only with \(a_{-\mathbf k}^\dagger\):

\begin{centeredalign}
b_{\mathbf k}=u_{\mathbf k}a_{\mathbf k}+v_{\mathbf k}a_{-\mathbf k}^\dagger,\qquad
\langle n_{\mathbf k}\rangle=|v_{\mathbf k}|^2,\qquad
u_{\mathbf k}^2-|v_{\mathbf k}|^2=1.
\end{centeredalign}

Quanta are created in pairs of opposite momentum, so the total momentum stays zero while the total number changes. A periodic modulation at angular frequency \(\Omega\) acts as a parametric drive and creates pairs whose frequencies add up to \(\Omega\). In cosmology the same transformation describes the creation of field quanta by the expansion of space.

In the Wilson experiment a coplanar microwave waveguide is terminated by a SQUID, whose inductance depends on the magnetic flux through it. Modulating the flux near 11 GHz moves the effective end of the line back and forth at a few percent of the speed of light in the line. A material mirror cannot approach such speeds, and for small amplitudes the number of created quanta grows as the square of the speed ratio. The circuit was cooled so that thermal photons were negligible. The emitted radiation was symmetric about half the drive frequency, and its two-mode squeezing showed that photons at \(\omega\) and \(\Omega-\omega\) were produced together, as pair creation requires.

The static attraction between conducting plates, which arises from the zero-point energy of the same modes, is treated under Quantum fluctuation. These calculations assume free fields. With interactions the vacuum itself becomes a correlated state, and Renormalization relates the parameters of the theory to measured masses and charges.

\subsection*{References For The Physical Derivation}
\begin{itemize}
\item \href{https://arxiv.org/abs/1105.4714}{C. M. Wilson and colleagues, Observation of the Dynamical Casimir Effect in a Superconducting Circuit; photon pairs and two-mode squeezing from a SQUID-modulated transmission line.}
\item \href{https://arxiv.org/abs/1103.0835}{P. D. Nation, J. R. Johansson, M. P. Blencowe and F. Nori, Stimulating uncertainty: Amplifying the quantum vacuum with superconducting circuits; Bogoliubov transformations and the dynamical Casimir effect.}
\end{itemize}





\clearpage

\section{Fermi–Dirac statistics}

\label{topic:fermi-dirac-statistics}

The conduction electrons of copper, about one per atom, contribute only about one percent of the heat capacity that a classical gas of free particles would have at room temperature. The classical Drude model could not explain this. The explanation lies in the occupation rule derived in Quantum statistical mechanics. A fermionic state holds zero or one particle, and at temperature \(T\) and chemical potential \(\mu\) its mean occupation is \(f(\epsilon)=[e^{(\epsilon-\mu)/k_BT}+1]^{-1}\).

At \(T=0\) the function \(f\) is a step. Every state below \(E_F=\mu(0)\) is filled and every state above it is empty. For free electrons of mass \(m\) with two spin states per wave vector, the filled states form a sphere of radius \(k_F\) in wave-vector space. Counting one wave vector per volume \((2\pi)^3/V\) gives the density \(n\), the Fermi energy and the density of states at the Fermi energy:

\begin{centeredalign}
n=\frac{k_F^3}{3\pi^2},\qquad
E_F=\frac{\hbar^2k_F^2}{2m}=\frac{\hbar^2}{2m}(3\pi^2n)^{2/3},\qquad
g(E_F)=\frac{3N}{2E_F}.
\end{centeredalign}

Copper has \(n=8.47\times10^{28}\,\mathrm{m^{-3}}\). This gives \(k_F=1.36\times10^{10}\,\mathrm{m^{-1}}\), \(E_F=7.0\) eV and a Fermi temperature \(T_F=E_F/k_B=8.2\times10^4\) K. Electrons at the Fermi surface move at \(1.6\times10^6\,\mathrm{m/s}\) even at absolute zero. The filled sphere exerts a degeneracy pressure \(\frac25nE_F\approx38\) GPa without any repulsion between the electrons. Room temperature is a low temperature for this gas: \(T/T_F\approx0.004\).

Heating can change occupations only within about \(k_BT\) of \(E_F\), because deeper electrons have no empty states nearby to move into. Roughly \(g(E_F)k_BT\) electrons each gain an energy of order \(k_BT\), so the energy rises as \(T^2\) and the heat capacity is linear in \(T\). The Sommerfeld expansion fixes the coefficient. For a function \(H(\epsilon)\) that is smooth near \(\mu\),

\begin{centeredalign}
\int_0^\infty H(\epsilon)f(\epsilon)\,d\epsilon=\int_0^\mu H(\epsilon)\,d\epsilon
+\frac{\pi^2}{6}(k_BT)^2H'(\mu)+O(T^4),\\
C_{\rm el}=\frac{\pi^2}{3}g(E_F)k_B^2T=\frac{\pi^2}{2}Nk_B\frac{T}{T_F}.
\end{centeredalign}

Applying the expansion to the particle number gives the small downward shift of \(\mu\) with temperature; applying it to the energy gives \(C_{\rm el}\). Per mole of copper, \(C_{\rm el}=\gamma T\) with \(\gamma=0.50\,\mathrm{mJ\,mol^{-1}K^{-2}}\). At 300 K this is \(0.018R\), where \(R\) is the gas constant, compared with the classical value \(3R/2\) for free particles and the lattice value \(3R\).

Below a few kelvin the lattice contribution falls as \(T^3\), and the electronic term dominates. A plot of \(C/T\) against \(T^2\) is then a straight line whose intercept is \(\gamma\). For copper the measured intercept is \(0.695\,\mathrm{mJ\,mol^{-1}K^{-2}}\), larger than the free-electron value by a factor of 1.38. The excess measures the density of states at the Fermi energy of the real band structure, enhanced by coupling to lattice vibrations. It is usually reported as a thermal effective mass.

The derivation assumes \(k_BT\ll E_F\) and a density of states that is smooth on the scale \(k_BT\). Electron-electron interactions change the coefficients but, in a Fermi liquid, keep the linear law. When \(T\gg T_F\) all occupations are small and \(f\) approaches the Maxwell--Boltzmann weight. The zero-or-one restriction comes from the antisymmetry discussed under Fermion. With a minus sign in the denominator instead, Bose--Einstein statistics allows macroscopic occupation of a single state.

\subsection*{References For The Physical Derivation}
\begin{itemize}
\item \href{https://arxiv.org/abs/0706.3360}{S. Giorgini, L. P. Pitaevskii and S. Stringari, Theory of ultracold Fermi gases; ideal Fermi gas, Fermi energy and low-temperature thermodynamics.}
\end{itemize}





\clearpage

\section{Dirac equation}

\label{topic:dirac-equation}

The \(n=2\) level of hydrogen is not a single level. Microwave and optical spectroscopy resolve components about 11 GHz apart, a fraction of about \(10^{-5}\) of the binding energy of that level. Dirac's relativistic equation for the electron accounts for this fine structure and shows that the electron carries spin one half.

A relativistic wave equation must reproduce \(E^2=p^2c^2+m^2c^4\). The Klein--Gordon equation does so with a second time derivative, and its conserved density can be negative. Dirac required instead an equation linear in \(\partial_t\), \(i\hbar\partial_t\psi=H\psi\), with a Hamiltonian linear in the momentum \(\mathbf p\):

\begin{centeredalign}
H=c\,\boldsymbol\alpha\cdot\mathbf p+\beta mc^2,\\
\alpha_i\alpha_j+\alpha_j\alpha_i=2\delta_{ij},\qquad
\alpha_i\beta+\beta\alpha_i=0,\qquad \beta^2=1.
\end{centeredalign}

Squaring gives \(H^2=c^2p^2+m^2c^4\) exactly when these anticommutation relations hold. The smallest matrices that satisfy them are \(4\times4\), so \(\psi\) has four components. The density \(\psi^\dagger\psi\) is positive and obeys a continuity equation. The free spectrum still contains negative energies \(-\sqrt{p^2c^2+m^2c^4}\); in Quantum field theory these solutions describe positrons.

Rotational symmetry then requires spin. The orbital angular momentum \(\mathbf L=\mathbf r\times\mathbf p\) does not commute with \(H\); instead \([H,\mathbf L]=-i\hbar c\,\boldsymbol\alpha\times\mathbf p\). Let \(\boldsymbol\Sigma\) be the \(4\times4\) matrix made of two copies of the Pauli matrices. It satisfies \([H,\hbar\boldsymbol\Sigma/2]=i\hbar c\,\boldsymbol\alpha\times\mathbf p\), so the conserved angular momentum is

\begin{centeredalign}
\mathbf J=\mathbf L+\frac{\hbar}{2}\boldsymbol\Sigma,\qquad [H,\mathbf J]=0,\qquad
\left(\frac{\hbar}{2}\boldsymbol\Sigma\right)^2=\frac34\hbar^2.
\end{centeredalign}

The internal part has the eigenvalue \(s(s+1)\hbar^2\) with \(s=1/2\). Spin one half is therefore a consequence of the linear relativistic equation, not an additional postulate. Coupling to a magnetic field in the same equation gives the magnetic moment with \(g=2\), as described under Electron.

For an electron bound to a point nucleus of charge \(Ze\), the Dirac equation with the Coulomb potential can be solved exactly. With \(\alpha\) the fine-structure constant, the energy depends on the principal quantum number \(n\) and on the total angular momentum \(j\), but not separately on the orbital quantum number \(l\):

\begin{centeredalign}
E_{nj}=mc^2\left[1+\left(\frac{Z\alpha}{n-j-\frac12+\sqrt{(j+\frac12)^2-Z^2\alpha^2}}\right)^2\right]^{-1/2}\\
\approx mc^2-\frac{mc^2(Z\alpha)^2}{2n^2}
-\frac{mc^2(Z\alpha)^4}{2n^4}\left(\frac{n}{j+\frac12}-\frac34\right).
\end{centeredalign}

The second term is the Bohr energy. The third is smaller by a factor of order \((Z\alpha)^2\approx5\times10^{-5}\) for hydrogen and gives the fine structure. In a nonrelativistic expansion it combines three effects: the relativistic correction to the kinetic energy, spin-orbit coupling with \(g=2\), and the Darwin term for S states. For \(n=2\) the levels with \(j=1/2\) lie below those with \(j=3/2\) by \(mc^2\alpha^4/32\). With the reduced-mass correction this is 10.943 GHz, compared with the measured \(2P_{3/2}\)--\(2P_{1/2}\) interval of 10.969 GHz.

The remaining 0.2 percent comes mainly from the anomalous magnetic moment of the electron, which multiplies the spin-orbit term by \(1+2a_e\). The formula also predicts that \(2S_{1/2}\) and \(2P_{1/2}\), which share \(n=2\) and \(j=1/2\), coincide. Their measured separation of 1057.8 MHz, the Lamb shift, is calculated in Quantum electrodynamics. For a point nucleus the solution with \(j=1/2\) fails when \(Z\alpha>1\). With a finite nucleus the lowest level reaches the negative-energy continuum near \(Z\approx170\), where a single-particle equation no longer applies.

\subsection*{References For The Physical Derivation}
\begin{itemize}
\item \href{https://arxiv.org/abs/hep-ph/0002158}{M. I. Eides, H. Grotch and V. A. Shelyuto, Theory of Light Hydrogenlike Atoms; Dirac energy levels of hydrogen and their radiative and recoil corrections.}
\end{itemize}





\clearpage

\section{Quantum electrodynamics}

\label{topic:quantum-electrodynamics}

In the Dirac theory of hydrogen, the \(2S_{1/2}\) and \(2P_{1/2}\) levels have exactly equal energies. In 1947 Lamb and Retherford drove microwave transitions between them and found \(2S_{1/2}\) higher by about 1000 MHz. The shift arises because the electron is coupled to the quantized electromagnetic field. Quantum electrodynamics treats that field as a dynamical quantum system rather than as a prescribed potential.

In natural units \(\hbar=c=1\), let \(\psi\) be the electron field with mass \(m\) and charge \(q=-e\), and let \(A_\mu\) be the electromagnetic potential. The theory is defined by the Lagrangian density

\begin{centeredalign}
\mathcal L_{\rm QED}=-\frac14F_{\mu\nu}F^{\mu\nu}+\bar\psi(i\gamma^\mu D_\mu-m)\psi,\\
D_\mu=\partial_\mu+iqA_\mu,\qquad F_{\mu\nu}=\partial_\mu A_\nu-\partial_\nu A_\mu.
\end{centeredalign}

Invariance under the local phase change \(\psi\to e^{-iq\chi}\psi\), together with \(A_\mu\to A_\mu+\partial_\mu\chi\), fixes the coupling term \(-qA_\mu\bar\psi\gamma^\mu\psi\). The current \(j^\mu=q\bar\psi\gamma^\mu\psi\) is conserved and is the source in Maxwell's equations, \(\partial_\mu F^{\mu\nu}=j^\nu\). Gauge invariance also implies the Ward--Takahashi identity, which makes charge renormalization the same for every charged particle. Perturbation theory in \(\alpha=e^2/4\pi\approx1/137\) describes emission, absorption and exchange of photons. For a bound electron, emission followed by reabsorption of a photon shifts the energy.

Welton's estimate shows where the shift comes from. Zero-point fluctuations of the field shake the electron, so it samples the Coulomb potential \(V\) over a small random displacement \(\delta\mathbf r\). Averaging \(V\) over that motion adds \(\frac16\langle\delta r^2\rangle\nabla^2V\). A field mode of angular frequency \(\omega\) drives a displacement proportional to \(1/\omega^2\), and the sum over modes gives a logarithm. With \(\hbar\) and \(c\) restored,

\begin{centeredalign}
\langle\delta r^2\rangle=\frac{2\alpha}{\pi}\left(\frac{\hbar}{mc}\right)^2
\ln\frac{\omega_{\max}}{\omega_{\min}},\\
\Delta E_{2S}=\frac16\langle\delta r^2\rangle\langle\nabla^2V\rangle
=\frac{\alpha^5mc^2}{6\pi}\ln\frac{\omega_{\max}}{\omega_{\min}}.
\end{centeredalign}

The last step uses \(\nabla^2V=4\pi\alpha\hbar c\,\delta^3(\mathbf r)\) for the Coulomb potential and \(|\psi_{2S}(0)|^2=1/8\pi a_0^3\), with \(a_0\) the Bohr radius. Only S states have a nonzero density at the nucleus, so \(2P_{1/2}\) is not shifted at this order. The upper limit is \(\hbar\omega_{\max}\approx mc^2\), where the nonrelativistic treatment fails. The lower limit is set by atomic binding, because the electron does not respond freely to slower fields. Bethe's 1947 calculation used an average excitation energy of about 17.8 Rydberg, or 242 eV.

With \(\alpha^5mc^2/6\pi h=135.6\) MHz and a logarithm of 7.65, the estimate is about 1040 MHz, with \(2S_{1/2}\) above \(2P_{1/2}\). The modern value, 1057.8 MHz, requires a relativistic treatment of the high-frequency modes and further terms. In vacuum polarization, virtual electron-positron pairs screen the proton's charge beyond the electron Compton wavelength, and this lowers \(2S\) by 27 MHz. The anomalous magnetic moment of the Electron and nuclear recoil also contribute.

Once the electron mass and charge are fixed by measurement, every term is finite; the procedure is described under Renormalization. The expansion assumes \(Z\alpha\ll1\). In heavy hydrogen-like ions the electron must be treated to all orders in the nuclear field, and the same theory is tested there in a strong-field regime.

\subsection*{References For The Physical Derivation}
\begin{itemize}
\item \href{https://arxiv.org/abs/hep-ph/0002158}{M. I. Eides, H. Grotch and V. A. Shelyuto, Theory of Light Hydrogenlike Atoms; self-energy, vacuum polarization and other contributions to the Lamb shift.}
\end{itemize}





\clearpage

\section{Boson}

\label{topic:boson}

Two photons that reach the two input faces of a balanced beam splitter at the same time leave together through one output face, never one through each. This holds provided they are identical in frequency, polarization and arrival time. Hong, Ou and Mandel observed it in 1987 as a dip to near zero in the coincidence rate of two detectors placed behind the outputs. The effect follows from the exchange symmetry of bosons.

The many-particle amplitude of identical bosons is unchanged when two particles are exchanged. In the occupation description of Fock space, creation operators of different modes commute, and one mode can hold any number \(n=0,1,2,\dots\) of quanta. Let \(a^\dagger\) and \(b^\dagger\) create a photon in the two input modes, and \(c^\dagger\) and \(d^\dagger\) in the two output modes. A lossless 50:50 beam splitter acts linearly on the modes; with one choice of phases,

\begin{centeredalign}
a^\dagger\mapsto\frac{c^\dagger+d^\dagger}{\sqrt2},\qquad
b^\dagger\mapsto\frac{c^\dagger-d^\dagger}{\sqrt2},\\
a^\dagger b^\dagger\ket{0}\mapsto\frac12\left[(c^\dagger)^2-(d^\dagger)^2\right]\ket{0}
=\frac{\ket{2,0}-\ket{0,2}}{\sqrt2}.
\end{centeredalign}

The cross terms \(-c^\dagger d^\dagger+d^\dagger c^\dagger\) cancel because the operators commute. In terms of paths, both photons transmitted and both photons reflected lead to the same final state with amplitudes of equal magnitude and opposite sign. For fermions, anticommutation makes the same cross terms add while \((c^\dagger)^2=0\), so two identical fermions always leave through different faces. For distinguishable particles the probabilities add, and a coincidence occurs with probability one half.

Real photons are wave packets, and only the overlapping parts of their states interfere. Let the second photon be delayed by \(\tau\), and let both have the normalized spectral intensity \(|\phi(\omega)|^2\). For a Gaussian spectrum with root-mean-square angular-frequency width \(\sigma\), the coincidence probability is

\begin{centeredalign}
P_c(\tau)=\frac12\left[1-\left|\int|\phi(\omega)|^2e^{i\omega\tau}d\omega\right|^2\right]
=\frac12\left(1-e^{-\sigma^2\tau^2}\right).
\end{centeredalign}

The width of the dip is set by the inverse bandwidth, not by the wavelength. For photons at 800 nm with a Gaussian spectrum of 5 nm full width at half maximum, \(\sigma=6.2\times10^{12}\,\mathrm{s^{-1}}\). The dip then has a full width of 270 fs at half depth, corresponding to about 80 micrometres of path. Translating one input mirror over this range traces the dip. Its depth, the visibility \(V=1-P_c(0)/P_c(\infty)\), equals the squared overlap of the two single-photon states.

Visibility measurements therefore test whether photons from separate emitters are identical, a requirement for photonic quantum logic. The calculation assumes exactly one photon in each input. Two independent laser pulses with Poisson photon numbers give at most 50 percent visibility, so a deeper dip excludes a description by classical fields. Loss and finite detector efficiency reduce the coincidence rate by the same factor at every delay and leave the visibility unchanged.

Elementary bosons include the photon, the gluons, the W and Z bosons and the Higgs boson, all with integer spin. Composite particles with an even number of fermion constituents, such as helium-4 atoms, behave as bosons when an experiment cannot resolve their structure. The same commutation rule produces stimulated emission and the condensation described under Bose--Einstein statistics.

\subsection*{References For The Physical Derivation}
\begin{itemize}
\item \href{https://arxiv.org/abs/2006.09335}{F. Bouchard and colleagues, Two-photon interference: the Hong-Ou-Mandel effect; beam-splitter transformation of two photons, dip width and visibility.}
\end{itemize}





\clearpage

\section{Quantum chromodynamics}

\label{topic:quantum-chromodynamics}

Electrons scattered from protons at large momentum transfer behave as if they struck nearly free point-like quarks. At low momentum transfer the same quarks are never observed in isolation. Both facts follow from the dependence of the strong coupling on the resolution scale. Quantum chromodynamics is an \(SU(3)\) gauge theory. Each quark field \(q_f\) of flavour \(f\) and mass \(m_f\) has three colour components. Eight gluon fields \(A^a_\mu\) act on them through generators \(T^a\) with \([T^a,T^b]=if^{abc}T^c\) and \(\operatorname{Tr}T^aT^b=\delta^{ab}/2\). In natural units \(\hbar=c=1\),

\begin{centeredalign}
\mathcal L=-\frac12\operatorname{Tr}F_{\mu\nu}F^{\mu\nu}
+\sum_f\bar q_f\big(i\gamma^\mu D_\mu-m_f\big)q_f,\\
D_\mu=\partial_\mu-ig_sA^a_\mu T^a,\qquad
F_{\mu\nu}=\partial_\mu A_\nu-\partial_\nu A_\mu-ig_s[A_\mu,A_\nu].
\end{centeredalign}

Squaring the commutator term gives three-gluon and four-gluon vertices, so gluons carry colour themselves. Virtual quarks and gluons make the effective coupling depend on the momentum scale \(\mu\) at which it is measured. At one loop,

\begin{centeredalign}
\mu\frac{dg_s}{d\mu}=-\frac{\beta_0}{16\pi^2}g_s^3,\qquad
\beta_0=\frac{11}{3}C_A-\frac43T_Fn_f=11-\frac23n_f,
\end{centeredalign}

with \(C_A=3\) the quadratic Casimir of the gluon representation, \(T_F=1/2\) the normalization of the quark generators, and \(n_f\) the number of quark flavours light compared with \(\mu\). The two signs have a physical origin. A quark loop polarizes the vacuum as an electron loop does in electrodynamics: virtual pairs screen a test charge. In a background colour-magnetic field, a massless field of spin \(s\) contributes in proportion to \((-1)^{2s}[(2s)^2-1/3]\). The term \(-1/3\) is the diamagnetic response of orbital motion, and \((2s)^2\) is the paramagnetic response of the spin. The factor \((-1)^{2s}\) is the sign of a fermion loop.

For gluons, \(s=1\) gives \(4-1/3=11/3\), and the spin term dominates. A Dirac quark gives \(-2/3\) per flavour. The vacuum then acts as a medium with magnetic permeability above one. Relativistic invariance makes its dielectric constant the inverse, below one, so a colour charge is antiscreened. In terms of \(\alpha_s=g_s^2/4\pi\) the equation integrates to

\begin{centeredalign}
\alpha_s(Q^2)=\frac{\alpha_s(\mu^2)}{1+\frac{\beta_0}{4\pi}\alpha_s(\mu^2)\ln(Q^2/\mu^2)}
=\frac{4\pi}{\beta_0\ln(Q^2/\Lambda^2)}.
\end{centeredalign}

With five light flavours \(\beta_0=23/3\). Starting from the measured \(\alpha_s=0.118\) at the \(Z\) mass, \(91.2\) GeV, the one-loop formula gives \(\alpha_s=0.173\) at 10 GeV and \(0.088\) at 1 TeV. Determinations from tau decays, quarkonia, jet rates and the \(Z\) width follow this decrease. The one-loop scale is \(\Lambda\approx90\) MeV. Higher orders and the renormalization scheme change \(\Lambda\), so its value is quoted with both specified. Asymptotic freedom requires \(\beta_0>0\), that is at most sixteen flavours.

The formula holds while \(\alpha_s\) is small. Its divergence as \(Q\) approaches \(\Lambda\) marks the breakdown of perturbation theory, not a measured infinity. Confinement and hadron masses are computed by lattice simulation. There the Wilson loop of a static quark--antiquark pair gives a potential that grows linearly with separation. The slope, the string tension, is close to 1 GeV per femtometre, until pair creation breaks the flux tube. A similar tension fixes the slope of meson Regge trajectories, which motivated the string models described under string theory. The underlying covariant derivative is derived under gauge theory.

\subsection*{References For The Physical Derivation}
\begin{itemize}
\item \href{https://arxiv.org/abs/1604.08082}{A. Deur, S. J. Brodsky and G. F. de Teramond, The QCD Running Coupling; beta function, running of alpha\_s and its measurements.}
\item \href{https://arxiv.org/abs/hep-ph/0502113}{F. Wilczek, Asymptotic Freedom: From Paradox to Paradigm; antiscreening by gluons and its physical interpretation.}
\item \href{https://arxiv.org/abs/hep-ph/0001312}{G. S. Bali, QCD forces and heavy quark bound states; static potential and string tension from Wilson loops.}
\end{itemize}





\clearpage

\section{Fermion}

\label{topic:fermion}

Atoms of helium-3, released from a trap and detected far away, arrive in pairs at nearby points less often than uncorrelated particles would. Atoms of helium-4 in the same apparatus arrive together more often. Electrons with parallel spins in a metal similarly avoid one another, even apart from their Coulomb repulsion. These effects follow from the sign that the amplitude of identical fermions acquires under exchange.

Let \(z=(\mathbf x,s)\) denote the position and spin projection of one particle. The amplitude of identical fermions changes sign when two complete coordinates are exchanged. With orthonormal one-particle states \(\varphi_a,\varphi_b\) for two particles, or \(\varphi_1,\dots,\varphi_N\) for \(N\) particles, the antisymmetric amplitudes are

\begin{centeredalign}
\Psi(z_1,z_2)=\frac{\varphi_a(z_1)\varphi_b(z_2)-\varphi_b(z_1)\varphi_a(z_2)}{\sqrt2},\\
\Psi(z_1,\dots,z_N)=\frac{1}{\sqrt{N!}}\det\left[\varphi_i(z_j)\right].
\end{centeredalign}

A determinant with two equal rows or two equal columns vanishes. Hence \(\Psi=0\) if two particles occupy the same one-particle state, which is the Pauli exclusion principle. The amplitude also vanishes whenever \(z_1=z_2\), so two particles with equal spin projection are never found at the same point. Opposite spin projections make the complete coordinates different. Two electrons can therefore share one spatial orbital in a spin singlet, as in the helium ground state or the chemical bond treated under Quantum chemistry.

The suppression extends over a finite distance. Consider the ground state of free electrons in a large volume, a Slater determinant of plane waves filling a Fermi sphere of radius \(k_F\). Let \(n_\sigma=k_F^3/6\pi^2\) be the density of one spin orientation. The probability density of finding a same-spin electron at distance \(r\) from a given one, divided by \(n_\sigma\), is

\begin{centeredalign}
g_{\uparrow\uparrow}(r)=1-\left[\frac{3j_1(k_Fr)}{k_Fr}\right]^2,\qquad
j_1(x)=\frac{\sin x-x\cos x}{x^2},\\
n_\sigma\int\left[g_{\uparrow\uparrow}(r)-1\right]d^3r=-1.
\end{centeredalign}

The bracket is the average of \(e^{i\mathbf k\cdot\mathbf r}\) over the filled sphere and comes from the exchange terms of the determinant. The function \(g_{\uparrow\uparrow}\) rises from zero at \(r=0\) to one half at \(k_Fr=1.81\) and approaches one with weak oscillations. The integral shows that this exchange hole contains exactly one missing electron of the same spin. Opposite spins remain uncorrelated in this state, \(g_{\uparrow\downarrow}=1\), until the Coulomb interaction is included.

In sodium, with one conduction electron per atom, \(k_F=0.92\times10^{10}\,\mathrm{m^{-1}}\), and the same-spin density falls to half within 0.20 nm of each electron. Each electron is thus surrounded by a deficit of negative charge, which lowers its electrostatic energy. For the uniform electron gas this exchange energy is \(-(3/4\pi)\,e^2k_F/4\pi\epsilon_0\) per electron. In Rydberg units it is \(-0.916/r_s\), with \(r_s\) the Wigner--Seitz radius in Bohr radii. Sodium has \(r_s=3.93\), which gives \(-3.2\) eV per electron, larger in magnitude than the mean kinetic energy of 1.9 eV.

The exchange sign is not an interaction and has no adjustable strength; it is fixed by the identity of the particles. Its connection with half-integer spin, the spin-statistics theorem, requires relativistic quantum field theory with locality and positive energy. In the occupation description of Fock space the same rule is the anticommutation \(\{a_i^\dagger,a_j^\dagger\}=0\), so that \((a_i^\dagger)^2=0\). Thermal occupations of fermionic states follow Fermi--Dirac statistics. Quasiparticles confined to two dimensions can have other exchange phases.

\subsection*{References For The Physical Derivation}
\begin{itemize}
\item \href{https://arxiv.org/abs/cond-mat/0612278}{T. Jeltes and colleagues, Comparison of the Hanbury Brown-Twiss effect for bosons and fermions; antibunching of helium-3 and bunching of helium-4 atoms.}
\item \href{https://arxiv.org/abs/0706.3360}{S. Giorgini, L. P. Pitaevskii and S. Stringari, Theory of ultracold Fermi gases; occupation and spin-resolved pair correlations of an ideal Fermi gas.}
\end{itemize}



\subsection*{Equation Sources}
\begin{itemize}
\item \href{https://arxiv.org/abs/cond-mat/0005069}{arXiv:cond-mat/0005069}
\item \href{https://arxiv.org/abs/quant-ph/0305150}{arXiv:quant-ph/0305150}
\item \href{https://arxiv.org/abs/hep-ph/0007343}{arXiv:hep-ph/0007343}
\end{itemize}

\clearpage

\section{Electron}

\label{topic:electron}

An electron held in a magnetic field precesses its spin slightly faster than it circulates in its cyclotron orbit. The small difference between the two frequencies fixes the electron's magnetic moment with a relative uncertainty of about \(10^{-13}\). Through the \(g\)-factor, it tests the Dirac equation and quantum electrodynamics together.

The electron has charge \(-e\), mass \(m=0.511\,\mathrm{MeV}/c^2\) and spin one half, with spin operator \(\mathbf S=\hbar\boldsymbol\sigma/2\) built from the Pauli matrices \(\boldsymbol\sigma\). Its magnetic moment is written \(\boldsymbol\mu=-g(e/2m)\mathbf S\). A classical rotating body with equal distributions of charge and mass has \(g=1\). In the nonrelativistic limit of the Dirac equation, the two upper spinor components obey a Schrödinger equation with kinetic energy \((\boldsymbol\sigma\cdot\boldsymbol\pi)^2/2m\). Here \(\boldsymbol\pi=\mathbf p+e\mathbf A\) is the kinetic momentum in a vector potential \(\mathbf A\):

\begin{centeredalign}
(\boldsymbol\sigma\cdot\boldsymbol\pi)^2=\boldsymbol\pi^2+e\hbar\,\boldsymbol\sigma\cdot\mathbf B,
\qquad
H=\frac{\boldsymbol\pi^2}{2m}+2\,\frac{e}{2m}\,\mathbf S\cdot\mathbf B.
\end{centeredalign}

The spin term arises because the components of \(\boldsymbol\pi\) do not commute in a magnetic field \(\mathbf B\): \([\pi_x,\pi_y]=-ie\hbar B_z\), and cyclically. Comparison with \(H=-\boldsymbol\mu\cdot\mathbf B\) gives \(g=2\). A Schrödinger equation with kinetic energy \(\boldsymbol\pi^2/2m\) contains no such term, and the Dirac value is not inserted by hand. With \(g=2\) the spin precession frequency \(g\,eB/2m\) equals the cyclotron frequency \(eB/m\), so the spin keeps a fixed angle to the velocity.

Coupling to the quantized electromagnetic field changes this value slightly. The electron emits and reabsorbs virtual photons, which modifies its response to a slowly varying external field. Schwinger's one-loop calculation of 1948 gave the anomaly \(a_e=(g-2)/2=\alpha/2\pi\approx0.0011614\), with \(\alpha\) the fine-structure constant. Higher orders form a series in \(\alpha/\pi\) whose coefficients are known through the fifth power:

\begin{centeredalign}
a_e=\frac12\left(\frac{\alpha}{\pi}\right)-0.328479\left(\frac{\alpha}{\pi}\right)^2\\
+1.181241\left(\frac{\alpha}{\pi}\right)^3-1.9122\left(\frac{\alpha}{\pi}\right)^4+\dots
\end{centeredalign}

In a Penning trap the anomaly is measured directly as a ratio of frequencies. The cyclotron frequency is \(\omega_c=eB/m\) and the spin precession frequency is \(\omega_s=g\,eB/2m\), so their difference isolates \(a_e\):

\begin{centeredalign}
\omega_a=\omega_s-\omega_c=a_e\,\omega_c,\qquad
\frac{g}{2}=1+\frac{\omega_a}{\omega_c}.
\end{centeredalign}

At \(B=5.36\) T, \(\omega_c/2\pi=150\) GHz and \(\omega_a/2\pi=174\) MHz. A single electron in a cryogenic trap occupies the lowest cyclotron levels, and quantum jumps between these levels and between spin states are detected one at a time. The 2023 measurement found \(a_e=0.00115965218059(13)\). The series summed through \((\alpha/\pi)^5\) gives \(0.00115965218\) with the value of \(\alpha\) from atom-recoil experiments. The remaining difference, about \(4\times10^{-12}\), is of the size of the calculated contributions from virtual muons and hadrons. At this level the comparison is limited by the independent measurements of \(\alpha\), which disagree among themselves by more than their stated uncertainties.

The Pauli description with \(g=2\) requires fields that vary slowly over the reduced Compton wavelength \(\hbar/mc=0.386\) pm and kinetic energies small compared with \(mc^2\). Corrections of relative order \((v/c)^2\) give the hydrogen fine structure derived under Dirac equation. The radiative effects that shift \(g\) also shift atomic levels, and the Lamb shift discussed in Quantum electrodynamics is their counterpart in hydrogen.

\subsection*{References For The Physical Derivation}
\begin{itemize}
\item \href{https://arxiv.org/abs/2209.13084}{X. Fan, T. G. Myers, B. A. D. Sukra and G. Gabrielse, Measurement of the Electron Magnetic Moment; Penning-trap determination of g/2 from cyclotron and anomaly frequencies.}
\item \href{https://arxiv.org/abs/1712.06060}{T. Aoyama, T. Kinoshita and M. Nio, Revised and Improved Value of the QED Tenth-Order Electron Anomalous Magnetic Moment; coefficients of the series in alpha/pi and comparison with experiment.}
\end{itemize}





\clearpage

\section{Spin foam}

\label{topic:spin-foam}

The amplitude for a particle to travel between two points is a sum over the paths joining them. Spin foam models apply the same rule to the geometry of space, described at each instant by a spin network. The history of a spin network is a two-dimensional complex whose faces carry spins and whose edges carry rotation-invariant couplings of those spins. In three spacetime dimensions the construction can be checked against classical gravity through the Ponzano--Regge model. Spacetime is triangulated by tetrahedra, and each edge \(e\) is labelled by a spin \(j_e\). The spin is read as a length \(l_e=j_e+1/2\) in units of the three-dimensional Planck length. Up to sign conventions the amplitude is

\begin{centeredalign}
Z=\sum_{\{j_e\}}\ \prod_e(2j_e+1)\ \prod_t\begin{Bmatrix}j_1&j_2&j_3\\ j_4&j_5&j_6\end{Bmatrix}_t .
\end{centeredalign}

Here the bracket is the Wigner \(6j\) symbol of the six spins on tetrahedron \(t\). It is the recoupling coefficient of three angular momenta known from atomic spectroscopy. The factor \(2j_e+1\) is the dimension of the spin-\(j_e\) representation. The \(6j\) symbol vanishes unless each triangle obeys the triangle inequality with an integer sum of spins. Only spins that can form real triangles contribute.

The connection with gravity appears at large spins. Ponzano and Regge found the following asymptotic form, and Roberts later proved it:

\begin{centeredalign}
\begin{Bmatrix}j_1&j_2&j_3\\ j_4&j_5&j_6\end{Bmatrix}\simeq
\frac{1}{\sqrt{12\pi V}}\cos\Big(\sum_{e=1}^6 l_e\theta_e+\frac\pi4\Big),\qquad l_e=j_e+\tfrac12.
\end{centeredalign}

Here \(V\) is the volume of the Euclidean tetrahedron with edge lengths \(l_e\). The angle \(\theta_e\) is the exterior dihedral angle at edge \(e\), between the outward normals of the two faces that share it. The sum \(\sum_el_e\theta_e\) is the Regge action of the tetrahedron. Regge calculus is the discrete form of general relativity, with curvature concentrated on edges as deficit angles. The amplitude is therefore a sum of \(e^{iS}\) and \(e^{-iS}\) for the two orientations. This is the form expected of a gravitational path integral in the semiclassical limit.

For a regular tetrahedron with all spins equal to \(N\), the volume is \(V=\frac{l^3}{6\sqrt2}\) with \(l=N+1/2\). Every exterior angle is \(\pi-\arccos(1/3)\approx1.911\). At \(N=20\) the exact \(6j\) symbol is \(-5.029\times10^{-3}\), and the asymptotic formula gives \(-5.034\times10^{-3}\). At \(N=40\) the two agree to three parts in \(10^4\). The sign changes with \(N\) follow the phase \(6l\theta\) of the cosine.

In three dimensions gravity has no local propagating degrees of freedom, and the Ponzano--Regge model is a quantization of it. The sum over spins diverges for triangulations with interior vertices. The divergence reflects a residual symmetry and is removed by gauge fixing. With a positive cosmological constant the model is replaced by the finite Turaev--Viro state sum. In four dimensions, models such as that of Engle, Pereira, Rovelli and Livine assign amplitudes to four-simplices. Their large-spin limit again contains the Regge action. Whether sums over many simplices reproduce general relativity, with its gravitational waves, is not established. The boundary states of a spin foam are the spin networks of loop quantum gravity, and each tetrahedron is the quantum tetrahedron described under quantum geometry.

\subsection*{References For The Physical Derivation}
\begin{itemize}
\item \href{https://arxiv.org/abs/gr-qc/9905087}{J. C. Baez, An Introduction to Spin Foam Models of Quantum Gravity and BF Theory; spin foams as histories of spin networks and the Ponzano-Regge model.}
\item \href{https://arxiv.org/abs/math-ph/9812013}{J. Roberts, Classical 6j-symbols and the tetrahedron; proof of the Ponzano-Regge asymptotic formula.}
\item \href{https://arxiv.org/abs/1205.2019}{A. Perez, The Spin Foam Approach to Quantum Gravity; three- and four-dimensional models and their semiclassical limits.}
\end{itemize}





\clearpage

\section{AdS/CFT correspondence}

\label{topic:ads-cft-correspondence}

A scalar field in anti-de Sitter space does not radiate away to infinity. Near the boundary its amplitude falls off as a definite power of a radial coordinate. The AdS/CFT correspondence identifies that power with the scaling dimension of an operator in a conformal field theory defined on the boundary. The relation follows from the bulk wave equation alone. In Poincaré coordinates the metric of \((d+1)\)-dimensional anti-de Sitter space with curvature radius \(L\) is

\begin{centeredalign}
ds^2=\frac{L^2}{z^2}\big(dz^2+\eta_{\mu\nu}dx^\mu dx^\nu\big),\qquad z>0,
\end{centeredalign}

where \(x^\mu\) are the \(d\) coordinates of the boundary, which is reached as \(z\to0\). A free scalar of mass \(m\) obeys \((\Box-m^2)\phi=0\), with \(\Box\phi=g^{-1/2}\partial_a(g^{1/2}g^{ab}\partial_b\phi)\). Near the boundary the derivatives along \(x\) carry an extra factor \(z^2\) relative to the radial terms. A power law \(\phi\propto z^\Delta\) is therefore tested against the radial part only. Using \(\sqrt g=(L/z)^{d+1}\) and \(g^{zz}=z^2/L^2\) gives

\begin{centeredalign}
\Box\,z^\Delta=\frac{z^{d+1}}{L^2}\,\partial_z\big(z^{1-d}\partial_z z^\Delta\big)
=\frac{\Delta(\Delta-d)}{L^2}\,z^\Delta
\qquad\Rightarrow\qquad \Delta(\Delta-d)=m^2L^2.
\end{centeredalign}

The two roots are \(\Delta_\pm=d/2\pm\sqrt{d^2/4+m^2L^2}\), and their sum is \(d\). A general solution behaves as \(\phi\to z^{d-\Delta}\phi_0(x)+z^{\Delta}A(x)\) with \(\Delta=\Delta_+\). The coefficient \(\phi_0\) of the slower falloff is fixed by a boundary condition. The rescaling \(z\to\lambda z\), \(x\to\lambda x\) is an isometry of the metric, and under it \(\phi_0\) transforms as a field of dimension \(d-\Delta\). The correspondence treats \(\phi_0\) as a source \(J\) coupled to a boundary operator \(\mathcal O\). The coupling \(\int d^dx\,J\mathcal O\) is dimensionless only if \(\mathcal O\) has dimension \(\Delta\). The dictionary is an equality of generating functionals:

\begin{centeredalign}
Z_{\rm CFT}[J]=\Big\langle \exp\int d^dx\,J\,\mathcal O\Big\rangle_{\rm CFT}
=Z_{\rm bulk}\big[\phi\to z^{d-\Delta}J\big].
\end{centeredalign}

Both roots are real only if \(m^2L^2\geq-d^2/4\). This is the Breitenlohner--Freedman bound. A field with slightly negative mass squared is stable in anti-de Sitter space, because the curvature confines it and its gradient energy outweighs the negative mass term. In the conformal field theory \(\Delta\) fixes the power law of the two-point function, \(\langle\mathcal O(x)\mathcal O(0)\rangle\propto|x|^{-2\Delta}\). The bulk calculation reproduces this from the propagation of \(\phi\) between two boundary points.

The best-tested case pairs string theory on \(AdS_5\times S^5\) with \(\mathcal N=4\) supersymmetric \(SU(N)\) gauge theory, so \(d=4\). The dilaton is massless in five dimensions, so \(\Delta=4\). It couples to the Lagrangian density \(\operatorname{Tr}F^2\) of the gauge theory, whose dimension is four. Kaluza--Klein modes on the five-sphere give scalars with \(m^2L^2=k(k-4)\) for integers \(k\geq2\). The relation gives \(\Delta=2+|k-2|=k\), which matches symmetrized traces of \(k\) scalar fields of the gauge theory. The \(k=2\) mode sits exactly at the bound, \(m^2L^2=-4\), and is stable. Supersymmetry protects these dimensions, so the free-field value can be compared with the classical bulk value.

The classical bulk calculation requires \(L\) to be large compared with the string length and the Planck length. For the \(\mathcal N=4\) theory this means large \(N\) and large 't Hooft coupling \(\lambda=g_{\rm YM}^2N\), since \(L^4/\alpha'^2=\lambda\). Away from that regime string and quantum corrections enter. The correspondence remains a conjecture: it has passed many quantitative checks of protected spectra and correlators, but it has no general proof. Because the bulk contains gravity, it is used to study quantum gravity in anti-de Sitter space through an ordinary gauge theory.

\subsection*{References For The Physical Derivation}
\begin{itemize}
\item \href{https://arxiv.org/abs/hep-th/9905111}{O. Aharony, S. S. Gubser, J. Maldacena, H. Ooguri and Y. Oz, Large N Field Theories, String Theory and Gravity; mass-dimension relation, field-operator dictionary and the AdS5 x S5 spectrum.}
\item \href{https://arxiv.org/abs/hep-th/9802150}{E. Witten, Anti De Sitter Space And Holography; boundary values of bulk fields as sources of boundary operators.}
\item \href{https://arxiv.org/abs/hep-th/9802109}{S. S. Gubser, I. R. Klebanov and A. M. Polyakov, Gauge Theory Correlators from Non-Critical String Theory; generating functional of boundary correlators.}
\end{itemize}





\clearpage

\section{Loop quantum gravity}

\label{topic:loop-quantum-gravity}

If space is quantized, the area of a surface should take discrete values, as the energy of an atom does. Loop quantum gravity gives a definite spectrum for area. It quantizes general relativity written with an \(SU(2)\) connection \(A^i_a\) and its conjugate momentum \(E^a_i\), a densitized triad that encodes the spatial metric. In units with \(c=1\), the Poisson bracket and the area of a surface \(S\) are

\begin{centeredalign}
\{A^i_a(x),E^b_j(y)\}=8\pi G\gamma\,\delta^b_a\delta^i_j\,\delta^3(x-y),\\
A(S)=\int_Sd^2\sigma\,\sqrt{n_aE^a_i\,n_bE^b_i}.
\end{centeredalign}

Here \(\gamma\) is the Barbero--Immirzi parameter, \(a\) a spatial index, \(i\) an internal index and \(n_a\) the normal to \(S\). The states are functions of holonomies, the path-ordered exponentials of \(A\) along the edges of a graph \(\Gamma\). Each edge carries a spin \(j\). Quantization turns \(E\) into \(-i\,8\pi\gamma l_P^2\,\delta/\delta A\), with \(l_P^2=\hbar G\). On an edge crossing \(S\), the flux of \(E\) through \(S\) acts as \(8\pi\gamma l_P^2\) times the \(SU(2)\) generator \(\mathbf J\) of that edge. The square root of \(\mathbf J^2=j(j+1)\) then gives a finite contribution from each crossing:

\begin{centeredalign}
\hat A(S)\,|\Gamma,\{j\}\rangle=8\pi\gamma l_P^2\sum_{p\in S\cap\Gamma}\sqrt{j_p(j_p+1)}\;|\Gamma,\{j\}\rangle.
\end{centeredalign}

The sum runs over the points where edges puncture the surface. A surface crossed by no edge has zero area. The spectrum is discrete, and its smallest nonzero eigenvalue comes from one spin-\(1/2\) puncture: \(a_0=4\sqrt3\pi\gamma l_P^2\). The formula holds for punctures by edges; nodes lying on the surface add further terms. The result is kinematical. The operator is defined before the Hamiltonian constraint is solved, and whether its spectrum survives in the full dynamics is debated.

The parameter \(\gamma\) is the main ambiguity. Classically it labels a family of canonical variables, \(A=\Gamma+\gamma K\), built from the spin connection \(\Gamma\) and the extrinsic curvature \(K\). Every value describes the same general relativity. After quantization the area spectrum scales with \(\gamma\), and different values give inequivalent quantum theories. The classical theory does not fix it.

The usual way to fix \(\gamma\) is to count horizon states. A black-hole horizon of area \(A\) is punctured by many edges. Counting the spin configurations compatible with that area gives an entropy proportional to \(A/\gamma\). Requiring the Bekenstein--Hawking value \(S=k_BA/4l_P^2\) yields an equation for \(\gamma\). In the counting of Domagala and Lewandowski, and of Meissner, it reads

\begin{centeredalign}
1=\sum_{j=1/2,\,1,\,3/2,\ldots}2\,e^{-2\pi\gamma\sqrt{j(j+1)}}
\qquad\Rightarrow\qquad \gamma=0.2375\ldots,\qquad a_0\approx5.2\,l_P^2.
\end{centeredalign}

With this value the smallest area is about \(1.4\times10^{-69}\) m\(^2\). An earlier count kept only spin-\(1/2\) punctures and gave \(\gamma=\ln2/(\pi\sqrt3)\approx0.127\). Other counting rules give other values, and the logarithmic correction to the entropy also depends on the rule. The procedure therefore matches one parameter; it does not predict the entropy. No experiment has resolved areas at this scale. Within the theory, the discreteness of area and volume for spin-network states is established, as described under spin network. The dynamics is formulated through spin foams and remains open.

\subsection*{References For The Physical Derivation}
\begin{itemize}
\item \href{https://arxiv.org/abs/gr-qc/0404018}{A. Ashtekar and J. Lewandowski, Background Independent Quantum Gravity: A Status Report; connection variables, the Barbero-Immirzi parameter and the area operator.}
\item \href{https://arxiv.org/abs/gr-qc/9411005}{C. Rovelli and L. Smolin, Discreteness of area and volume in quantum gravity; discrete spectra of geometric operators on spin-network states.}
\item \href{https://arxiv.org/abs/gr-qc/0407051}{M. Domagala and J. Lewandowski, Black hole entropy from Quantum Geometry; corrected horizon state counting and the resulting value of the Barbero-Immirzi parameter.}
\end{itemize}





\clearpage

\section{Spin network}

\label{topic:spin-network}

In a lattice gauge theory, a line of electric flux gives a definite flux through every surface it pierces and none through surfaces it misses. Loop quantum gravity gives space the same structure. Area belongs to lines that pierce a surface, and volume to points where several lines meet. The states with this property are spin networks, built like lattice gauge fields but with the group \(SU(2)\) on an arbitrary graph. Each edge \(e\) of a graph \(\Gamma\) carries a spin \(j_e\). Each node \(n\) carries an intertwiner \(\iota_n\), a tensor that couples the spins meeting there to total spin zero. The state depends on the connection \(A\) through the holonomies \(h_e\) along the edges:

\begin{centeredalign}
\Psi_{\Gamma,j,\iota}[A]=\Big(\bigotimes_n\iota_n\Big)\cdot\Big(\bigotimes_e D^{(j_e)}\big(h_e[A]\big)\Big),\qquad
h_e=\mathcal P\exp\int_eA.
\end{centeredalign}

Here \(D^{(j)}\) is the spin-\(j\) matrix of a group element, and the dot contracts each matrix index with an index of an intertwiner. Gauge transformations act at the nodes and rotate the indices of all edges meeting there. The state is invariant because each intertwiner is a rotation-invariant tensor. Whether such a tensor exists, and how many independent ones exist, follows from the addition of angular momenta.

The simplest graph is a closed loop with one two-valent node. Two spins couple to total spin zero only if they are equal, and then in exactly one way. The intertwiner is the identity map between the two edges, so the node can be removed. The state is the Wilson loop \(\operatorname{Tr}D^{(j)}(h_\gamma)\) of the spin-\(j\) representation around the loop \(\gamma\). An edge of spin \(1/2\) cannot join an edge of spin \(1\) at a two-valent node: no invariant exists. The spin along a line therefore changes only where a third edge meets it. At a three-valent node the spins must satisfy \(|j_1-j_2|\leq j_3\leq j_1+j_2\) with an integer sum. The intertwiner is again unique, a Clebsch--Gordan coefficient. Four spin-\(1/2\) edges admit two independent intertwiners.

Geometric operators act on these states locally. For a surface \(S\) crossed by edges at points \(p\), with no node on the surface, the area is diagonal. The volume of a region \(R\) is a sum of contributions from the nodes inside it:

\begin{centeredalign}
\hat A(S)\,\Psi=8\pi\gamma l_P^2\sum_{p\in S\cap\Gamma}\sqrt{j_p(j_p+1)}\;\Psi,\qquad
\hat V(R)\,\Psi=\sum_{n\in R}\hat V_n\Psi.
\end{centeredalign}

Here \(\gamma\) is the Barbero--Immirzi parameter and \(l_P\) the Planck length. Each crossing contributes the square root of the Casimir operator of its edge. Consider a loop of spin \(j\) that crosses a closed surface twice. The area eigenvalue is \(16\pi\gamma l_P^2\sqrt{j(j+1)}\). For \(j=1/2\) and \(\gamma=0.2375\) this is \(8\sqrt3\pi\gamma l_P^2\approx10.3\,l_P^2\). The volume of the loop is zero. A node contributes volume only if at least four edges meet there with tangents not all in one plane. For three edges, the triple product of their generators vanishes on the unique invariant state.

The loop therefore carries area but no volume. It is a flux line of geometry, in the same way that a Wilson line in lattice gauge theory is a line of electric flux. The triad measured by the area plays the part of the electric field, and its square gives the Casimir. These statements are kinematical. They hold on the state space before the Hamiltonian constraint is imposed. Diffeomorphism invariance makes only the combinatorics and knotting of the graph physical, not its coordinates. The four-valent node is worked out as the quantum tetrahedron under quantum geometry. The area spectrum and the value of \(\gamma\) are discussed under loop quantum gravity.

\subsection*{References For The Physical Derivation}
\begin{itemize}
\item \href{https://arxiv.org/abs/gr-qc/9505006}{C. Rovelli and L. Smolin, Spin Networks and Quantum Gravity; spin-network basis of loop quantum gravity.}
\item \href{https://arxiv.org/abs/gr-qc/9411007}{J. C. Baez, Spin Network States in Gauge Theory; holonomies, intertwiners and gauge-invariant states on graphs.}
\item \href{https://arxiv.org/abs/gr-qc/9411005}{C. Rovelli and L. Smolin, Discreteness of area and volume in quantum gravity; area and volume spectra on spin-network states.}
\end{itemize}





\clearpage

\section{Quantum chemistry}

\label{topic:quantum-chemistry}

A proton and a hydrogen atom bind into the molecular ion \(\mathrm H_2^+\). The bond is held by a single electron, and the depth of the potential well is 2.79 eV. The binding comes from the wave nature of the electron: its amplitudes near the two nuclei interfere. The simplest calculation that shows this, the linear combination of atomic orbitals, is also the starting point of molecular-orbital theory.

Protons are 1836 times heavier than electrons and move much more slowly. In the Born--Oppenheimer approximation they are held at a fixed separation \(R\). The electronic energy at each \(R\), including the proton repulsion, then serves as the potential for nuclear motion. In atomic units the electronic Hamiltonian is \(H=-\tfrac12\nabla^2-1/r_a-1/r_b+1/R\), with \(r_a\) and \(r_b\) the distances to the two protons. The trial states combine the hydrogen 1s orbitals \(\phi_a\) and \(\phi_b\) centred on each proton:

\begin{centeredalign}
\psi_\pm=\frac{\phi_a\pm\phi_b}{\sqrt{2(1\pm S)}},\qquad
E_\pm(R)=-\frac12+\frac1R-\frac{J\pm K}{1\pm S}.
\end{centeredalign}

The variational principle selects these combinations. The \(2\times2\) matrices of \(H\) and of the overlap in the basis \(\phi_a,\phi_b\) are unchanged when the protons are interchanged, so their eigenvectors are the even and odd combinations. Each \(E_\pm\) is an upper bound to the exact energy of that symmetry. The three integrals are the overlap \(S=\langle\phi_a|\phi_b\rangle\), the attraction \(J=\langle\phi_a|r_b^{-1}|\phi_a\rangle\) of one atomic cloud to the other proton, and the resonance integral \(K=\langle\phi_a|r_a^{-1}|\phi_b\rangle\). The last exists only because the orbitals overlap:

\begin{centeredalign}
S=e^{-R}\left(1+R+\frac{R^2}{3}\right),\qquad
J=\frac1R-e^{-2R}\left(1+\frac1R\right),\qquad
K=e^{-R}(1+R).
\end{centeredalign}

In the even, bonding state the electron density accumulates between the protons, where both attract it. In the odd, antibonding state it has a node on the midplane. Minimizing \(E_+(R)\) gives an equilibrium separation of 2.49 bohr, or 0.132 nm, and a well depth of 1.76 eV below a hydrogen atom and a free proton. The exact Born--Oppenheimer solution gives 2.00 bohr, or 0.106 nm, and 2.79 eV. Letting the orbital exponent vary with \(R\) already improves the result to 2.00 bohr and 2.35 eV.

The antibonding curve \(E_-(R)\) is repulsive at all separations. Light that promotes the electron from the bonding to the antibonding state therefore dissociates the ion into a proton and a hydrogen atom. The two-orbital problem is a Two-state quantum system: the coupling between two localized states splits them into symmetric and antisymmetric combinations, as at an avoided crossing. The Born--Oppenheimer separation fails where two electronic energy surfaces approach each other, as in many photochemical reactions.

In \(\mathrm H_2\) two electrons occupy the bonding orbital with opposite spins, in a singlet. The antisymmetry discussed under Fermion forces additional electrons into the antibonding orbital. With four electrons, as in \(\mathrm{He}_2\), bonding and antibonding contributions cancel and no chemical bond forms. Electron repulsion makes the single-orbital description of \(\mathrm H_2\) inaccurate at large separation, where it gives equal weight to ionic and covalent configurations.

Configuration interaction and coupled-cluster methods add this correlation systematically. Density-functional theory includes it through an approximate functional of the electron density. Each computes the same Born--Oppenheimer energy surface. The cost of an exact calculation grows exponentially with the number of orbitals, so strongly correlated molecules, such as many transition-metal complexes, remain difficult. Computing such energies is one proposed application of quantum computers; demonstrations so far concern small molecules.

\subsection*{References For The Physical Derivation}
\begin{itemize}
\item \href{https://arxiv.org/abs/1808.10402}{S. McArdle, S. Endo, A. Aspuru-Guzik, S. Benjamin and X. Yuan, Quantum computational chemistry; the Born-Oppenheimer electronic structure problem, basis sets, correlation methods and quantum algorithms.}
\end{itemize}





\clearpage

\section{Standard Model}

\label{topic:standard-model}

The masses of the W and Z bosons, 80.37 and 91.19 GeV in units with \(c=1\), are not independent quantities. Their ratio is fixed by one angle, which also determines how the Z boson couples to electrons and neutrinos. This constraint, tested at electron-positron and hadron colliders, follows from the gauge structure and particle content of the Standard Model.

The Standard Model is a gauge theory with group \(SU(3)_c\times SU(2)_L\times U(1)_Y\). The factor \(SU(3)_c\) acts on colour and is the subject of Quantum chromodynamics. The electroweak factors act differently on the two chiralities of each fermion: left-handed fields form \(SU(2)_L\) doublets and right-handed fields are singlets. One generation, listed with its representation of \(SU(3)_c\times SU(2)_L\) and hypercharge \(Y\) as a subscript, is

\begin{centeredalign}
Q_L=\begin{pmatrix}u_L\\d_L\end{pmatrix}:(3,2)_{1/6},\qquad
u_R:(3,1)_{2/3},\qquad d_R:(3,1)_{-1/3},\\
L_L=\begin{pmatrix}\nu_L\\e_L\end{pmatrix}:(1,2)_{-1/2},\qquad
e_R:(1,1)_{-1},\qquad H:(1,2)_{1/2}.
\end{centeredalign}

Two further generations repeat the pattern with larger masses. The electric charge is \(Q=T_3+Y\), with \(T_3=\pm\frac12\) the weak isospin of a doublet member. This gives \(2/3\) and \(-1/3\) for the quarks and \(0\) and \(-1\) for the leptons. The hypercharges are tightly constrained. Counting colours and writing right-handed fields as left-handed antiparticles, the sums of \(Y\) and of \(Y^3\) over one generation both vanish. These cancellations of gauge anomalies, required for a consistent quantum theory, work only between quarks and leptons together.

The Higgs doublet \(H\) has a vacuum expectation value \(v/\sqrt2\) in its neutral component. Its kinetic energy then gives mass to three combinations of the gauge fields, the \(W^\pm\) and the \(Z\), and leaves the photon massless. With \(g\) and \(g'\) the \(SU(2)_L\) and \(U(1)_Y\) couplings and \(\theta_W\) the weak mixing angle,

\begin{centeredalign}
m_W=\frac{gv}{2},\qquad m_Z=\frac{\sqrt{g^2+g'^2}\,v}{2},\qquad
\tan\theta_W=\frac{g'}{g},\\
e=g\sin\theta_W,\qquad m_W=m_Z\cos\theta_W.
\end{centeredalign}

Muon decay measures the Fermi constant \(G_F\), and W exchange gives \(G_F/\sqrt2=g^2/8m_W^2\). Hence \(v=(\sqrt2G_F)^{-1/2}=246.2\) GeV. Eliminating \(g\) with \(e^2=4\pi\alpha\) gives a prediction for the W mass:

\begin{centeredalign}
m_W^2=\frac{\pi\alpha}{\sqrt2\,G_F\sin^2\theta_W},\qquad
m_Z=\frac{m_W}{\cos\theta_W}.
\end{centeredalign}

Take \(G_F=1.1664\times10^{-5}\,\mathrm{GeV^{-2}}\), \(\sin^2\theta_W=0.231\) from Z-pole asymmetries, and the coupling at the Z mass, \(\alpha\approx1/128\). The formula then gives \(m_W\approx80.3\) GeV and \(m_Z\approx91.5\) GeV, both within half a percent of the measured values. With the low-energy value \(\alpha=1/137\) the same formula gives 77.5 GeV, so the running described under Renormalization is required. The remaining difference is a loop correction dominated by the top quark and proportional to \(G_Fm_t^2\). Its size indicated the top-quark mass before the quark was observed in 1995.

With 19 parameters, the model accounts for collider measurements of electroweak and strong processes. Its tree-level relations hold at the percent level, and tests at the per-mille level require loop corrections, so they test the quantum theory itself. Neutrino oscillations show that neutrinos have mass, which the minimal model with only left-handed neutrinos lacks. Gravity, dark matter and the excess of matter over antimatter also lie outside it.

\subsection*{References For The Physical Derivation}
\begin{itemize}
\item \href{https://arxiv.org/abs/0901.0241}{P. Langacker, Introduction to the Standard Model and Electroweak Physics; gauge group, fermion representations, gauge-boson masses and precision tests.}
\end{itemize}





\clearpage

\section{Klein–Gordon equation}

\label{topic:klein-gordon-equation}

Nuclear forces fall to negligible strength within a few femtometres, whereas the Coulomb force has unlimited range. In 1935 Yukawa attributed the short range to the mass of an exchanged particle, later identified as the pion. The relation between mass and range follows from the Klein--Gordon equation, the relativistic wave equation for a particle without spin.

Replacing \(E\) by \(i\hbar\partial_t\) and \(\mathbf p\) by \(-i\hbar\nabla\) in \(E^2=p^2c^2+m^2c^4\) gives an equation for a field \(\phi(\mathbf x,t)\). Plane waves solve it for either sign of the energy:

\begin{centeredalign}
\left(\frac{1}{c^2}\frac{\partial^2}{\partial t^2}-\nabla^2+\frac{m^2c^2}{\hbar^2}\right)\phi=0,\\
\phi=e^{i(\mathbf p\cdot\mathbf x-Et)/\hbar},\qquad E=\pm\sqrt{p^2c^2+m^2c^4}.
\end{centeredalign}

The negative-energy solutions have no lower bound. A particle described by them could, once coupled to the electromagnetic field, radiate without limit. A second difficulty concerns probability. The equation implies a continuity equation \(\partial_t\rho+\nabla\cdot\mathbf j=0\), with

\begin{centeredalign}
\rho=\frac{i\hbar}{2mc^2}\left(\phi^*\partial_t\phi-\phi\,\partial_t\phi^*\right),\qquad
\mathbf j=\frac{\hbar}{2mi}\left(\phi^*\nabla\phi-\phi\nabla\phi^*\right),\\
\rho=\frac{E}{mc^2}|\phi|^2\quad\text{for a plane wave}.
\end{centeredalign}

For negative-energy plane waves \(\rho\) is negative. More generally, because the equation is second order in time, \(\phi\) and \(\partial_t\phi\) can be chosen independently at an initial instant, and \(\rho\) can be made negative anywhere. It is therefore not a probability density. Pauli and Weisskopf resolved this in 1934 by treating \(\phi\) as a quantized field. Multiplied by the charge, \(\rho\) becomes a charge density whose sign distinguishes particles from antiparticles. Positive-frequency solutions describe \(\pi^+\) mesons, and negative-frequency solutions describe \(\pi^-\) mesons with positive energy.

For a static point source of strength \(g\) the field obeys \((\nabla^2-\mu^2)\phi=-g\,\delta^3(\mathbf r)\), with \(\mu=mc/\hbar\). The solution falls off exponentially, and a second nucleon coupled to the same field acquires a potential energy proportional to it:

\begin{centeredalign}
\phi(r)=\frac{g\,e^{-\mu r}}{4\pi r},\qquad
V(r)\propto-\frac{e^{-r/\lambda}}{r},\qquad \lambda=\frac{\hbar}{mc}.
\end{centeredalign}

The range is the reduced Compton wavelength of the exchanged particle. With the charged-pion mass \(139.57\,\mathrm{MeV}/c^2\) and \(\hbar c=197.33\) MeV fm, \(\lambda=1.41\) fm. Nucleon-nucleon scattering in partial waves of high angular momentum, which probe large separations, follows the one-pion-exchange potential. Because the pion is pseudoscalar, that potential also contains spin and isospin factors and a tensor part, but its central part has this Yukawa form. Below about 1 fm, heavier mesons and the quark structure of nucleons take over.

The Klein--Gordon equation with a Coulomb potential also describes pionic atoms, in which a \(\pi^-\) replaces an electron. Their X-ray transition energies are used to determine the charged-pion mass. As a single-particle equation it applies when potentials change slowly over \(\lambda\) and energies remain well below the pair-creation threshold \(2mc^2\). For particles of spin one half it is replaced by the Dirac equation, although each component of a free Dirac spinor still satisfies it. In Quantum field theory every free relativistic field obeys it, as the dispersion relation of each mode.

\subsection*{References For The Physical Derivation}
\begin{itemize}
\item \href{https://arxiv.org/abs/1105.2919}{R. Machleidt and D. R. Entem, Chiral effective field theory and nuclear forces; one-pion exchange as the long-range part of the nucleon-nucleon interaction.}
\item \href{https://arxiv.org/abs/1605.03300}{M. Trassinelli and colleagues, Measurement of the charged pion mass using X-ray spectroscopy of exotic atoms; pionic-atom transitions and the charged-pion mass.}
\end{itemize}





\clearpage

\section{Bose–Einstein statistics}

\label{topic:bose-einstein-statistics}

Rubidium atoms cooled in a magnetic trap to a few hundred nanokelvin collect abruptly in the lowest quantum state of the trap. After the trap is switched off, a narrow peak appears in the velocity distribution above a broad thermal cloud. The temperature at which this happens follows from counting how many atoms the excited states can hold. For noninteracting bosons at temperature \(T\) and chemical potential \(\mu\), Quantum statistical mechanics gives the mean occupation of a single-particle state of energy \(\epsilon\) as \(\overline n_B=[e^{\beta(\epsilon-\mu)}-1]^{-1}\), with \(\beta=1/k_BT\).

Occupations cannot be negative, so \(\mu\) must stay below the lowest single-particle energy, set here to zero. At a fixed \(\mu<0\) every excited state holds a finite mean number of particles. For a uniform gas of spinless particles of mass \(m\), the excited states can therefore hold at most the density reached as \(\mu\to0\):

\begin{centeredalign}
n_{\rm ex}\leq\int\frac{d^3p}{(2\pi\hbar)^3}\,\frac{1}{e^{\beta p^2/2m}-1}
=\frac{\zeta(3/2)}{\lambda_T^3},\qquad
\lambda_T=\sqrt{\frac{2\pi\hbar^2}{mk_BT}}.
\end{centeredalign}

Here \(\zeta(3/2)\approx2.612\) is a value of the Riemann zeta function and \(\lambda_T\) is the thermal de Broglie wavelength. If the density \(n\) exceeds this bound, the surplus cannot be accommodated in excited states. It accumulates in the ground state, whose occupation \((e^{-\beta\mu}-1)^{-1}\) grows without limit as \(\mu\) approaches zero from below. The condition \(n\lambda_T^3=\zeta(3/2)\) defines the transition temperature, and the same bound gives the condensed fraction below it:

\begin{centeredalign}
k_BT_c=\frac{2\pi\hbar^2}{m}\left(\frac{n}{\zeta(3/2)}\right)^{2/3},\qquad
\frac{N_0}{N}=1-\left(\frac{T}{T_c}\right)^{3/2}.
\end{centeredalign}

Condensation thus begins when the thermal wavelength becomes comparable with the spacing \(n^{-1/3}\) between particles. No attraction between the atoms is involved. The cause is the counting of symmetric states: a bosonic mode can hold any number of particles, and permuting identical particles produces no new state. With Maxwell--Boltzmann weights the corresponding relation is \(n\lambda_T^3=e^{\beta\mu}\), which has no upper bound, so excited states accommodate any density.

London applied this result to liquid helium-4 in 1938. A mass density of \(0.145\,\mathrm{g\,cm^{-3}}\) gives \(n=2.18\times10^{28}\,\mathrm{m^{-3}}\) and \(T_c\approx3.1\) K, close to the observed superfluid transition at 2.17 K. The difference reflects the strong interactions in the liquid, for which ideal-gas counting is only a guide. Neutron scattering finds that less than ten percent of the atoms occupy the zero-momentum state even near absolute zero.

Dilute alkali gases approach the ideal limit much more closely. In a harmonic trap with geometric-mean angular frequency \(\bar\omega\), the density of states grows as \(\epsilon^2\) instead of \(\sqrt\epsilon\). The same counting then gives

\begin{centeredalign}
k_BT_c=\hbar\bar\omega\left(\frac{N}{\zeta(3)}\right)^{1/3}
\approx0.94\,\hbar\bar\omega N^{1/3},\qquad
\frac{N_0}{N}=1-\left(\frac{T}{T_c}\right)^3.
\end{centeredalign}

For \(N=10^6\) atoms and \(\bar\omega=2\pi\times100\,\mathrm{s^{-1}}\) this gives \(T_c\approx450\) nK. Condensation was observed in rubidium and sodium in 1995 at transition temperatures between about 0.2 and 2 microkelvin. The results assume many particles and weak interactions. A finite particle number lowers \(T_c\) by a relative amount of order \(N^{-1/3}\), and repulsive interactions shift it by a few percent in traps. In a uniform two-dimensional gas the corresponding integral diverges as \(\mu\to0\), and no condensation occurs at finite temperature.

Photons and phonons in thermal equilibrium are not conserved in number. Their chemical potential is therefore zero, and the same occupation gives Planck's law instead of a condensate. Reversing the sign in the denominator gives Fermi--Dirac statistics, which fills states one by one up to a Fermi energy. The symmetric exchange rule itself is discussed under Boson.

\subsection*{References For The Physical Derivation}
\begin{itemize}
\item \href{https://arxiv.org/abs/cond-mat/9806038}{F. Dalfovo, S. Giorgini, L. P. Pitaevskii and S. Stringari, Theory of Bose-Einstein condensation in trapped gases; ideal-gas transition temperature and condensate fraction in uniform and harmonic traps, with finite-size and interaction shifts.}
\end{itemize}





\clearpage

\section{String theory}

\label{topic:string-theory}

Mesons of increasing spin \(J\) lie close to straight lines when \(J\) is plotted against mass squared. A rotating relativistic string with its ends moving at the speed of light gives this relation. String theory begins from the quantum mechanics of such a string, and its best-established results concern the spectrum. A string of tension \(T\) sweeps out a worldsheet \(X^\mu(\tau,\sigma)\) in \(D\)-dimensional spacetime. With \(\hbar=c=1\), the Nambu--Goto action is proportional to the area of that worldsheet:

\begin{centeredalign}
S=-T\int d\tau\,d\sigma\,\sqrt{-\det\big(\partial_\alpha X^\mu\,\partial_\beta X_\mu\big)},\qquad
T=\frac{1}{2\pi\alpha'}.
\end{centeredalign}

Consider a straight open string of length \(2R\) rotating rigidly, with free ends moving at the speed of light. Integrating the relativistic energy and angular momentum densities gives \(E=\pi TR\) and \(J=\pi TR^2/2\), so \(J=\alpha'E^2\). The mesons \(\rho(775)\), \(a_2(1320)\) and \(\rho_3(1690)\) fit \(J\approx0.46+0.89M^2\), with \(M\) in GeV. The slope corresponds to a tension of about \(0.9\) GeV per femtometre, close to the flux-tube tension of quantum chromodynamics.

Quantization is simplest in light-cone gauge. It removes the two longitudinal directions and leaves \(D-2\) transverse oscillators \(\alpha^i_n\) for each harmonic \(n\). For an open string with free ends the mass is

\begin{centeredalign}
\alpha'M^2=N+a,\qquad
N=\sum_{n\geq1}\sum_{i=1}^{D-2}\alpha^i_{-n}\alpha^i_n,\qquad
a=\frac{D-2}{2}\sum_{n\geq1}n\ \to\ -\frac{D-2}{24}.
\end{centeredalign}

The constant \(a\) is the zero-point energy of the oscillators, half a quantum per mode, summed over harmonics. With a cutoff, \(\sum_nne^{-\epsilon n}=\epsilon^{-2}-1/12+O(\epsilon^2)\). The divergent term is removed by a local counterterm on the worldsheet. The finite remainder \(-1/12\) equals the value \(\zeta(-1)\) of the Riemann zeta function. The same subtraction gives the Casimir energy of a field between two plates.

The lowest excited level, \(N=1\), consists of the states \(\alpha^i_{-1}|0;k\rangle\), with one vector index taking \(D-2\) values. A massive vector in \(D\) dimensions has \(D-1\) polarizations; only a massless one has \(D-2\). Lorentz invariance, which light-cone gauge does not make manifest, therefore requires \(1+a=0\):

\begin{centeredalign}
\frac{D-2}{24}=1\;\Rightarrow\;D=26,\qquad
\alpha'M^2=-1\ \ (N=0),\qquad \alpha'M^2=0\ \ (N=1).
\end{centeredalign}

A complete check confirms this: the Lorentz generators built from the oscillators close only for \(D=26\) and \(a=-1\). At level \(N=2\) the \(324\) states form a symmetric traceless tensor of \(SO(25)\), the rotation group of a massive particle at rest. The ground state is a tachyon, which signals that the bosonic string in flat space is unstable.

Established results are of this kind: consistency conditions and spectra of a perturbative quantum theory. Closed strings contain a massless spin-two state whose low-energy interactions are those of general relativity. That is why the theory is studied as a theory of quantum gravity. With worldsheet fermions the critical dimension becomes ten, and a projection of the spectrum removes the tachyon. For certain supersymmetric black holes, counting string states reproduces the Bekenstein--Hawking entropy. No string excitation or extra dimension has been detected, and the string length is unknown. The choice among many four-dimensional solutions is open. The AdS/CFT correspondence supplies a nonperturbative formulation in anti-de Sitter backgrounds.

\subsection*{References For The Physical Derivation}
\begin{itemize}
\item \href{https://arxiv.org/abs/0908.0333}{D. Tong, Lectures on String Theory; Nambu-Goto action, light-cone quantization and the critical dimension.}
\item \href{https://arxiv.org/abs/hep-th/9411028}{J. Polchinski, What is String Theory?; open and closed string spectra and the zero-point energy.}
\end{itemize}





\clearpage

\section{Quantum statistical mechanics}

\label{topic:quantum-statistical-mechanics}

Thermal equilibrium assigns probabilities to quantum energy states while respecting the quantities that can be exchanged with the surroundings. For a system that exchanges energy with a large reservoir at temperature \(T\) but has fixed particle number, the canonical state is determined by the Hamiltonian \(H\):

\begin{centeredalign}
\rho_\beta=\frac{e^{-\beta H}}{Z},\qquad
Z=\operatorname{Tr}e^{-\beta H},\qquad \beta=(k_BT)^{-1}.
\end{centeredalign}

The trace sums over the actual many-particle state space, including exchange statistics and boundary conditions. It is not an independent classical distribution imposed on particle positions. Derivatives of the same partition function connect average energy and fluctuations:

\begin{centeredalign}
\langle H\rangle=-\partial_\beta\log Z,\qquad
\operatorname{Var}H=\partial_\beta^2\log Z,\qquad
C=\frac{\operatorname{Var}H}{k_BT^2}.
\end{centeredalign}

The last equality holds for a temperature-independent Hamiltonian in the canonical ensemble. It explains why equilibrium energy fluctuations measure the heat capacity: both are responses of the same probability weights to temperature. A different ensemble constrains different fluctuations and must be analyzed separately.

For independent modes of energy \(\epsilon\) exchanging particles with a reservoir at chemical potential \(\mu\), set \(z=\exp[-\beta(\epsilon-\mu)]\). A fermionic mode contributes \(1+z\) to the grand partition function because only occupations zero and one are allowed. A bosonic mode contributes the geometric sum \((1-z)^{-1}\) when \(z<1\). Differentiation gives

\begin{centeredalign}
\overline n_F=\frac{1}{e^{\beta(\epsilon-\mu)}+1},\qquad
\overline n_B=\frac{1}{e^{\beta(\epsilon-\mu)}-1}.
\end{centeredalign}

The sign difference follows from the allowed state occupations. It produces Fermi filling and Bose enhancement before a specific interaction is added. In the dilute limit both expressions approach the Maxwell--Boltzmann weight. Interactions generally prevent factorization into independent mode contributions, although an effective quasiparticle description may recover it approximately.

For a single harmonic mode, the same sum gives \(Z=[2\sinh(\beta\hbar\omega/2)]^{-1}\) and mean energy \((\hbar\omega/2)\coth(\beta\hbar\omega/2)\). The mean approaches the zero-point energy as temperature tends to zero; quadrature fluctuations persist while the energy variance vanishes. At high temperature the mean energy approaches \(k_BT\). This links the oscillator algebra to thermodynamic response through the choice of preparation.

Equilibrium statistics do not specify the rate of approach to equilibrium. Relaxation requires dynamics, reservoir coupling and sometimes additional conserved quantities. A Gibbs state can describe stationary observables while giving no information about a transport coefficient or memory time. Constructing those predictions requires extending the equilibrium description by the corresponding physical evolution.

\subsection*{References For The Physical Derivation}
\begin{itemize}
\item \href{https://arxiv.org/abs/0706.3360}{S. Giorgini, L. P. Pitaevskii and S. Stringari, Theory of ultracold atomic Fermi gases; ideal and interacting quantum gases.}
\end{itemize}

\subsection*{Relations In The Original Papers}

\href{https://arxiv.org/html/1110.3234v1#S2.E27}{arXiv:1110.3234v1, S2.E27}. A thermal harmonic mode has a geometric distribution of number-state populations. The single-mode bosonic example, not the interacting Fermi-gas discussion.



\clearpage

\section{Quantum gravity}

\label{topic:quantum-gravity}

A black hole of one solar mass is predicted to emit thermal radiation at about sixty nanokelvin. That prediction comes from quantum fields on a classical curved spacetime. A theory in which the geometry itself is quantized does not yet exist in tested form. The scale at which it should matter follows from dimensional analysis. Planck's constant \(\hbar\), Newton's constant \(G\) and the speed of light \(c\) form only one length, one time and one mass:

\begin{centeredalign}
l_P=\sqrt{\frac{\hbar G}{c^3}}=1.62\times10^{-35}\,{\rm m},\qquad
t_P=\frac{l_P}{c}=5.39\times10^{-44}\,{\rm s},\\
m_P=\sqrt{\frac{\hbar c}{G}}=2.18\times10^{-8}\,{\rm kg},\qquad
m_Pc^2=1.22\times10^{19}\,{\rm GeV}.
\end{centeredalign}

Near the Planck mass the Compton wavelength \(\hbar/mc\) of a particle equals its Schwarzschild radius \(2Gm/c^2\), up to a factor of order one. A probe localized to the Planck length would therefore form a black hole around itself. Gravitational corrections to a scattering amplitude at energy \(E\) have relative size \((E/m_Pc^2)^2\). At the 13 TeV of the Large Hadron Collider this ratio is about \(10^{-30}\).

The Hawking temperature can be derived in a few steps. Outside a mass \(M\) the Schwarzschild metric has \(g_{tt}=-(1-r_s/r)c^2\) with \(r_s=2GM/c^2\). In a thermal state the imaginary time \(\tau=it\) is periodic with period \(\hbar/k_BT\). This follows from the Gibbs operator \(e^{-\beta H}\) of quantum statistical mechanics, which evolves states in imaginary time. Near the horizon set \(r=r_s+\rho^2/4r_s\). The Euclidean metric becomes

\begin{centeredalign}
ds^2\simeq\Big(\frac{\rho}{2r_s}\Big)^2c^2d\tau^2+d\rho^2,\qquad
\frac{\hbar}{k_BT_H}=\frac{4\pi r_s}{c}
\;\Rightarrow\; T_H=\frac{\hbar c^3}{8\pi GMk_B}.
\end{centeredalign}

This is a plane in polar coordinates with angle \(c\tau/2r_s\). It is smooth at \(\rho=0\) only if that angle has period \(2\pi\), which fixes the period of \(\tau\) and hence the temperature. Hawking obtained the same result from the mixing of positive and negative frequencies of a quantum field during collapse. For a solar mass \(T_H=6.2\times10^{-8}\) K. This is far below the 2.7 K of the cosmic microwave background, so such holes absorb more radiation than they emit. Only holes lighter than about \(4.5\times10^{22}\) kg, less than a lunar mass, are hotter than the background. The associated Bekenstein--Hawking entropy is \(S=k_BA/4l_P^2\), with \(A\) the horizon area, about \(10^{77}k_B\) for a solar mass.

The calculation treats geometry as classical and matter fields as quantum. It is controlled when the horizon curvature is small on the Planck scale, that is for \(M\gg m_P\). Other established results share this semiclassical status. Colella, Overhauser and Werner measured the gravitational phase shift of neutrons in an interferometer. Treating general relativity as an effective field theory gives calculable quantum corrections to the Newtonian potential at energies far below \(m_Pc^2\), too small to measure.

A theory valid at the Planck scale itself is not established. Perturbative quantization of general relativity needs new counterterms at each loop order, so it loses predictive power there. Open questions include the final stage of evaporation and whether the radiation carries the information of what fell in. The microscopic origin of the entropy is also open in general. String theory, loop quantum gravity and the AdS/CFT correspondence give partial answers in restricted settings. None of them has been tested by experiment.

\subsection*{References For The Physical Derivation}
\begin{itemize}
\item \href{https://arxiv.org/abs/gr-qc/0108040}{S. Carlip, Quantum Gravity: a Progress Report; Planck scales, black hole thermodynamics and the status of the approaches.}
\item \href{https://arxiv.org/abs/gr-qc/0311082}{C. P. Burgess, Quantum Gravity in Everyday Life: General Relativity as an Effective Field Theory; controlled low-energy quantum corrections.}
\end{itemize}





\clearpage

\section{Quantum geometry}

\label{topic:quantum-geometry}

A tetrahedron in flat space is fixed by its six edge lengths. Once its four face areas are known, two further numbers, such as two dihedral angles, specify its shape. Quantum geometry asks what happens when areas and angles become operators. For a single tetrahedron the answer is exact. The face areas can be sharp, but two angles cannot be sharp together when the volume differs from zero.

Describe face \(i\) by its area vector \(\mathbf A_i\), normal to the face with length equal to its area. For a tetrahedron the four outward vectors sum to zero. Take edge vectors \(\mathbf a,\mathbf b,\mathbf c\) from one vertex, so that three faces have \(\mathbf A_1=\frac12\mathbf a\times\mathbf b\), \(\mathbf A_2=\frac12\mathbf b\times\mathbf c\) and \(\mathbf A_3=\frac12\mathbf c\times\mathbf a\). Vector algebra gives \(\mathbf A_1\cdot(\mathbf A_2\times\mathbf A_3)=\frac18[\mathbf a\cdot(\mathbf b\times\mathbf c)]^2\), and the volume is \(V=|\mathbf a\cdot(\mathbf b\times\mathbf c)|/6\). Hence

\begin{centeredalign}
\sum_{i=1}^4\mathbf A_i=0,\qquad
V^2=\frac29\,\big|\mathbf A_1\cdot(\mathbf A_2\times\mathbf A_3)\big|.
\end{centeredalign}

In loop quantum gravity and in spin-foam models each area vector becomes an operator proportional to an angular momentum, \(\mathbf A_i=8\pi\gamma l_P^2\,\mathbf J_i\). Here \(l_P\) is the Planck length and \(\gamma\) the Barbero--Immirzi parameter. Face \(i\) carries a spin \(j_i\), so its area \(8\pi\gamma l_P^2\sqrt{j_i(j_i+1)}\) is sharp. The condition \(\sum_i\mathbf J_i=0\) selects the rotation-invariant states of four spins. Their number is the number of ways to couple the four spins to total spin zero. The products \(\mathbf J_1\cdot\mathbf J_2\) and \(\mathbf J_1\cdot\mathbf J_3\) measure angles between face normals. With \([J^a,J^b]=i\epsilon^{abc}J^c\) and \(\hbar=1\), their commutator is

\begin{centeredalign}
[\mathbf J_1\cdot\mathbf J_2,\,\mathbf J_1\cdot\mathbf J_3]=i\,\mathbf J_1\cdot(\mathbf J_2\times\mathbf J_3)\equiv iQ,\\
V=(8\pi\gamma l_P^2)^{3/2}\sqrt{\tfrac29\,|Q|}.
\end{centeredalign}

Two angles commute only on states of zero volume. A quantum tetrahedron with volume therefore has an uncertainty relation between its dihedral angles, as a spin has between its components. Its classical shape exists only as a mean value. The relative spread of the angles decreases as the spins grow.

The smallest case has four faces of spin \(1/2\). Two pairs of spin-\(1/2\) faces couple to total spin zero in two independent ways, so the shape space is two-dimensional. Let faces 1 and 2 couple to spin \(k\in\{0,1\}\). Then \(\mathbf J_1\cdot\mathbf J_2=[k(k+1)-3/2]/2\) takes the values \(-3/4\) and \(1/4\). The operator \(Q\) has eigenvalues \(\pm\sqrt3/4\), one for each orientation. Both give the volume \(V\approx0.31\,(8\pi\gamma)^{3/2}l_P^3\), about \(4.5\,l_P^3\) for \(\gamma=0.2375\). In either volume eigenstate \(\mathbf J_1\cdot\mathbf J_2\) has mean \(-1/4\) and variance \(1/4\): the angle is maximally uncertain.

The numerical prefactor of the volume depends on how the operator is regularized on a graph. The vanishing of \(Q\) at nodes with fewer than four faces does not. The calculation is kinematical: it concerns one tetrahedron at one instant, before any dynamics of general relativity is imposed. Gluing quantum tetrahedra gives the four-valent nodes of a spin network, and their histories are summed in spin foam models. Whether these discrete spectra describe physical space depends on that dynamics, which remains open, and on the value of \(\gamma\) discussed under loop quantum gravity.

\subsection*{References For The Physical Derivation}
\begin{itemize}
\item \href{https://arxiv.org/abs/gr-qc/9707010}{A. Barbieri, Quantum tetrahedra and simplicial spin networks; area vectors as spin operators, volume formula and the angle uncertainty relation.}
\item \href{https://arxiv.org/abs/gr-qc/9903060}{J. C. Baez and J. W. Barrett, The Quantum Tetrahedron in 3 and 4 Dimensions; Hilbert space of a quantum tetrahedron from geometric quantization.}
\end{itemize}





\clearpage

\section{Quantum spacetime}

\label{topic:quantum-spacetime}

To locate an event within a distance \(\Delta x\), a probe must carry momentum of order \(\hbar/\Delta x\). Its energy also curves spacetime. When \(\Delta x\) approaches the Planck length, the probe would collapse into a black hole larger than the region probed. Heuristic arguments of this kind, together with estimates from string scattering, suggest a smallest resolvable length. Quantum spacetime models encode it by deforming commutators of positions and momenta. None has been derived from a complete theory or confirmed by experiment. They are proposals whose consequences can be calculated and bounded.

A one-dimensional model of Kempf, Mangano and Mann replaces the canonical commutator by a momentum-dependent one. The Robertson relation \(\Delta x\,\Delta p\geq|\langle[x,p]\rangle|/2\), with \(\langle p^2\rangle=(\Delta p)^2+\langle p\rangle^2\), then gives

\begin{centeredalign}
[x,p]=i\hbar\big(1+\beta p^2\big),\qquad
\Delta x\,\Delta p\geq\frac\hbar2\Big(1+\beta(\Delta p)^2+\beta\langle p\rangle^2\Big),
\end{centeredalign}

where \(\beta>0\) has dimensions of inverse momentum squared. For \(\langle p\rangle=0\), dividing by \(\Delta p\) shows that the bound on \(\Delta x\) is smallest at \(\Delta p=1/\sqrt\beta\). No state can have a smaller position uncertainty than

\begin{centeredalign}
\Delta x_{\min}=\hbar\sqrt\beta,\qquad
\beta=\frac{\beta_0}{(m_Pc)^2}\;\Rightarrow\;\Delta x_{\min}=\sqrt{\beta_0}\,l_P,
\end{centeredalign}

with \(\beta_0\) dimensionless and \(m_P\) the Planck mass. Position eigenstates no longer exist; the most localized states are packets of width \(\Delta x_{\min}\). Momentum can still be sharp, and a momentum-space wave function describes the state. As an effective description the model is meant for momenta well below \(m_Pc/\sqrt{\beta_0}\), where the deformation is small.

A harmonic oscillator of mass \(m\) and frequency \(\omega\) shows a concrete consequence. In the variable \(u=\arctan(\sqrt\beta\,p)\) the position acts as \(i\hbar\sqrt\beta\,\partial_u\). The eigenvalue problem becomes a Pöschl--Teller equation with the exact spectrum

\begin{centeredalign}
E_n=\hbar\omega\Big[\big(n+\tfrac12\big)\sqrt{1+\tfrac{\kappa^2}{4}}+\frac\kappa2\big(n^2+n+\tfrac12\big)\Big],
\qquad \kappa=\beta m\hbar\omega.
\end{centeredalign}

For small \(\kappa\) each level shifts by \(\beta m\hbar^2\omega^2(2n^2+2n+1)/4\), and the spacing grows linearly with \(n\). The size is the issue. For an electron with \(\hbar\omega=1\) meV and \(\beta_0=1\), \(\kappa\approx3\times10^{-54}\). Experiments therefore bound \(\beta_0\) rather than measure it, and the bounds lie many orders of magnitude above one. How the deformation applies to composite bodies is unresolved, which complicates tests with massive oscillators.

Other proposals act on spacetime coordinates directly. Snyder in 1947 made coordinates fail to commute, with commutators proportional to the Lorentz generators and to a length squared. Lorentz symmetry is kept, and the spatial coordinates acquire discrete spectra. Doplicher, Fredenhagen and Roberts derived spacetime uncertainty relations such as \(\Delta x^0\sum_i\Delta x^i\gtrsim l_P^2\). Their condition is that localizing an event must not create a horizon. Loop quantum gravity instead quantizes areas and volumes of space, and string theory limits resolution through the size of the string. The scales involved are set out under quantum gravity.

\subsection*{References For The Physical Derivation}
\begin{itemize}
\item \href{https://arxiv.org/abs/hep-th/9412167}{A. Kempf, G. Mangano and R. B. Mann, Hilbert Space Representation of the Minimal Length Uncertainty Relation; deformed commutator, minimal length and the oscillator spectrum.}
\item \href{https://arxiv.org/abs/1203.6191}{S. Hossenfelder, Minimal Length Scale Scenarios for Quantum Gravity; arguments for a minimal length, Snyder-type models and experimental bounds.}
\item \href{https://arxiv.org/abs/hep-th/0303037}{S. Doplicher, K. Fredenhagen and J. E. Roberts, The quantum structure of spacetime at the Planck scale and quantum fields; spacetime uncertainty relations from the absence of horizons.}
\end{itemize}





\clearpage

\section{Fock space}

\label{topic:fock-space}

An excited atom in an empty cavity can decay to its ground state, leaving one photon in a cavity mode. The light of a laser contains no definite number of photons at all. A description with a fixed number of particles cannot represent either situation. Fock space collects the states of zero, one, two and more identical particles in one Hilbert space, so that a single operator can change the particle number.

Let \(\mathcal H\) be the one-particle Hilbert space with orthonormal modes \(\varphi_1,\varphi_2,\dots\). The \(N\)-particle space is the symmetric part of \(\mathcal H^{\otimes N}\) for bosons and the antisymmetric part for fermions, written \(S_\pm\mathcal H^{\otimes N}\). Fock space is their direct sum. Its basis states are labeled by occupation numbers \(n_i\), the number of particles in mode \(i\):

\begin{centeredalign}
\mathcal F_\pm=\bigoplus_{N=0}^{\infty}S_\pm\mathcal H^{\otimes N},\qquad
\ket{n_1,n_2,\dots},\qquad \hat N=\sum_ia_i^\dagger a_i.
\end{centeredalign}

For bosons the creation and annihilation operators change one occupation by one unit, with factors fixed by the commutation relations:

\begin{centeredalign}
a_i^\dagger\ket{\dots,n_i,\dots}=\sqrt{n_i+1}\,\ket{\dots,n_i+1,\dots},\\
a_i\ket{\dots,n_i,\dots}=\sqrt{n_i}\,\ket{\dots,n_i-1,\dots},\qquad
[a_i,a_j^\dagger]=\delta_{ij}.
\end{centeredalign}

For fermions \(n_i\) is zero or one, the commutator becomes an anticommutator, and each operator carries a sign \((-1)^{\sum_{k<i}n_k}\) fixed by an ordering of the modes. Exchange symmetry is then built in: a state specifies how many particles occupy each mode, not which particle is where. A one-body operator with matrix elements \(h_{ij}=\langle\varphi_i|h|\varphi_j\rangle\) becomes \(\sum_{ij}h_{ij}a_i^\dagger a_j\), the same expression for every particle number. Processes that create or absorb particles contain unequal numbers of creation and annihilation operators.

The emission of light shows what the factor \(\sqrt{n+1}\) does. Consider a two-level atom with ground state \(g\) and excited state \(e\), coupled to one field mode of angular frequency \(\omega\) that contains \(n\) photons. Near resonance, in the rotating-wave approximation, the coupling is \(\hbar\lambda(\sigma_+a+\sigma_-a^\dagger)\). Here \(\sigma_-\) takes the atom from \(e\) to \(g\), \(\sigma_+\) does the reverse, and \(\lambda\) is a coupling frequency. The ladder relations give

\begin{centeredalign}
\langle g,n+1|H_{\rm int}|e,n\rangle=\hbar\lambda\sqrt{n+1},\qquad
\langle e,n-1|H_{\rm int}|g,n\rangle=\hbar\lambda\sqrt{n}.
\end{centeredalign}

Transition rates are proportional to squared matrix elements. Emission into the mode therefore proceeds at a rate proportional to \(n+1\), and absorption at a rate proportional to \(n\). The part proportional to \(n\) is stimulated emission; the remaining term is spontaneous emission, present even in an empty mode. This is Einstein's relation between his A and B coefficients, here a consequence of the operator algebra. With equilibrium populations \(N_e/N_g=e^{-\hbar\omega/k_BT}\), the balance \(N_e(n+1)=N_gn\) gives \(n=(e^{\hbar\omega/k_BT}-1)^{-1}\), the Planck occupation.

The ratio of stimulated to spontaneous emission into a mode equals its occupation. For visible light at 600 nm in surroundings at 300 K, the thermal occupation is \(n\approx2\times10^{-35}\), so thermal sources emit spontaneously. A 1 GHz microwave mode at the same temperature has \(n\approx6\times10^3\), and stimulated processes dominate. A laser maintains a large \(n\) in one mode by population inversion, which is why its emission is coherent and directional.

A fixed set of modes describes noninteracting particles exactly. Interactions mix occupation states, so the occupation basis remains a basis but its states are no longer stationary. Each mode of a free field is a Quantum harmonic oscillator, and Quantum field theory uses Fock space as the state space of free particles. Bose--Einstein statistics and Fermi--Dirac statistics give the equilibrium distributions of the occupation numbers.

\subsection*{References For The Physical Derivation}
\begin{itemize}
\item \href{https://arxiv.org/abs/1812.10732}{H. Tasaki, Introduction to the second quantization formalism for non-relativistic quantum mechanics; Fock space and creation and annihilation operators for bosons and fermions.}
\end{itemize}





\part{Historical And Interpretive Annotations}

\chapter{Annotations: history, interpretations, and popular frames}

History, textbooks and interpretations explain how the theory was found and how its probabilities are read.

Each entry points to the treatments in which the corresponding physics is worked out.

\section*{Topics}

\begin{longtable}{p{0.92\linewidth}}
\toprule
Page \\
\midrule
\endhead
\phantomsection\label{page:introduction-to-quantum-mechanics-book}Introduction to Quantum Mechanics (book) \emph{(annotation)} \\
\phantomsection\label{page:quantum-mind}Quantum mind \emph{(annotation)} \\
\phantomsection\label{page:erwin-schr-dinger}Erwin Schrödinger \emph{(annotation)} \\
\phantomsection\label{page:history-of-quantum-mechanics}History of quantum mechanics \emph{(annotation)} \\
\phantomsection\label{page:introduction-to-quantum-mechanics}Introduction to quantum mechanics \emph{(annotation)} \\
\phantomsection\label{page:history-of-quantum-field-theory}History of quantum field theory \emph{(annotation)} \\
\phantomsection\label{page:david-hilbert}David Hilbert \emph{(annotation)} \\
\phantomsection\label{page:quantum-computing-a-gentle-introduction}Quantum Computing: A Gentle Introduction \emph{(annotation)} \\
\phantomsection\label{page:quantum-mysticism}Quantum mysticism \emph{(annotation)} \\
\phantomsection\label{page:interpretations-of-quantum-mechanics}Interpretations of quantum mechanics \emph{(annotation)} \\
\phantomsection\label{page:quantum-theory-concepts-and-methods}Quantum Theory: Concepts and Methods \emph{(annotation)} \\
\phantomsection\label{page:qbism}QBism \emph{(annotation)} \\
\phantomsection\label{page:relational-quantum-mechanics}Relational quantum mechanics \emph{(annotation)} \\
\phantomsection\label{page:modern-quantum-mechanics}Modern Quantum Mechanics \emph{(annotation)} \\
\phantomsection\label{page:applications-of-quantum-mechanics}Applications of quantum mechanics \emph{(annotation)} \\
\phantomsection\label{page:werner-heisenberg}Werner Heisenberg \emph{(annotation)} \\
\phantomsection\label{page:the-quantum-universe}The Quantum Universe \emph{(annotation)} \\
\bottomrule
\end{longtable}

\clearpage

\section{Introduction to Quantum Mechanics (book)}

\label{topic:introduction-to-quantum-mechanics-book}

\par\smallskip\noindent\textit{Historical, bibliographic or interpretive entry.}

A student who meets quantum mechanics through David Griffiths' Introduction to Quantum Mechanics starts with a wave function on a line, not with abstract vectors. The book appeared in 1995 and is widely used in undergraduate courses. Its third edition (2018) was written with Darrell Schroeter. The opening chapter states the statistical interpretation: \(|\Psi(x,t)|^2\,dx\) is the probability of finding the particle between \(x\) and \(x+dx\) at time \(t\).

The route then solves the time-independent Schrödinger equation for a sequence of one-dimensional potentials. These are the infinite square well, the harmonic oscillator, the free particle, the delta-function well and the finite square well. The oscillator is solved twice, once with ladder operators and once by a power series, so the algebra and the differential equation are seen to give the same spectrum. For a well of width \(a\) with impenetrable walls the energies and wave functions are

\begin{centeredalign}
E_n=\frac{n^2\pi^2\hbar^2}{2ma^2},\qquad
\psi_n(x)=\sqrt{\frac{2}{a}}\,\sin\frac{n\pi x}{a},\qquad n=1,2,3,\dots
\end{centeredalign}

For an electron in a well of width \(a=1\) nm the lowest energy is \(0.38\) eV. The transition from \(n=2\) to \(n=1\) releases \(1.13\) eV, a photon of wavelength near \(1.1\,\mu\mathrm{m}\). This is the order of magnitude by which confinement shifts optical transitions in semiconductor quantum wells and dots, although an effective mass and finite barriers change the numbers.

Only after these examples does the book introduce the formalism: Hilbert space, Hermitian operators, discrete and continuous spectra, the generalized uncertainty principle and Dirac notation. Three dimensions follow, with the hydrogen atom, angular momentum and spin, and then identical particles. The second half covers approximation methods. These include time-independent perturbation theory, the variational principle, the WKB approximation, time-dependent perturbation theory with emission and absorption of radiation, the adiabatic approximation with Berry's phase, and scattering. A closing afterword discusses the Einstein--Podolsky--Rosen argument, Bell's theorem, Schrödinger's cat and the interpretation of measurement.

The pedagogical choice is to let calculus on familiar potentials carry the reader before the linear algebra is made explicit. Sakurai's Modern Quantum Mechanics takes the opposite route, starting from spin and Dirac notation. Readers can continue with the treatments Particle in a box, Quantum harmonic oscillator, Schrödinger equation, Perturbation theory (quantum mechanics) and Scattering.







\clearpage

\section{Quantum mind}

\label{topic:quantum-mind}

\par\smallskip\noindent\textit{Historical, bibliographic or interpretive entry.}

A neuron either fires or stays quiet, and the difference involves the motion of many ions across its membrane. Proposals grouped under the name quantum mind suggest that quantum superpositions of such brain states, or of molecular states inside neurons, take part in cognition or consciousness. The best known is the orchestrated objective reduction model of Roger Penrose and Stuart Hameroff. It places superpositions in microtubules and links their collapse to gravity. Matthew Fisher proposed in 2015 that phosphorus nuclear spins in calcium phosphate clusters could stay coherent for long times.

Any such proposal must survive decoherence. A superposition of two states that differ in the positions of charges is recorded by the surrounding ions and water molecules, which scatter differently from the two configurations. Max Tegmark estimated in 2000 how quickly these collisions destroy interference. He compared the decoherence time \(\tau_{\rm dec}\) with the time \(\tau_{\rm dyn}\) over which neural states change:

\begin{centeredalign}
\tau_{\rm dec}\approx10^{-20}\text{--}10^{-13}\ \mathrm{s}\ \ll\
\tau_{\rm dyn}\approx10^{-3}\text{--}10^{-1}\ \mathrm{s}.
\end{centeredalign}

The shorter value applies to a neuron in a superposition of firing and not firing, the longer to a kink-like polarization excitation in a microtubule. Even the longer time is ten orders of magnitude below the dynamical time. On this estimate, the variables that enter neural computation behave classically. Hagan, Hameroff and Tuszyński replied that different assumptions raise the microtubule estimate to \(10^{-5}\)--\(10^{-4}\) s. Their changes concern the separation of the superposed states, charge versus dipole coupling and the dielectric constant. They argued that ordered water near microtubules could extend the time further.

What is established: quantum mechanics determines chemical bonding and the structure of ion channels, and tunnelling and spin-dependent radical-pair reactions occur in some biological molecules. What is not established: no experiment has shown a superposition of neural states, entanglement between neurons, or any contribution of quantum coherence to perception or consciousness. Fisher's proposal points to testable consequences, for example different effects of the lithium isotopes 6 and 7, whose nuclear spins differ; these remain open. Readers can continue with Quantum decoherence, Quantum biology, Measurement problem and Wave function collapse.

\subsection*{References For The Physical Derivation}
\begin{itemize}
\item \href{https://arxiv.org/abs/quant-ph/9907009}{M. Tegmark, The importance of quantum decoherence in brain processes; decoherence times of 1e-20 to 1e-13 s against dynamical times of 1e-3 to 1e-1 s.}
\item \href{https://arxiv.org/abs/quant-ph/0005025}{S. Hagan, S. R. Hameroff and J. A. Tuszyński, Quantum Computation in Brain Microtubules? Decoherence and Biological Feasibility; revised microtubule estimate of 1e-5 to 1e-4 s.}
\end{itemize}





\clearpage

\section{Erwin Schrödinger}

\label{topic:erwin-schr-dinger}

\par\smallskip\noindent\textit{Historical, bibliographic or interpretive entry.}

De Broglie's proposal of 1924 assigned an electron the wavelength \(h/p\) but did not say how such a wave evolves in a potential. Erwin Schrödinger (1887--1961) supplied the equation in four papers of 1926 titled Quantisierung als Eigenwertproblem. Energy levels became eigenvalues of a differential operator, and boundary conditions replaced the quantum rules of the older theory. The equation and the matrix elements that connect it with Heisenberg's arrays are

\begin{centeredalign}
i\hbar\frac{\partial\psi}{\partial t}=\Bigl[-\frac{\hbar^2}{2m}\nabla^2+V(\mathbf r)\Bigr]\psi,
\qquad A_{mn}=\int\psi_m^*\,A\,\psi_n\,d^3r .
\end{centeredalign}

The opening paper obtained the hydrogen spectrum \(E_n\propto-1/n^2\) from the requirement that the wave function be finite and single-valued. Later papers treated the oscillator, the rigid rotor, the Stark effect by perturbation theory, and finally the time-dependent equation. In 1926 he also constructed oscillator wave packets that move like a classical particle without spreading, now called coherent states.

In the same year Schrödinger showed how matrix mechanics follows from wave mechanics. Expanding in the energy eigenfunctions \(\psi_n\) turns each operator \(A\) into the matrix \(A_{mn}\) above. The matrices of \(x\) and \(p=-i\hbar\,\partial_x\) then satisfy \(xp-px=i\hbar\) and oscillate at the frequencies \((E_m-E_n)/\hbar\). For the oscillator the only nonzero position elements are \(x_{n,n+1}=\sqrt{(n+1)\hbar/(2m\omega)}\), which are Heisenberg's transition amplitudes. Dirac and Jordan's transformation theory and von Neumann's Hilbert space later gave the equivalence its general form.

In 1935, responding to Einstein, Podolsky and Rosen, Schrödinger analysed composite systems whose joint state is not a product of states of the parts. He called this entanglement and noted that a measurement on one part can steer the other into different sets of states. The pair state \((\ket{\uparrow\downarrow}-\ket{\downarrow\uparrow})/\sqrt2\) gives each spin the density operator \(I/2\), yet outcomes along any common axis are perfectly anticorrelated. In the same year he described a cat whose survival is correlated with a decaying atom, to show the difficulty of applying an indefinite state to a macroscopic object.

He shared the 1933 Nobel Prize in Physics with Dirac. In 1944 he wrote What is Life?, which discussed how genetic information could be stored in an aperiodic molecular structure. Readers can continue with Schrödinger equation, Schrödinger picture, Quantum entanglement, Schrödinger's cat and Einstein--Podolsky--Rosen paradox.

\subsection*{References For The Physical Derivation}
\begin{itemize}
\item \href{https://arxiv.org/abs/quant-ph/0612147}{H. M. Wiseman, S. J. Jones and A. C. Doherty, Steering, Entanglement, Nonlocality, and the EPR Paradox; operational definition of the steering introduced by Schrödinger in 1935.}
\end{itemize}





\clearpage

\section{History of quantum mechanics}

\label{topic:history-of-quantum-mechanics}

\par\smallskip\noindent\textit{Historical, bibliographic or interpretive entry.}

In 1900 the measured spectrum of radiation from a heated cavity disagreed with classical physics at high frequencies. Classical equipartition gives each electromagnetic mode the mean energy \(k_BT\), so the energy density grows without bound with frequency. Max Planck found a formula that fitted the data. He then derived it by assuming that material oscillators exchange energy with radiation of frequency \(\nu\) in units \(h\nu\):

\begin{centeredalign}
u(\nu,T)=\frac{8\pi h\nu^3}{c^3}\,\frac{1}{e^{h\nu/k_BT}-1}\ \longrightarrow\
\begin{cases}\dfrac{8\pi\nu^2}{c^3}\,k_BT, & h\nu\ll k_BT,\\[6pt]
\dfrac{8\pi h\nu^3}{c^3}\,e^{-h\nu/k_BT}, & h\nu\gg k_BT.\end{cases}
\end{centeredalign}

The low-frequency limit is the classical Rayleigh--Jeans law, and the high-frequency limit is Wien's law. Fitting the formula to the spectra gave Planck values of \(h\) and \(k_B\) within a few percent of the modern ones. From \(k_B\) he obtained Avogadro's number and the elementary charge.

Einstein in 1905 treated light itself as quanta of energy \(h\nu\). This explained why the photoelectron energy \(h\nu-W\), with \(W\) the work function, depends on frequency and not on intensity. In 1907 he applied quantization to lattice vibrations and explained the fall of the heat capacity of solids at low temperature. Bohr's model of 1913 quantized the electron's angular momentum in hydrogen and reproduced the Balmer lines. Sommerfeld extended it to elliptical orbits and fine structure. This old quantum theory failed for helium and needed ad hoc rules for the anomalous Zeeman effect.

Between 1923 and 1927 a new mechanics replaced the orbits. Compton's scattering of X-rays by electrons (1923) confirmed photon momentum \(h/\lambda\). De Broglie (1924) assigned matter the wavelength \(\lambda=h/p\). In 1925 Heisenberg replaced orbits by arrays of transition amplitudes, and Born and Jordan recognized them as matrices. In 1926 Schrödinger wrote a wave equation for the same physics:

\begin{centeredalign}
pq-qp=\frac{h}{2\pi i}\,I,\qquad i\hbar\frac{\partial\psi}{\partial t}=H\psi .
\end{centeredalign}

Pauli derived the hydrogen spectrum from matrix mechanics early in 1926, and Schrödinger showed the two forms equivalent. Born then interpreted \(|\psi|^2\) as a probability density. Pauli had stated the exclusion principle in 1925, and Uhlenbeck and Goudsmit had proposed electron spin. In 1927 Davisson and Germer, and independently G. P. Thomson, observed electron diffraction, and Heisenberg stated the uncertainty relation. The Solvay conference of October 1927 closed this formative period. Readers can continue with Old quantum theory, Wave--particle duality, Canonical commutation relation, Schrödinger equation and Uncertainty principle.







\clearpage

\section{Introduction to quantum mechanics}

\label{topic:introduction-to-quantum-mechanics}

\par\smallskip\noindent\textit{Historical, bibliographic or interpretive entry.}

A hydrogen lamp viewed through a grating shows sharp lines at fixed wavelengths, the brightest red one at \(656\) nm. Classical mechanics allows an electron orbit of any energy and so predicts a continuous spectrum. Quantum mechanics accounts for the lines with three facts that recur throughout the subject: quantized energies, superposition of states, and the Born rule for probabilities.

Quantization. A bound system has a discrete set of energies, the eigenvalues of its Hamiltonian \(H\). For the electron in hydrogen, with reduced mass \(\mu\), elementary charge \(e\) and principal quantum number \(n\), the bound energies and the emitted photon frequency \(\nu\) are

\begin{centeredalign}
E_n=-\frac{\mu e^4}{2(4\pi\varepsilon_0)^2\hbar^2}\,\frac{1}{n^2}
\approx-\frac{13.6\ \mathrm{eV}}{n^2},\qquad h\nu=E_m-E_n .
\end{centeredalign}

A photon carries away the energy difference between two levels. The transition from \(n=3\) to \(n=2\) gives \(13.6\times(1/4-1/9)=1.89\) eV, a wavelength of \(656\) nm, the observed red Balmer line. Transitions ending on \(n=1\) lie in the ultraviolet, beginning at \(121.6\) nm. Fine structure, the Lamb shift and hyperfine splitting appear only when relativity, the quantized radiation field and the nuclear spin are added.

Superposition. If \(\ket{\psi_1}\) and \(\ket{\psi_2}\) are possible states, so is \(\alpha\ket{\psi_1}+\beta\ket{\psi_2}\) for complex numbers \(\alpha\) and \(\beta\). Superposition is not a statement of ignorance: the relative phase of \(\alpha\) and \(\beta\) changes the outcome of later interference experiments.

Born rule. A measurement whose outcomes correspond to orthonormal states \(\ket{a}\) gives outcome \(a\) with probability \(|\langle a|\psi\rangle|^2\). For a spin-1/2 particle prepared along \(z\) and sent through a Stern--Gerlach magnet tilted by the angle \(\theta\),

\begin{centeredalign}
P(\uparrow_\theta)=|\langle\uparrow_\theta|\uparrow_z\rangle|^2=\cos^2\frac{\theta}{2},
\qquad P(\downarrow_\theta)=\sin^2\frac{\theta}{2}.
\end{centeredalign}

At \(\theta=90^\circ\) the two outcomes are equally likely, and each atom still lands in one of only two spots. Repeated runs reproduce the frequencies, not the individual outcomes. Between preparation and measurement the state changes deterministically by the Schrödinger equation \(i\hbar\,\partial_t\ket{\psi}=H\ket{\psi}\). Uncertainty relations and tunnelling follow from these rules, and the tensor-product rule for composite systems adds entanglement. The treatments Schrödinger equation, Quantum superposition, Born rule, Spin (physics) and Spectral line develop each step.







\clearpage

\section{History of quantum field theory}

\label{topic:history-of-quantum-field-theory}

\par\smallskip\noindent\textit{Historical, bibliographic or interpretive entry.}

The quantum mechanics of 1926 treated electrons quantum mechanically but left the electromagnetic field classical, so it could not derive the rate of spontaneous emission. In 1927 Dirac quantized the radiation field by treating each mode, with wave vector \(\mathbf k\) and polarization \(\lambda\), as a harmonic oscillator:

\begin{centeredalign}
H_{\rm rad}=\sum_{\mathbf k,\lambda}\hbar\omega_k\Bigl(a^\dagger_{\mathbf k\lambda}a_{\mathbf k\lambda}+\tfrac12\Bigr),
\qquad [a_{\mathbf k\lambda},a^\dagger_{\mathbf k'\lambda'}]=\delta_{\mathbf k\mathbf k'}\delta_{\lambda\lambda'} .
\end{centeredalign}

The number of photons in a mode is the excitation level of its oscillator. The matrix element \(\langle n+1|a^\dagger|n\rangle=\sqrt{n+1}\) makes the emission rate into a mode proportional to \(n+1\). The term \(n\) is stimulated emission and the term \(1\) is spontaneous emission, in Einstein's ratio. Jordan and Wigner (1928) introduced anticommutators for fermion fields, and Heisenberg and Pauli (1929) formulated a general quantized field theory. Dirac's relativistic electron equation (1928) had negative-energy solutions, interpreted in 1931 as antiparticles; Anderson observed the positron in 1932. Fermi's theory of beta decay (1934) created and destroyed electrons and neutrinos, so particle number was no longer conserved.

Calculations beyond lowest order gave infinite results, for example for the electron self-energy. In 1947 Lamb and Retherford measured a splitting of about 1 GHz between the \(2S_{1/2}\) and \(2P_{1/2}\) levels of hydrogen, which Dirac's equation predicts to be degenerate. Bethe obtained about 1040 MHz by absorbing the divergent part into the observed electron mass. Tomonaga, Schwinger, Feynman and Dyson made this renormalization systematic in quantum electrodynamics. Schwinger's 1948 result for the anomaly of the electron magnetic moment is

\begin{centeredalign}
a_e=\frac{g-2}{2}=\frac{\alpha}{2\pi}\approx0.001161,\qquad a_e^{\rm exp}\approx0.0011597 .
\end{centeredalign}

Higher orders in the fine-structure constant \(\alpha\) account for the remaining difference. Yang and Mills (1954) extended gauge invariance to non-abelian symmetry groups. Glashow, Weinberg and Salam built the electroweak theory between 1961 and 1968; Weinberg and Salam used the Higgs mechanism of 1964 to give the W and Z bosons their masses. 't Hooft and Veltman proved in 1971--72 that such theories are renormalizable. Gross, Wilczek and Politzer found asymptotic freedom in 1973, which led to quantum chromodynamics as the theory of the strong interaction. Wilson's renormalization group explained renormalization as the change of effective couplings with length scale. The W and Z bosons were observed in 1983 and the Higgs boson in 2012. Readers can continue with Quantum field theory, Quantum electrodynamics, Renormalization, Gauge theory and Standard Model.

\subsection*{References For The Physical Derivation}
\begin{itemize}
\item \href{https://arxiv.org/abs/hep-th/9812203}{G. 't Hooft, The Glorious Days of Physics: Renormalization of Gauge Theories; recollections of the proof of renormalizability and the establishment of the Standard Model.}
\end{itemize}





\clearpage

\section{David Hilbert}

\label{topic:david-hilbert}

\par\smallskip\noindent\textit{Historical, bibliographic or interpretive entry.}

A wave function on a line and an infinite list of amplitudes describe the same state, provided the squared magnitudes have a finite sum. This identification underlies the equivalence of wave and matrix mechanics. It grew from mathematics developed by David Hilbert (1862--1943). Between 1904 and 1910 Hilbert studied integral equations by reducing them to quadratic forms in infinitely many variables. He called the set of their characteristic values a spectrum, before any link to atomic spectra was known, and separated discrete points from continuous parts.

The space of square-summable sequences, and its equivalence with square-integrable functions proved by Riesz and Fischer in 1907, are now called Hilbert space. For an orthonormal basis \(\phi_n\), such as the eigenfunctions of a Hamiltonian, the expansion coefficients \(c_n\) satisfy

\begin{centeredalign}
\psi(x)=\sum_nc_n\phi_n(x),\qquad c_n=\int\phi_n^*\psi\,dx,\qquad
\int|\psi|^2dx=\sum_n|c_n|^2 .
\end{centeredalign}

The last identity, Parseval's relation, states that the norm, and therefore every Born probability, is the same in both descriptions. A standard check expands \(f(x)=x\) on \((-\pi,\pi)\) in sines, with coefficients \(2(-1)^{n+1}/n\); Parseval's relation then gives \(\sum_n1/n^2=\pi^2/6\).

In a paper of 1928, Hilbert, John von Neumann and Lothar Nordheim discussed the foundations of quantum mechanics. Von Neumann then defined abstract Hilbert space by axioms. In 1929 he proved the spectral theorem for unbounded self-adjoint operators, the case needed for position, momentum and energy. Every such observable becomes \(A=\int\lambda\,dE(\lambda)\) for a projection-valued measure \(E\). A state \(\rho\) then gives an outcome set \(B\) the probability \(\operatorname{Tr}[\rho E(B)]\). His 1932 book Mathematische Grundlagen der Quantenmechanik collected these results.

The 1924 Methoden der mathematischen Physik by Richard Courant and Hilbert supplied eigenvalue methods that physicists applied to wave mechanics in 1926. Hilbert's sixth problem of 1900 had asked for an axiomatic treatment of physics, which von Neumann's work partly carried out for quantum mechanics. Readers can continue with Hilbert space, Spectral theory, Self-adjoint operator, Projection-valued measure and Mathematical formulation of quantum mechanics.







\clearpage

\section{Quantum Computing: A Gentle Introduction}

\label{topic:quantum-computing-a-gentle-introduction}

\par\smallskip\noindent\textit{Historical, bibliographic or interpretive entry.}

Three polarizing filters can transmit light that two of them block. A horizontal and a vertical polarizer extinguish an unpolarized beam. A third polarizer at \(45^\circ\) inserted between them lets one eighth of the original intensity through, \(\tfrac12\cos^245^\circ\cos^245^\circ=\tfrac18\). Eleanor Rieffel and Wolfgang Polak introduce the qubit through photon polarization and this experiment in Quantum Computing: A Gentle Introduction (MIT Press, 2011). The authors had earlier written a survey of quantum computing for non-physicists.

The book is written for readers trained in computer science or mathematics rather than physics. Its opening part covers single and multiple qubits, measurement, entanglement, the BB84 key-distribution protocol, unitary gates and reversible versions of classical functions. The middle part treats quantum algorithms, from Deutsch's and Simon's problems to Shor's factoring algorithm and Grover's search. The final part covers entangled subsystems, quantum error correction, fault tolerance, and alternative models such as measurement-based and adiabatic computation.

Grover's algorithm shows the style of argument. A search over \(N\) items with one marked item \(\ket{w}\) starts from the uniform superposition \(\ket{s}\). Each iteration flips the sign of \(\ket{w}\) and then reflects about \(\ket{s}\). Two reflections make a rotation by \(2\theta\) in the plane of \(\ket{w}\) and \(\ket{s}\), so after \(k\) iterations

\begin{centeredalign}
\sin\theta=\frac{1}{\sqrt N},\qquad P_k=\sin^2\bigl((2k+1)\theta\bigr),\qquad
k_{\rm opt}\approx\frac{\pi}{4}\sqrt N .
\end{centeredalign}

For \(N=4\) a single iteration gives \(3\theta=90^\circ\) and finds the marked item with certainty. For \(N=10^6\) about 785 iterations replace the half million queries that classical search needs on average. The quadratic gain assumes that the oracle acts coherently on a superposition, and it is known to be optimal for unstructured search. Readers can continue with Qubit, Quantum logic gate, Quantum circuit, Quantum algorithm and Quantum error correction.

\subsection*{References For The Physical Derivation}
\begin{itemize}
\item \href{https://arxiv.org/abs/quant-ph/9809016}{E. G. Rieffel and W. Polak, An Introduction to Quantum Computing for Non-Physicists; the authors' earlier survey for non-physicists, covering basic principles, cryptography, teleportation, Shor's and Grover's algorithms and error correction.}
\end{itemize}





\clearpage

\section{Quantum mysticism}

\label{topic:quantum-mysticism}

\par\smallskip\noindent\textit{Historical, bibliographic or interpretive entry.}

Since the 1970s a body of popular writing, often called quantum mysticism, has connected quantum theory with consciousness, Eastern philosophy or healing practices. Fritjof Capra's The Tao of Physics (1975) is an early and influential example. Some founders of the theory also wrote on philosophy. Bohr chose the yin-yang symbol for his coat of arms to express complementarity, and Schrödinger wrote about Vedanta. The physical content of such claims can be compared with what quantum mechanics states.

A frequent claim is that an observer's mind creates the outcome of a measurement. In the theory, a measurement is a physical interaction that correlates the system with a detector. Decoherence then suppresses interference between the detector's pointer states within times far shorter than any human reaction. Von Neumann and Wigner discussed placing the collapse in the observer's consciousness. This is one interpretive proposal among several, not a consequence of the formalism.

A second claim is that entanglement connects everything instantly, so that intentions could act at a distance. Entangled particles do show correlations that no local hidden-variable model reproduces, as Bell tests confirm. The same theory forbids using them to send a signal. Whatever operation Bob performs on his particle, described by operators \(M_k\) on his system \(B\), Alice's reduced state is unchanged:

\begin{centeredalign}
\rho_A'=\operatorname{Tr}_B\Bigl[\sum_kK_k\,\rho_{AB}\,K_k^\dagger\Bigr]=\rho_A,\qquad
K_k=I_A\otimes M_k,\qquad \sum_kM_k^\dagger M_k=I_B .
\end{centeredalign}

This holds for unitary operations and for measurements whose outcomes are not communicated. Alice's statistics carry no information about Bob's choice until a classical message arrives. A third claim reads the uncertainty principle as subjective doubt. In the theory it bounds the spreads of position and momentum in any prepared state, \(\sigma_x\sigma_p\geq\hbar/2\).

Questions of meaning, value and spiritual experience lie outside physics, which neither supports nor refutes them. Where a claim makes a physical prediction, such as a mental influence on a random quantum outcome or a healing effect of entanglement, it can be tested by standard methods. No such effect has been established. Readers can continue with Quantum entanglement, Measurement in quantum mechanics, Quantum decoherence and Uncertainty principle.







\clearpage

\section{Interpretations of quantum mechanics}

\label{topic:interpretations-of-quantum-mechanics}

\par\smallskip\noindent\textit{Historical, bibliographic or interpretive entry.}

Every interpretation of quantum mechanics must reproduce the same laboratory statistics, from atomic spectra to Bell-test correlations. All therefore share two rules. Between interventions a closed system evolves unitarily under its Hamiltonian \(H\), and outcome \(a\) occurs with the Born probability, where \(P_a\) is the projector onto the outcome:

\begin{centeredalign}
\ket{\psi(t)}=e^{-iHt/\hbar}\ket{\psi(0)},\qquad p(a)=\langle\psi|P_a|\psi\rangle .
\end{centeredalign}

They also share decoherence. When a system becomes correlated with many environmental degrees of freedom, interference between macroscopically distinct alternatives becomes unobservable in practice. Decoherence explains why a detector shows no interference between its pointer positions. It does not by itself select which outcome occurs. That question, the measurement problem, is where interpretations differ.

The Copenhagen view of Bohr and Heisenberg describes the apparatus classically and treats the state as a tool for predicting results of specified experiments. Everett's relative-state formulation (1957) keeps unitary evolution without exception: all outcomes occur in decohered branches, and probabilities must be recovered from branch weights. The de Broglie--Bohm theory adds particle positions guided by the wave function and reproduces the Born rule when positions are initially distributed as \(|\psi|^2\). QBism treats the state as an agent's degrees of belief, and relational quantum mechanics treats outcomes as relative to the interacting system. These approaches agree with standard quantum mechanics in all predicted statistics and differ in what they take to exist.

Objective-collapse models change the dynamics itself. In the model of Ghirardi, Rimini and Weber (GRW, 1986), each particle undergoes spontaneous localizations at rate \(\lambda\) to a region of width \(r_C\). For one nucleon these are too rare to observe. In a rigid body of \(10^{23}\) nucleons the rates add. A superposition of two positions of the body separated by more than \(r_C\) then disappears within about \(10^{-7}\) s. Each localization also gives a random momentum kick, so every particle gains energy at the rate

\begin{centeredalign}
\frac{dE}{dt}=\frac{3\lambda\hbar^2}{4m\,r_C^2}\qquad\text{for a particle of mass } m .
\end{centeredalign}

With the GRW values \(\lambda=10^{-16}\ \mathrm{s^{-1}}\) and \(r_C=10^{-7}\) m, a nucleon gains about \(5\times10^{-44}\) W. This spontaneous heating, and the radiation emitted by randomly kicked charges, are absent in standard quantum mechanics. Searches with X-ray detectors, cold solids and mechanical oscillators have excluded part of the parameter range, but not the original GRW values. An underground radiation measurement ruled out the parameter-free version of the gravity-related Diósi--Penrose model. Readers can continue with Measurement problem, Quantum decoherence, Wave function collapse, Born rule and Bell's theorem.

\subsection*{References For The Physical Derivation}
\begin{itemize}
\item \href{https://arxiv.org/abs/2203.04231}{M. Carlesso, S. Donadi, L. Ferialdi, M. Paternostro, H. Ulbricht and A. Bassi, Present status and future challenges of non-interferometric tests of collapse models; bounds from X-ray emission, bulk heating and mechanical oscillators.}
\item \href{https://arxiv.org/abs/2111.13490}{S. Donadi, K. Piscicchia, C. Curceanu, L. Diósi, M. Laubenstein and A. Bassi, Underground test of gravity-related wave function collapse; exclusion of the parameter-free Diósi--Penrose model.}
\end{itemize}





\clearpage

\section{Quantum Theory: Concepts and Methods}

\label{topic:quantum-theory-concepts-and-methods}

\par\smallskip\noindent\textit{Historical, bibliographic or interpretive entry.}

Asher Peres's Quantum Theory: Concepts and Methods (1993) treats quantum theory as a set of rules for computing the probabilities of outcomes of tests that follow specified preparations. It does not assign properties to a system between preparation and test. The book joins foundational results to techniques later central to quantum information: composite systems, Bell's theorem, contextuality, information and thermodynamics, semiclassical methods, chaos and the measuring process.

A central tool is the generalized measurement, or positive operator-valued measure (POVM). Its outcomes \(k\) correspond to positive operators \(E_k\) that need not be projectors; they only have to sum to the identity. Neumark's theorem, discussed in the book, shows that every such measurement is a projective measurement on the system combined with an auxiliary system. In 1988 Peres applied this freedom to two non-orthogonal states \(\ket{\psi_1}\) and \(\ket{\psi_2}\) prepared with equal probability. No measurement identifies them every time without error. A three-outcome POVM can do so part of the time: two outcomes name the state with certainty, and the third is inconclusive. The probabilities and the largest conclusive fraction are

\begin{centeredalign}
p_k=\operatorname{Tr}(\rho E_k),\qquad E_k\geq0,\qquad \sum_kE_k=I,\qquad
P_{\rm conclusive}=1-|\langle\psi_1|\psi_2\rangle| .
\end{centeredalign}

This limit was found independently by Ivanovic, Dieks and Peres. It sets how often an eavesdropper can identify non-orthogonal signal states without error, a quantity used in analyses of quantum key distribution. For orthogonal states the conclusive fraction is one, and for identical states it is zero.

The chapters on inseparability derive the Clauser--Horne--Shimony--Holt inequality for local hidden variables, with \(\pm1\) outcomes \(A_0,A_1\) on one side and \(B_0,B_1\) on the other. The singlet state violates it:

\begin{centeredalign}
|\langle A_0B_0\rangle+\langle A_0B_1\rangle+\langle A_1B_0\rangle-\langle A_1B_1\rangle|
\leq2,\qquad \text{quantum maximum } 2\sqrt2 .
\end{centeredalign}

The chapter on contextuality gives a proof of the Kochen--Specker theorem with 33 directions in three-dimensional space. No rule assigns the value 1 to exactly one direction of every orthogonal triad among them. Readers can continue with POVM, Bell's theorem, Quantum nonlocality, Projection-valued measure and Measurement in quantum mechanics.







\clearpage

\section{QBism}

\label{topic:qbism}

\par\smallskip\noindent\textit{Historical, bibliographic or interpretive entry.}

Two experimenters with different information about how a spin was prepared may assign it different density operators. QBism, the quantum-Bayesian interpretation developed by Carlton Caves, Christopher Fuchs and Rüdiger Schack from about 2002, takes this as the general case. A quantum state is an agent's personal degrees of belief about the outcomes of that agent's own future actions on a system. The outcomes are the agent's experiences. The formalism is unchanged; what changes is the meaning of a state assignment.

QBists read the Born rule as a consistency condition linking probabilities for different experiments, comparable to the rules that keep a gambler's bets coherent. Fuchs and Schack made this explicit with a symmetric informationally complete measurement (SIC). In dimension \(d\) it has \(d^2\) outcomes \(H_i=\Pi_i/d\), whose rank-one projectors satisfy \(\operatorname{Tr}(\Pi_i\Pi_j)=1/(d+1)\) for \(i\neq j\). The probabilities \(p(i)\) of this reference measurement fix the state completely. For any other measurement with outcomes \(D_j\), the Born rule becomes

\begin{centeredalign}
q(j)=\sum_{i=1}^{d^2}\Bigl[(d+1)\,p(i)-\frac1d\Bigr]\,r(j|i),\qquad
r(j|i)=\operatorname{Tr}(\Pi_iD_j) .
\end{centeredalign}

Here \(r(j|i)\) is the probability of \(D_j\) if the reference measurement had been performed and had given \(i\). A classical agent would use the law of total probability, \(q(j)=\sum_ip(i)\,r(j|i)\). The quantum rule differs by the factor \(d+1\) and the shift \(1/d\), and QBism takes these numbers as the empirical content of quantum theory. For a qubit the four SIC states point to the corners of a regular tetrahedron on the Bloch sphere.

On this reading, a measurement on one member of an entangled pair updates only the measuring agent's expectations about the distant member. No physical influence travels, and QBists describe quantum mechanics as local. Critics reply that this leaves open why the Born rule succeeds and what the world is like apart from agents. QBism predicts no measured statistics different from standard quantum mechanics. Readers can continue with Born rule, POVM, Quantum state and Measurement problem.

\subsection*{References For The Physical Derivation}
\begin{itemize}
\item \href{https://arxiv.org/abs/1311.5253}{C. A. Fuchs, N. D. Mermin and R. Schack, An Introduction to QBism with an Application to the Locality of Quantum Mechanics; the interpretation and its account of entangled correlations.}
\item \href{https://arxiv.org/abs/1301.3274}{C. A. Fuchs and R. Schack, Quantum-Bayesian Coherence: The No-Nonsense Version; the Born rule written with a SIC reference measurement.}
\end{itemize}





\clearpage

\section{Relational quantum mechanics}

\label{topic:relational-quantum-mechanics}

\par\smallskip\noindent\textit{Historical, bibliographic or interpretive entry.}

Wigner's friend measures the spin of an electron inside a sealed laboratory and sees a definite result. Wigner, outside, describes the friend and the electron together by unitary evolution and assigns them an entangled state. Carlo Rovelli's relational quantum mechanics (1996) accepts both descriptions as correct. The values of physical variables are relative to the system that interacts with them, as simultaneity is relative to an observer in special relativity.

Let the spin be prepared as \(\alpha\ket{\uparrow}+\beta\ket{\downarrow}\), and let \(\ket{F_\uparrow}\) and \(\ket{F_\downarrow}\) be the friend's recording states. Relative to Wigner the measurement produces the state \(\ket{\Psi}\) below. Let \(P_\Phi\) project onto \(\ket{\Phi}=(\ket{\uparrow}\ket{F_\uparrow}+\ket{\downarrow}\ket{F_\downarrow})/\sqrt2\). For \(\alpha=\beta=1/\sqrt2\),

\begin{centeredalign}
\ket{\Psi}=\alpha\ket{\uparrow}\ket{F_\uparrow}+\beta\ket{\downarrow}\ket{F_\downarrow},\qquad
\langle\Psi|P_\Phi|\Psi\rangle=1 .
\end{centeredalign}

An equal mixture of the two recorded outcomes would give probability \(1/2\) for the same projector. In principle Wigner can tell the two descriptions apart by this interference measurement on the whole laboratory. For a macroscopic friend, decoherence makes it impossible in practice, which is why everyday observers agree. Relational quantum mechanics holds that the friend's outcome is a fact relative to the friend, while the entangled state is correct relative to Wigner. When Wigner asks the friend, the answer is a further interaction, and its result is correlated with the spin exactly as \(\ket{\Psi}\) requires.

The interpretation removes observer-independent states but keeps the Born rule and unitary evolution, so it makes no new experimental prediction. Discussion has centred on whether facts relative to different observers can be compared. In 2022 Adlam and Rovelli added a postulate that information held by one observer is stored in physical variables that other observers can measure. Readers can continue with Measurement problem, Quantum decoherence, Quantum entanglement and Wave function collapse.

\subsection*{References For The Physical Derivation}
\begin{itemize}
\item \href{https://arxiv.org/abs/quant-ph/9609002}{C. Rovelli, Relational Quantum Mechanics; values of physical quantities taken as relative to other systems, illustrated by one observer measuring a system while a second observer describes both.}
\item \href{https://arxiv.org/abs/2203.13342}{E. Adlam and C. Rovelli, Information is Physical: Cross-Perspective Links in Relational Quantum Mechanics; the postulate that observers' information is stored in measurable variables.}
\end{itemize}





\clearpage

\section{Modern Quantum Mechanics}

\label{topic:modern-quantum-mechanics}

\par\smallskip\noindent\textit{Historical, bibliographic or interpretive entry.}

A beam of silver atoms passing a non-uniform magnetic field splits into two spots, not a continuous smear. J. J. Sakurai's Modern Quantum Mechanics begins with this Stern--Gerlach experiment instead of the wave function. Sakurai died in 1982; San Fu Tuan completed the book from his notes, and it was published in 1985. Later editions, from 2011, were revised with Jim Napolitano.

The opening argument uses sequential Stern--Gerlach filters. Atoms selected with \(S_z=+\hbar/2\) pass a second apparatus oriented along \(x\). The \(S_x=+\hbar/2\) beam then enters a third apparatus along \(z\), and both \(S_z\) outcomes reappear. The \(S_x\) selection has destroyed the earlier \(S_z\) information. A consistent description requires a two-dimensional complex vector space, in which

\begin{centeredalign}
\ket{S_x;\pm}=\frac{1}{\sqrt2}\bigl(\ket{+}\pm\ket{-}\bigr),\qquad
\ket{S_y;\pm}=\frac{1}{\sqrt2}\bigl(\ket{+}\pm i\ket{-}\bigr).
\end{centeredalign}

Here \(\ket{\pm}\) are the \(S_z\) eigenstates. Each later filter passes half of the incoming atoms, so one quarter of the atoms leaving the initial filter emerge in the \(S_z=-\hbar/2\) spot. The factor \(i\) in the \(S_y\) states is forced, because real coefficients cannot give probability \(1/2\) between every pair of the three axes. From here Sakurai introduces kets, bras, operators, matrix representations, the density operator and the commutator \([S_i,S_j]=i\hbar\,\epsilon_{ijk}S_k\) in the simplest system that has them.

The later chapters follow the same order of ideas. Quantum dynamics comes next, in the Schrödinger and Heisenberg pictures, with propagators, Feynman's path integral, gauge transformations and the Aharonov--Bohm effect. Angular momentum is developed as the theory of rotations, with Clebsch--Gordan coefficients, the Wigner--Eckart theorem, and spin correlations with Bell's inequality. Symmetries including parity and time reversal, approximation methods, scattering and identical particles complete the core.

The book is a common graduate text. Its starting point makes the incompatibility of measurements the organizing fact, whereas Griffiths' Introduction to Quantum Mechanics starts from the wave equation. Readers can continue with Spin (physics), Pauli matrices, Angular momentum operator, Heisenberg picture and Path integral formulation.







\clearpage

\section{Applications of quantum mechanics}

\label{topic:applications-of-quantum-mechanics}

\par\smallskip\noindent\textit{Historical, bibliographic or interpretive entry.}

Several widespread technologies work only because energy levels are discrete and transitions between them follow quantum rules. Lasers, magnetic resonance imaging, atomic clocks and transistors are the best known. Two of them show the physics in a single relation each.

Laser. In 1917 Einstein showed that radiation can stimulate an atom in its upper level \(2\) to emit a photon of the same frequency, direction and phase as the incident one. Requiring thermal equilibrium with Planck's radiation law fixes the ratio of spontaneous to stimulated emission. Amplification then requires a positive gain coefficient \(g(\nu)\):

\begin{centeredalign}
\frac{A_{21}}{B_{21}}=\frac{8\pi h\nu^3}{c^3},\qquad
g(\nu)=\sigma(\nu)\Bigl(N_2-\frac{g_2}{g_1}N_1\Bigr)>0 .
\end{centeredalign}

Here \(A_{21}\) is the spontaneous emission rate, \(B_{21}\) the stimulated emission coefficient per unit spectral energy density, \(\sigma(\nu)\) the transition cross section, \(N_{1,2}\) the level populations and \(g_{1,2}\) their degeneracies. In thermal equilibrium \(N_2/g_2<N_1/g_1\), so gain requires a population inversion produced by pumping through other levels. In a helium--neon laser, collisions with excited helium atoms pump the neon level whose decay emits at \(632.8\) nm. The \(\nu^3\) factor makes short-wavelength lasers harder to build, since spontaneous emission drains the inversion faster.

Magnetic resonance. A proton has spin \(1/2\) and a magnetic moment. In a field \(B\) its two spin states differ in energy by \(\hbar\gamma B\), where \(\gamma\) is the gyromagnetic ratio. The spin precesses at the Larmor frequency, and thermal equilibrium at temperature \(T\) gives a small net polarization:

\begin{centeredalign}
\omega_L=\gamma B,\qquad \frac{\gamma}{2\pi}=42.58\ \mathrm{MHz\,T^{-1}},\qquad
\frac{N_\uparrow-N_\downarrow}{N}\approx\frac{\hbar\gamma B}{2k_BT}.
\end{centeredalign}

In a \(1.5\) T scanner the resonance lies at \(63.9\) MHz, and at body temperature the net polarization is about five protons per million. The signal is still measurable because water contains about \(6.7\times10^{28}\) protons per cubic metre. A magnetic field gradient makes the frequency depend on position, which forms the image. Relaxation times of the precessing magnetization differ between tissues and provide contrast.

Atomic clocks use the same resonance principle with a hyperfine transition; the SI second is defined by the caesium-133 frequency \(9\,192\,631\,770\) Hz. Transistors depend on the \(1.12\) eV band gap of silicon and on the Fermi--Dirac occupation of electron states. Readers can continue with Quantum optics, Spin (physics), Two-state quantum system, Quantum metrology and Fermi--Dirac statistics.







\clearpage

\section{Werner Heisenberg}

\label{topic:werner-heisenberg}

\par\smallskip\noindent\textit{Historical, bibliographic or interpretive entry.}

In 1925 the energies of atoms were known from spectral lines, but electron orbits had never been observed. Werner Heisenberg (1901--1976) built a mechanics from observable quantities alone: transition frequencies \(\omega_{nm}=(E_n-E_m)/\hbar\) and radiation amplitudes between pairs of states. A coordinate became an array \(x_{nm}e^{i\omega_{nm}t}\). The product of two arrays had to respect the combination rule of spectral frequencies:

\begin{centeredalign}
(xy)_{nm}=\sum_k x_{nk}\,y_{km},\qquad \omega_{nk}+\omega_{km}=\omega_{nm} .
\end{centeredalign}

Born recognized this rule as matrix multiplication, which does not commute. With Jordan he showed that the quantum condition takes the form \(xp-px=i\hbar\), and the 1926 paper of Born, Heisenberg and Jordan developed the theory systematically. Heisenberg's July 1925 paper already obtained the oscillator energies \((n+\tfrac12)\hbar\omega\), including the zero-point term absent from the old quantum theory. An operator \(A\) evolves by \(dA/dt=(i/\hbar)[H,A]\), the equation of the Heisenberg picture.

In 1926 Heisenberg explained the two spectra of helium by the exchange symmetry of its two electrons, and in 1928 he traced ferromagnetism to the same exchange interaction. His 1927 paper on uncertainty argued, with a gamma-ray microscope, that position and momentum cannot both be sharply defined. Kennard proved the precise form for standard deviations the same year:

\begin{centeredalign}
\sigma_x\sigma_p\geq\frac{\hbar}{2},\qquad
\langle E_{\rm kin}\rangle\geq\frac{\sigma_p^2}{2m}\geq\frac{\hbar^2}{8m\sigma_x^2} .
\end{centeredalign}

An electron confined to \(\sigma_x=0.1\) nm therefore has a mean kinetic energy of at least \(0.95\) eV, which sets the electronvolt scale of atomic binding. Heisenberg received the 1932 Nobel Prize in Physics for the creation of quantum mechanics. He later introduced isospin (1932) and the S-matrix (1943), and during the Second World War he worked in the German nuclear fission programme. Readers can continue with Heisenberg picture, Uncertainty principle, Canonical commutation relation, Commutator and S-matrix.

\subsection*{References For The Physical Derivation}
\begin{itemize}
\item \href{https://arxiv.org/abs/quant-ph/0404009}{I. J. R. Aitchison, D. A. MacManus and T. M. Snyder, Understanding Heisenberg's 'Magical' Paper of July 1925; reconstruction of the calculations of the 1925 paper for an anharmonic oscillator.}
\end{itemize}





\clearpage

\section{The Quantum Universe}

\label{topic:the-quantum-universe}

\par\smallskip\noindent\textit{Historical, bibliographic or interpretive entry.}

Brian Cox and Jeff Forshaw's The Quantum Universe (2011) explains quantum mechanics to general readers without the wave equation. Following Feynman's popular account of quantum electrodynamics, it assigns a small clock to every path a particle could take. The hand of the clock turns as the particle moves, and the clocks of all paths arriving at a point are added like arrows. The squared length of the resulting arrow gives the probability of finding the particle there.

In equations, the clock is the phase \(e^{iS/\hbar}\), where \(S\) is the classical action of the path. For a free particle of mass \(m\) that covers a distance \(D\) in time \(T\), the action, the amplitude \(A\) and the probability \(P\) are

\begin{centeredalign}
S=\frac{mD^2}{2T},\qquad A\propto\sum_{\rm paths}e^{iS/\hbar},\qquad P\propto|A|^2 .
\end{centeredalign}

Changing \(D\) by a small amount \(\delta D\) changes the phase by \(p\,\delta D/\hbar\), with \(p=mD/T\) the momentum. Two slits a distance \(d\) apart, with a screen at distance \(L\), give paths to the screen position \(y\) that differ in length by about \(yd/L\). The arrows realign each time this phase difference grows by \(2\pi\), so the bright fringes are spaced by

\begin{centeredalign}
\Delta y=\frac{2\pi\hbar L}{pd}=\frac{\lambda L}{d},\qquad \lambda=\frac{h}{p}.
\end{centeredalign}

The de Broglie wavelength thus follows from the clock rule. Contributions of neighbouring paths cancel except near paths of stationary action, which recovers classical trajectories when \(S\gg\hbar\). The book uses the same arithmetic for the uncertainty principle, atomic energy levels, the Pauli exclusion principle, band structure and the transistor. Its closing chapters use electron degeneracy pressure to estimate the largest mass of a white dwarf star, about 1.4 solar masses. Readers can continue with Path integral formulation, Wave interference, Wave--particle duality and Fermi--Dirac statistics.







\backmatter


\chapter{Notation Of The Mechanism Descriptions}\label{chap:notation}
The opening spin calculation gives a concrete meaning to the notation
used in the companion manuscripts. The operation is commutator
evolution, and the carrier is the two-spin operator space. Together
they specify the kind of motion being considered. The Hamiltonian
and its domain, the measured polarization and the preparation then
determine a particular prediction. A physical pair of spins supplies
the coupling and the duration of the experiment. These descriptions
are written
\begin{centeredalign}
M=(\Omega,\Xi),\qquad I_{\mathrm{op}}=(M;C,R,P),\qquad
I_{\mathrm{real}}=(I_{\mathrm{op}};A).
\end{centeredalign}
The nesting expresses dependence, not a sequence of instructions.
The later description contains the earlier one together with the
information needed to calculate a more specific consequence. It
also permits shared inner descriptions to appear in several physical
realizations.

For the spins, \(\Xi\) contains density operators on
\(\mathbb C^2\otimes\mathbb C^2\). The operation \(\Omega\)
is the Hamiltonian-commutator class, realized here by
\(L\rho=-i[\hbar gZ_1Z_2,\rho]/\hbar\).
\(C\) includes positivity, unit trace and the stated domain.
\(R\) specifies the polarization
\(\operatorname{Tr}(\rho X_1)\); \(P\) fixes the initial
joint state, the evolution interval and any applied pulses.
The particular coupling \(g\) and apparatus belong to
\(A\). Changing \(g\) within this family changes the
oscillation frequency. Replacing the interaction by another
Hamiltonian changes \(L\), while measuring \(Z_1\) changes
\(R\) and the set of equations needed for the prediction.

For a particle, \(\Omega\) can identify a differential
operation such as the kinetic Laplacian, but the generator
is that operation on the domain fixed by \(\Xi\) and
\(C\). Dirichlet and Neumann conditions therefore give
different generators even when the bulk formula matches.
Exchange symmetry and gauge constraints similarly change
the admissible state space. A physical modification that
changes these choices belongs in the corresponding inner
description, as well as in the description of the device.

Comparing two mechanisms requires maps of like mathematical
objects. For linear state evolutions, a map \(\alpha\)
from source states to target states gives
\begin{centeredalign}
\Delta_\alpha=L_t\alpha-\alpha L_s,\qquad
R_t\alpha=\beta R_s.
\end{centeredalign}
The first relation measures failure to preserve evolution.
The second compares the measured quantities, with
\(\beta\) translating their units or outcome labels.
The two-spin reduction to \(x\) retains the immediate
measurement but loses the rate determined by \(c\).
It thus illustrates why the two comparisons are distinct.

For a more general relation \(\Omega_s:X_s\to Y_s\),
the source and target can have different input and output
spaces. The general operator defect is then
\begin{centeredalign}
\Delta_{\alpha,\beta}=\Omega_t\alpha-\beta\Omega_s,
\qquad \alpha:X_s\to X_t,\quad \beta:Y_s\to Y_t.
\end{centeredalign}
For the preceding linear evolution comparison, both output
spaces are state spaces and \(\beta=\alpha\); the defect
reduces to \(\Delta_\alpha\). A nonlinear coordinate change
uses a derivative or a pullback on test functions instead
of ordinary multiplication by a linear state map. These
choices give the notation its physical meaning: the
comparison is made between specified evolutions or
relations, with their domains and measured consequences.


\chapter{Sources And Reproduction}

Equation Sources links corpus-aligned relations for 6 topics. Relations In The Original Papers links 16 displays from 7 articles across 15 topics. The cited relations and their conditions appear beside the relevant treatment. The mechanism tree and cross-topic comparisons are the authors' explanatory organization of established theory. The edition manifest records the individual equation and article identifiers. Independent scientific review of every chapter remains to be completed.

\end{document}